\documentclass[9pt,twoside]{extarticle}
\usepackage{Breviary-Appendix}
%%
\begin{document}
\renewcommand{\thepage}{\arabic{page}**}
\begin{center} \includegraphics[width=1 \textwidth]{20251031_131444495} \end{center}
%\frontmatter
\subsection*{\center \textred{\textsc{DIE 12 DECEMBRIS}}}
\section*{\center \textred{B. M}ariæ \textred{V}irg. de \textred{G}uadalupe}
\subsubsection*{\center \textred{DUPLEX MAJUS I CLASSIS}}
    \fancyhead[C]{\textsc{12 Dec. - B. M. V. de Guadalupe.}}
\begin{multicols}{2}
\begin{tightlist}{}
\item {\small \textred{Omnia de Communi Festorum B. Mariæ Virginis [161], præter ea quæ hic habentur propria.}}
\end{tightlist}

\subsubsection*{\center \textsc{IN I VESPERIS}}
\begin{tightlist}{}
\item {\small \textred{Antiphonæ de Laudibus, Psalmi ut in I Vesperis de Communi [161]; Capitulum, Responsorium et Hymnus item de Communi.}}
\item \textred{℣.} Non fecit táliter omni natióni. \textred{℟.} Et judícia sua non manifestávit eis, allelúja.
\item \textred{Ad Magnif. Ant.} Elégi et sanctificávi * locum istum, ut sit ibi nomen meum, et permáneant óculi mei et cor meum ibi cunctis diébus.
\end{tightlist}

\subsubsection*{\center \textred{Oratio}}
\lettrine{D}eus, qui sub beatíssimæ Vírginis Maríæ singulári patrocínio constitútos perpétuis benefíciis nos cumulári voluísti: præsta supplícibus tuis; ut, cujus hódie commemoratióne lætámur in terris, ejus conspéctu perfruámur in cælis. Per Dóminum.

\begin{tightlist}{}
\item \smallred{Fit Commemoratio Feriæ.}
\item \smallred{Completorium de Dominica.}
\item \smallred{In fine Hymnorum Completorii et Horarum dicitur:}
\item \textred{G}lória tibi, Dómine,
\item Qui natus es de Vírgine,
\item Cum Patre, et Sancto Spíritu,
\item In sempitérna sǽcula.
\item {\small \textred{Ad Nunc dimittis Ant.} Allelúja, Sancta.}
\end{tightlist}

\subsubsection*{\center \textsc{AD MATUTINUM}}
\begin{tightlist}{}
\item \textred{Invitat.} In honórem beatíssimæ Maríæ Vírginis * Jubilémus Dómino.
\textred{Ps. 94} Veníte, exsultémus. 2.
\end{tightlist}

\subsubsection*{\center \textred{Hymnus}}
\lettrine{P}ræclára custos Vírginum,\\
Intácta Mater Núminis,\\
Cæléstis aulæ jánua,\\
Spes nostra, cæli gáudium.\\
\\
\textred{I}nter rubéta lílium,\\
Colúmba formosíssima,\\
Virga e radíce gérminans\\
Nostro medélam vúlneri,\\
\\
\textred{T}urris dracóni impérvia,\\
Amíca stella náufragis,\\
Tuére nos a fráudibus,\\
Tuáque luce dírige.\\
\\
\textred{E}rróris umbras díscute,\\
Syrtes dolósas ámove,\\
Fluctus tot inter, déviis\\
Tutam reclúde sémitam.\\
\\
\textred{G}lória tibi, Dómine,\\
Qui natus es de Vírgine,\\
Cum Patre, et Sancto Spíritu,\\
In sempitérna sǽcula.\\
\textred{A}men.

\subsubsection*{\center \textsc{IN I NOCTURNO}}
\begin{tightlist}{}
\item \textred{Ant.} A solis ortu * usque ad occásum, magnum est nomen meum in géntibus.
\item {\small \textred{Psalmi de Communi [165].}}
\item \textred{Ant.} Ego in altíssimis * habitávi, et thronus meus in colúmna nubis.
\item \textred{Ant.} Invénta sum * a non quæréntibus me; palam appárui iis, qui me non interrogábant.
\item \textred{℣.} Notum fecit Dóminus salutáre suum. \textred{℟.} In conspéctu géntium revelávit glóriam suam, allelúja.
\item {\small \textred{Bened.} Alma Virgo vírginum intercédat pro nobis ad Dóminum.}
\end{tightlist}

\subsubsection*{\center \normalfont De libro Ecclesiástici}
\subsubsection*{\textred{Lectio j \hfill Cap. 24, 5-31}}
\lettrineThree{E}go ex ore Altíssimi prodívi, primogénita ante omnem creatúram:
ego feci in cælis ut orirétur lumen indefíciens, et sicut nébula texi omnem terram:
ego in altíssimis habitávi, et thronus meus in colúmna nubis.
Gyrum cæli circuívi sola, et profúndum abýssi penetrávi, in flúctibus maris ambulávi,
et in omni terra steti: et in omni pópulo,
et in omni gente primátum hábui:
et ómnium excelléntium et humílium corda virtúte calcávi: et in his ómnibus réquiem quæsívi, et in hereditáte Dómini morábor.
Tunc præcépit, et dixit mihi Creátor ómnium: et qui creávit me, requiévit in tabernáculo meo,
et dixit mihi: In Jacob inhábita, et in Israël hereditáre, et in eléctis meis mitte radíces.

\begin{tightlist}{}
\item \textred{℟.} Vidi speciósam sicut colúmbam ascendéntem désuper rivos aquárum: cujus inæstimábilis odor erat nimis in vestiméntis ejus: \textred{*} Et sicut dies verni circúmdabant eam flores rosárum, et lília convállium.\\ \textred{℣.} Quæ est ista, quæ ascéndit per desértum sicut vírgula fumi ex aromátibus myrrhæ et thuris? \textred{E}t.
\item {\small \textred{Bened.} Nos cum prole pia benedícat Virgo María.}
\end{tightlist}

\subsubsection*{\center \textred{Lectio ij}}
\lettrine{A}b inítio, et ante sǽcula creáta sum, et usque ad futúrum sǽculum non désinam, et in habitatióne sancta coram ipso ministrávi.
Et sic in Sion firmáta sum, et in civitáte sanctificáta simíliter requiévi, et in Jerúsalem potéstas mea.
Et radicávi in pópulo honorificáto, et in parte Dei mei heréditas illíus, et in plenitúdine sanctórum deténtio mea.
Quasi cedrus exaltáta sum in Líbano, et quasi cypréssus in monte Sion:
quasi palma exaltáta sum in Cades, et quasi plantátio rosæ in Jéricho:
quasi olíva speciósa in campis, et quasi plátanus exaltáta sum juxta aquam in platéis.
Sicut cinnamómum, et bálsamum aromatízans odórem dedi: quasi myrrha elécta dedi suavitátem odóris:
et quasi storax, et gálbanus, et úngula, et gutta, et quasi Líbanus non incísus vaporávi habitatiónem meam, et quasi bálsamum non mixtum odor meus.

\begin{tightlist}{}
\item \textred{℟.} Quæ est ista, quæ ascéndit sicut auróra consúrgens, pulchra ut luna, elécta ut sol, \textred{*} Terríbilis ut castrórum ácies ordináta? 
\textred{℣.} Fília Sion, tota formósa et suávis es, pulchra ut luna, elécta ut sol. \textred{T}erríbilis.
\item {\small \textred{Bened.} Sancta Dei Génitrix sit nobis auxiliátrix.}
\end{tightlist}

\subsubsection*{\center \textred{Lectio iij}}
\lettrine{E}go quasi terebínthus exténdi ramos meos, et rami mei honóris et grátiæ.
Ego quasi vitis fructificávi suavitátem odóris: et flores mei fructus honóris et honestátis.
Ego mater pulchræ dilectiónis, et timóris, et agnitiónis, et sanctæ spei.
In me grátia omnis viæ et veritátis, in me omnis spes vitæ et virtútis.
Transíte ad me omnes qui concupíscitis me, et a generatiónibus meis implémini:
spíritus enim meus super mel dulcis, et heréditas mea super mel et favum:
memória mea in generatiónes sæculórum.
Qui edunt me, adhuc esúrient: et qui bibunt me, adhuc sítient.
Qui audit me, non confundétur: et qui operántur in me, non peccábunt.
Qui elúcidant me, vitam ætérnam habébunt.

\begin{tightlist}{}
\item \textred{℟.} Quæ est ista, quæ procéssit sicut sol, et formósa tamquam Jerúsalem? \textred{*} Vidérunt eam fíliæ Sion, et beátam dixérunt, et regínæ laudavérunt eam. \textred{℣.} Et sicut dies verni circúmdabant eam flores rosárum, et lília convállium. \textred{V}idérunt. Glória Patri. \textred{V}idérunt.
\end{tightlist}

\subsubsection*{\center \textsc{IN II NOCTURNO}}
\begin{tightlist}{}
\item \textred{Ant.} Signum magnum * appáruit in cælo: múlier amícta sole, et luna sub pédibus ejus.
\item \textred{Ant.} Pópulus qui ambulábat * in ténebris, vidit lucem magnam: habitántibus in regióne umbræ mortis, lux orta est eis.
\item \textred{Ant.} Vidérunt * omnes términi terræ salutáre Dei nostri.
\item \textred{℣.} Introíbimus in tabernáculum ejus. \textred{℟.} Adorábimus in loco ubi stetérunt pedes ejus, allelúja.
\item {\small \textred{Bened.} Précibus suæ Matris benedícat nos Fílius Dei Patris.}
\end{tightlist}

\subsubsection*{\center \textred{Lectio iv}}
\lettrineThree{A}nno a reparáta salúte millésimo quingentésimo trigésimo primo, Deípara Virgo Joánni Dídaco, pio rudíque neóphyto, Méxici in colle Tepejacénsi, uti antíqua et constánti traditióne mandátur, sese vidéndam prǽbuit, eúmque peramánter álloquens, epíscopum adíre, ipsíque ædem ínibi sibi sacram excitándam nuntiáre jussit. Rei veritátem solérter exploratúrus, responsiónem dístulit Joánnes de Zumárraga, loci antístes: ac neóphyto, beatíssimæ Vírginis aspéctu imperióque rursus permóto, legationémque lácrimis et précibus iteránti, præcépit, ut signum, quo magnæ Dei Matris ostenderétur volúntas, eníxe péteret.

\begin{tightlist}{}
\item \textred{℟.} Signum magnum appáruit in cælo, múlier amícta sole et luna sub pédibus ejus, \textred{*} Et in cápite ejus coróna duódecim stellárum.
\textred{℣.} Fundaménta ejus in móntibus sanctis, díligit Dóminus portas Sion super ómnia tabernácula Jacob. \textred{E}t.
\item {\small \textred{Bened.} Jesus Maríæ Fílius sit nobis clemens et propítius.}
\end{tightlist}

\subsubsection*{\center \textred{Lectio v}}
\lettrine{D}um viam a colle Tepejacénsi remotiórem neóphytus arríperet, et Méxicum, ne pátruo letháli morbo corrépto extréma deéssent sacraménta, sacerdótem advocatúrus conténderet, eídem benigníssima Virgo tértio occúrrit; mæréntem de pátrui sanitáte certiórem facit; rosásque pulchérrimas, contra loci hiemísque asperitátem, recens obórtas, in ejus pállio compósitas, epíscopo deférri injúngit. Mandátis obséquitur Dídacus, cujus in pállio, rosis coram epíscopo effúsis, Maríæ sanctæ imágo qua prorsus spécie in colle prope urbem sese exhibúerat, mirum in modum depícta conspícitur. Tanto percúlsi prodígio religiósam cives effígiem in episcopáli sacéllo rite asservándam curant: quæ, paulo post, solémni pompa ad ædem in collo Tepejacénsi exstrúctam transláta, insígni géntium cunctárum veneratióne præfúlget.

\begin{tightlist}{}
\item \textred{℟.} Quæ est ista, quæ progréditur quasi auróra consúrgens: \textred{*} Pulchra ut luna, elécta ut sol?
\textred{℣.} Quasi arcus refúlgens inter nébulas glóriæ, et quasi flos rosárum in diébus vernis. \textred{P}ulchra.
\item {\small \textred{Bened.} Fílius Vírginis Maríæ det nobis gáudia vitæ.}
\end{tightlist}

\subsubsection*{\center \textred{Lectio vj}}
\lettrine{M}agnífico deínceps excépta templo, cui canonicórum collégium Románi Pontífices ad divíni cultus splendórem addidérunt, Mexicánæ gentis pietátem in Deíparam summópere adáuxit; ingentíque abhinc quátuor sǽculis cólitur populórum ac miraculórum frequéntia. Quaprópter eam uti præsentíssimum advérsus públicas privatásque calamitátes præsídium archiepíscopus Mexicánus ceteríque illárum regiónum antístites, ómnium órdinum consensióne, in primáriam adlegérunt univérsæ Mexicánæ natiónis patrónam, ritéque eléctam Benedíctus décimus quartus auctoritáte apostólica declarávit, atque Offícium et Missam sub título Beátæ Vírginis Guadalupénsis recitári concéssit. Leo vero décimus tértius, iterátis Mexicanórum prǽsulum petitiónibus benígne ánnuens, novíssimum hoc Offícium, ex sacrórum Rítuum Congregatiónis consúlto, recitári indúlsit, Virginísque effígiem prodígiis atque cultu célebrem áurea coróna, suo nómine et jussu, solémni ritu condecorári decrévit. Quam quidem quarto Idus Octóbris anno millésimo octingentésimo nonagésimo quinto, coram epíscopis fere ómnibus reipúblicæ et áliis bene multis ex réliquis Américæ pártibus, máxima incolárum advenarúmque frequéntia et gratulatióne archiepíscopus Mexicánus pretiosíssimam impósuit.

\begin{tightlist}{}
\item \textred{℟.} Elégi et sanctificávi locum istum: \textred{*} Ut sit ibi nomen meum, et permáneant óculi mei, et cor meum ibi cunctis diébus.
\textred{℣.} Non fecit táliter omni natióni, et judícia sua non manifestávit eis. \textred{U}t. Glória Patri. \textred{U}t.
\end{tightlist}

\subsubsection*{\center \textsc{IN III NOCTURNO}}
\begin{tightlist}{}
\item \textred{Ant.} Tabernáculum Dei * María, in médio civitátis ejus, et non commovébitur.
\item \textred{Ant.} Egréssa es * in salútem pópuli tui, in salútem cum Christo tuo.
\item \textred{Ant.} Gloriósa dicta sunt * de te, cívitas Dei: fundávit te Dóminus in móntibus sanctis.
\item \textred{℣.} Fílii tui de longe vénient. \textred{℟.} Et fíliæ tuæ de látere surgent, allelúja.
\item {\small \textred{Bened.} Fons Evangélii répleat nos dógmate cæli.}
\end{tightlist}

\subsubsection*{\center  \normalfont Léctio sancti Evangélii secúndum Lucam}
\subsubsection*{\center \textred{Lectio vij \hfill Cap. I, 39-47}}
\lettrine{I}n illo témpore: Exsúrgens María ábiit in montána cum festinatióne, in civitátem Juda. Et intrávit in domum Zacharíæ, et salutávit Elísabeth. Et réliqua.

\begin{center}Homilía sancti Bernárdi Abbátis.\\
\textred{De verbis Apoc. cap. 12, ad finem}\end{center}
\lettrineThree{M}irabátur Elísabeth persónam veniéntis, dicens: Unde hoc mihi, ut véniat Mater Dómini mei ad me? Commendábat et vocem salutántis, adjíciens: Ut facta est vox salutatiónis tuæ in áuribus meis, exsultávit in gáudio infans in útero meo: et beatificábat fidem credéntis: Beáta, ínquiens, quæ credidísti. Magna quidem præcónia; sed et devóta humílitas nihil sibi passa retinére, in eum magis univérsa refúdit, cujus in se benefícia laudabántur.

\begin{tightlist}{}
\item \textred{℟.} Felix namque es, sacra Virgo María, et omni laude digníssima: \textred{*} Quia ex te ortus est sol justítiæ, Christus, Deus noster.
\textred{℣.} Ora pro pópulo, intérveni pro clero, intercéde pro devóto femíneo sexu: séntiant omnes tuum juvámen, quicúmque célebrant tuam sanctam festivitátem. \textred{Q}uia.
\item {\small \textred{Bened.} Ab hoste malígno deféndat nos Dei Génitrix Virgo.}
\end{tightlist}

\subsubsection*{\center \textred{Lectio viij}}
\lettrine{T}u, inquit, magníficas Matrem Dómini, sed magníficat ánima mea Dóminum. In voce mea fílium pérhibes exsultásse in gáudio, sed exsultávit spíritus meus in Deo, salutári meo: et ipse quoque tamquam amícus sponsi gaudet ad vocem sponsi. Beátam me esse dicis, quæ crédidi; sed credulitátis et beatitúdinis causa respéctus est supérnæ pietátis; ut ex hoc magis beátam me dicant omnes generatiónes, quia ancíllam húmilem et exíguam respéxit Deus. Verúmtamen numquid putámus, fratres, Elísabeth sanctam in eo, quod per Spíritum útique loquebátur errásse? Absit. Beáta plane quam respéxit Deus, et beáta quæ crédidit. Hic enim magnus divínæ respectiónis éxstitit fructus.

\begin{tightlist}{}
\item \textred{℟.} Beátam me dicent omnes generatiónes: \textred{*} Quia fecit mihi Dóminus magna qui potens est, et sanctum nomen ejus.
\textred{℣.} Et misericórdia ejus a progénie in progénies timéntibus eum. \textred{Q}uia.
\item {\small \textred{Bened.} In omni tribulatióne et angústia succúrrat nobis pia Virgo María.
\item \textred{Si dicenda sit ix Lectio de Homilia Dominicæ, erit:}
\item \textred{Bened.} Verba Redemptóris pia sint suffrágia nobis.}
\end{tightlist}

\subsubsection*{\center \textred{Lectio ix}}
\lettrine{I}neffábili síquidem artifício Spíritus supervéniens tantæ humilitáti, magnanímitas tanta in secretário virgínei cordis accéssit, ut quod vidélicet nec humílitas tanta minúerit magnanimitátem, nec magnanímitas tanta humilitátem. Sed cum in sua æstimatióne tam húmilis esset, nihilóminus et in promissiónis credulitáte magnánimis, ut quæ nihil áliud quam exíguam sese reputábat ancíllam, ad inscrutábile sacraméntum nullátenus se dubitáret eléctam, et veram Dei et hóminis Genitrícem créderet mox futúram. Jam te, Mater misericórdiæ, per ipsum sinceríssimæ tuæ mentis afféctum tuis jacens provolúta pédibus luna, mediatrícem sibi apud solem justítiæ constitútam, devótis supplicatiónibus interpéllat, ut in lúmine tuo vídeat lumen, et solis grátiam tuo mereátur obténtu, quam vere amávit præ ómnibus, et ornávit stola glóriæ índuens, et corónam pulchritúdinis ponens in cápite tuo.

\begin{tightlist}{}
\item \textred{℟.} Christi Virgo dilectíssima, virtútum operátrix, opem fer míseris: \textred{*} Súbveni, Dómina, clamántibus ad te júgiter. \textred{℣.} Quóniam peccatórum mole prémimur, et non est qui ádjuvet. \textred{S}úbveni. Glória Patri. \textred{S}úbveni.
\item Te Deum laudámus. 6.
\item \textred{Sacerd. ℣.} Ora pro nobis, sancta Dei Génitrix. \textred{℟.} Ut digni efficiámur promissiónibus Christi, allelúja.
\end{tightlist}

\subsubsection*{\center \textsc{AD LAUDES}}
\begin{center}\textred{et per Horas}\end{center}
\begin{tightlist}{}
\item \textred{Ant. 1} Quæ est ista * quæ procéssit sicut sol, et formósa tamquam Jerúsalem.
\item {\small \textred{Psalmi de Dominica I loco.}}
\item \textred{2} Vidérunt eam * fíliæ Sion vernántem in flóribus rosárum, et beatíssimam prædicavérunt.
\item \textred{3} Flores apparuérunt * in terra nostra: te laudámus, sancta Dei Génitrix.
\item \textred{4} Cantáte ei * cánticum novum: annuntiáte inter gentes glóriam ejus.
\item \textred{5} Tu glória * Jerúsalem, tu lætítia Israel, tu honorificéntia pópuli tui.
\end{tightlist}

\subsubsection*{\center \textred{Capitulum \hfill Eccli. 24, 23-24}}
\lettrine{E}go quasi vitis fructificávi suavitátem odóris: et flores mei fructus honóris et honestátis. Ego mater pulchræ dilectiónis, et timóris, et agnitiónis, et sanctæ spei.

\subsubsection*{\center \textred{Hymnus}}
\lettrine{A}lma Parens, omni carens\\
Corruptélæ mácula,\\
Quam elégit qui confrégit\\
Diræ mortis víncula.\\
\\
\textred{C}lemens, præsta, ut qui festa\\
Tua gaudent cólere,\\
Gratuléntur et læténtur\\
Veræ lucis lúmine.\\
\\
\textred{P}rimum quidem nobis fidem\\
Tuis auge précibus,\\
Spe labéntes firma mentes,\\
Caritáte róbora.\\
\\
\textred{M}ater Dei, nostræ spei\\
Causa et exórdium,\\
Aufer bella et flagélla,\\
Famem, pestem, gládium.\\
\\
\textred{E}genórum et vinctórum\\
Consoláre gémitum:\\
Fave votis, dans ægrótis\\
Optátum remédium.\\
\\
\textred{P}ax et quies nostros dies\\
Fáciant lætíssimos:\\
Inimícos fac amícos,\\
Seda malos hómines.\\
\\
\textred{O} María, Mater pia,\\
Tuum da subsídium,\\
Ut regnémus et laudémus\\
Tuum semper Fílium.\\
\\
\textred{P}ater Deus, Fili Deus,\\
Deus alme Spíritus,\\
Per ætérna nos gubérna\\
Sǽcula, Deus Trínitas.\\
\textred{A}men.

\begin{tightlist}{}
\item \textred{℣.} Omnis terra adóret te, et psallat tibi. \textred{℟.} Psalmum dicat nómini tuo, allelúja.
\item \textred{Ad Bened. Ant.} Quasi arcus * refúlgens inter nébulas glóriæ, et quasi flos rosárum in diébus vernis.
\end{tightlist}

\subsubsection*{\center \textred{Oratio}}
\lettrine{D}eus, qui sub beatíssimæ Vírginis Maríæ singulári patrocínio constitútos perpétuis benefíciis nos cumulári voluísti: præsta supplícibus tuis; ut, cujus hódie commemoratióne lætámur in terris, ejus conspéctu perfruámur in cælis. Per Dóminum.

\begin{tightlist}{}
\item {\small \textred{Et non fit Commemoratio Octavæ.}
\item \textred{Ad Horas Antiphonæ de Laudibus, Psalmi de Dominica, ad Primam tamen ut in Festis et in Responsorio ℣.} Qui de Vírgine dignátus es nasci. \textred{Reliqua ut in Communi [177].}}
\end{tightlist}

\subsubsection*{\center \textsc{IN II VESPERIS}}
\begin{tightlist}{}
\item {\small \textred{Antiphonæ de Laudibus, Psalmi, Capitulum et Hymnus} Ave, maris stella, \textred{de Communi [177].}}
\item \textred{℣.} Non fecit táliter omni natióni. \textred{℟.} Et judícia sua non manifestávit eis, allelúja.
\item \textred{Ad Magnif. Ant.} Leva in circúitu * óculos tuos et vide: omnes isti congregáti sunt, venérunt tibi: fílii tui de longe vénient, et fíliæ tuæ de látere surgent.
\item {\small \textred{Fit Commemoratio sequentis:}}
\item \textred{Ant.} In patiéntia tua * possedísti ánimam tuam, Lúcia, sponsa Christi: odísti quæ in mundo sunt, et corúscas cum Angelis: sánguine próprio inimícum vicísti.
\item \textred{℣.} Diffúsa est grátia in lábiis tuis. \textred{℟.} Proptérea benedíxit te Deus in ætérnum, allelúja.
\end{tightlist}

\subsubsection*{\center \textred{Oratio}}
\lettrine{E}xáudi nos, Deus, salutáris noster: ut, sicut de beátæ Lúciæ Vírginis et Mártyris tuæ festivitáte gaudémus; ita piæ devotiónis erudiámur afféctu. Per Christum.

\begin{tightlist}{}
\item {\small \textred{Et fit Commemoratio Feriæ.}}
\item {\small \textred{Completorium de Dominica. Ad Nunc dimittis Ant.} Allelúja, Sancta.}
\end{tightlist}

\end{multicols}

\begin{center}\textcolor{red}{\rule{\linewidth}{0.5pt}} \end{center}

\vfill
\begin{center} IMPRIMATUR.\\
Tornaci, die 11 Februarii 1939.\\
\textsc{J. Lecouvet}, Vic. Gen.\end{center}
\vfill

\fancyfoot[C]{{\small Imprimé en Belgique par la Société S. Jean l'Evangéliste,
Desclée & Cie, Tournai. - 5767}}
\fancyhead[C]{}
\fancyfoot[C]{}
\cleardoublepage
\pagenumbering{arabic}
\renewcommand{\thepage}{\arabic{page}*}
%\mainmatter
\section*{\center \textsc{Appendix}\\ \large \textit{\textred{Ex Breviario Carmelitarum 1866 - Pars Hiemalis}}}
\begin{multicols}{2} % Start the two-column layout
\subsection*{\center \textred{DIE XXIV. NOVEMBRIS.}}
\subsubsection*{\center \textred{INFRA OCTAVAM}}
\subsection*{\center \textsc{S. Joannis a Cruce.}}
%\subsubsection*{\center \textred{Duplex minus.}}
%\begin{center}\textred{Vide festum, supra, 535.}\end{center}
\fancyhead[C]{\textsc{24 Nov. - Infra Octav. S. Joannis a Cruce}}

\begin{tightlist}{}
\item \textred{Per Octavam S. Joannis, fit de ea comm. in Laud. et Vesp. per Añam ad Bened. et Añam ad Magnif. de ij. Vesperis, excepto Sabbato, in quo ad Vesp. fit comm. S. Joannis per Añam ad Magnificat de primis Vesperis Festi:} 
\item \textred{Aña ad Magnif. j. Vesp.} Collaudábunt multi sapiéntiam ejus, et usque in sǽculum non delébitur.
\item \textred{℣.} Ora pro nobis, sancte pater Joánnes. \textred{℟.} Ut digni efficiámur promissiónibus Christi, allelúja.
\item Orémus. \hfill \textred{Oratio.}
\end{tightlist}
\lettrine{D}eus, qui sanctum Joánnem Confessórem tuum, atque Doctórem, Patrem nostrum perféctæ sui abnegatiónis, et crucis amatórem exímium effecísti: concéde; ut ejus imitatióni júgiter inhæréntes, glóriam assequámur ætérnam. Per Dóminum.

\begin{tightlist}{}
\item \textred{Aña ad Bened.} Dedit illi sciéntiam Sanctórum, honestávit illum in labóribus, et complévit labóres illíus.
\item \textred{℣.} Justus germinábit sicut lílium. \textred{℟.} Et florébit in ætérnum ante Dóminum.
\item \textred{Aña ad Magnif. ij. Vesp.} Quæsívit verba utília, et conscrípsit sermónes rectíssimos, et veritáte plenos.
\item \textred{℣.} Ora pro nobis, sancte pater Joánnes. \textred{℟.} Ut digni efficiámur promissiónibus Christi.
\end{tightlist}

\begin{center}\textcolor{red}{\rule{\linewidth}{0.5pt}} \end{center}

\subsection*{\center \textsc{\textred{DIE XXVIII. NOV.}}}
\subsubsection*{\center \textred{Dies Octava}}
\subsection*{\center \textsc{Præsentationis B.M.V.}}
\subsubsection*{\center \textred{Duplex minus.}}
    \fancyhead[C]{\textsc{28 Nov - In Octava Præsentationis B.M.V.}}
\begin{tightlist}{}
\item \textred{Omnia ut in Commune Beatæ Mariæ Virginis} [161], \textred{exceptis hic signatis.}
\end{tightlist}
\subsubsection*{\center \textsc{AD I. VESPERAS.}}
\begin{tightlist}{}
\item \textred{Aña. prima de Laudibus} Pulchra es et decóra, fília Jerúsalem, terríbilis ut castrórum ácies ordináta. \textred{quæ sola dicitur super Ps. præc.}
\end{tightlist}
\subsubsection*{\center \textred{Oratio.}}
\lettrine{B}eátæ Maríæ semper Vírginis, tríbue nos Dómine supplicatióne tuéri: ut cujus venerábilem Præsentatiónem celebrámus obséquiis, ejus intercessiónibus et méritis commendémur. Per Dóminum.

%\begin{tightlist}{}
%\item \textred{℣.} Ora pro nobis sancte pater Joánnes. \textred{℟.} Ut digni efficiámur promissiónibus Christi, allelúja. 
%\item \textred{Ad Magnificat, Aña.} Quæsívit verba utília, et conscrípsit sermónes rectíssimos, et veritáte plenos.


\begin{tightlist}{}
\item \textred{Comm. Oct. S. Joannis. Aña.} Quæsívit, \textred{ut supra.} 
%Quæsívit verba utília, et conscrípsit sermónes rectíssimos, et veritáte plenos. 
%\item \textred{℣.} Ora pro nobis sancte Pater Joánnes. \textred{℟.} Ut digni efficiámur promissiónibus Christi, allelúja.
%\item Orémus. \hfill \textred{Oratio.}
\end{tightlist}
%\lettrine{D}eus, qui sanctum Joánnem Confessórem tuum atque Doctórem, Patrem nostrum, perféctæ sui abnegatiónis, et crucis amatórem exímium effecísti: concéde, ut ejus imitatióni júgiter inhæréntes, glóriam assequámur ætérnam. Per Christum.

\begin{tightlist}{}
\item \textred{Et fit Comm. Feriæ.}
%\item \textred{Comm. Oct. S. Joannis, Aña.} Quæsívit verba utília, et conscrípsit sermónes rectíssimos, et veritáte plenos.
%\item \textred{℣.} Ora pro nobis sancte Pater Joánnes. \textred{℟.} Ut digni efficiámur promissiónibus Christi.
\end{tightlist}
%\subsubsection*{\center \textred{Oratio.}}
%\lettrine{D}eus, qui sanctum Joánnem Confessórem tuum atque Doctórem, Patrem nostrum, perféctæ sui abnegatiónis, et crucis amatórem exímium effecísti: concéde, ut ejus imitatióni júgiter inhæréntes, glóriam assequámur ætérnam. Per Dóminum.

\begin{tightlist}{}
\item \textred{In fine Hymnorum Completorii et Horarum dicatur ℣.} Glória tibi Dómine, Qui natus es de Vírgine.
\item \textred{Ad Nunc dimittis, Aña.} Sub tuum præsídium.
\end{tightlist}

\subsubsection*{\center \textsc{AD MATUTINUM.}}
\begin{tightlist}{}
\item \textred{In primo Nocturno, Lectiones de Scriptura occurrente.} 
\item \textred{In secundo Nocturno, Lectiones de Sermone} Hílarem vídeo, \textred{ut in Commune.} [171]
\end{tightlist}

\subsubsection*{\center \textsc{IN TERTIO NOCTURNO.}}
\begin{tightlist}{}
\item \textred{In ultima Aña, loco} Festivitátem, \textred{dicitur} Præsentatiónem.
\item {\small \textred{Bened.} Fons Evangélii répleat nos dógmate cæli.}
\end{tightlist}

\subsubsection*{\center Léctio sancti Evangélii secúndum Matthǽum.}
\begin{center}\textred{\textbf{Lectio vij}.} \hfill \textred{Cap. I.}\end{center}
\lettrine{L}iber generatiónis Jesu Christi fílii David, fílii Abraham. Abraham génuit Isaac. Isaac autem génuit Jacob. Et réliqua. 

\begin{center}Homilía sancti Joánnis Damascéni.\\
\textred{Lib 4. de fide orthod.}\end{center}

\lettrineThree{Q}uod Joseph ex Davídica tribu descénderit, Matthǽus et Lucas sacratíssimi Evangelístæ dilúcide demonstrárunt. Sanctæ vero Vírginis generatiónem ambo siluérunt. Decet autem scire non fuísse morem Hebrǽis, neque divínæ Scriptúræ, genealogías téxere mulíerum. Nam lex fuit, tribum non desponsáre ex ália tribu. Atqui Joseph ex Davídica descéndit tribu, et justus erat: hoc enim de ipso divínum testátur Evangélium; non ígitur præter legem, sanctam Vírginem ad desponsatiónem duxísset, nisi ex ejus sceptro descendísset. Cum ígitur osténderit Evangelísta descénsum Joseph, satisfáctum esse putándum est.

\begin{tightlist}{}
\item \textred{℟.} Beátam me dicent omnes generatiónes: \textred{*} Quia fecit mihi Dóminus magna, qui potens est: et sanctum nomen ejus. \textred{℣.} Et misericórdia ejus a progénie in progénies timéntibus eum. \textred{Q}uia.
\item {\small \textred{Bened.} Ab hoste malígno deféndat nos Dei Génitrix Virgo.}
\end{tightlist}

\subsubsection*{\center\textred{Lectio viij.}}
\lettrine{E}x série ígitur Nathan fílii David, Levi génuit Melchi; et Panthérem: Panther génuit Barpanthérem, sic ab eo denominátum. Hic Barpánther génuit Jóachim : Jóachim vero génuit sanctam Deíparam. Jóachim ergo venerándam et dignam laude Annam in matrimónium duxit. Sed quemádmodum vetus Anna stérilis facta per votum, et repromissiónem hábuit Samuélem, hunc in modum et hæc per supplicatiónem, et repromissiónem a Deo Deíparam profert; ut et in hoc nulla illústrium, et gloriosárum postérior habeátur.

\begin{tightlist}{}
\item \textred{℟.} Felix namque es, sacra Virgo María, et omni laude digníssima: \textred{*} Quia ex te ortus est Sol justítiæ, Christus, Deus noster. \textred{℣.} Ora pro pópulo, intérveni pro clero, intercéde pro devóto femíneo sexu, séntiant omnes tuum juvámen, quicúmque célebrant tuam sanctam Præsentatiónem. \textred{Q}uia.
\item {\small \textred{Bened.} In omni tribulatióne et angústia succúrrat nobis pia Virgo María.}
\end{tightlist}

\subsubsection*{\center\textred{Lectio ix.}}
\lettrine{P}arit ergo grátia (Anna enim hoc interpretátur) Dóminam: quod profécto índicat Maríæ nomen. Nam revéra Dómina facta est ómnium creaturárum, cum Conditóris ómnium effecta est mater. Editur partu in domo ovílis Jóachim, et addúcitur in templum. Deínde in domo Dómini plantáta, et impinguáta Spíritu : véluti olíva fructífera, omnis virtútis habitáculum facta est, ab omni sæculári vita et carnáli concupiscéntia procul mentem sequéstrans et sic vírginem ánimam cum córpore consérvans, ut decébat eam, quæ in sinu Deum susceptúra esset. Sanctus enim exsístens, in sanctis requiéscit. Sic ígitur sanctificatiónem omnem exércet, et templum sanctum, et admirábile altíssimi Dei digne demonstrátur.

\begin{tightlist}{}
\item \textred{℟.} Christi Virgo dilectíssima, virtútum operátrix, opem fer míseris: \textred{*} Súbveni Dómina clamántibus ad te júgiter. \textred{℣.} Quóniam peccatórum mole prémimur, et non est qui ádjuvet. \textred{S}úbveni. Glória Patri. \textred{S}úbveni.
\item Te Deum laudámus. 6.
\item \textred{Sacerd. ℣.} Ora pro nobis, sancta Dei Génitrix. \textred{℟.} Ut digni efficiámur promissiónibus Christi, allelúja.
\item \textred{In Laud. fit commem. S. Joannis, Aña.} Dedit illi sciéntiam sanctórum, honestávit illum in labóribus, et complévit labóres illíus.
\item \textred{℣.} Justus germinábit sicut lílium. \textred{℟.} Et florébit in ætérnum ante Dóminum.
\item Orémus \hfill \textred{Oratio.}
\end{tightlist}
\lettrine{D}eus, qui sanctum Joánnem Confessórem tuum, atque Doctórem, Patrem nostrum perféctæ sui abnegatiónis, et crucis amatórem exímium effecísti: concéde, ut ejus imitatióni júgiter inhæréntes, glóriam assequámur ætérnam. Per Christum.

\begin{tightlist}{}
\item \textred{Deinde fit Comm. Feriæ.}
\end{tightlist}

\subsubsection*{\center \textsc{IN II. VESPERIS.}}
\begin{tightlist}{}
\item \textred{Ad Magnif. Aña.} Beáta Dei Génitrix, María, Virgo perpétua, templum Dómini, Sacrárium Spíritus sancti; tu sola sine exémplo placuísti Dómino nostro Jesu Christo. Ora pro pópulo, intérveni pro clero, intercéde pro devóto femíneo sexu.
\item \textred{Commem. Bb Dionysii et Redempti, Aña.} Isti sunt sancti, qui pro Dei amóre minas hóminum contempsérunt, sancti Mártyres in regnum ætérnum exsúltant cum Angelis. O quam pretiósa est mors sanctórum, qui assídue assístunt ante Dóminum, et ab ínvicem non sunt separati!
\item \textred{℣.} Lætámini in Dómino, et exsultáte justi. \textred{℟.} Et gloriámini omnes recti corde, allelúja.
\item Orémus. \hfill \textred{Oratio.}
\end{tightlist}
\lettrine{D}eus, qui mirábili dispositióne beátos Dionýsium et Redémptum per maris perícula ad palmam martýrii perduxísti: eórum intercessióne concéde; ut, inter mundánas varietátes et sæculária desidéria, fidéles usque ad mortem in confessióne tui nóminis maneámus. Per Dóminum.

\begin{tightlist}{}
\item \textred{Et Commem. Oct. S. Joan. per Añam.} Quæsívit verba utília, et conscrípsit sermónes rectíssimos, et veritáte plenos. 
\item \textred{℣.} Ora pro nobis sancte Pater Joánnes. \textred{℟.} Ut digni efficiámur promissiónibus Christi.
\item \textred{Oratio} Deus, qui sanctum Joánnem \textred{ut supra.}
\item \textred{Et fit Comm. Feriæ,} 
\item \textred{et deinde comm. S. Saturnini Martyr, Aña.} Iste Sanctus pro lege Dei sui certávit usque ad mortem, et a verbis impiórum non tímuit; fundátus enim erat supra firmam petram.
\item \textred{℣.} Glória et honóre coronásti eum Dómine. \textred{℟.} Et constituísti eum super ópera mánuum tuárum.
\item Orémus. \hfill \textred{Oratio.}
\end{tightlist}
\lettrine{D}eus, qui nos beáti Saturníni Mártyris tui concédis natalítio pérfrui: ejus nos tríbue méritis adjuvári. Per Dóminum.

\begin{center}\textcolor{red}{\rule{\linewidth}{0.5pt}} \end{center}


\subsection*{\center \textsc{\textred{DIE XXIX. NOVEMB.}}}
\subsubsection*{\center \textred{Quinta die infra Octav.}}
\subsection*{\center \textsc{S. Joannis a Cruce}}
\fancyhead[C]{\textsc{29 Nov - Die V Infra Octav. S. Joannis}}
\subsubsection*{\center \textred{Semiduplex.}}
\subsubsection*{\center \textsc{AD MATUTINUM.}}
\begin{tightlist}{}
\item \textred{Invitatorium, Hymnus, prima Aña j. Nocturni, et Psalmi ut infra in die Oct. ejusdem. ℣. et ℟℟. ut ibidem secundum ordinem diei et Nocturni.}
\item \rubric{Id est: Quando in propriis locis aliud non assignatur, Versus, et tria Responsoria sunt de Festo secundum ordinem dierum et Nocturnorum: ita ut Feria secunda et quinta ℣. et ℟℟. sint de primo Noct., Feria tertia, et sexta de secundo Noct., Feria quarta et Sabbato de tertio Noct.\\ Benedictiones etiam ante lectiones dicantur secundum ordinem Nocturnorum:}
\end{tightlist}
\subsubsection*{\textred{Feriis II et V:}}
\begin{tightlist}{}
\item \textred{℣.} Amávit eum Dóminus, et ornávit eum. \textred{℟.} Stolam glóriæ índuit eum.
\item \textred{℟. j.} Euge serve bone et fidélis, quia super pauca fuísti fidélis, supra multa te constítuam: \textred{*} Intra in gáudium Dómini tui. \textred{℣.} Dómine quinque talénta tradidísti mihi, ecce ália quinque superlucrátus sum. \textred{I}ntra.
\item \textred{℟. ij.} Justus germinábit sicut lílium: \textred{*} Et florébit in ætérnum ante Dóminum. \textred{℣.} Plantátus in domo Dómini, in átriis domus Dei nostri. \textred{E}t.
\item \textred{℟. iij.} Iste cognóvit justítiam, et vidit mirabília magna, et exorávit Altíssimum. \textred{*} Et invéntus est in número sanctórum. \textred{℣.} Iste est, qui contémpsit vitam mundi, et pervénit ad cœléstia regna. \textred{E}t. Glória Patri. \textred{E}t.
\end{tightlist}

\subsubsection*{\textred{Feriis III et VI:}}
\begin{tightlist}{}
\item \textred{℣.} Justum dedúxit Dóminus per vias rectas. \textred{℟.} Et osténdit illi regnum Dei. 
\item \textred{℟. j.} Desidérium ánimæ ejus tribuísti ei Dómine. \textred{*} Et voluntáte labiórum ejus non fraudásti eum. \textred{℣.} Posuísti Dómine super caput ejus corónam de lápide pretióso. \textred{E}t.
\item \textred{℟. ij.} Amávit eum Dóminus, et ornávit eum: stolam glóriæ índuit eum: \textred{*} Et ad portas paradísi coronávit eum. \textred{℣.} Induit eum Dóminus lorícam fídei, et ornávit eum. \textred{E}t.
\item \textred{℟. iij.} Iste sanctus digne in memóriam vértitur hóminum: qui ad gáudium tránsiit Angelórum: quóniam in hac peregrinatióne solo córpore constitútus, cogitatióne, et aviditáte, \textred{*} In illa ætérna pátria conversátus est. \textred{℣.} Vínculis carnis absolútus, taléntum sibi créditum Dómino suo duplicátum reportávit. \textred{I}n. Glória Patri. \textred{I}n.
\end{tightlist}

\subsubsection*{\textred{Feria IV et Sabbato:}}
\begin{tightlist}{}
\item \textred{℣.} Justus ut palma florébit in domo Dómini. \textred{℟.} Sicut cedrus Líbani multiplicábitur.
\item \textred{℟. j.} Magnificávit eum Dóminus in conspéctu regum: \textred{*} Et dedit illi corónam glóriæ. \textred{℣.} Státuit illi testaméntum sempitérnum. \textred{E}t.
\item \textred{℟. ij.} Vir Israelíta, gaude cohéres Christi: \textred{*} Ora pro nobis. \textred{℣.} Ut précibus tuis a Deo véniam sperémus, devóti ad te confúgimus. \textred{O}ra.
\item \textred{℟. iij.} Sancte Joánnes, Christi Conféssor, audi rogántes sérvulos, \textred{*} Et impetrátam nobis cœ́litus tu defer indulgéntiam. \textred{℣.} O sancte Joánnes, sidus áureum Dómini grátia, servórum gémitus sólita súscipe cleméntia. \textred{E}t. Glória Patri. \textred{E}t. 
\item Te Deum laudámus. 6.
\item \textred{Sacerd. ℣.} Os justi meditábitur sapiéntiam. \textred{℟.} Et lingua ejus loquétur judícium, allelúja.
\end{tightlist}

\begin{tightlist}{}
\item \textred{Lectio vero prima et secunda erit de Scriptura occurr. tertia de Octava, ut sequitur.}
\end{tightlist}

\subsubsection*{\center Sermo sancti Bernárdi Abbátis.}
\begin{center}\textred{In festo S. Malach. Serm. 2.}\\
\textred{\textbf{Lectio iii}.}\end{center}
\lettrine{Q}uasi unus ómnium parens vivébat ómnibus Joánnes, quasi gallína pullos suos, sic fovébat omnes, et in velaménto alárum suárum protegébat. Non sexus, non ætas, non condítio discernebátur, aut persóna : déerat némini, expánso ómnibus grémio pietátis. De quacúmque tribulatióne clamarétur ad eum, própriam reputábat, nisi quod in sua pátiens, in aliéna compátiens erat, plerúmque et impátiens. Non tamen ita ómnibus intendébat, ut se solum expóneret, solum cura excíperet generáli. Erat et sui sollícitus, se ipsum custodiébat. Ita benígne totus suus, et totus ómnium erat, ut nec cáritas a custódia sui, nec propríetas ab utilitáte commúni eum impedíre, vel retardáre in áliquo viderétur.

\begin{tightlist}{}
\item \textred{Ad Laudes, sola prima Aña.} Addúcam eum in montem sanctum meum, et lætificábo eum in domo oratiónis meæ. 
\item \textred{Reliqua omnia ut infra in die Oct.}
\item \textred{Comm. Feriæ.}
\item \textred{Et fit comm. S. Saturnini Mart. Aña.} Nisi granum fruménti cadens in terram mórtuum fúerit, ipsum solum manet.
\item \textred{℣.} Magna est glória ejus in salutári tuo. \textred{℟.} Glóriam et magnum decórem impónes super eum.
\item Orémus. \hfill \textred{Oratio.}
\end{tightlist}
\lettrine{D}eus, qui nos beáti Saturníni Mártyris tui concédis natalítio pérfrui: ejus nos tríbue méritis adjuvári. Per Dóminum.

\begin{tightlist}{}
\item \textred{De Vigilia S. Andreæ nihil nisi in Missa.}
\item \textred{Vesp. S. Andreæ Apostoli.} 644.
\end{tightlist}

\begin{center}\textcolor{red}{\rule{\linewidth}{0.5pt}} \end{center}


%\subsection*{\center \textsc{DIE XXX. NOVEMBR.}\\ 
%In Festo S. Andreæ Apostoli.}
%\fancyhead[C]{\textsc{30 Nov - S. Andreæ Apostoli}}
%\begin{center}\textred{Duplex minus P.C.}\end{center}
%\begin{tightlist}{}
%\item \textred{Omnia de Communi Apostolorum [3] præter hic signata.}
%\end{tightlist}
%\subsubsection*{\center Ad Vesperas.}
%\begin{tightlist}{}
%\item \textred{Aña.} Unus ex duóbus, qui secúti fúerant Jesum, erat Andréas frater Simónis Petri, allelúja.
%\item \textred{Capit.} Jam non estis. [3]
%\item \textred{℟.} Vir iste in pópulo suo mitíssimus appáruit, sanctitáte autem et grátia plenus: iste est, qui assídue orat * Pro pópulo, et pro civitáte ista.
%\item \textred{℣.} Pro eo, ut me dilígerent, detrahébant mihi: ego autem orábam. Pro.
%\item Glória Patri. Pro.
%\item \textred{Hymn} Exsúltet. [3]
%\item \textred{℣.} Diléxit Andréam Dóminus. \textred{℟.} In odórem suavitátis, allelúja.
%\item \textred{Ad Magnif. Aña.} Ambulans Jesus juxta mare Galilææ, vidit Petrum, et Andréam, fratrem ejus, et ait illis: Veníte post me, at fáciam vos fíeri piscatóres hóminum. At illi contínuo relíctis rétibus, secúti sunt eum.
%\end{tightlist}
%\subsubsection*{\center\textred{Oratio.}}
%\lettrine{Q}uæsumus, omnípotens Deus, ut beátus Andréas Apóstolus, cujus prævenímus festivitátem, tuum pro nobis implóret auxílium; ut a nostris reátibus absolúti, a cunctis étiam perículis eruámur. Per Dóminum.
%
%\begin{tightlist}{}
%\item \textred{Si hoc Festum transferatur, dicatur in j. Vesp. Oratio.} Majestátem, \textred{ut infra ad Laudes.}
%\end{tightlist}
%
%\subsubsection*{\center \textsc{AD MATUTINUM.}}
%\begin{tightlist}{}
%\item \textred{Invitat.} Christum Regem regum adorémus Dóminum: * Qui martyrio crucis beátum glorificávit Andréam Apóstolum.
%\item \textred{Ps.} Veníte. 2.
%\item \textred{Hymn.} Exsúltet cœlum.
%\end{tightlist}
%
%\subsubsection*{\center \textsc{IN PRIMO NOCTURNO.}}
%\begin{tightlist}{}
%\item \textred{Aña.} Vidit Dóminus Petrum et Andréam, et vocávit eos.
%\item \textred{Ps.} Cœli enárrant. [4]
%\item \textred{Aña.} Veníte post me, dicit Dóminus, et fáciam vos fíeri piscatóres hóminum.
%\item \textred{Ps.} Benedícam.
%\item \textred{Aña.} Relíctis rétibus suis, secúti sunt Dóminum Redemptórem. 
%\item \textred{Ps.} Eructávit.
%\item \textred{℣.} In omnem terram. [7]
%\end{tightlist}
%
%\subsubsection*{\center De Epístola beáti Pauli Apóstoli ad Romános.}
%\textred{\textbf{Lectio j}.} \hfill \textred{Cap. 10.}\\
%\lettrineThree{F}ratres, finis legis, Christus, ad justítiam omni credénti. Móyses enim scripsit, quóniam justítiam, quæ ex lege est, qui fécerit homo, vivet in ea. Quæ autem ex fide est justítia, sic dicit: Ne díxeris in corde tuo: Quis ascéndet in cœlum? idest, Christum dedúcere. Aut quis descéndet in abýssum? hoc est, Christum a mórtuis revocáre. Sed quid dicit Scriptúra? Prope est verbum in ore tuo, et in corde tuo: hoc est verbum fídei, quod prædicámus. Quia si confiteáris in ore tuo Dóminum Jesum, et in corde tuo credíderis, quod Deus illum suscitávit a mórtuis, salvus eris.
%
%\begin{tightlist}{}
%\item \textred{℟.} Ambulans Jesus juxta mare Galilææ vidit duos fratres, Simónem, qui vocátur Petrus, et Andréam fratrem ejus, mitténtes rétia in mare, et ait illis: * Veníte post me, et fáciam vos fíeri piscatóres hóminum. 
%\item \textred{℣.} Erant enim piscatóres, et ait illis. Veníte.
%\end{tightlist}
%
%\subsubsection*{\center\textred{Lectio ij.}}
%\lettrine{C}orde enim créditur ad justítiam, ore autem conféssio fit ad salútem. Dicit enim Scriptúra: Omnis, qui credit in illum, non confundétur. Non enim est distínctio Judæi et Græci: nam idem Dóminum ómnium, dives in omnes, qui ínvocant illum. Omnis enim, quicúmque invocáverit nomen Dómini, salvus erit. Quómodo ergo invocábunt, in quem non credidérunt? Aut quómodo credent ei, quem non audiérunt? Quómodo autem áudient sine prædicánte? Quómodo vero prædicábunt, nisi mittántur? sicut scriptum est: Quam speciósi pedes evangelizántium pacem, evangelizántium bona!
%
%\begin{tightlist}{}
%\item \textred{℟.} Mox ut vocem Dómini prædicántis audívit beátus Andréas, relíctis rétibus, quorum usu, actúque vivébat, ætérnæ vitæ secútus est * Præmia largiéntem.
%\item \textred{℣.} Ad uníus jussiónis vocem Petrus et Andréas secúti sunt Redemptórem. Præmia.
%\end{tightlist}
%
%\subsubsection*{\center\textred{Lectio iij.}}
%\lettrine{S}ed non omnes obédiunt Evangélio. Isaías enim dicit: Dómine, quis crédidit audítui nostro? Ergo fides ex audítu, audítus autem per verbum Christi. Sed dico: Numquid non audiérunt? Et quidem in omnem terram exívit sonus eórum, et in fines orbis terræ verba eórum.
%
%\begin{tightlist}{}
%\item \textred{℟.} Homo Dei ducebátur, ut crucifígerent eum: pópulus autem clamábat voce magna, dicens: * Innocens ejus sanguis sine causa damnátur.
%\item \textred{℣.} Cumque carnífices dúcerent eum, ut crucifigerétur: factus est concúrsus populórum clamántium, et dicéntium. Innocens.
%Glória Patri. Innocens.
%\end{tightlist}
%
%\subsubsection*{\center \textsc{IN SECUNDO NOCTURNO.}}
%\begin{tightlist}{}
%\item \textred{Aña.} Unus ex duóbus, qui secúti fúerant Jesum, erat Andréas frater Simónis Petri, allelúja.
%\item \textred{Ps.} Omnes Gentes. [8]
%\item \textred{Aña.} Diléxit Andréam Dóminus in odórem suavitátis, allelúja.
%\item \textred{Ps.} Exáudi Deus.
%\item \textred{Aña.} Dignus sibi Dóminus computávit Mártyrem, quem vocávit Apóstolum, dum esset in mari, allelúja.
%\item \textred{Ps.} Exáudi Deus.
%\item \textred{℣.} Constítues eos. [10]
%\end{tightlist}
%
%
%\subsubsection*{\center \textred{Lectio iv.}}
%\begin{center}\textred{Ex ejus pass. quam scribunt Presb. et Diac. Achajæ.}\end{center}
%\lettrineThree{A}ndréas Apóstolus Bethsáidæ natus, qui est Galilææ vicus, frater Petri, discípulus Joánnis Baptístæ, cum eum de Christo dicéntem audísset: Ecce Agnus Dei, secútus Jesum, fratrem quoque suum ad eúmdem perdúxit. Cum póstea una cum fratre piscarétur in mari Galilææ, ambo a prætereúnte Christo Dómino ante álios Apóstolos vocáti illis verbis: Veníte post me, et fáciam vos fíeri piscatóres hóminum; nullam interponéntes moram, et relíctis rétibus secúti sunt eum. Post cujus passiónem, et resurrectiónem Andréas, cum in Scythiam Európæ, quæ ei província ad Christi fidem disseminándam obtígerat, venísset, deínde Epírum ac Thráciam peragrásset, doctrína et miráculis innumerábiles hómines ad Christum convértit. Post Patras Achájæ proféctus, et in ea urbe plúrimis ad veritátem Evangélicam perdúctis, Ægéam procónsulem, prædicatióni Evangélicæ resisténtem, libérrime increpávit, quod, qui judex hóminum habéri vellet, Christum Deum ómnium júdicem a dæmónibus elúsus non agnósceret.
%
%\begin{tightlist}{}
%\item \textred{℟.} Doctor bonus, et amícus Dei Andréas dúcitur ad crucem. Adspíciens a longe, vidit crucem, et ait: Salve crux: * Súscipe discípulum ejus, qui pepéndit in te magíster meus Christus. 
%\item \textred{℣.} Salve crux, quæ in córpore Christi dedicáta es, et ex membris ejus tamquam margarítis ornáta. Súscipe.
%\end{tightlist}
%
%\subsubsection*{\center\textred{Lectio v.}}
%\lettrine{T}um Ægéas irátus: Désine, inquit, Christum jactáre, cui simília verba nihil profuérunt, quo minus a Judæis crucifigerétur. Andréam vero de Christo nihilóminus líbere prædicántem, quod pro salúte humáni géneris se crucifigéndum obtulísset, ímpia oratióne interpéllat, ac demum hortátur, ut sibi cónsulens, diis velit immoláre. Cui Andréas: Ego omnipoténti Deo, qui unus et verus est, ímmolo quotídie, non taurórum carnes, nec hircórum sánguinem, sed immaculátum Agnum in altári: cujus carnem posteáquam omnis pópulus credéntium manducáverit, Agnus qui sacrificátus est, ínteger persevérat et vivus. Quámobrem ira accénsus Ægéas, jubet eum in cárcerem detrúdi: unde pópulus Andréam fácile liberásset, nisi ipse sedásset multitúdinem, veheméntius rogans, ne se ad optatíssimam martyrii corónam properántem impedírent. 
%
%\begin{tightlist}{}
%\item \textred{℟.} O bona crux, quæ decórem et pulchritúdinem ex membris Dómini suscepísti: * Accipe me ab homínibus, et redde me magístro meo Christo: ut per te me recípiat, qui per te me redémit.
%\item \textred{℣.} Salve crux, quæ in córpore Christi dedicáta es, et ex membris ejus tamquam margarítis ornáta. Accipe me.
%\end{tightlist}
%
%\subsubsection*{\center\textred{Lectio vi.}}
%\lettrine{I}gitur paulo post in tribúnal prodúctum, cum Ægéas crucis extolléntem mystéria, sibíque suam impietátem exprobrántem, diútius ferre non posset, in crucem tolli, et Christi mortem imitári jussit. Addúctus Andréas ad locum martyrii, cum crucem vidísset, longe exclamáre cœpit: O bona crux, quæ decórem ex membris Dómini suscepísti, diu desideráta, sollícite amáta, sine intermissióne quæsíta, et aliquándo cupiénti ánimo præparáta; áccipe me ab homínibus, et redde me Magístro meo: ut per te me recípiat, qui per te me redémit. Itaque cruci affíxus est: in qua bíduum vivus pendens, et Christi fidem prædicáre nunquam intermíttens, ad eum migrávit, cujus mortis similitúdinem concupíerat. Quæ ómnia Presbyteri et Diáconi Achájae, qui ejus passiónem scripsérunt, se ita ut commemoráta sunt, audísse et vidísse testántur. Ejus ossa primum Constantíno Imperatóre Constantinópolim, deínde Amálphim transláta sunt. Caput, Pio Secúndo Pontífice, Romam allátum, in Basílica sancti Petri collocátum est.
%
%\begin{tightlist}{}
%\item \textred{℟.} Diléxit Andréam Dóminus in odórem suavitátis: dum pendéret in cruce, dignum sibi computávit mártyrem, quem vocávit Apóstolum dum esset in mari: * Et ídeo amícus Dei appellátus est.
%\item \textred{℣.} Elégit eum Dóminus, et excélsum fecit illum. Et ídeo. 
%\item Glória Patri. Et ídeo.
%\end{tightlist}
%
%\subsubsection*{\center \textsc{IN TERTIO NOCTURNO}}
%\begin{tightlist}{}
%\item \textred{Aña.} Ego, si patíbulum crucis expavéscerem, crucis glóriam non prædicárem.
%\item \textred{Ps.} Confitébimur. [10]
%\item \textred{Aña.} Andréas vero rogábat pópulum, ut non impedíret passiónem ipsíus.
%\item \textred{Ps.} Dóminus regnávit.
%\item \textred{Aña.} Bíduo vivens pendébat in cruce pro Christi nómine beátus Andréas, et docébat pópulum.
%\item \textred{Ps.} Dóminus regnávit.
%\item \textred{℣.} Nimis honoráti. [12]
%\end{tightlist}
%
%\subsubsection*{\center Léctio sancti Evangélii secúndum Matthæum.}
%\textred{\textbf{Lectio vij}.} \hfill \textred{Cap. 14.}\\
%\lettrine{I}n illo témpore: Ambulans Jesus juxta mare Galilææ, vidit duos fratres, Simónem, qui vocátur Petrus, et Andréam fratrem ejus. Et réliqua.
%
%\begin{center}Homilía sancti Gregórii Papæ.\\
%\textred{Homilia 5. in Evang.}\end{center}
%
%\lettrineThree{A}udístis, fratres caríssimi, quia ad uníus jussiónis vocem Petrus et Andréas relíctis rétibus secúti sunt Redemptórem. Nulla vero hunc fácere adhuc mirácula víderant: nihil ab eo de præmio ætérnæ retributiónis audíerant: et tamen ad unum DÓmini præcéptum, hoc, quod possidére videbántur, oblíti sunt. Quanta nos ejus mirácula vidémus, quot flagéllis afflígimur, quantis minárum asperitátibus deterrémur, et tamen vocántem sequi contémnimus?
%
%\begin{tightlist}{}
%\item \textred{℟.} Orávit sanctus Andréas, et dum respíceret in cœlum, voce magna clamávit, et dixit: Tu es Deus meus, quem vidi: * Ne me patiáris ab ímpio júdice depóni: quia virtútem sanctæ crucis agnóvi. 
%\item \textred{℣.} Tu es magíster meus Christus, quem diléxi, quem cognóvi, quem conféssus sum: tantúmmodo in ista voce me exáudi. Ne me patiáris.
%\end{tightlist}
%\subsubsection*{\center\textred{Lectio viij.}}
%\lettrine{I}n cœlo jam sedet, qui de conversióne nos ádmonet: jam jugo fídei colla Géntium súbdidit: jam mundi glóriam stravit: jam ruínis ejus crebrescéntibus distrícti sui judícii diem propinquántem denúntiat: et tamen supérba mens nostra adhuc non vult hoc sponte desérere, quod quotídie perdit invíta. Quid ergo, fratres caríssimi, quid in ejus judício dictúri sumus, qui ab amóre præséntis sæculi nec præcéptis fléctimur, nec verbéribus emendámur?
%
%\begin{tightlist}{}
%\item \textred{℟.} Expándi manus meas tota die in cruce ad pópulum non credéntem; sed contradicéntem mihi: * Qui ámbulant vias non bonas, sed post peccáta sua.
%\item \textred{℣.} Deus ultiónum Dóminus: Deus ultiónum líbere egit: exaltáre, qui júdicas terram; redde retributiónem supérbis. Qui.
%\end{tightlist}
%\subsubsection*{\center\textred{Lectio ix.}}
%\lettrine{S}ed fortásse áliquis tácitis sibi cogitatiónibus dicat: Ad vocem Domínicam utérque iste piscátor, quid, aut quantum dimísit, qui pene nihil hábuit? Sed hac in re, fratres caríssimi, afféctum debémus pótius pensáre, quam censum. Multum relíquit, qui sibi nihil retínuit; multum relíquit, qui quantúmlibet parum, totum deséruit. Certe nos et hábita cum amóre possidémus, et ea, quæ mínime habémus, ex desidério quærimus. Multum ergo Petrus et Andréas dimísit, quando utérque étiam desidéria habéndi derelíquit.
%
%\begin{tightlist}{}
%\item \textred{℟.} Vir iste in pópulo suo mitíssimus appáruit: sanctitáte autem et grátia plenus: iste est, qui assídue orat * Pro pópulo, et pro civitáte ista.
%\item \textred{℣.} Pro eo, ut me dilígerent, detrahébant mihi: ego autem orábam. Pro.
%\item Glória Patri. Pro pópulo.
%\item Te Deum laudámus.
%\item \textred{Sacerdotalis ℣.} Dedísti hereditátem. \textred{℟.} Timéntibus nomen tuum Dómine, allelúja.
%\end{tightlist}
%\subsubsection*{\center \textsc{AD LAUDES}}
%\begin{center} \textred{et per Horas, Aña.}\end{center}
%\lettrine{S}alve crux pretiósa, súscipe discípulum ejus, qui pepéndit in te magíster meus Christus.
%
%\begin{tightlist}{}
%\item \textred{Ps.} Dóminus regnávit, \textred{cum reliquis.} 30.
%\item \textred{Aña.} Beátus Andréas orábat, dicens: Dómine Rex ætérnæ glóriæ, súscipe me pendéntem in patíbulo.
%\item \textred{Aña.} Maximílla Christo amábilis, tulit corpus Apóstoli, et óptimo loco cum aromátibus sepelívit.
%\item \textred{Aña.} Qui persequebántur justum, demersísti eos, Dómine, in profúndum, et in ligno crucis dux justi fuísti.
%\item \textred{Capit.} Jam non estis.
%\item \textred{Hymn.} Ætérna Christi múnera.
%\item \textred{℣.} Annuntiavérunt, \textred{de Communi. [13]}
%\item \textred{Ad Benedictus, Aña.} Cum pervenísset beátus Andréas ad locum, ubi crux paráta erat, exclamávit, et dixit: O bona crux, diu desideráta, et jam concupiscénti ánimo præparáta: secúrus et gaudens vénio ad te, ita et tu exsúltans suscípias me discípulum ejus, qui pepéndit in te.
%\end{tightlist}
%
%\subsubsection*{\center\textred{Oratio.}}
%\lettrine{M}ajestátem tuam, Dómine, supplíciter exorámus: ut sicut Ecclésiæ tuæ beátus Andréas Apóstolus éxstitit prædicátor et rector; ita apud te sit pro nobis perpétuus intercéssor. Per Dóminum.
%
%\begin{tightlist}{}
%\item \textred{Ad Horas, Capitula, ℟℟. et ℣℣. de Communi. [13]}
%\end{tightlist}
%
%\subsubsection*{\center \textsc{AD VESPERAS}}
%\begin{tightlist}{}
%\item \textred{Añæ Laudum.}
%\item \textred{Ps.} Dixit Dominus, \textred{cum reliquis de Comm. [14]}
%\item \textred{℣.} Diléxit Andréam Dóminus. \textred{℟.} In odórem suavitátis, allelúja.
%\item \textred{Ad Magnificat, Aña.} Concéde nobis hóminem justum, redde nobis hóminem sanctum, ne interfícias hóminem Deo carum, justum, mansuétum, et pium.
%\item \textred{Et fit comm. seq.}
%\end{tightlist}

%\begin{center}\textcolor{red}{\rule{\linewidth}{0.5pt}} \end{center}

%\section*{\center \textsc{\textred{Festa Decembris.}}}
%\begin{center}\textcolor{red}{\rule{\linewidth}{0.5pt}} \end{center}

\subsection*{\center \textred{DIE I. DECEMBRIS.}}
\subsubsection*{\center \textred{Dies Octava.}}
\subsection*{\center \textsc{S. Joannis a Cruce.}}
\subsubsection*{\center \textred{Conf. et Doct., Patris Nostri.\\ Duplex minus.}}
%\begin{center}\textred{Vide festum, supra, 535.}\end{center}
\fancyhead[C]{\textsc{1 Dec. - Dies Octava S. Joannis a Cruce.}}

\begin{tightlist}{}
\item \textred{Omnia de Comm. unius Conf. non Pontif.} [85] \textred{præter hic signata.}
\end{tightlist}
\subsubsection*{\center \textsc{AD I. VESPERAS.}}
\begin{tightlist}{}
\item \textred{Nisi præcedat Festum duplex.}
\item \textred{Aña j. Laud.} Addúcam eum in montem sanctum meum, et lætificábo eum in domo oratiónis meæ. \textred{quæ sola dicitur super Psalm. præc.}
\end{tightlist}

\subsubsection*{\center\textred{Hymnus.}}
\lettrine{R}egis ætérni generóse miles,\\ Ordinis nostri cólumen Joánnes,\\ Quos tuis læti méritis dicámus,\\ Súscipe plausus.\\
\\ 
O satis felix! Speculátor alti\\ Núminis, votis animóque Martyr,\\ Pœ́nitens Virgo, memoránde vates,\\ Mystice Doctor.\\
\\
Sæpe cum Christo, et Genitríce pura\\ Verba commísces, superísque vivens:\\ Hinc fluit candor, tua quo refúlgent\\ Scripta, perénnis.\\
\\
Mente perfúsa rádiis ab alto,\\ Montis ascénsum, tenebrásque noctis,\\ Et facem vivam récolens amóris,\\ Alta revélas.\\
\\
Te sacri Verbi retegénte sensus,\\ Visa calígo procul exsuláre,\\ Dum facis lumen rútilans oríri\\ Noctis ab umbris.\\
\\
Pérsonent laudes Tríadi Joánnis,\\ Cui Pater lumen, tribuítque dona\\ Spíritus, vitæ réserat volúmen\\ Cœ́licus Agnus.\\
Amen.

\begin{tightlist}{}
\item \textred{℣.} Ora pro nobis sancte Pater Joánnes. \textred{℟.} Ut digni efficiámur promissiónibus Christi, allelúja. 
\item \textred{Ad Magnificat, Aña.} Collaudábunt multi sapiéntiam ejus, et usque in sǽculum non delébitur.
\end{tightlist}

\subsubsection*{\center \textred{Oratio.}}
\lettrine{D}eus, qui sanctum Joánnem Confessórem tuum, atque Doctórem, Patrem nostrum perféctæ sui abnegatiónis, et crucis amatórem exímium effecísti: concéde, ut ejus imitatióni júgiter inhæréntes, glóriam assequámur ætérnam. Per Christum.

\subsubsection*{\center \textsc{AD MATUTINUM.}}
\subsubsection*{\center\textred{Hymnus.}}
\lettrine{D}iem Joánnes ádvehit,\\ Crucis micántem glória,\\ Cui fixus hæsit spíritu,\\ Fixámque gessit córpori.\\
\\
Non probra, nec ludíbria,\\ Fames, flagélla, víncula,\\ Sedáre jam valent sitim,\\ Qua dura fortis áppetit.\\
\\
Hæc poscit ille gáudia,\\ Hæc vult labóris prǽmia\\ Hæc vota, confíxi Ducis\\ Reférre signum pérpetim.\\
\\
Noctem volútans mysticam\\ In nocte lucem pércipit,\\ De luce flammam súscipit,\\ Carméli ut agmen ínstruat.\\
\\
Te, Christe, laudent Cœ́lites,\\ Tibíque corda dédita\\ Currant, Joánne prǽvio,\\ Ad vera vitæ gáudia.\\
Amen.

\begin{tightlist}{}
\item \textred{In primo Nocturno, Lectiones de Scriptura occurrente.}
\end{tightlist}
\subsubsection*{\center \textsc{IN SECUNDO NOCTURNO.}}
\subsubsection*{\center Sermo sancti Bernárdi Abbátis.}
\begin{center}\textred{In nat. S. Vict. Conf. Serm. 1.\\
\textbf{Lectio iv}.}\end{center}
\lettrineThree{J}oánnis vita, et glória speciális, non tam ad glóriam, quam ad virtútem próvocat omnes, qui recto sunt corde. Non recti plane, sed pervérsi est ánimi, ante quǽrere glóriam, quam exercére virtútem, et velle coronári, qui legítime non certáverit. Vanum est vobis, inquit, ante lucem súrgere. Ita est. Frustra ad celsitúdinem nítitur glóriæ qui prius non cláruit virtúte. Absit mihi gloriári, nisi in glória illórum, quibus Prophéta congrátulans, ait: Dómine, in lúmine vultus tui ambulábunt, et in nómine tuo exsultábunt tota die, et in justítia tua exaltabúntur. Et infert: Quóniam glória virtútis eórum tu es. Pulchre non eórum, sed virtútis eórum glória commendátur.

\begin{tightlist}{}
\item \textred{℟.} Desidérium ánimæ ejus tribuísti ei Dómine. \textred{*} Et voluntáte labiórum ejus non fraudásti eum. \textred{℣.} Posuísti Dómine super caput ejus corónam de lápide pretióso. \textred{E}t.
\end{tightlist}

\subsubsection*{\center\textred{Lectio v.}}
\lettrine{H}abémus, delictíssimi, in vita Joánnis, et quod digne mirémur, et quod salúbrius imitémur. Quis non omni admiratióne stupéndum ducat, carne circúmdatum, cœlis paténtibus, carnis óculos supérno infígere lúmini, et vidére visiónes Dei? Hæc, atque simília in sancto hómine venerámur. Mérito quidem quando et sine salútis perículo possunt non fíeri, et sine salútis perículo néqueunt usurpári. Tútius sane æmulánda solidióra, quam sublimióra, et quæ magis virtútem redóleant, glóriam minus. Studeámus proínde móribus conformári.

\begin{tightlist}{}
\item \textred{℟.} Amávit eum Dóminus, et ornávit eum: stolam glóriæ índuit eum: \textred{*} Et ad portas paradísi coronávit eum. \textred{℣.} Induit eum Dóminus lorícam fídei, et ornávit eum. \textred{E}t ad.
\end{tightlist}

\subsubsection*{\center\textred{Lectio vi.}}
\lettrine{Æ}mulémur in viro sóbrium victum, devótum afféctum, æmulémur mansuetúdinem, spíritus, castimóniam córporis, oris custódiam, ánimi puritátem, pónere frænum iræ, et modum linguæ, dormíre párcius, oráre frequéntius, commonére nosmetípsos psalmis, hymnis, cánticis spirituálibus: diébus júngere noctes, et divínis láudibus occupári. Æmulémur charísmata melióra: discámus ab ipso, quod mitis, et húmilis corde fuit. Æmulémur, inquam, quod éxstitit jucúndus ad hóspites, pátiens ad peccántes, benígnus ad omnes. Hoc enim mélius: in his forma est, cui imprimámur: in miráculis glória, a qua reprimámur: Illa lætíficent, ista ædíficent; móveant illa, ista promóveant.

\begin{tightlist}{}
\item \textred{℟.} Iste sanctus digne in memóriam vértitur hóminum: qui ad gáudium tránsiit Angelórum: quóniam in hac peregrinatióne solo córpore constitútus, cogitatióne, et aviditáte, \textred{*} In illa ætérna pátria conversátus est. \textred{℣.} Vínculis carnis absolútus, taléntum sibi créditum Dómino suo duplicátum reportávit. \textred{I}n illa. Glória Patri. \textred{I}n illa.
\end{tightlist}

\subsubsection*{\center \textsc{IN TERTIO NOCTURNO.}}
\subsubsection*{\center Léctio sancti Evangélii secúndum Lucam.}
\textred{\textbf{Lectio vij}. \hfill Cap. 11.}\\

\lettrine{I}n illo témpore: Dixit Jesus turbis: Nemo lucérnam accéndit, et in abscóndito ponit, neque sub módio, sed supra candelábrum; ut, qui ingrediúntur lumen vídeant. Et réliqua.

\begin{center}Homilía sancti Joánnis Chrysóstomi.\\
\textred{Hom. 10. in cap. 5. Matth.}\end{center}

\lettrineThree{P}er hanc comparatiónem vult osténdere, quare Sanctos suos ipse Christus maniféstat, et non sinit eos esse abscónsos, id est, quia nec qui accéndunt lucérnam, ídeo eam accéndunt, ut ponant sub módio, sed super candelábrum, ut lúceat ómnibus, qui in domo sunt. Qui sunt accensóres lucérnæ? Pater, et Fílius. Quæ est illa lucérna? Verbum divínum, de quo dictum est: Lucérna pédibus meis verbum tuum. Quod est candelábrum? Ecclésia, quæ bájulat verbum vitæ; consequénter dicit et Paulus: inter quos lucétis, sicut luminária in sǽculo, verbum vitæ continéntes. Consequénter et omnis ecclesiásticus vir habens verbum Dei, dícitur candelábrum.

\begin{tightlist}{}
\item \textred{℟.} Magnificávit eum Dóminus in conspéctu regum: \textred{*} Et dedit illi corónam glóriæ. \textred{℣.} Státuit illi testaméntum sempitérnum. \textred{E}t dedit.
\end{tightlist}

\subsubsection*{\center\textred{Lectio viij.}}
\lettrine{M}odii autem sunt hómines mundiáles, vácui a Deo, et ab ómnibus, quæ sunt Dei. Dicúntur autem módii, quóniam sicut módii désuper quidem vácui sunt, desúbtus autem pleni: sic omnes mundi amatóres, et carnis, in spirituálibus quidem et divínis, quæ superióra recte esse dicúntur, vácui sunt, et insensáti: in inferióribus autem, id est, in mundiálibus, et terrénis aliquátenus pleni vidéntur, et sapiéntes. Proptérea sicut in módio si volúeris áliquid repónere, necésse est, ut cadat deórsum ad inferióra ejus, et tenebrósa: sic hómini carnáli, et mundiáli, et vácuo, si díxeris áliquid, quod est Dei, non illud tenet.

\begin{tightlist}{}
\item \textred{℟.} Vir Israelíta, gaude cohéres Christi: \textred{*} Ora pro nobis. \textred{℣.} Ut précibus tuis a Deo véniam sperémus, devoti ad te confúgimus. \textred{O}ra.
\end{tightlist}

\subsubsection*{\center\textred{Lectio ix.}}
\lettrine{N}unc autem sciens, quóniam ille est laudábilis, qui facit quod docet, addit, et dicit; sic lúceat lux vestra coram homínibus, ut vídeant ópera vestra bona, et magníficent Patrem vestrum, qui in cœlis est: id est, sic illumináte, non ut verba vestra tantúmmodo áudiant hómines, sed ut ópera vestra bona vídeant, ut quos illuminavéritis per verbum, condiátis per exémpla óperum vestrórum: quóniam qui docet, et facit quod docet, vere ille docet: qui autem non facit quod docet, non álium docet, sed seípsum condémnat. Et mélius est fácere, et non docére, quam docére, et non fácere; quóniam qui facit, et si tacúerit, áliquos córrigit suo exémplo.\\

\begin{tightlist}{}
\item \textred{℟.} Sancte Joánnes Christi Conféssor, audi rogántes sérvulos, \textred{*} Et impetrátam nobis cœ́litus tu defer indulgéntiam. \textred{℣.} O sancte Joánnes sidus áureum Dómini grátia, servórum gémitus sólita súscipe cleméntia. \textred{E}t. Glória Patri. \textred{E}t.
\item Te Deum laudámus. 6. 
\item \textred{Sacerd. ℣.} Os justi meditábitur sapiéntiam. \textred{℟.} Et lingua ejus loquétur judícium, allelúja.
\end{tightlist}

\subsubsection*{\center \textsc{AD LAUDES.}}
\begin{center}\textred{et per Horas, Aña.}\end{center}
\lettrine{A}ddúcam eum in montem sanctum meum, et lætificábo eum in domo oratiónis meæ.

\begin{tightlist}{}
\item \textred{Ps. de Dominica.} 30.
\item \textred{Aña.} Apéruit os suum in oratióne, et Dóminus spíritu intelligéntiæ replévit illum.
\item \textred{Aña.} Cibábit illum Dóminus pane vitæ, et intelléctus, et aqua sapiéntiæ salutáris potávit illum.
\item \textred{Aña.} Vir sápiens plebem suam érudit, et fructus sensus illíus fidéles sunt.
\item \textred{Aña.} Ipse tamquam imbres mittet elóquia sapiéntiæ suæ.
\end{tightlist}

\subsubsection*{\center \textred{Capitulum. \hfill Eccl. 31.}}
\lettrineThree{B}eátus vir, qui invéntus est sine mácula, et qui post aurum non ábiit, nec sperávit in pecúnia et thesáuris. Quis est hic, et laudábimus eum? fecit enim mirabília in vita sua.

\subsubsection*{\center\textred{Hymnus.}}
\lettrine{D}um Crucem gestat Dóminus, Joánni\\ Alloquens offert méritas corónas,\\ Is crucis tantum cálices amáros,\\ Súmere poscit.\\
\\
Sunt pati ac sperni sua vota, merces;\\ Flamma sed nunquam saturátur ingens:\\ Emori nulli sub honóre notus,\\ Instat, habétque.\\
\\
Pérgraves tandem Crucis hos labóres\\ Morte conclúdit, meritúmque pandit\\ Finis extrémi globus igne flagrans\\ Lapsus Olympo.\\
\\
Fit micans æther rádiis, cubíle,\\ Lámpadum lumen tenebréscit omne,\\ Testis est fragrans odor ipse vitæ\\ Córpore sparsus.\\
\\
Summa laus Patri, genitóque Verbo,\\ Et tibi compar, utriúsque Flamen,\\ Quo Beatórum pia turba semper\\ Páscitur igne.\\
Amen.

\begin{tightlist}{}
\item \textred{℟.} Justus germinábit sicut lílium. \textred{℟.} Et florébit in ætérnum ante Dóminum, allelúja.
\item \textred{Ad Benedictus Aña.} Dedit illi sciéntiam sanctórum, honestávit illum in labóribus, et complévit labóres illíus.
\end{tightlist}

\subsubsection*{\center\textred{Oratio.}}
\lettrine{D}eus, qui sanctum Joánnem Confessórem tuum, atque Doctórem, Patrem nostrum perféctæ sui abnegatiónis, et crucis amatórem exímium effecísti: concéde, ut ejus imitatióni júgiter inhæréntes, glóriam assequámur ætérnam. Per Christum.

\subsubsection*{\center \textsc{AD VESPERAS.}}
\begin{tightlist}{}
\item \textred{Añæ Laud. Ps. et capit. de Communi. [93]}
%\item \textred{Hymn.} Regis ætérni. \textred{ut in j. Vesperis.}
\end{tightlist}
\subsubsection*{\center\textred{Hymnus.}}
\lettrine{R}egis ætérni generóse miles,\\ Ordinis nostri cólumen Joánnes,\\ Quos tuis læti méritis dicámus,\\ Súscipe plausus.\\
\\ 
O satis felix! Speculátor alti\\ Núminis, votis animóque Martyr,\\ Pœ́nitens Virgo, memoránde vates,\\ Mystice Doctor.\\
\\
Sæpe cum Christo, et Genitríce pura\\ Verba commísces, superísque vivens:\\ Hinc fluit candor, tua quo refúlgent\\ Scripta, perénnis.\\
\\
Mente perfúsa rádiis ab alto,\\ Montis ascénsum, tenebrásque noctis,\\ Et facem vivam récolens amóris,\\ Alta revélas.\\
\\
Te sacri Verbi retegénte sensus,\\ Visa calígo procul exsuláre,\\ Dum facis lumen rútilans oríri\\ Noctis ab umbris.\\
\\
Pérsonent laudes Tríadi Joánnis,\\ Cui Pater lumen, tribuítque dona\\ Spíritus, vitæ réserat volúmen\\ Cœ́licus Agnus.\\
Amen.
\begin{tightlist}{}
\item \textred{℣.} Ora pro nobis sancte Pater Joánnes. \textred{℟.} Ut digni efficiámur promissiónibus Christi, allelúja. 
\item \textred{Ad Magnificat, Aña.} Quæsívit verba utília, et conscrípsit sermónes rectíssimos, et veritáte plenos.
\item \textred{Et fit comm. S. Bibianæ, Aña.} Accínxit fortitúdine lumbos suos, et roborávit bráchium suum: ideóque lucérna ejus non exstinguétur in sempitérnum.
\item \textred{℣.} Diffúsa est grátia in lábiis tuis. \textred{℟.} Proptérea benedíxit te Deus in ætérnum.
\item Orémus. \hfill \textred{Oratio.}
\end{tightlist}
\lettrine{D}eus, ómnium largítor bonórum, qui in fámula tua Bibiána cum virginitátis flore martýrii palmam conjunxísti: mentes nostras ejus intercessióne tibi caritáte conjúnge; ut, amótis perículis, prǽmia consequámur ætérna. Per Christum.

\begin{tightlist}{}
\item \textred{Et fit Commemoratio Feriæ.}
\end{tightlist}

\begin{center}\textcolor{red}{\rule{\linewidth}{0.5pt}} \end{center}

\subsection*{\center \textred{DIE II. JANUARII.}}
\subsubsection*{\center \textred{In Octava}\\ \textsc{Sancti Stephani.}}
\subsubsection*{\center \textred{Duplex minus.}}

\begin{tightlist}{}
\item \textred{Omnia fiant sicut in die Festi.} 387. \textred{præter hic signata.}
\end{tightlist}
\fancyhead[C]{\textsc{2 Jan. - Dies Octava S. Stephani.}}

\subsubsection*{\center \textsc{IN PRIMO NOCTURNO.}}
\subsubsection*{\center De Actibus Apostolórum.}
\begin{center}\textred{\textbf{Lectio j}. \hfill Cap. 7.}\end{center}
\lettrineThree{D}ura cervíce, et incircumcísis córdibus et áuribus, vos semper Spirítui sancto resístitis, sicut patres vestri, ita et vos.
Quem prophetárum non sunt persecúti patres vestri? et occidérunt eos, qui prænuntiábant de advéntu Justi, cujus vos nunc proditóres, et homicídæ fuístis:
qui accepístis legem in dispositióne Angelórum, et non custodístis.
Audiéntes autem hæc dissecabántur córdibus suis, et stridébant déntibus in eum.

\begin{tightlist}{}
\item \textred{℟.} Stéphanus autem. 388.
\end{tightlist}

\subsubsection*{\center \textred{Lectio ij.}}
\lettrine{C}um autem esset plenus Spíritu sancto, inténdens in cælum, vidit glóriam Dei, et Jesum stantem a dextris Dei, et ait: Ecce vídeo cælos apértos, et Fílium hóminis stantem a dextris Dei.
Exclamántes autem voce magna continuérunt aures suas, et ímpetum fecérunt unanímiter in eum.
Et ejiciéntes eum extra civitátem lapidábant: et testes deposuérunt vestiménta sua secus pedes adolescéntis, qui vocabátur Saulus.
Et lapidábant Stéphanum invocántem, et dicéntem: Dómine Jesu, súscipe spíritum meum.

\begin{tightlist}{}
\item \textred{℟.} Vidébant. 389.
\end{tightlist}

\subsubsection*{\center \textred{Lectio iij.}}
\lettrine{P}ósitis autem génibus, clamávit voce magna, dicens: Dómine ne státuas illis hoc peccátum. Et cum hoc dixísset, obdormívit in Dómino. Saulus autem erat conséntiens neci ejus.
Facta est autem in illa die persecútio magna in Ecclésia, quæ erat Jerosólymis, et omnes dispérsi sunt per regiónes Judǽæ et Samaríæ præter Apóstolos.
Curavérunt autem Stéphanum viri timoráti, et fecérunt planctum magnum super eum.

\begin{tightlist}{}
\item \textred{℟.} Lapidavérunt. 390.
\end{tightlist}

\subsubsection*{\center \textsc{IN SECUNDO NOCTURNO.}}
\subsubsection*{\center Sermo sancti Augustíni Epíscopi.}
\begin{center}\textred{Sem. 2. de S. Steph.\\ \textbf{Lectio iv.}}\end{center}
\lettrineThree{P}ost hestérnum festivíssimum diem, quo Salvatóris nostri Christi nobis Natívitas illúxit, étiam hodiérnus dies beáti mártyris Stéphani coróna illustrátur. Mártyris illíus mérita nulla pars orbis ignórat: passus enim est in ipso princípio Ecclésiæ, id est, in ipsa urbe Jerosólyma. Ibi enim Diáconus ministrávit, et in ipso juventútis flore decórem ætátis suæ sánguine purpurávit. Pássio ejus insígnis est, multúmque mirábilis: hanc modo de libro Actuum Apostolórum, cum legerétur, non solum audívimus, sed étiam óculis spectávimus.

\begin{tightlist}{}
\item \textred{℟.} Vidébant. 389.
\end{tightlist}

\subsubsection*{\center \textred{Lectio v.}}
\lettrine{C}hristus ergo caput Mártyrum prior passus est pro nobis, relínquens vobis exémplum, ut sequámini vestígia ejus. Cujus passiónis vestígia prior secútus beátus Stéphanus, confiténdo Christum lapidátus a Judǽis, corónam méruit, tamquam suo sibi nómine pósitam. Stéphanus enim Græce, Latíne coróna appellátur. Jam coróna nomen habébat, et ídeo palmam martýrii suo nómine præferébat. 

\begin{tightlist}{}
\item \textred{℟.} Stéphanus servus. 390.
\end{tightlist}

\subsubsection*{\center \textred{Lectio vj.}}
\lettrine{Q}ui cum lapidarétur, non solum non exspectábat de persecutóribus reportáre vindíctam, sed eis pótius a Deo véniam postulábat. Memínerat enim dixísse Dóminum: Mihi vindíctam, ego retríbuam, dicit Dóminus. Et íterum: Ne díxeris, Ulcíscar me de inimícis meis: sed exspécta Dóminum, ut tibi auxílio sit. Exspectáre nos jubet Dóminus Deus, ut in die futúræ retributiónis cum sanctis Martýribus vindicémur. 

\begin{tightlist}{}
\item \textred{℟.} Impetum. 391.
\end{tightlist}

\subsubsection*{\center \textsc{IN TERTIO NOCTURNO.}}
\begin{tightlist}{}
\item \textred{Lect. ut in die.} 391.
\end{tightlist}

\subsubsection*{\center \textsc{AD LAUDES.}}
\begin{tightlist}{}
\item \textred{Omnia ut in die, cum Oratione.} Omnípotens, \textred{et Comm. ut infra.}
\end{tightlist}
\subsubsection*{\center \textred{Oratio.}}
\lettrine{O}mnípotens sempitérne Deus, qui primítias Mártyrum in beáti Levítæ Stéphani sánguine dedicásti: tríbue, quǽsumus, ut pro nobis intercéssor exsístat, qui pro suis étiam persecutóribus exorávit Dóminum nostrum Jesum Christum Fílium tuum. Qui tecum vivit.

\begin{tightlist}{}
\item \textred{Comm. S. Euphrosynæ, Aña.} Símile est regnum cœlórum sagénæ missæ in mare, et ex omni génere píscium congregánti: quam cum impléta esset, educéntes, et secus littus sedéntes, elegérunt bonos in vasa sua, malos autem foras misérunt.
\item \textred{℣.} Elégit eam Deus, et præelégit eam. \textred{℟.} Et habitáre eam facit in tabernáculo suo.
\item Orémus. \hfill \textred{Oratio.}
\end{tightlist}
\lettrine{D}eus, qui sanctam Euphrósynam Vírginem ad sanctitátis stúdium mirabíliter excitásti, et in eódem júgiter direxísti: fac nos, ejus méritis et intercessióne, perfécta caritáte fervéntes. Per Christum.

\begin{tightlist}{}
\item \textred{Et comm. Oct. S. Joannis, Ana.} Hic est discípulus ille, qui testimónium perhíbuit: et scimus, quia verum est testimónium ejus.
\item \textred{℣.} Virgo est eléctus a Dómino. \textred{℟.} Atque inter céteros magis diléctus, allelúja.
\item Orémus. \hfill \textred{Oratio.}
\end{tightlist}
\lettrine{E}cclésiam tuam, Dómine, benígnus illústra: ut beáti Joánnis Apóstoli tui et Evangelístæ illumináta doctrínis, ad dona pervéniat sempitérna. Per Christum.


\begin{tightlist}{}
\item \textred{Et comm. Oct. SS. Innocentium, Ana.} Heródes irátus occídit multos púeros in Béthlehem Judæ civitáte David.
\item \textred{℣.} Ex ore infántium et lacténtium. \textred{℟.} Perfecísti laudem, Dómine Deus noster, allelúja.
\item Orémus. \hfill \textred{Oratio.}
\end{tightlist}
\lettrine{D}eus, cujus hodiérna die præcónium Innocéntes mártyres non loquéndo, sed moriéndo conféssi sunt: ómnia in nobis vitiórum mala mortífica; ut fidem tuam, quam lingua nostra lóquitur, étiam móribus vita fateátur. Qui vivis.


\begin{tightlist}{}
\item \textred{Ad Horas, ut in die, cum Oratione} Omnípotens.
\end{tightlist}

\subsubsection*{\center \textsc{AD VESPERAS.}}
\begin{tightlist}{}
\item \textred{Aña.} Tecum princípium, \textred{et cetera omnia ut in ij. Vesp. Festi,} 395. \textred{cum Oratione} Omnípotens.
\item \textred{Fit Comm. S. Euphrosynæ, Aña.} Quinque prudéntes Vírgines accepérunt óleum in vasis suis cum lampádibus. Média autem nocte clamor factus est: Ecce sponsus venit, exíte óbviam ei.
\item \textred{℣.} Diffúsa est grátia in lábiis tuis. \textred{℟.} Proptérea benedíxit te Deus in ætérnum.
\item \textred{Oratio} Deus, qui sanctam \textred{ut supra.}
\item \textred{Deinde pro Commem. S. Joannis, Aña.} Ecce ego Joánnes vidi óstium apértum in cœlo: et ecce sedes pósita erat in eo, et in médio sedis et in circúitu ejus quátuor animália plena óculis ante et retro, et dabant glóriam, et honórem, et benedictiónem sedénti super thronum, vivénti in sǽcula sæculórum.
\item \textred{℣.} Valde honorándus est beátus Joánnes. \textred{℟.} Qui supra pectus Dómini in cœna recúbuit, allelúja.
\item \textred{Oratio} Ecclésiam tuam \textred{ut supra.}
\item \textred{Pro Oct. SS. Innoc. Aña.} A bimátu et infra occídit multos púeros Heródes propter Dóminum.
\item \textred{℣.} Hi sunt qui cum muliéribus non sunt coinquináti. \textred{℟.} Vírgines enim sunt, et sequúntur Agnum quocúmque íerit, allelúja.
\item \textred{Oratio} Deus, cujus hodiérna \textred{ut supra.}
\end{tightlist}

\begin{center}\textcolor{red}{\rule{\linewidth}{0.5pt}} \end{center}

\subsection*{\center \textred{DIE III. JANUARII.}}
\subsubsection*{\center \textred{In Octava}\\ \textsc{S. Joannis Evangelistæ}}
\subsubsection*{\center \textred{Duplex minus.}}

\begin{tightlist}{}
\item \textred{Omnia fiant sicut in die Festi.} 396. \textred{exceptis hic signatis.}
\end{tightlist}
\fancyhead[C]{\textsc{3 Jan. - Dies Octava S. Joannis.}}
\subsubsection*{\center \textsc{IN PRIMO NOCTURNO.}}
\subsubsection*{\center De libro Apocalýpsis beáti Joánnis Apóstoli.}
\begin{center}\textred{\textbf{Lectio j}. \hfill Cap. 4.}\end{center}
\lettrineThree{P}ost hæc vidi: et ecce óstium apértum in cælo: et vox prima, quam audívi, tamquam tubæ loquéntis mecum, dicens: Ascénde huc, et osténdam tibi, quæ opórtet fíeri post hæc.
Et statim fui in spíritu: et ecce sedes pósita erat in cælo, et supra sedem sedens.
Et qui sedébat, símilis erat aspéctui lápidis jáspidis, et sardínis: et iris erat in circúitu sedis, símilis visióni smarágdinæ.
Et in circúitu sedis sedília vigíntiquátuor, et super thronos vigíntiquátuor senióres sedéntes, circumamícti vestiméntis albis, et in capítibus eórum corónæ áureæ.

\begin{tightlist}{}
\item \textred{℟.} Valde honorándus. 398.
\end{tightlist}

\subsubsection*{\center \textred{Lectio ij.}}
\lettrine{E}t de throno procedébant fúlgura, et voces, et tonítrua: et septem lámpades ardéntes ante thronum, qui sunt septem Spíritus Dei.
Et in conspéctu sedis tamquam mare vítreum símile crystállo: et in médio sedis, et in circúitu sedis quátuor animália plena óculis ante et retro.
Et ánimal primum símile leóni, et secúndum ánimal símile vítulo, et tértium ánimal habens fáciem quasi hóminis, et quartum ánimal símile áquilæ volánti.
Et quátuor animália, síngula eórum habébant alas senas: et in circúitu, et intus plena sunt óculis.

\begin{tightlist}{}
\item \textred{℟.} In illa die. 398.
\end{tightlist}

\subsubsection*{\center \textred{Lectio iij.}}
\lettrine{E}t réquiem non habébant die ac nocte, dicéntia: Sanctus, Sanctus, Sanctus Dóminus Deus omnípotens, qui erat, et qui est, et qui ventúrus est.
Et cum darent illa animália glóriam, et honórem, et benedictiónem sedénti super thronum, vivénti in sǽcula sæculórum,
procidébant vigínti quátuor senióres ante sedéntem in throno, et adorábant vivéntem in sǽcula sæculórum, et mittébant corónas suas ante thronum, dicéntes:
Dignus es Dómine Deus noster accípere glóriam, et honórem, et virtútem: quia tu creásti ómnia, et propter voluntátem tuam erant, et creáta sunt.

\begin{tightlist}{}
\item \textred{℟.} Qui vícerit. 399.
\end{tightlist}

\subsubsection*{\center \textsc{IN SECUNDO NOCTURNO.}}
\subsubsection*{\center Ex Tractátu sancti Augustíni Epíscopi in Joánnem.}
\begin{center}\textred{\textbf{Lectio iv}. \hfill Tract. 36.}\end{center}
\lettrineThree{I}n quátuor Evangéliis, vel pótius quátuor libris uníus Evangélii, sanctus Joánnes Apóstolus non immérito secúndum intelligéntiam spiritálem áquilæ comparátus, áltius multóque sublímius áliis tribus eréxit prædicatiónem suam; et in ejus erectióne étiam corda nostra érigi vóluit. Nam céteri tres Evangelístæ, tamquam cum hómine Dómino in terra ambulántes, de divinitáte ejus pauca dixérunt: istum autem quasi pigúerit in terra ambuláre, sicut ipso exórdio sui sermónis intónuit, eréxit se non solum super terram, et super omnem ámbitum áeris, et cœli, sed super omnem étiam exércitum Angelórum, omnémque constitutiónem invisibílium potestátum; et pervénit ad eum, per quem facta sunt ómnia, dicéndo: In princípio erat Verbum, et Verbum erat apud Deum, et Deus erat Verbum.

\begin{tightlist}{}
\item \textred{℟.} Hic est. 399.
\end{tightlist}

\subsubsection*{\center \textred{Lectio v.}}
\lettrine{H}uic tantæ sublimitáti princípii étiam cétera cóngrua prædicávit: et de Dómini divinitáte, quómodo nullus álius, est locútus: hoc eructábat, quod bíberat. Non enim sine causa de illo in isto ipso Evangélio narrátur, quia et in convívio super pectus Dómini discumbébat. De illo ergo péctore in secréto bíberat: sed quod in secréto bibit, in manifésto eructávit: ut pervéniat ad omnes Gentes non solum incarnátio Fílii Dei, pássio, et resurréctio: sed étiam quod erat ante incarnatiónem, únicus Patri, Verbum Patris, coætérnus generánti, æquális ei, a quo missus est.

\begin{tightlist}{}
\item \textred{℟.} Hic est discípulus. 400.
\end{tightlist}

\subsubsection*{\center \textred{Lectio vj.}}
\lettrine{A}quila ipse est Joánnes, sublímium prædicátor, et lucis intérnæ atque ætérnæ fixis óculis contemplátor. Dicúntur enim et pulli aquilárum a paréntibus sic probári: patris scílicet ungue suspéndi, et rádiis solis oppóni: qui firme contemplátus fúerit, fílius agnóscitur: si ácie palpitáverit, tamquam adulterínus ab ungue dimíttitur. Jam ergo vidéte, quam sublímia loqui débuit, qui est áquilæ comparátus: et tamen étiam nos humi repéntes, infírmi, et vix ullíus moménti inter hómines, audémus tractáre ista, et ista expónere, et putámus nos aut cápere posse, cum cogitámus, aut capi, dum dícimus.

\begin{tightlist}{}
\item \textred{℟.} Diligébat. 400.
\end{tightlist}

\subsubsection*{\center \textsc{IN TERTIO NOCTURNO.}}
\begin{tightlist}{}
\item \textred{Omnia ut in die.} 400.
\end{tightlist}

\subsubsection*{\center \textsc{AD LAUDES.}}
\begin{tightlist}{}
\item \textred{Omnia ut in die.} 402. \textred{sed fiat commem. SS. Innocentium, Aña.} Heródes irátus occídit multos púeros in Béthlehem Judæ civitáte David.
\item \textred{℣.} Ex ore infántium et lacténtium. \textred{℟.} Perfecísti laudem, Dómine Deus noster, allelúja.
\item Orémus. \hfill \textred{Oratio.}
\end{tightlist}
\lettrine{D}eus, cujus hodiérna die præcónium Innocéntes mártyres non loquéndo, sed moriéndo conféssi sunt: ómnia in nobis vitiórum mala mortífica; ut fidem tuam, quam lingua nostra lóquitur, étiam móribus vita fateátur. Qui vivis.


\subsubsection*{\center \textsc{AD VESPERAS.}}
\begin{tightlist}{}
\item \textred{Aña.} Tecum princípium, \textred{et cetera omnia ut in die.} 404.
\item \textred{Pro comm. SS. Innocentium. Aña.} Innocéntes pro Christo infántes occísi sunt, ab iníquo rege lacténtes interfécti sunt; ipsum sequúntur Agnum sine mácula: et dicunt semper, Glória tibi Dómine.
\item \textred{℣.} Hi sunt, qui cum muliéribus non sunt coinquináti. \textred{℟.} Vírgines enim sunt: et sequúntur Agnum quocúmque íerit, allelúja.
\item \textred{Oratio} Deus, cujus hodiérna, \textred{ut supra.}
\end{tightlist}

\begin{center}\textcolor{red}{\rule{\linewidth}{0.5pt}} \end{center}

\subsection*{\center \textred{DIE IV. JANUARII.}}
\subsubsection*{\center \textred{In Octava}\\ \textsc{SS. Innocentium.}}
\subsubsection*{\center \textred{Duplex minus.}}

\begin{tightlist}{}
\item \textred{Omnia ut in Festo,} 405. \textred{præter hic signata.}
\end{tightlist}

\subsubsection*{\center \textsc{IN PRIMO NOCTURNO.}}
\subsubsection*{\center De Epístola beáti Pauli Apóstoli ad Romános.}
\begin{center}\textred{\textbf{Lectio j}. \hfill Cap. 5.}\end{center}
\lettrineThree{J}ustificáti ergo ex fide, pacem habeámus ad Deum per Dóminum nostrum Jesum Christum:
per quem et habémus accéssum per fidem in grátiam istam, in qua stamus, et gloriámur in spe glóriæ filiórum Dei.
Non solum autem, sed et gloriámur in tribulatiónibus: sciéntes quod tribulátio patiéntiam operátur:
patiéntia autem probatiónem, probátio vero spem,
spes autem non confúndit: quia cáritas Dei diffúsa est in córdibus nostris per Spíritum sanctum, qui datus est nobis.

\begin{tightlist}{}
\item \textred{℟.} Sub altáre. 406.
\end{tightlist}

\subsubsection*{\center \textred{Lectio ij.}}
\lettrine{U}t quid enim Christus, cum adhuc infírmi essémus, secúndum tempus pro ímpiis mórtuus est?
Vix enim pro justo quis móritur: nam pro bono fórsitan quis áudeat mori?
Comméndat autem caritátem suam Deus in nobis: quóniam cum adhuc peccatóres essémus, secúndum tempus
Christus pro nobis mórtuus est: multo ígitur magis nunc justificáti in sánguine ipsíus, salvi érimus ab ira per ipsum.

\begin{tightlist}{}
\item \textred{℟.} Effudérunt. 407.
\end{tightlist}

\subsubsection*{\center \textred{Lectio iij.}}
\lettrine{S}i enim cum inimíci essémus, reconciliáti sumus Deo per mortem Fílii ejus: multo magis reconciliáti, salvi érimus in vita ipsíus.
Non solum autem, sed et gloriámur in Deo per Dóminum nostrum Jesum Christum, per quem nunc reconciliatiónem accépimus.
Proptérea sicut per unum hóminem peccátum in hunc mundum intrávit, et per peccátum mors, et ita in omnes hómines mors pertránsiit, in quo omnes peccavérunt.

\begin{tightlist}{}
\item \textred{℟.} Sub throno. 407.
\end{tightlist}

\subsubsection*{\center \textsc{IN SECUNDO NOCTURNO.}}
\subsubsection*{\center Sermo sancti Augustíni Epíscopi.}
\begin{center}\textred{Serm. 1. de Innocent.\\ qui est 8. de SS.\\ \textbf{Lectio iv.}}\end{center}
\lettrineThree{N}ascénte Dómino, luctus cœpit, non cœlo, sed mundo: indícitur mátribus lamentátio, Angelis exsultátio, infántibus transmigrátio. Deus est, qui natus est: innocéntes illi debéntur víctima, qui venit damnáre mundi malítiam. Agnélli debent immolári, quia Agnus futúrus est crucifígi, qui tollit peccáta mundi. Sed oves úlulant matres, quia agnos perdunt sine voce balántes. Grande martýrium, crudéle spectáculum! Exímitur machǽra, et nulla intervénit causa: sola stridet invídia, cum qui natus est, nulli fáciat violéntiam: sed oves cérnimus matres, quæ super agnus lugent. Vox in Rama audíta est: plorátus et ululátus magnus: pígnora sunt, non crédita, sed creáta; non depósita, sed expósita.

\begin{tightlist}{}
\item \textred{℟.} Adoravérunt. 408.
\end{tightlist}

\subsubsection*{\center \textred{Lectio v.}}
\lettrine{T}estis est ipsa natúra, quæ pugnábat in mártyres, cum tyránnus jactáret infántes. Mater crines cápitis dissipábat, quæ ornaméntum cápitis amittébat. Quantis modis infántem volébat abscóndere, et ipse se infántulus publicábat! Nesciébat tacére, quia necdum didícerat formidáre. Pugnábat mater, et cárnifex: ille trahébat, illa tenébat. Ad carníficem mater clamábat: Quid séparas a me, quem génui ex me? Uterus génuit, non mansit ille, cum vixit, úbera mea frustra lacte contórquens; caute portávi, quem a te vídeo manu crudéli jactári. Modo eum effudérunt víscera mea, et tu elídis ad terram.

\begin{tightlist}{}
\item \textred{℟.} Isti sunt. 408.
\end{tightlist}

\subsubsection*{\center \textred{Lectio vj.}}
\lettrine{A}lia acclamábat mater, cum exáctor latro non compélleret simul occídi cum párvulo matrem: Ut quid me dimíttis inánem? Si culpa est, mea est: si non est crimen, junge mortem, et líbera matrem. Alia dicébat: Quid quǽritis? unum quǽritis, et multos occíditis: et ad unum, qui unus est, attíngere non potéstis. Alia contra clamábat: Veni jam, veni Salvátor mundi, quámdiu quæréris? nullum times, vídeat te miles, et nostros non occídat infántes. Miscebátur lamentátio matrum, et ad cœlum transíbat oblátio parvulórum.

\begin{tightlist}{}
\item \textred{℟.} Cantábant. 408.
\end{tightlist}

\subsection*{\center \textsc{IN TERTIO NOCTURNO.}}
\begin{tightlist}{}
\item \textred{Omnia ut in die.} 408. \textred{et dicitur} Te Deum laudámus. 6.
\item \textred{Ad Laudes, Horas, et ij. Vesp. omnia ut in die, sine comm. Octavarum.}
\end{tightlist}


\begin{center}\textcolor{red}{\rule{\linewidth}{0.5pt}} \end{center}


\section*{\center \textsc{Officia Passionis D.N.J.C.}}

\begin{center} \textsc{\textred{Mobilibus diebus post Septuagesimam et non ultra Quadragesimam recitanda.\\
Ex Decr. S. R. C. 31. Januarii 1879.}}\end{center}

\begin{center}\textcolor{red}{\rule{\linewidth}{0.5pt}} \end{center}
\fancyhead[C]{\textsc{Orationis D.N.J.C. in Monte Oliveti}} % Book title in header center for even pages
\subsection*{\center \textred{\textsc{FERIA TERTIA}}}
\subsubsection*{\center \textred{post Domin. Septuages.}}
\subsection*{\center \textsc{Orationis D.N.J.C.\\ in Monte Oliveti.}}
\subsubsection*{\center \textred{Duplex minus. P. C.}}
%\begin{center}\textcolor{red}{\rule{\linewidth}{0.5pt}} \end{center}

\subsubsection*{\center \textsc{AD VESPERAS.}}
\begin{tightlist}{}
\item \textred{Aña.} Venit Jesus cum discípulis suis in villam, quæ dícitur Gethsémani, et dixit illis: Sedéte hic, donec vadam illuc et orem.
\item \textred{Ps. de Offic. præced.}
\end{tightlist}

\subsubsection*{\textred{Capitulum. \hfill Matth. 26.}}
\lettrineThree{V}enit Jesus ad discípulos suos, et ínvenit eos dormiéntes, et dicit Petro: Sic non potuístis una hora vigiláre mecum? Vigiláte et oráte ut non intrétis in tentatiónem.

\begin{tightlist}{}
\item \textred{℟.} Prócidit Jesus in fáciem suam orans et dicens: Pater mi, si possíbile est, tránseat a me calix iste. \textred{*} Et factus in agonía prolíxius orábat. 
\item \textred{℣.} Et factus est sudor ejus sicut guttæ sánguinis decurréntis in terram. \textred{E}t factus. Glória Patri. \textred{E}t factus. 
\end{tightlist}

\subsubsection*{\center\textred{Hymnus.}}
\lettrine{A}dspice ut Verbum Patris a supérnis,\\ Sédibus clemens, et amóre flagrans,\\ Pérditis culpa génitis medéri\\ Pergit Adámi.\\
\\
Flébilem mundi míserans ruínam,\\ Et volens nostros reparáre casus,\\ Orat, et prona véniam precátur\\ Fronte Magíster.\\
\\
Flúctuat secum tot acérba volvens:\\ Hunc, ait supplex, cálicem dolóris,\\ Mi Pater, transfer; tua sed volúntas,\\ Non mea, fiat.\\
\\
Cum premat tristis pavor ima cordis,\\ Déficit languens Dóminus; per artus\\ Sánguinis sudor fluit, atque guttis\\ Terra madéscit.\\
\\
At celer summo véniens Olympo,\\ Angelus Jesum récreat jacéntem:\\ Córpori vires rédeunt novóque\\ Róbore surgit.\\
\\
Laus, honor Patri genitǽque Proli,\\ Cui datum nomen super omne nomen,\\ Et Paracléto decus atque virtus\\ Omne per ævum.\\
Amen.

\begin{tightlist}{}
\item \textred{℣.} Tristis est ánima mea. \textred{℟.} Usque ad mortem.
\item \textred{Ad Magnif. Aña.} Pósitis génibus orábat dicens: Pater, si vis, transfer cálicem istum a me; verúmtamen non mea volúntas sed tua fiat.
\end{tightlist}

\subsubsection*{\center \textred{Oratio.}}
\lettrine{D}omine Jesu Christe, qui in horto, verbo et exémplo nos oráre docuísti ad tentatiónum perícula superánda, concéde propítius, ut nos oratióni semper inténti, ejus copiósum fructum cónsequi mereámur. Qui vivis.

\begin{tightlist}{}
\item \textred{In fine Hymn. Completorii et Horarum dic. ℣.} Glória tibi Dómine. Qui nos salvásti in Cruce.
\item \textred{Ad Nunc dimittis, Aña.} Hoc signum Crucis.
\end{tightlist}

\subsubsection*{\center \textsc{AD MATUTINUM.}}
\begin{tightlist}{}
\item \textred{Invitat.} Christum Jesum in monte Olivárum orántem * Veníte adorémus.
\item \textred{Ps.} Veníte, exsultémus. 2.
\item \textred{Hymn. ut supra ad Vesp.}
\end{tightlist}

\subsubsection*{\center \textsc{IN PRIMO NOCTURNO.}}
\begin{tightlist}{}
\item \textred{Aña.} Ante oratiónem prǽpara ánimam tuam, et noli esse quasi homo, qui tentat Deum.
\item \textred{Ps. 16.} Exáudi, Dómine, 67.
\end{tightlist}
%\subsubsection*{\center Psalmus xvj}
%\lettrineThree{E}xáudi, Dómine, justítiam meam: * inténde deprecatiónem meam.
%
%Auribus pércipe oratiónem meam, * non in lábiis dolósis.
%
%De vultu tuo judícium meum pródeat: * óculi tui vídeant æquitátes.
%
%Probásti cor meum, et visitásti nocte: * igne me examinásti, et non est invénta in me iníquitas.
%
%Ut non loquátur os meum ópera hóminum: * propter verba labiórum tuórum ego custodívi vias duras.
%
%Pérfice gressus meos in sémitis tuis: * ut non moveántur vestígia mea.
%
%Ego clamávi, quóniam exaudísti me, Deus: * inclína aurem tuam mihi, et exáudi verba mea.
%
%Mirífica misericórdias tuas, * qui salvos facis sperántes in te.
%
%A resisténtibus déxteræ tuæ custódi me, * ut pupíllam óculi.
%
%Sub umbra alárum tuárum prótege me: * a fácie impiórum qui me afflixérunt.
%
%Inimíci mei ánimam meam circumdedérunt, ádipem suum conclusérunt : * os eórum locútum est supérbiam.
%
%Projiciéntes me nunc circumdedérunt me: * óculos suos statuérunt declináre in terram.
%
%Suscepérunt me sicut leo parátus ad prædam: * et sicut cátulus leónis hábitans in ábditis.
%
%Exsúrge, Dómine, prǽveni eum, et supplánta eum: * éripe ánimam meam ab ímpio, frámeam tuam ab inimícis manus tuæ.
%
%Dómine, a paucis de terra dívide eos in vita eórum: * de abscónditis tuis adimplétus est venter eórum.
%
%Saturáti sunt fíliis: * et dimisérunt relíquias suas párvulis suis.
%
%Ego autem in justítia apparébo conspéctui tuo: * satiábor cum apparúerit glória tua.
%
%Glória Patri, et Fílio, * et Spirítui Sancto.
%
%Sicut erat in princípio, et nunc, et semper, * et in sǽcula sæculórum. Amen.
\begin{tightlist}{}
%\item \textred{Aña.} Ante oratiónem præpara ánimam tuam, et noli esse quasi homo, qui tentat Deum.
\item \textred{Aña.} Non impediáris oráre semper, et ne vereáris usque ad mortem justificári.
\item \textred{Ps. 27.} Ad te, Dómine, 86.
\end{tightlist}

%\subsubsection*{\center Psalmus xvij}
%\lettrine{A}d te, Dómine, clamábo, Deus meus, ne síleas a me: * nequándo táceas a me, et assimilábor descendéntibus in lacum.
%
%Exáudi, Dómine, vocem deprecatiónis meæ dum oro ad te: * dum extóllo manus meas ad templum sanctum tuum.
%
%Ne simul trahas me cum peccatóribus: * et cum operántibus iniquitátem ne perdas me.
%
%Qui loquúntur pacem cum próximo suo, * mala autem in córdibus eórum.
%
%Da illis secúndum ópera eórum, * et secúndum nequítiam adinventiónum ipsórum.
%
%Secúndum ópera mánuum eórum tríbue illis, * redde retributiónem eórum ipsis.
%
%Quóniam non intellexérunt ópera Dómini, et in ópera mánuum ejus * déstrues illos, et non ædificábis eos.
%
%Benedíctus Dóminus, * quóniam exaudívit vocem deprecatiónis meæ.
%
%Dóminus adjútor meus, et protéctor meus : * in ipso sperávit cor meum et adjútus sum.
%
%Et reflóruit caro mea: * et ex voluntáte mea confitébor ei.
%
%Dóminus fortitúdo plebis suæ : * et protéctor salvatiónum Christi sui est.
%
%Salvum fac pópulum tuum, Dómine, et bénedic hereditáti tuæ, * et rege eos, et extólle illos usque in ætérnum.
%
%Glória Patri, et Fílio, * et Spirítui Sancto.
%
%Sicut erat in princípio, et nunc, et semper, * et in sǽcula sæculórum. Amen.
\begin{tightlist}{}
%\item \textred{Aña.} Non impediáris oráre semper, et ne vereáris usque ad mortem justificári.
\item \textred{Aña.} Omnia quæcúmque petiéritis in oratióne credéntes, accipiétis.
\item \textred{Ps. 141.} Voce mea, 238.
\end{tightlist}
%\subsubsection*{\center \textred{Psalmus cxlj}
%\lettrine{V}oce mea ad Dóminum clamávi: * voce mea ad Dóminum deprecátus sum:
%
%Effúndo in conspéctu ejus oratiónem meam, * et tribulatiónem meam ante ipsum pronúntio.
%
%In deficiéndo ex me spíritum meum: * et tu cognovísti sémitas meas.
%
%In via hac, qua ambulábam, * abscondérunt láqueum mihi.
%
%Considerábam ad déxteram, et vidébam: * et non erat qui cognósceret me.
%
%Périit fuga a me: * et non est qui requírat ánimam meam.
%
%Clamávi ad te, Dómine, * dixi: Tu es spes mea, pórtio mea in terra vivéntium.
%
%Inténde ad deprecatiónem meam: * quia humiliátus sum nimis.
%
%Líbera me a persequéntibus me: * quia confortáti sunt super me.
%
%Educ de custódia ánimam meam ad confiténdum nómini tuo: * me exspéctant justi, donec retríbuas mihi.
%
%Glória Patri, et Fílio, * et Spirítui Sancto.
%
%Sicut erat in princípio, et nunc, et semper, * et in sǽcula sæculórum. Amen.
\begin{tightlist}{}
%\item \textred{Aña.} Omnia quæcúmque petiéritis in oratióne credéntes, accipiétis.
\item \textred{℣.} Tristis est ánima mea. \textred{℟.} Usque ad mortem.
\end{tightlist}

\subsubsection*{\center De libro Tobíæ.}
\begin{center}\textred{\textbf{Lectio j}. \hfill Cap. 12.}\end{center}
\lettrineThree{B}ona est orátio cum jejúnio, et eleemósyna magis quam thesáuros auri recóndere; quóniam eleemósyna a morte líberat; et ipsa est quæ purgat peccáta et facit inveníre misericórdiam et vitam ætérnam. Qui autem fáciunt peccátum et iniquitátem, hostes sunt ánimæ suæ. Manifésto ergo vobis veritátem, et non abscóndam a vobis occúltum sermónem. Quando orábas cum lácrymis, et sepeliébas mórtuos, et derelinquébas prándium tuum, et mórtuos abscondébas per diem in domo tua, et nocte sepeliébas eos, ego óbtuli oratiónem tuam Dómino. Et quia accéptus eras Deo, necésse fuit ut tentátio probáret te.

\begin{tightlist}{}
\item \textred{℟.} In ómnibus grátias ágite: \textred{*} Sine intermissióne oráte. \textred{℣.} Hæc est enim volúntas Dei in Christo Jesu in ómnibus vobis. \textred{S}ine.
\end{tightlist}

\subsubsection*{\center De Epístola beáti Jacóbi Apóstoli.}
\begin{center}\textred{\textbf{Lectio ij}. \hfill Cap. 5.} \end{center}
\lettrine{O}ráte pro ínvicem, ut salvémini: multum enim valet deprecátio justi assídua. Elías homo erat símilis nobis passíbilis, et oratióne orávit, ut non plúeret super terram, et non pluit annos tres et menses sex. Et rursum orávit, et cœlum dedit plúviam, et terra dedit fructum suum. Fratres mei, si quis ex vobis erráverit a veritáte, et convérterit quis eum, scire debet, quóniam qui convérti fécerit peccatórem ab erróre viæ suæ, salvábit ánimam ejus a morte, et opériet multitúdinem peccatórum.

\begin{tightlist}{}
\item \textred{℟.} Cum stábitis ad orándum, \textred{*} Dimíttite si quid habétis advérsus áliquem. \textred{℣.} Ut et Pater vester, qui in cœlis est, dimíttat vobis peccáta vestra. \textred{D}imíttite.
\end{tightlist}
\subsubsection*{\center De Epístola beáti Pauli Apóstoli ad Hebrǽos.}
\begin{center}\textred{\textbf{Lectio iij}. \hfill Cap. 5.}\end{center}
\lettrine{C}hristus non semetípsum clarificávit, ut Póntifex fíeret: sed qui locútus est ad eum: Fílius meus es tu; ego hódie génui te. Quemádmodum et in álio loco dicit: Tu es Sacérdos in ætérnum secúndum órdinem Melchísedech. Qui in diébus carnis suæ preces, supplicationésque ad eum qui possit illum salvum fácere a morte, cum clamóre válido et lácrymis ófferens, exaudítus est pro sua reveréntia. Et quidem cum esset Fílius Dei, dídicit ex iis, quæ passus est, obediéntiam, et consummátus, factus est ómnibus obtemperántibus sibi causa salútis ætérnæ, appellátus a Deo Póntifex juxta órdinem Melchísedech.

\begin{tightlist}{}
\item \textred{℟.} Exaudívit Dóminus deprecatiónem meam, \textred{*} Dóminus oratiónem meam suscépit. \textred{℣.} Clamávi in toto corde meo, exáudi me, Dómine. \textred{D}óminus. Glória Patri. \textred{D}óminus.
\end{tightlist}

\subsubsection*{\center \textsc{IN SECUNDO NOCTURNO.}}
\begin{tightlist}{}
\item \textred{Aña.} Cum orátis non éritis sicut hypócritæ, qui amant in synagógis, et in ángulis plateárum stantes oráre, ut videántur ab homínibus.
\item \textred{Ps. 42.} Júdica me, 108.
\end{tightlist}

%\subsubsection*{\center Psalmus xlij}
%\lettrineThree{J}údica me, Deus, et discérne causam meam de gente non sancta, * ab hómine iníquo, et dolóso érue me.
%
%Quia tu es, Deus, fortitúdo mea: * quare me repulísti? et quare tristis incédo, dum afflígit me inimícus?
%
%Emítte lucem tuam et veritátem tuam: * ipsa me deduxérunt, et adduxérunt in montem sanctum tuum, et in tabernácula tua.
%
%Et introíbo ad altáre Dei: * ad Deum, qui lætíficat juventútem meam.
%
%Confitébor tibi in cíthara, Deus, Deus meus: * quare tristis es, ánima mea? et quare contúrbas me?
%
%Spera in Deo, quóniam adhuc confitébor illi: * salutáre vultus mei, et Deus meus.
%
%Glória Patri, et Fílio, * et Spirítui Sancto.
%
%Sicut erat in princípio, et nunc, et semper, * et in sǽcula sæculórum. Amen.

\begin{tightlist}{}
%\item \textred{Aña.} Cum orátis non éritis sicut hypócritæ, qui amant in synagógis, et in ángulis plateárum stantes oráre, ut videántur ab homínibus.
\item \textred{Aña.} Tu autem, cum oráveris, intra in cubículum tuum, et clauso óstio, ora Patrem tuum in abscóndito.
\item \textred{Ps. 63.} Exáudi, Deus, 253.
\end{tightlist}

%\subsubsection*{\center Psalmus lxiij}
%\lettrine{E}xáudi, Deus, oratiónem meam cum déprecor: * a timóre inimíci éripe ánimam meam.
%
%Protexísti me a convéntu malignántium: * a multitúdine operántium iniquitátem.
%
%Quia exacuérunt ut gládium linguas suas: * intendérunt arcum rem amáram, ut sagíttent in occúltis immaculátum.
%
%Súbito sagittábunt eum, et non timébunt: * firmavérunt sibi sermónem nequam.
%
%Narravérunt ut abscónderent láqueos: * dixérunt: Quis vidébit eos?
%
%Scrutáti sunt iniquitátes: * defecérunt scrutántes scrutínio.
%
%Accédet homo ad cor altum: * et exaltábitur Deus.
%
%Sagíttæ parvulórum factæ sunt plagæ eórum: * et infirmátæ sunt contra eos linguæ eórum.
%
%Conturbáti sunt omnes qui vidébant eos: * et tímuit omnis homo.
%
%Et annuntiavérunt ópera Dei, * et facta ejus intellexérunt.
%
%Lætábitur justus in Dómino, et sperábit in eo: * et laudabúntur omnes recti corde.
%
%Glória Patri, et Fílio, * et Spirítui Sancto.
%
%Sicut erat in princípio, et nunc, et semper, * et in sǽcula sæculórum. Amen.

\begin{tightlist}{}
%\item \textred{Aña.} Tu autem, cum oráveris, intra in cubículum tuum, et clauso óstio, ora Patrem tuum in abscóndito.
\item \textred{Aña.} Factum est autem in illis diébus, éxiit Jesus in montem oráre, et erat pernóctans in oratióne Dei.
\item \textred{Ps. 120.} Levávi óculos, 93.
\end{tightlist}

%\subsubsection*{\center Psalmus cxx}
%\lettrine{L}evávi óculos meos in montes, * unde véniet auxílium mihi.
%
%Auxílium meum a Dómino, * qui fecit cælum et terram.
%
%Non det in commotiónem pedem tuum: * neque dormítet qui custódit te.
%
%Ecce, non dormitábit neque dórmiet, * qui custódit Israël.
%
%Dóminus custódit te, Dóminus protéctio tua, * super manum déxteram tuam.
%
%Per diem sol non uret te: * neque luna per noctem.
%
%Dóminus custódit te ab omni malo: * custódiat ánimam tuam Dóminus.
%
%Dóminus custódiat intróitum tuum, et éxitum tuum: * ex hoc nunc, et usque in sǽculum.
%
%Glória Patri, et Fílio, * et Spirítui Sancto.
%
%Sicut erat in princípio, et nunc, et semper, * et in sǽcula sæculórum. Amen.

\begin{tightlist}{}
%\item \textred{Aña.} Factum est autem in illis diébus, éxiit Jesus in montem oráre, et erat pernóctans in oratióne Dei.
\item \textred{℣.} Pater mi, si possíbile est. \textred{℟.} Tránseat a me calix iste.
\end{tightlist}

\subsubsection*{\center Ex Tractátu sancti Cypriáni Epíscopi et Mártyris.}
\begin{center}\textred{De Oratione Dom.}\\
\textred{\textbf{Lectio iv}.}
\end{center}
\lettrineThree{N}ec verbis tantum, sed et factis Dóminus oráre nos dócuit ipse orans frequénter et déprecans, et quid fácere nos oportéret exémpli sui contestatióne demónstrans sicut scriptum est: Ipse autem fuit secédens in solitúdinem, et adórans. Et íterum: Exívit in montem oráre, et fuit pernóctans in oratióne Dei. Quod si ille orábat, qui sine peccáto erat, quanto magis peccatóres opórtet oráre? Et si ille per totam noctem júgiter vígilans contínuis précibus orábat, quanto nos magis in frequentánda oratióne debémus nocte vigiláre? Orábat autem Dóminus, et rogábat non pro se, quid enim pro se ínnocens precarétur? sed pro delíctis nostris, sicut et ipse declárat, cum dicit ad Petrum: Ecce sátanas expetívit, ut vos vexáret quómodo tríticum; ego autem rogávi pro te ne defíciat fides tua. Et póstmodum pro ómnibus Patrem deprecátur dicens: Non pro his autem rogo solis, sed et pro illis, qui creditúri sunt per verbum ipsórum in me, ut omnes unum sint, sicut tu, Pater, in me, et ego in te, ut et ipsi in nobis unum sint.

\begin{tightlist}{}
\item \textred{℟.} Simon, Simon, ecce sátanas expetívit vos, ut cribráret sicut tríticum: \textred{*} Ego autem rogávi pro te, ut non defíciat fides tua. \textred{℣.} Et tu aliquándo convérsus, confírma fratres tuos. \textred{E}go.
\end{tightlist}

\subsubsection*{\center Ex commentáriis sancti Ansélmi Epíscopi.}
\begin{center}\textred{In Epistolam ad Hebrǽos.\\
\textbf{Lectio v}. \hfill Cap. 5.}\end{center}
\lettrine{S}icut verus Póntifex óbtulit preces: nam in Evangélio sæpe légitur orásse, et máxime apud Lucam, qui Sacerdótis in eo descrípsit persónam. Sed et quidquid ipse egit in carne, preces et supplicatiónes fuérunt pro homínibus. Tota vita sua Patrem orávit de resurrectióne carnis suæ, ac de nostra salúte, et instánte jam Passióne óbtulit supplicatiónes, id est humíllimas et instantíssimas oratiónes cum summa devotióne cordis et affectióne, quando factus in agonía prolíxius orábat, et factus est sudor ejus sicut guttæ sánguinis decurréntis in terram. Et has preces, supplicationésque óbtulit, id est óbviam tulit Patri per exauditiónem sibi obviánti. Obtulit eas ad eum, qui posset illum a morte salvum fácere, id est resuscitáre, ad eum, quem scivit esse sufficiéntem ut salváret eum, id est immortálem et impassíbilem fáceret, eripiéndo a morte, ut nec ánima remanéret in inférno, nec caro putrésceret in sepúlcro.

\begin{tightlist}{}
\item \textred{℟.} Orátio humiliántis se nubes penetrábit: \textred{*} Et donec propínquet, non consolábitur. \textred{℣.} Et non discédet donec Altíssimus adspíciat. \textred{E}t.
\end{tightlist}

\subsubsection*{\center\textred{Lectio vj.}}
\lettrine{E}t óbtulit has cum clamóre válido, id est cum intentióne piæ devotiónis vehementíssima et efficacíssima, ut quando prolíxius orábat, et étiam cum lácrymis; quia in illa prolíxa oratióne credéndus est lácrymas effudísse, cum et guttæ sánguinis pro sudóre decúrrerent ab ejus córpore, et exaudítus est, quia quod quæsíerat, accépit in resurrectióne. Exaudítus est, id est, extra númerum aliórum ómnium audítus est, quóniam super omnem creatúram post labóris sui milítiam, a Patre sublimátus est, et hoc pro sui reveréntia, id est secúndum hoc, quod ipse sicut Dei Fílius dignus est reveréntia et veneratióne. Vel pro sua reveréntia in Deum, id est propter quod ipse super ómnia Patrem revéritus est, et honorávit. Sive pro sua reveréntia, id est, sicut sua relígio méruit. Sánguinis quoque ejus effúsio potest intélligi clamor válidus, in quo exaudítus est pro reveréntia ejúsdem passiónis. Reveréntia est, quod sine peccáto passus est sola caritáte.

\begin{tightlist}{}
\item \textred{℟.} Fiant aures tuæ auscultántes, et óculi tui apérti, \textred{*} Ut áudias oratiónem servi tui. \textred{℣.} Quam ego oro coram te hódie nocte ac die. \textred{U}t. Glória Patri. \textred{U}t.
\end{tightlist}

\subsubsection*{\center \textsc{IN TERTIO NOCTURNO.}}
\begin{tightlist}{}
\item \textred{Aña.} Tribulátio et angústia invenérunt me, mandáta tua meditátio mea est.
\item \textred{Ps. 3.} Dómine, quid multiplicáti, 25.
\end{tightlist}

%\subsubsection*{\center Psalmus iij}
%\lettrineThree{D}ómine, quid multiplicáti sunt qui tríbulant me? * multi insúrgunt advérsum me.
%
%Multi dicunt ánimæ meæ : * Non est salus ipsi in Deo ejus.
%
%Tu autem, Dómine, suscéptor meus es, * glória mea, et exáltans caput meum.
%
%Voce mea ad Dóminum clamávi: * et exaudívit me de monte sancto suo.
%
%Ego dormívi, et soporátus sum: * et exsurréxi, quia Dóminus suscépit me.
%
%Non timébo míllia pópuli circumdántis me: * exsúrge, Dómine, salvum me fac, Deus meus.
%
%Quóniam tu percussísti omnes adversántes mihi sine causa: * dentes peccatórum contrivísti.
%
%Dómini est salus: * et super pópulum tuum benedíctio tua.
%
%Glória Patri, et Fílio, * et Spirítui Sancto.
%
%Sicut erat in princípio, et nunc, et semper, * et in sǽcula sæculórum. Amen.

\begin{tightlist}{}
%\item \textred{Aña.} Tribulátio et angústia invenérunt me, mandáta tua meditátio mea est.
\item \textred{Aña.} Sustínui qui simul contristarétur, et non fuit, et qui consolarétur et non invéni.
\item \textred{Ps. 12.} Usquequo, Dómine, 130.
\end{tightlist}

%\subsubsection*{\center Psalmus xij}
%\lettrine{U}squequo, Dómine, obliviscéris me in finem? * úsquequo avértis fáciem tuam a me?
%
%Quámdiu ponam consília in ánima mea, * dolórem in corde meo per diem?
%
%Usquequo exaltábitur inimícus meus super me? * réspice, et exáudi me, Dómine, Deus meus.
%
%Illúmina óculos meos ne umquam obdórmiam in morte : * nequándo dicat inimícus meus : præválui advérsus eum.
%
%Qui tríbulant me, exsultábunt si motus fúero : * ego autem in misericórdia tua sperávi.
%
%Exsultábit cor meum in salutári tuo : cantábo Dómino qui bona tríbuit mihi : * et psallam nómini Dómini altíssimi.
%
%Glória Patri, et Fílio, * et Spirítui Sancto.
%
%Sicut erat in princípio, et nunc, et semper, * et in sǽcula sæculórum. Amen.

\begin{tightlist}{}
%\item \textred{Aña.} Sustínui qui simul contristarétur, et non fuit, et qui consolarétur et non invéni.
\item \textred{Aña.} Adhǽsit paviménto ánima mea: vivífica me secúndum verbum tuum.
\item \textred{Ps. 87.} Dómine, Deus salútis, 277.
\end{tightlist}

%\subsubsection*{\center Psalmus lxxxvij}
%\lettrine{D}ómine, Deus salútis meæ : * in die clamávi, et nocte coram te.
%
%Intret in conspéctu tuo orátio mea : * inclína aurem tuam ad precem meam :
%
%Quia repléta est malis ánima mea : * et vita mea inférno appropinquávit.
%
%Æstimátus sum cum descendéntibus in lacum : * factus sum sicut homo sine adjutório, inter mórtuos liber.
%
%Sicut vulneráti dormiéntes in sepúlcris, quorum non es memor ámplius : * et ipsi de manu tua repúlsi sunt.
%
%Posuérunt me in lacu inferióri : * in tenebrósis, et in umbra mortis.
%
%Super me confirmátus est furor tuus : * et omnes fluctus tuos induxísti super me.
%
%Longe fecísti notos meos a me : * posuérunt me abominatiónem sibi.
%
%Tráditus sum, et non egrediébar : * óculi mei languérunt præ inópia.
%
%Clamávi ad te, Dómine, tota die : * expándi ad te manus meas.
%
%Numquid mórtuis fácies mirabília : * aut médici suscitábunt, et confitebúntur tibi?
%
%Numquid narrábit áliquis in sepúlcro misericórdiam tuam, * et veritátem tuam in perditióne?
%
%Numquid cognoscéntur in ténebris mirabília tua, * et justítia tua in terra obliviónis?
%
%Et ego ad te, Dómine, clamávi : * et mane orátio mea prævéniet te.
%
%Ut quid, Dómine, repéllis oratiónem meam : * avértis fáciem tuam a me?
%
%Pauper sum ego, et in labóribus a juventúte mea : * exaltátus autem, humiliátus sum et conturbátus.
%
%In me transiérunt iræ tuæ : * et terróres tui conturbavérunt me.
%
%Circumdedérunt me sicut aqua tota die : * circumdedérunt me simul.
%
%Elongásti a me amícum et próximum : * et notos meos a miséria.
%
%Glória Patri, et Fílio, * et Spirítui Sancto.
%Sicut erat in princípio, et nunc, et semper, * et in sǽcula sæculórum. Amen.

\begin{tightlist}{}
%\item \textred{Aña.} Adhæsit paviménto ánima mea: vivífica me secúndum verbum tuum.
\item \textred{℣.} Vigiláte et oráte. \textred{℟.} Ut non intrétis in tentatiónem.
\end{tightlist}

\subsubsection*{\center Léctio sancti Evangélii secúndum Lucam.}
\begin{center}\textred{\textbf{Lectio vij}. \hfill Cap. 22.}\end{center}
\lettrine{I}n illo témpore: Egréssus Jesus ibat secúndum consuetúdinem in montem Olivárum. Secúti sunt autem illum et discípuli. Et réliqua.

\begin{center} Homilía sancti Ambrósii Epíscopi.\\
\textred{Lib. 10. Comment. in Luc. Cap. 22.}\end{center}
\lettrineThree{T}ransfer a me cálicem istum, quasi homo mortem recúsans, quasi Deus senténtiam suam servans. Opórtet enim mori nos sǽculo, ut resurgámus Deo; ut juxta divínam senténtiam lex maledictiónis resolúta in terræ limum natúræ fine solvátur. Quod autem ait: Non mea volúntas, sed tua fiat; suam ad hóminem rétulit, Patris ad Divinitátem. Volúntas enim hóminis, temporális: volúntas Divinitátis, ætérna. Non ália volúntas Patris, ália Fílii. Una enim volúntas, ubi una Divínitas. Disce tamen Deo esse subjéctus, ut non quod ipse vis éligas, sed quod Deo scias esse placitúrum.

\begin{tightlist}{}
\item \textred{℟.} Veni in altitúdinem maris, \textred{*} Et tempéstas demérsit me. \textred{℣.} Salvum me fac, Deus, quóniam intravérunt aquæ usque ad ánimam meam. \textred{E}t.
\end{tightlist}

\subsubsection*{\center\textred{Lectio viij.}}
\lettrine{D}eínde verbórum ipsórum proprietátem considerémus. Tristis est, inquit, ánima mea. Et álibi: Nunc ánima mea turbáta est valde. Non ergo suscípiens, sed suscépta turbátur. Anima enim obnóxia passiónibus, Divínitas líbera. Dénique spíritus promptus, caro autem infírma. Tristis autem est non ipse, sed ánima. Non est tristis sapiéntia, non divína substántia, sed ánima. Suscépit enim ánimam meam, suscépit corpus meum. Non me feféllit, ut áliud esset, et áliud viderétur. Tristis videbátur, et tristis erat, non pro sua passióne, sed pro nostra dispersióne.

\begin{tightlist}{}
\item \textred{℟.} Intret in conspéctu tuo orátio mea, \textred{*} Inclína aurem tuam ad precem meam. \textred{℣.} Quia repléta est malis ánima mea, et vita mea inférno appropinquávit. \textred{I}nclína.
\end{tightlist}

\subsubsection*{\center\textred{Lectio ix.}}
\lettrine{D}enique ait: Percútiam pastórem, et dispergéntur oves gregis. Tristis erat, quia nos párvulos relinquébat. Céterum, quam constánter se morti obtúlerit, consequéntia declárant: quandóquidem quæréntibus occúrrit, turbátos confirmávit, trépidos provocávit, proditórem ósculi dignatióne suscépit. Nec illud distat a vero, si tristis erat pro persecutóribus suis, quos sciébat immánis sacrilégii pœnas datúros. Et ídeo dixit: Transfer hunc cálicem a me. Non quia Deus, Dei Fílius, mortem timébat, sed quia nec malos pro se períre volébat.

\begin{tightlist}{}
\item \item \textred{℟.} Prócidit Jesus in fáciem suam orans et dicens: Pater mi si possíbile est, tránseat a me calix iste; \textred{*} Et factus in agonía prolíxius orábat. \textred{℣.} Et factus est sudor ejus sicut guttæ sánguinis decurréntis in terram. \textred{E}t. Glória Patri. \textred{E}t.
\item Te Deum laudámus. 6.
\item \textred{Sacerd. ℣.} Verba mea áuribus pércipe Dómine. \textred{℟.} Et intéllige clamórem meum.
\end{tightlist}

\subsubsection*{\center \textsc{AD LAUDES}}
\begin{center}\textred{et per horas, Aña.}\end{center}
\lettrine{V}enit Jesus cum discípulis suis in villam, quæ dícitur Gethsémani, et dixit illis: Sedéte hic, donec vadam illuc, et orem.

\begin{tightlist}{}
\item \textred{Ps. de Dominica.} 30.
\item \textred{Aña.} Et assúmpto Petro, et duóbus fíliis Zebedǽi, cœpit contristári et mœstus esse.
\item \textred{Aña.} Tunc ait illis: Tristis est ánima mea usque ad mortem: sustinéte hic, et vigiláte mecum.
\item \textred{Aña.} Et progréssus pusíllum, prócidit in fáciem suam, orans et dicens: Pater mi, si possíbile est, tránseat a me calix iste.
\item \textred{Aña.} Pater mi, si non potest hic calix transíre, nisi bibam illum, fiat volúntas tua.
\end{tightlist}
\subsubsection*{\textred{Capitulum. \hfill Matth. 26.}}
\lettrineThree{V}enit Jesus ad discípulos suos, et ínvenit eos dormiéntes, et dicit Petro: Sic non potuístis una hora vigiláre mecum? Vigiláte et oráte, ut non intrétis in tentatiónem.

\subsubsection*{\center\textred{Hymnus.}}
\lettrine{V}enit e cœlo Mediátor alto,\\
Quem sacri dudum cecinére vates:\\
Parce mœróri, lacrymísque amáris,\\
Fília Sion.\\
\\
Attulit mortem vetus hortus, unde\\
Culpa prodívit: novus iste vitam\\
Hortus en affert, ubi nocte Jesus\\
Pérmanet orans.\\
\\
Víndicem placat Genitóris iram:\\
Réprimit forti jaciénda dextra\\
Fúlmina: occúrrit scéleri expiándo\\
Sponte Redémptor.\\
\\
Sic teret duros láqueos avérni,\\
Et diu clausam reserábit aulam,\\
Nos ad ætérni révocans beáta\\
Gáudia regni.\\
\\
Laus, honor Patri, genitǽque Proli\\
Cui datum nomen super omne nomen,\\
Et Paracléto decus atque virtus,\\
Omne per ævum.\\
Amen.

\begin{tightlist}{}
\item \textred{℣.} Factus est sudor ejus. \textred{℟.} Sicut guttæ sánguinis.
\item \textred{Ad Benedictus, Aña.} Factus in agonía prolíxius orábat: Et factus est sudor ejus sicut guttæ sánguinis decurréntis in terram.
\end{tightlist}
\subsubsection*{\center \textred{Oratio.}}
\lettrine{D}omine Jesu Christe, qui in horto, verbo et exémplo nos oráre docuísti ad tentatiónum perícula superánda, concéde propítius, ut nos oratióni semper inténti, ejus copiósum fructum cónsequi mereámur. Qui vivis.

\begin{tightlist}{}
\item \textred{Ad Horas Psalmi de Dominica, ad Primam tamen ut in Festis.}
\end{tightlist}

\subsubsection*{\center \textsc{AD PRIMAM.}}
\begin{tightlist}{}
\item \textred{℣.} Qui nos salvásti in Cruce.
\end{tightlist}

\subsubsection*{\center \textsc{AD TERTIAM.}}
\subsubsection*{\textred{Capitulum. \hfill Matth. 26.}}
\lettrine{V}enit Jesus ad discípulos suos, et ínvenit eos dormiéntes, et dicit Petro: Sic non potuístis una hora vigiláre mecum? Vigiláte et oráte, ut non intrétis in tentatiónem.

\begin{tightlist}{}
%\item \textred{Capit.} Venit Jesus, \textred{ut supra.}
\item \textred{℟.} Tristis est \textred{*} Anima mea. \textred{T}ristis. \textred{℣.} Usque ad mortem. \textred{A}nima. Glória Patri. \textred{T}ristis.
\item \textred{℣.} Pater mi, si possíbile est. \textred{℟.} Tránseat a me calix iste.
\end{tightlist}
\subsubsection*{\center \textsc{AD SEXTAM.}}
\subsubsection*{\textred{Capitulum. \hfill Luc. 21.}}
\lettrine{V}igiláte ítaque, omni témpore orántes, ut digni habeámini fúgere ista ómnia quæ futúra sunt, et stare ante Fílium hóminis.

\begin{tightlist}{}
\item \textred{℟.} Pater mi, \textred{*} Si possíbile est. \textred{P}ater. \textred{℣.} Tránseat a me calix iste. \textred{S}i. Glória Patri. \textred{P}ater.
\item \textred{℣.} Vigiláte et oráte. \textred{℟.} Ut non intrétis in tentatiónem.
\end{tightlist}

\subsubsection*{\center \textsc{AD NONAM.}}
\subsubsection*{\textred{Capitulum. \hfill Matth. 5.}}
\lettrine{O}ráte pro persequéntibus et calumniántibus vos, ut sitis fílii Patris vestri, qui in cœlis est.

\begin{tightlist}{}
\item \textred{℟.} Vigiláte, \textred{*} Et oráte. \textred{V}igiláte. \textred{℣.} Ut non intrétis in tentatiónem. \textred{E}t. Glória Patri. \textred{V}igiláte.
\item \textred{℣.} Factus est sudor ejus. \textred{℟.} Sicut guttæ sánguinis.
\end{tightlist}

\subsubsection*{\center \textsc{AD VESPERAS.}}
\begin{tightlist}{}
\item \textred{Añæ Laudum, Psalmi de Dominica,} 56. \textred{et loco ultimi, Ps. 115.} Crédidi. 93.
\item \textred{Capit. Hymnus, et ℣. ut in primis Vesperis.}
\item \textred{Ad Magnif. Aña.} Ecce appropinquávit hora, et Fílius hóminis tradétur in manus peccatórum.
\item \textred{Completorium de Dominica.}
\end{tightlist}

\begin{center}\textcolor{red}{\rule{\linewidth}{0.5pt}} \end{center}
\subsection*{\center \textred{\textsc{FERIA TERTIA}\\
post Domin. Sexages.}}
\subsection*{\center \textsc{Commem. Passionis D.N.J.C.}}
\subsubsection*{\center \textred{Duplex minus. P. C.}}
%\begin{center}\textcolor{red}{\rule{\linewidth}{0.5pt}} \end{center}

\subsubsection*{\center \textsc{AD VESPERAS.}}
\fancyhead[C]{\textsc{Comm. Passionis D.N.J.C.}}
\begin{tightlist}{}
\item \textred{Aña.} Cálicem salutáris accípiam et nomen Dómini invocábo.
\item \textred{Psalmi de Offic. præc.}
\end{tightlist}
\subsubsection*{\textred{Capitulum. \hfill Philip. 2.}}
\lettrineThree{F}ratres, Hoc enim sentíte in vobis, quod et in Christo Jesu, qui, cum in forma Dei esset, non rapínam arbitrátus est, esse se æquálem Deo: sed semetípsum exinanívit, formam servi accípiens, in similitúdinem hóminum factus, et hábitu invéntus ut homo. Humiliávit semetípsum factus obédiens usque ad mortem, mortem autem crucis.

\begin{tightlist}{}
\item \textred{℟.} Crucem tuam adorámus, Dómine, sanctámque recólimus passiónem: \textred{*} Ecce enim propter lignum venit gáudium in univérso mundo. 
\item \textred{℣.} Deus misereátur nostri, et benedícat nobis, illúminet vultum suum super nos, et misereátur nostri. \textred{E}cce. Glória Patri. \textred{E}cce.
\end{tightlist}
\subsubsection*{\center\textred{Hymnus.}}
\lettrine{M}œréntes óculi spárgite lácrymas,\\
Et luctu résonent íntima córdium,\\
Illátas réfero Númini ab ímpiis\\
Pœnas, et fera vúlnera.\\
\\
Accíncta heu! gládiis turba satéllitum\\
Arréptum Dóminum fústibus ímpetit;\\
Nunc cædit cólaphis, nunc quatit hórridis\\
Divínum caput íctibus.\\
\\
Haud finis scéleri: tráditur ímprobo\\
Christus carnífici; nec mora, bárbarus\\
In Regem súperum non timet ímpio\\
Ausu, vértere déxteram.\\
\\
Audíte, o pópuli: Numen amábile,\\
Manánte ex húmeris úndique sánguine,\\
Lictóris rábidi sústinet ímpetum,\\
Et vocem premit ínnocens.\\
\\
Quis non illácrymet? Jam nova cóndidit\\
Torménta indómitæ gentis iníquitas;\\
Infígit cérebro proh dolor! éfferis\\
Sertum vépribus ásperum.\\
\\
Heu crimen! tráhitur fúnibus ímprobis\\
Funésti Dóminus supplícii ad locum;\\
Illic oppétiit fúnere, spíritum\\
Patri restítuens suo.\\
\\
Passo pro míseris tam fera vúlnera,\\
In terris résonet débita glória;\\
Sacrúmque assídue Nomen in æthera\\
Humánum genus éfferat.\\
Amen.

\begin{tightlist}{}
\item \textred{℣.} Oblátus est quia ipse vóluit. \textred{℟.} Et livóre ejus sanáti sumus.
\item \textred{Ad Magnificat, Aña.} O vos omnes, qui transítis per viam, atténdite, et vidéte, si est dolor sicut dolor meus.
\end{tightlist}

\subsubsection*{\center\textred{Oratio.}}
\lettrine{O}mnípotens sempitérne Deus, qui humáno géneri, ad imitándum humilitátis exémplum, Salvatórem nostrum carnem súmere, et crucem subíre fecísti: concéde propítius; ut, sicut solémnem commemoratiónem Passiónis ejus celebrámus, ita et patiéntiæ ipsíus habére documénta, et resurrectiónis consórtia mereámur. Per eúmdem Dóminum.

\begin{tightlist}{}
\item \textred{Ad Complet. et per Horas in fine Hymn. dicitur ℣.} Glória tibi Dómine, Qui nos salvásti in Cruce.
\item \textred{Ad Nunc dimittis, Aña.} Hoc signum Crucis * erit in cælo, cum Dóminus ad judicándum vénerit.
\end{tightlist}

\subsubsection*{\center \textsc{AD MATUTINUM.}}
\begin{tightlist}{}
\item \textred{Invit.} Christum Regem crucifíxum * Veníte, adorémus.
\item \textred{Ps.} Veníte, exsultémus. 2. 
\end{tightlist}
\subsubsection*{\center\textred{Hymnus.}}
\lettrine{A}dspice, infámi Deus ipse ligno\\
Pendet, effúso mádidus cruóre,\\
Adspice, immíti manus alma clavo\\
Fínditur alte.\\
\\
Hunc, velut pravi scéleris minístrum,\\
Inter indígnos médium latrónes\\
Cerne: crudélis fuit ista gentis\\
Dira volúntas.\\
\\
Pallet heu! vultus; caput ecce lassum\\
Fléctitur: clausis óculis Redémptor\\
Spíritum sacro méritis onústum\\
Fundit ab ore.\\
\\
O cor, æs durum súperas, geméndo\\
Ni scelus tergas; tua culpa Christum\\
Stípiti affíxit: tua culpa morti\\
Súbdidit atræ.\\
\\
Sit Deo ætérnum decus omne in ævum\\
Qui pius nostri géneris Redémptor,\\
Críminis labem míseris nocéntem\\
Sánguine tersit.\\
Amen.

\subsubsection*{\center \textsc{IN PRIMO NOCTURNO.}}
\begin{tightlist}{}
\item \textred{Aña.} Adstitérunt reges terræ, et príncipes convenérunt in unum advérsus Dóminum, et advérsus Christum ejus.
\item \textred{Ps. 2.} Quare fremuérunt, 24.
\end{tightlist}

%\subsubsection*{\center Psalmus ij}
%\lettrineThree{Q}uare fremuérunt Gentes: * et pópuli meditáti sunt inánia?
%
%Astitérunt reges terræ, et príncipes convenérunt in unum * advérsus Dóminum, et advérsus Christum ejus.
%
%Dirumpámus víncula eórum: * et projiciámus a nobis jugum ipsórum.
%
%Qui hábitat in cælis, irridébit eos: * et Dóminus subsannábit eos.
%
%Tunc loquétur ad eos in ira sua, * et in furóre suo conturbábit eos.
%
%Ego autem constitútus sum Rex ab eo super Sion montem sanctum ejus, * prǽdicans præcéptum ejus.
%
%Dóminus dixit ad me: * Fílius meus es tu, ego hódie génui te.
%
%Póstula a me, et dabo tibi Gentes hereditátem tuam, * et possessiónem tuam términos terræ.
%
%Reges eos in virga férrea, * et tamquam vas fíguli confrínges eos.
%
%Et nunc, reges, intellígite: * erudímini, qui judicátis terram.
%
%Servíte Dómino in timóre: * et exsultáte ei cum tremóre.
%
%Apprehéndite disciplínam, nequándo irascátur Dóminus, * et pereátis de via justa.
%
%Cum exárserit in brevi ira ejus: * beáti omnes qui confídunt in eo.
%
%Glória Patri, et Fílio, * et Spirítui Sancto.
%
%Sicut erat in princípio, et nunc, et semper, * et in sǽcula sæculórum. Amen.
\begin{tightlist}{}
%\item \textred{Aña.} Adstitérunt reges terræ, et príncipes convenérunt in unum advérsus Dóminum, et advérsus Christum ejus.
\item \textred{Aña.} Multiplicáti sunt, qui tríbulant me, multi insúrgunt advérsum me.
\item \textred{Ps. 3.} Dómine, quid multiplicáti, 25.
\end{tightlist}

%\subsubsection*{\center Psalmus iij}
%\lettrine{D}ómine quid multiplicáti sunt qui tríbulant me? * multi insúrgunt advérsum me.
%
%Multi dicunt ánimæ meæ : * Non est salus ipsi in Deo ejus.
%
%Tu autem, Dómine, suscéptor meus es, * glória mea, et exáltans caput meum.
%
%Voce mea ad Dóminum clamávi: * et exaudívit me de monte sancto suo.
%
%Ego dormívi, et soporátus sum: * et exsurréxi, quia Dóminus suscépit me.
%
%Non timébo míllia pópuli circumdántis me: * exsúrge, Dómine, salvum me fac, Deus meus.
%
%Quóniam tu percussísti omnes adversántes mihi sine causa: * dentes peccatórum contrivísti.
%
%Dómini est salus: * et super pópulum tuum benedíctio tua.
%
%Glória Patri, et Fílio, * et Spirítui Sancto.
%
%Sicut erat in princípio, et nunc, et semper, * et in sǽcula sæculórum. Amen.

\begin{tightlist}{}
%\item \textred{Aña.} Multiplicáti sunt, qui tríbulant me, multi insúrgunt advérsum me.
\item \textred{Aña.} Divisérunt sibi vestiménta mea, et super vestem meam misérunt sortem.
\item \textred{Ps. 21.} Deus, Deus meus, 222.
\end{tightlist}

%\subsubsection*{\center Psalmus xxj}
%\lettrine{D}eus, Deus meus, réspice in me : quare me dereliquísti? * longe a salúte mea verba delictórum meórum.
%
%Deus meus, clamábo per diem, et non exáudies: * et nocte, et non ad insipiéntiam mihi.
%
%Tu autem in sancto hábitas: * laus Israël.
%
%In te speravérunt patres nostri: * speravérunt, et liberásti eos.
%
%Ad te clamavérunt, et salvi facti sunt: * in te speravérunt, et non sunt confúsi.
%
%Ego autem sum vermis, et non homo: * oppróbrium hóminum, et abjéctio plebis.
%
%Omnes vidéntes me derisérunt me: * locúti sunt lábiis, et movérunt caput.
%
%Sperávit in Dómino, erípiat eum: * salvum fáciat eum, quóniam vult eum.
%
%Quóniam tu es, qui extraxísti me de ventre: * spes mea ab ubéribus matris meæ. In te projéctus sum ex útero.
%
%De ventre matris meæ Deus meus es tu, * ne discésseris a me :
%
%Quóniam tribulátio próxima est: * quóniam non est qui ádjuvet.
%
%Circumdedérunt me vítuli multi: * tauri pingues obsedérunt me.
%
%Aperuérunt super me os suum: * sicut leo rápiens et rúgiens.
%
%Sicut aqua effúsus sum: * et dispérsa sunt ómnia ossa mea.
%
%Factum est cor meum tamquam cera liquéscens: * in médio ventris mei.
%
%Aruit tamquam testa virtus mea, et lingua mea adhǽsit fáucibus meis: * et in púlverem mortis deduxísti me.
%
%Quóniam circumdedérunt me canes multi: * concílium malignántium obsédit me.
%
%Fodérunt manus meas et pedes meos: * dinumeravérunt ómnia ossa mea.
%
%Ipsi vero consideravérunt et inspexérunt me: * divisérunt sibi vestiménta mea, et super vestem meam misérunt sortem.
%
%Tu autem, Dómine, ne elongáveris auxílium tuum a me: * ad defensiónem meam cónspice.
%
%Erue a frámea, Deus, ánimam meam: * et de manu canis únicam meam.
%
%Salva me ex ore leónis: * et a córnibus unicórnium humilitátem meam.
%
%Narrábo nomen tuum frátribus meis: * in médio Ecclésiæ laudábo te.
%
%Qui timétis Dóminum, laudáte eum: * univérsum semen Jacob, glorificáte eum.
%
%Tímeat eum omne semen Israël: * quóniam non sprevit, neque despéxit deprecatiónem páuperis :
%
%Nec avértit fáciem suam a me: * et cum clamárem ad eum, exaudívit me.
%
%Apud te laus mea in ecclésia magna: * vota mea reddam in conspéctu timéntium eum.
%
%Edent páuperes, et saturabúntur: et laudábunt Dóminum qui requírunt eum: * vivent corda eórum in sǽculum sǽculi.
%
%Reminiscéntur et converténtur ad Dóminum * univérsi fines terræ :
%
%Et adorábunt in conspéctu ejus * univérsæ famíliæ Géntium.
%
%Quóniam Dómini est regnum: * et ipse dominábitur Géntium.
%
%Manducavérunt et adoravérunt omnes pingues terræ : * in conspéctu ejus cadent omnes qui descéndunt in terram.
%
%Et ánima mea illi vivet: * et semen meum sérviet ipsi.
%
%Annuntiábitur Dómino generátio ventúra: * et annuntiábunt cæli justítiam ejus pópulo qui nascétur, quem fecit Dóminus.
%
%Glória Patri, et Fílio, * et Spirítui Sancto.
%
%Sicut erat in princípio, et nunc, et semper, * et in sǽcula sæculórum. Amen.

\begin{tightlist}{}
%\item \textred{Aña.} Divisérunt sibi vestiménta mea, et super vestem meam misérunt sortem.
\item \textred{℣.} Oblátus est, quia ipse vóluit. \textred{℟.} Et livóre ejus sanáti sumus.
\end{tightlist}

\subsubsection*{\center De Epístola beáti Pauli Apóstoli ad Romános}
\begin{center}\textred{\textbf{Lectio j}. \hfill Cap. 5. et 6.}\end{center}
\lettrineThree{C}omméndat caritátem suam Deus in nobis: quóniam cum adhuc peccatóres essémus, secúndum tempus, Christus pro nobis mórtuus est: multo ígitur magis nunc justificáti in sánguine ipsíus, salvi érimus ab ira per ipsum. Si enim, cum inimíci essémus, reconciliáti sumus Deo per mortem Fílii ejus, multo magis reconciliáti, salvi érimus in vita ipsíus. Non solum autem, sed et gloriámur in Deo per Dóminum nostrum Jesum Christum, per quem nunc reconciliatiónem accépimus. Si enim, uníus delícto mors regnávit per unum: multo magis abundántiam grátiæ, et donatiónis, et justítiæ accipiéntes, in vita regnábunt per unum Jesum Christum.

\begin{tightlist}{}
\item \textred{℟.} Vínea mea elécta, ego te plantávi, dicit Dóminus, et tu facta es mihi nimis amára; \textred{*} Quia parásti crucem Salvatóri tuo. \textred{℣.} Pópule meus, quid feci tibi, aut in quo contristávi te? respónde mihi. \textred{Q}uia.
\end{tightlist}

\subsubsection*{\center\textred{Lectio ij.}}
\lettrine{I}gitur, sicut per uníus delíctum, in omnes hómines in condemnatiónem, sic et per uníus justítiam, in omnes hómines in justificatiónem vitæ. Sicut enim, per inobediéntiam uníus hóminis, peccatóres constitúti sunt multi: ita et per uníus obeditiónem justi constituéntur multi. Lex autem subintrávit, ut abundáret delíctum: ubi autem abundávit delíctum, superabundávit grátia. Ut sicut regnávit peccátum in mortem, ita et grátia regnet per justítiam in vitam ætérnam, per Jesum Christum Dóminum nostrum.

\begin{tightlist}{}
\item \textred{℟.} Ego edúxi te de Ægýpto, demérso Pharaóne in mare rubrum, et ante te præívi in colúmna nubis. \textred{*} Et tu me tradidísti princípibus sacerdótum, et me duxísti ad prætórium Piláti. \textred{℣.} Pópule meus, quid feci tibi, aut in quo contristávi te? respónde mihi. \textred{E}t.
\end{tightlist}

\subsubsection*{\center\textred{Lectio iij.}}
\lettrine{Q}uid ergo dicémus? Permanébimus in peccáto, ut grátia abúndet? Absit: qui enim mórtui sumus peccáto, quómodo adhuc vivémus in illo? An ignorátis, quia quicúmque baptizáti sumus in Christo Jesu, in morte ipsíus baptizáti sumus? Consepúlti enim sumus cum illo, per Baptísmum, in mortem: ut quómodo Christus surréxit a mórtuis per glóriam Patris; ita et nos in novitáte vitæ ambulémus. Si enim complantáti facti sumus similitúdini mortis ejus; simul et resurrectiónis érimus. Hoc sciéntes, quia vetus homo noster simul crucifíxus est, ut destruátur corpus peccáti, et ultra non serviámus peccáto. Qui enim mórtuus est, justificátus est a peccáto. Si autem mórtui sumus cum Christo, crédimus quia simul étiam vivémus cum Christo.

\begin{tightlist}{}
\item \textred{℟.} Ego propter te flagellávi Ægýptum cum primogénitis suis; \textred{*} Et tu me flagellátum tradidísti, qui tamquam agnus, coram tondénte se, obmútui, et non apérui os meum. \textred{℣.} Pópule meus, quid feci tibi, aut in quo contristávi te? respónde mihi. \textred{E}t. Glória Patri. \textred{E}t.
\end{tightlist}

\subsubsection*{\center \textsc{IN SECUNDO NOCTURNO.}}
\begin{tightlist}{}
\item \textred{Aña.} Inimíci mei dixérunt mala mihi, quando moriétur, et períbit nomen ejus?
\item \textred{Ps. 40.} Beátus qui, 122.
\end{tightlist}

%\subsubsection*{\center Psalmus xl}
%\lettrineThree{B}eátus qui intélligit super egénum, et páuperem: * in die mala liberábit eum Dóminus.
%
%Dóminus consérvet eum, et vivíficet eum, et beátum fáciat eum in terra: * et non tradat eum in ánimam inimicórum ejus.
%
%Dóminus opem ferat illi super lectum dolóris ejus: * univérsum stratum ejus versásti in infirmitáte ejus.
%
%Ego dixi : Dómine, miserére mei: * sana ánimam meam, quia peccávi tibi.
%
%Inimíci mei dixérunt mala mihi: * Quando moriétur, et períbit nomen ejus?
%
%Et si ingrediebátur ut vidéret, vana loquebátur: * cor ejus congregávit iniquitátem sibi.
%
%Egrediebátur foras, * et loquebátur in idípsum.
%
%Advérsum me susurrábant omnes inimíci mei: * advérsum me cogitábant mala mihi.
%
%Verbum iníquum constituérunt advérsum me: * Numquid qui dormit non adjíciet ut resúrgat?
%
%Etenim homo pacis meæ, in quo sperávi: * qui edébat panes meos, magnificávit super me supplantatiónem.
%
%Tu autem, Dómine, miserére mei, et resúscita me: * et retríbuam eis.
%
%In hoc cognóvi quóniam voluísti me: * quóniam non gaudébit inimícus meus super me.
%
%Me autem propter innocéntiam suscepísti: * et confirmásti me in conspéctu tuo in ætérnum.
%
%Benedíctus Dóminus, Deus Israël, a sǽculo et usque in sǽculum: * fiat, fiat.
%
%Glória Patri, et Fílio, * et Spirítui Sancto.
%
%Sicut erat in princípio, et nunc, et semper, * et in sǽcula sæculórum. Amen.

\begin{tightlist}{}
%\item \textred{Aña.} Inimíci mei dixérunt mala mihi, quando moriétur, et períbit nomen ejus?
\item \textred{Aña.} Aliéni insurrexérunt in me, et fortes quæsiérunt ánimam meam.
\item \textred{Ps. 53.} Deus, in nómine, 41.
\end{tightlist}

%\subsubsection*{\center Psalmus liij}
%\lettrine{D}eus, in nómine tuo salvum me fac: * et in virtúte tua júdica me.
%
%Deus, exáudi oratiónem meam: * áuribus pércipe verba oris mei.
%
%Quóniam aliéni insurrexérunt advérsum me, et fortes quæsiérunt ánimam meam: * et non proposuérunt Deum ante conspéctum suum.
%
%Ecce enim Deus ádjuvat me: * et Dóminus suscéptor est ánimæ meæ.
%
%Avérte mala inimícis meis: * et in veritáte tua dispérde illos.
%
%Voluntárie sacrificábo tibi, * et confitébor nómini tuo, Dómine: quóniam bonum est :
%
%Quóniam ex omni tribulatióne eripuísti me: * et super inimícos meos despéxit óculus meus.
%
%Glória Patri, et Fílio, * et Spirítui Sancto.
%
%Sicut erat in princípio, et nunc, et semper, * et in sǽcula sæculórum. Amen.

\begin{tightlist}{}
%\item \textred{Aña.} Aliéni insurrexérunt in me, et fortes quæsiérunt ánimam meam.
\item \textred{Aña.} Conculcavérunt me inimíci mei tota die, advérsum me omnes cogitatiónes eórum in malum.
\item \textred{Ps. 55.} Miserére mei, 158.
\end{tightlist}

%\subsubsection*{\center Psalmus lv}
%\lettrine{M}iserére mei, Deus, quóniam conculcávit me homo: * tota die impúgnans tribulávit me.
%
%Conculcavérunt me inimíci mei tota die: * quóniam multi bellántes advérsum me.
%
%Ab altitúdine diéi timébo: * ego vero in te sperábo.
%
%In Deo laudábo sermónes meos, in Deo sperávi: * non timébo quid fáciat mihi caro.
%
%Tota die verba mea exsecrabántur: * advérsum me omnes cogitatiónes eórum in malum.
%
%Inhabitábunt, et abscóndent: * ipsi calcáneum meum observábunt.
%
%Sicut sustinuérunt ánimam meam, pro níhilo salvos fácies illos: * in ira pópulos confrínges.
%
%Deus, vitam meam annuntiávi tibi: * posuísti lácrimas meas in conspéctu tuo.
%
%Sicut et in promissióne tua: * tunc converténtur inimíci mei retrórsum :
%
%In quacúmque die invocávero te: * ecce cognóvi quóniam Deus meus es.
%
%In Deo laudábo verbum, in Dómino laudábo sermónem: * in Deo sperávi, non timébo quid fáciat mihi homo.
%
%In me sunt, Deus, vota tua, * quæ reddam, laudatiónes tibi.
%
%Quóniam eripuísti ánimam meam de morte, et pedes meos de lapsu: * ut pláceam coram Deo in lúmine vivéntium.
%
%Glória Patri, et Fílio, * et Spirítui Sancto.
%
%Sicut erat in princípio, et nunc, et semper, * et in sǽcula sæculórum. Amen.

\begin{tightlist}{}
%\item \textred{Aña.} Conculcavérunt me inimíci mei tota die, advérsum me omnes cogitatiónes eórum in malum.
\item \textred{℣.} Próprio Fílio suo non pepércit Deus. \textred{℟.} Sed pro ómnibus trádidit illum.
\end{tightlist}

\subsubsection*{\center Sermo sancti Augustíni Epíscopi.}
\begin{center}\textred{Sermo 41. de Passione Domini.\\ \textbf{Lectio iv}.}\end{center}
\lettrineThree{A}dmonet nos, fratres caríssimi, ad solemnitátem Domínicæ Passiónis dies ipse, in quo eam nec muta eleménta tacuérunt. Célebret eam lux fídei, linguis hóminum, quam conclamavérunt étiam siléntia tenebrárum. Hódie Dóminus noster in statéra crucis prétium nostræ salútis appéndit, et una morte univérsum mundum, sicut ómnium Cónditor, ita ómnium Reparátor absólvit. Indubitánter enim credámus, quod totum mundum redémit, qui plus dedit, quam totus mundus valéret. Méritum enim redémptæ mercédis dígnitas insígnis prétii supergréssa est. Inter redémptum et rediméntem dispensátio fuit, compensátio non fuit. Qui ergo non habébat peccáta própria, digne delévit aliéna. Solus hic pia víctima pro ómnibus cécidit, ut omnes leváret. Et quia débitum solus non hábuit, recte fœnus misericórdiæ pro debitóribus erogávit.

\begin{tightlist}{}
\item \textred{℟.} Ego propter te Chananæórum reges percússi, et dedi tibi sceptrum regále; \textred{*} Et tu dedísti cápiti meo spíneam corónam, et percussísti arúndine caput meum. \textred{℣.} Pópule meus, quid feci tibi, aut in quo contristávi te? respónde mihi. \textred{E}t.
\end{tightlist}

\subsubsection*{\center\textred{Lectio v.}}
\lettrine{P}erpénde inter hæc, qui talem pro nobis dedit pecúniam, qualem a nobis sit exactúrus usúram. In hac ítaque die fides prophéticæ annuntiatiónis impléta est, ita dicéntis: Corpus meum dedi percutiéntibus: et genas meas velléntibus; fáciem meam non avérti a fœditáte sputórum. Suscépit mala nostra, ut tribúeret bona sua. Hinc intelligámus, quantum hóminem dilígere dignátus sit ante culpam, quem sic díligit post ruínam. Agnósce, homo, quantum váleas et quantum débeas: et dum tantam redemptiónis tuæ pérspicis dignitátem, ipse tibi indícito peccándi pudórem. Ecce, pro ímpio píetas flagellátur, pro stulto sapiéntia illúditur, pro mendáce véritas necátur, damnátur justítia pro iníquo, misericórdia affícitur pro crudéli, pro mísero replétur sincéritas acéto, inebriátur felle dulcédo, addícitur innocéntia pro reo, móritur vita pro mórtuo.

\begin{tightlist}{}
\item \textred{℟.} Ego duxi te per desértum quadragínta annis, et manna cibávi te; \textred{*} Et tu me cecidísti álapis, et flagéllis. \textred{℣.} Pópule meus, quid feci tibi, aut in quo contristávi te? respónde mihi. \textred{E}t.
\end{tightlist}

\subsubsection*{\center\textred{Lectio vj.}}
\lettrine{E}xpávit scelus hóminum natúra rerum, et quem creatúra rebéllis non agnóscit, eum mundi Dóminum tremens terra testátur, et cœli Regem sol fúgiens confitétur. Clámyde coccínea indúitur, quia sánguine Mártyrum suórum Ecclésiæ corpus ornátur. Coróna spínea cápiti ejus impónitur, quia púnctio peccatórum nostrórum, quorum remissióne Redemptóris glória strúitur, áridis tríbulis comparátur. Studeámus nunc e divérso, ut membrórum vita cápitis sit coróna. Quod vero sitíre se in cruce pósitus dicit, fidem incrédulæ gentis concupíscit. Sed e contra acétum malítiæ pórrigunt, quia vinum sapiéntiæ, quod a Deo accéperant, peccándo corrúperant. Velum templi scínditur, quia Synagóga honóre nudátur, observátio antíqua dissólvitur, Ecclésiæ únitas præmonstrátur. Monuménta aperiúntur, quia mortis jura jure superántur.

\begin{tightlist}{}
\item \textred{℟.} Ego exaltávi te magna virtúte: \textred{*} Et tu me suspendísti in patíbulo crucis, in qua expándi manus meas ad pópulum non credéntem et contradicéntem. \textred{℣.} Pópule meus, quid feci tibi, aut in quo contristávi te? respónde mihi. \textred{E}t. Glória Patri. \textred{E}t.
\end{tightlist}
\subsubsection*{\center \textsc{IN TERTIO NOCTURNO.}}
\begin{tightlist}{}
\item \textred{Aña.} Fílii hóminum dentes eórum arma, et sagíttæ, et lingua eórum gládius acútus. 
\item \textred{Ps. 56.} Miserére mei, 159.
\end{tightlist}

%\subsubsection*{\center Psalmus lvi}
%\lettrineThree{M}iserére mei, Deus, miserére mei: * quóniam in te confídit ánima mea.
%
%Et in umbra alárum tuárum sperábo: * donec tránseat iníquitas.
%
%Clamábo ad Deum altíssimum: * Deum, qui benefécit mihi.
%
%Misit de cælo, et liberávit me: * dedit in oppróbrium conculcántes me.
%
%Misit Deus misericórdiam suam, et veritátem suam, et erípuit ánimam meam de médio catulórum leónum: * dormívi conturbátus.
%
%Fílii hóminum dentes eórum arma et sagíttæ : * et lingua eórum gládius acútus.
%
%Exaltáre super cælos, Deus: * et in omnem terram glória tua.
%
%Láqueum paravérunt pédibus meis: * et incurvavérunt ánimam meam.
%
%Fodérunt ante fáciem meam fóveam: * et incidérunt in eam.
%
%Parátum cor meum, Deus, parátum cor meum: * cantábo, et psalmum dicam.
%
%Exsúrge, glória mea, exsúrge psaltérium et cíthara: * exsúrgam dilúculo.
%
%Confitébor tibi in pópulis, Dómine: * et psalmum dicam tibi in Géntibus :
%
%Quóniam magnificáta est usque ad cælos misericórdia tua, * et usque ad nubes véritas tua.
%
%Exaltáre super cælos, Deus: * et super omnem terram glória tua.
%
%Glória Patri, et Fílio, * et Spirítui Sancto.
%
%Sicut erat in princípio, et nunc, et semper, * et in sǽcula sæculórum. Amen.

\begin{tightlist}{}
%\item \textred{Aña.} Fílii hóminum dentes eórum arma, et sagíttæ, et lingua eórum gládius acútus. 
\item \textred{Aña.} Intendérunt arcum rem amáram, ut sagíttent in occúltis immaculátum.
\item \textred{Ps. 63.} Exáudi, Deus, 253.
\end{tightlist}

%\subsubsection*{\center Psalmus lxiij}
%\lettrine{E}xáudi, Deus, oratiónem meam cum déprecor: * a timóre inimíci éripe ánimam meam.
%
%Protexísti me a convéntu malignántium: * a multitúdine operántium iniquitátem.
%
%Quia exacuérunt ut gládium linguas suas: * intendérunt arcum rem amáram, ut sagíttent in occúltis immaculátum.
%
%Súbito sagittábunt eum, et non timébunt: * firmavérunt sibi sermónem nequam.
%
%Narravérunt ut abscónderent láqueos: * dixérunt: Quis vidébit eos?
%
%Scrutáti sunt iniquitátes: * defecérunt scrutántes scrutínio.
%
%Accédet homo ad cor altum: * et exaltábitur Deus.
%
%Sagíttæ parvulórum factæ sunt plagæ eórum: * et infirmátæ sunt contra eos linguæ eórum.
%
%Conturbáti sunt omnes qui vidébant eos: * et tímuit omnis homo.
%
%Et annuntiavérunt ópera Dei, * et facta ejus intellexérunt.
%
%Lætábitur justus in Dómino, et sperábit in eo: * et laudabúntur omnes recti corde.
%
%Glória Patri, et Fílio, * et Spirítui Sancto.
%
%Sicut erat in princípio, et nunc, et semper, * et in sǽcula sæculórum. Amen.

\begin{tightlist}{}
%\item \textred{Aña.} Intendérunt arcum rem amáram, ut sagíttent in occúltis immaculátum.
\item \textred{Aña.} Factus sum sicut homo sine adjutório inter mórtuos liber.
\item \textred{Ps. 87.} Dómine, Deus, 277.
\end{tightlist}{}

%\subsubsection*{\center Psalmus lxxxvij}
%\lettrine{D}ómine, Deus salútis meæ : * in die clamávi, et nocte coram te.
%
%Intret in conspéctu tuo orátio mea : * inclína aurem tuam ad precem meam :
%
%Quia repléta est malis ánima mea : * et vita mea inférno appropinquávit.
%
%Æstimátus sum cum descendéntibus in lacum : * factus sum sicut homo sine adjutório, inter mórtuos liber.
%
%Sicut vulneráti dormiéntes in sepúlcris, quorum non es memor ámplius : * et ipsi de manu tua repúlsi sunt.
%
%Posuérunt me in lacu inferióri : * in tenebrósis, et in umbra mortis.
%
%Super me confirmátus est furor tuus : * et omnes fluctus tuos induxísti super me.
%
%Longe fecísti notos meos a me : * posuérunt me abominatiónem sibi.
%
%Tráditus sum, et non egrediébar : * óculi mei languérunt præ inópia.
%
%Clamávi ad te, Dómine, tota die : * expándi ad te manus meas.
%
%Numquid mórtuis fácies mirabília : * aut médici suscitábunt, et confitebúntur tibi?
%
%Numquid narrábit áliquis in sepúlcro misericórdiam tuam, * et veritátem tuam in perditióne?
%
%Numquid cognoscéntur in ténebris mirabília tua, * et justítia tua in terra obliviónis?
%
%Et ego ad te, Dómine, clamávi : * et mane orátio mea prævéniet te.
%
%Ut quid, Dómine, repéllis oratiónem meam : * avértis fáciem tuam a me?
%
%Pauper sum ego, et in labóribus a juventúte mea : * exaltátus autem, humiliátus sum et conturbátus.
%
%In me transiérunt iræ tuæ : * et terróres tui conturbavérunt me.
%
%Circumdedérunt me sicut aqua tota die : * circumdedérunt me simul.
%
%Elongásti a me amícum et próximum : * et notos meos a miséria.
%
%Glória Patri, et Fílio, * et Spirítui Sancto.
%
%Sicut erat in princípio, et nunc, et semper, * et in sǽcula sæculórum. Amen.

\begin{tightlist}{}
%\item \textred{Aña.} Factus sum sicut homo sine adjutório inter mórtuos liber.
\item \textred{℣.} Factus est obédiens usque ad mortem. \textred{℟.} Mortem autem crucis.
\end{tightlist}

\subsubsection*{\center Léctio sancti Evangélii secúndum Joánnem.}
\begin{center}\textred{\textbf{Lectio vij}. \hfill Cap. 19.}\end{center}
\lettrine{I}n illo témpore: Sciens Jesus, quia ómnia consummáta sunt, ut consummarétur scriptúra, dixit: Sítio. Et réliqua.

\begin{center}Homilía sancti Joánnis Chrysóstomi.\\
\textred{Hom. 84. in c. 19. Joann.}\end{center}
\lettrineThree{S}ciens Jesus ómnia consummáta, ait, Sítio: étiam hoc verbo prophetíam implens. Tu autem adstántium pravitátem intéllige; nam licet innúmeros habeámus inimícos, qui nos gravíssime persecúti sint, cum tamen intérimi eos vidémus, fléctimur misericórdia. Hi vero nil moti, magis efferántur, et irrídent, acetúmque in spóngia pórrigunt, ut bibat tamquam reus: nam propter hoc hyssópum ádditur, quo gustáto, mórtuus est. Vides quam non perturbáte, et quam poténter ómnia transigéntem? Quod et prósequens osténdit; cum enim ómnia consummáta essent, inclináto cápite, quod affíxum non erat, trádidit spíritum, hoc est exspirávit.

\begin{tightlist}{}
\item \textred{℟.} Ego te potávi aqua salútis de petra: \textred{*} Et tu me potásti felle, et acéto. \textred{℣.} Pópule meus, quid feci tibi, aut in quo contristávi te? respónde mihi. \textred{E}t.
\end{tightlist}

\subsubsection*{\center\textred{Lectio viij.}}
\lettrine{A}tqui non post inclinátum caput exspirátur; sed post exspiratiónem caput inclinátur. Hic autem contra áccidit, per quod osténdit eum Evangelísta esse ómnium Dóminum. Judǽi autem camélum devorántes, et cúlicem excolántes, tantum nefas ausi, de die consúlant: Quóniam Parascéve erat, ut non remanérent in cruce córpora sábbato, rogavérunt Pilátum, ut frangeréntur eórum crura. Vides quanti sit véritas? Per ea, quæ illi fáciunt, prophetía implétur; ália enim ex hoc manifestáta est. Venérunt ergo mílites, et aliórum fregérunt crura, non Christi; sed ad Judæórum grátiam conciliándam latus ejus láncea aperuérunt, et mórtuo adhuc insúltant. O péssimam voluntátem, ac scelestíssimam!

\begin{tightlist}{}
\item \textred{℟.} Ego ante te apérui mare: \textred{*} Et tu aperuísti láncea latus Salvatóri tuo, qui redémit mundum in sánguine suo. \textred{℣.} Pópule meus, quid feci tibi, aut in quo contristávi te? respónde mihi. \textred{E}t.
\end{tightlist}

\subsubsection*{\center\textred{Lectio ix.}}
\lettrine{N}oli tamen perturbári, dilectíssime: quæ enim malo illi ánimo patrárunt veritáti consentiébant. Prophetía hinc impléta est, ínquiens: Vidébunt, in quem transfixérunt. Neque hoc tantum, sed non creditúris fuit fídei arguméntum, ut Thomæ, et qui cum eo erant. Ad hæc et arcánum mystérium consummátum est, exívit enim sanguis, et aqua. Non casu, et simplíciter hi fontes scaturiérunt; sed quóniam ex ambóbus Ecclésia constitúta est. Sciunt hoc initiáti, per aquam enim regeneráti, sánguine, et carne nutríti sunt. Hinc mystéria ortum habent, ut quóties ad admirándum cálicem accédis, tamquam ab ipso látere háuriens accédas.

\begin{tightlist}{}
\item \textred{℟.} Crucem tuam adorámus, Dómine, sanctámque recólimus passiónem: \textred{*} Ecce enim propter lignum venit gáudium in univérso mundo. \textred{℣.} Deus misereátur nostri, et benedícat nobis, illúminet vultum suum super nos, et misereátur nostri. \textred{E}cce. Glória Patri. \textred{E}cce.
\item Te Deum laudámus. 6.
\item \textred{Sacerd. ℣.} Adorámus te, Christe, et benedícimus tibi. \textred{℟.} Quia per mortem tuam redemísti mundum.
\end{tightlist}

\subsubsection*{\center \textsc{AD LAUDES.}}
\begin{center}\textred{et per Horas, Aña.}\end{center}
\lettrine{F}ui flagellátus tota die, et castigátio mea in matutínis.

\begin{tightlist}{}
\item \textred{Ps. de Dominica.} 30.
\item \textred{Aña.} Veni ad montem myrrhæ et ad collem thuris; tamquam agnus ad occisiónem ductus obmútui, et non apérui os meum.
\item \textred{Aña.} Fodérunt manus meas, et pedes meos, dinumeravérunt ómnia ossa mea.
\item \textred{Aña.} Consolántem me quæsívi et non invéni: dedérunt in escam meam fel, et in siti mea potavérunt me acéto.
\item \textred{Aña.} Cum accepísset Jesus acétum dixit: Consummátum est, et inclináto cápite, emísit spíritum.
\end{tightlist}

\subsubsection*{\textred{Capitulum. \hfill Philipp. 2.}}
\lettrineThree{F}ratres, Hoc enim sentíte in vobis quod et in Christo Jesu, qui cum in forma Dei esset, non rapínam arbitrátus est esse se æquálem Deo, sed semetípsum exinanívit, formam servi accípiens in similitúdinem hóminum factus, et hábitu invéntus ut homo. Humiliávit semetípsum factus obédiens usque ad mortem, mortem autem crucis.

\subsubsection*{\center\textred{Hymnus.}}
\lettrine{S}ævo dolórum túrbine\\
Jactátur, atris óbrutus\\
Pœnis, acérba sústinens,\\
Redémptor affíxus cruci.\\
\\
Pedes, manúsque horríbili\\
Clavi cruéntant vúlnere,\\
Cor, vultus, artus, péctora\\
Sacro redúndant sánguine.\\
\\
Flet, orat, et clamans obit:\\
Cor Matris ictum cóncidit:\\
Heu, Mater! Heu Fili! dolor\\
Ingráta frangat péctora.\\
\\
Montes, sepúlcra, sáxaque\\
Scindúntur; arva, flúmina,\\
Rupes et æquor cóntremunt,\\
Templíque velum scínditur.\\
\\
Sol, luna, cœlum, sídera\\
Plangunt, et orbis íngemit:\\
O vos viri, vos párvuli,\\
Nuptæ, puéllæ plángite.\\
\\
Adstáre mœréntes cruci,\\
Pedes beátos úngite,\\
Laváte fletu, térgite\\
Comis, et ore lámbite.\\
\\
Tu caritátis víctima,\\
Ut nostra tollas crímina,\\
Nobis salúbri pérficis\\
Adoptiónem sánguine.\\
\\
Nostra ergo pax, et gáudium\\
Sis vita, Jesu, et prǽmium;\\
Sis ductor, et lux in via,\\
Merces, coróna in Pátria.\\
Amen.

\begin{tightlist}{}
\item \textred{℣.} Ipse passus est propter iniquitátes nostras. \textred{℟.} Et attrítus est propter scélera nostra.
\item \textred{Ad Benedictus, Aña.} Ad Jesum autem cum veníssent, ut vidérunt eum jam mórtuum, non fregérunt ejus crura, sed unus mílitum láncea latus ejus apéruit, et contínuo exívit sanguis et aqua.
\end{tightlist}
\subsubsection*{\center\textred{Oratio.}}
\lettrine{O}mnípotens sempitérne Deus, qui humáno géneri, ad imitándum humilitátis exémplum, Salvatórem nostrum carnem súmere, et crucem subíre fecísti: concéde propítius; ut, sicut solémnem commemoratiónem Passiónis ejus celebrámus, ita et patiéntiæ ipsíus habére documénta, et resurrectiónis consórtia mereámur. Per eúmdem Dóminum.

\begin{tightlist}{}
\item \textred{Ad Horas Psalmi de Dominica, ad Primam tamen ut in Festis.}
\end{tightlist}

\subsubsection*{\center \textsc{AD PRIMAM.}}
\begin{tightlist}{}
\item \textred{℣.} Qui nos salvásti in cruce.
\end{tightlist}

\subsubsection*{\center \textsc{AD TERTIAM.}}
\subsubsection*{\textred{Capitulum. \hfill Philipp. 2.}}
\lettrine{F}ratres, Hoc enim sentíte in vobis quod et in Christo Jesu, qui cum in forma Dei esset, non rapínam arbitrátus est esse se æquálem Deo, sed semetípsum exinanívit, formam servi accípiens, in similitúdinem hóminum factus, et hábitu invéntus ut homo. Humiliávit semetípsum factus obédiens usque ad mortem, mortem autem crucis.

\begin{tightlist}{}
%\item \textred{Capit. Fratres, ut supra.}
\item \textred{℟.} Oblátus est, \textred{*} Quia ipse vóluit. \textred{O}blátus. \textred{℣.} Et livóre ejus sanáti sumus. \textred{Q}uia. Glória Patri. \textred{O}blátus. 
\item \textred{℣.} Próprio Fílio suo non pepércit Deus. \textred{℟.} Sed pro nobis ómnibus trádidit illum.
\end{tightlist}

\subsubsection*{\center \textsc{AD SEXTAM.}}
\subsubsection*{\textred{Capitulum. \hfill Zach. 13.}}
\lettrine{Q}uid sunt plagæ istæ in médio mánuum tuárum? Et dicet? His plagátus sum in domo eórum, qui diligébant me.

\begin{tightlist}{}
\item \textred{℟.} Próprio Fílio suo \textred{*} Non pepércit Deus. \textred{P}róprio. \textred{℣.} Sed pro nobis ómnibus trádidit illum. \textred{N}on. Glória Patri. \textred{P}róprio.
\item \textred{℣.} Factus est obédiens usque ad mortem. \textred{℟.} Mortem autem crucis.
\end{tightlist}

\subsubsection*{\center \textsc{AD NONAM.}}
\subsubsection*{\textred{Capitulum. \hfill Is. 53.}}
\lettrine{G}eneratiónem ejus quis enarrábit? Quia abscíssus est de terra vivéntium, propter scelus pópuli mei percússi eum.

\begin{tightlist}{}
\item \textred{℟.} Factus est obédiens \textred{*} Usque ad mortem. \textred{F}actus. \textred{℣.} Mortem autem crucis. \textred{U}sque. Glória Patri. \textred{F}actus.
\item \textred{℣.} Ipse passus est propter iniquitátes nostras. \textred{℟.} Et attrítus est propter scélera nostra.
\end{tightlist}

\subsubsection*{\center \textsc{AD VESPERAS.}}
\begin{tightlist}{}
%\item \textred{Añæ Laudum.}
\item \textred{Aña.} Fui flagellátus tota die, et castigátio mea in matutínis.
\item \textred{Ps. 115.} Crédidi. 93.
\item \textred{Aña.} Veni ad montem myrrhæ et ad collem thuris; tamquam agnus ad occisiónem ductus obmútui, et non apérui os meum.
\item \textred{Ps. 119.} Ad Dóminum. 93. 
\item \textred{Aña.} Fodérunt manus meas, et pedes meos, dinumeravérunt ómnia ossa mea.
\item \textred{Ps. 139.} Eripe me. 236. 
\item \textred{Aña.} Consolántem me quæsívi et non invéni: dedérunt in escam meam fel, et in siti mea potavérunt me acéto.
\item \textred{Ps. 140.} Dómine, clamávi. 237. 
\item \textred{Aña.} Cum accepísset Jesus acétum dixit: Consummátum est, et inclináto cápite, emísit spíritum.
\item \textred{Ps. 141.} Voce mea. 238.
\item \textred{Capit. Hymnus et ℣. ut in j. Vesperis.}
\item \textred{Ad Magnif. Aña.} Depónens Joseph Corpus Jesu invólvit síndone, et pósuit eum in monuménto excíso: fuit planctus magnus in illa die. Mulíeres autem sedéntes contra sepúlcrum dolébant, sicut doléri solet in morte primogéniti.
\item \textred{Completorium de Dominica.}
\item \textred{Ad Nunc dimittis, Aña.} Hoc signum Crucis * erit in cælo, cum Dóminus ad judicándum vénerit.
\end{tightlist}

\begin{center}\textcolor{red}{\rule{\linewidth}{0.5pt}} \end{center}

\subsection*{\center \textsc{DIE IV. FEBRUARII.}}
\subsubsection*{\center \textred{INFRA OCTAVAM}}
\subsection*{\center \textsc{S. Andreæ Corsini.}}
%\subsubsection*{\center \textred{Duplex minus.}}
\fancyhead[C]{\textsc{4 Feb - Infra Octav. S. Andreæ Corsini}}

\begin{tightlist}{}
\item \textred{Infra Octavam S. Andreæ Corsini, fit de ea com. in Laud. per Añam, quæ habetur ad Benedictus, et in Vesp. per Añam quæ habetur ad Magnificat de ij. Vesp. Festi, exceptis Vesp. Sabbati, in quibus dicitur Aña, quæ habetur ad Magnificat in primis Vesp.}
\end{tightlist}

\begin{tightlist}{}
\item \textred{Ad j. Vesp. Magnif. Aña.} Natus est firmaméntum gentis, rector fratrum stabiliméntum pópuli.
\item \textred{℣.} Ora pro nobis sancte Andréa. \textred{℟.} Ut digni efficiámur promissiónibus Christi, (allelúja).
\end{tightlist}
\subsubsection*{\center \textred{Oratio.}}
\lettrine{D}eus, qui in Ecclésia tua nova semper instáuras exémpla virtútum: da pópulo tuo; beáti Andréæ Confessóris tui, atque Pontíficis ita sequi vestígia, ut assequátur et prǽmia. Per Dóminum.

\begin{tightlist}{}
\item \textred{Ad Benedictus, Aña.} Fecit quod rectum erat in óculis Dómini, et diréxit eum in ómnibus, et adjúvit eum Deus, et divulgátum est nomen ejus.
\item \textred{℣.} Justus germinábit sicut lílium. \textred{℟.} Et florébit in ætérnum ante Dóminum.
\item \textred{Ad ij. Vesp. Magnif. Aña.} Hic est Fratrum amátor, et pópuli: hic est qui multum orat pro pópulo, et univérsa civitáte.
\item \textred{℣.} Ora pro nobis sancte Andréa. \textred{℟.} Ut digni efficiámur promissiónibus Christi.
\end{tightlist}

\begin{center}\textcolor{red}{\rule{\linewidth}{0.5pt}} \end{center}

%\subsection*{\center \textsc{DIE VI. FEBRUARII.\\
%Beatæ Archangelæ Girlani}}
%\fancyhead[C]{\textsc{6 Feb - Beatæ Archangelæ Girlani}}
%\begin{center}\textred{Virginis, Ordinis Nostri. Duplex minus.}\end{center}
%\begin{tightlist}{}
%\item \textred{Omnia de Comm. unius Virg. [106] præter sequentia.}
%\end{tightlist}
%\subsubsection*{\center \textred{Oratio.}}
%\lettrine{D}eus, qui beátam Archángelam Vírginem in benedictiónibus dulcédinis prævéntam heróicis ditásti virtútibus; ipsa interveniénte concéde, ut dono grátiæ tuæ muníti in terris, Beatórum choris sociári mereámur in cœlis. Per Dóminum.
%
%\begin{tightlist}{}
%\item \textred{In j. Nocturno, Lectiones de Script. occurr. Sed si in Quadragesima celebretur, Lectiones de Communi: quod per totam Quadrag. observabitur in Festis proprias Lectiones non habentibus.}
%\end{tightlist}
%
%\subsubsection*{\center \textsc{IN SECUNDO NOCTURNO.}}
%\subsubsection*{\center\textred{Lectio iv.}}
%\lettrineThree{A}rchángela in sæculo Eleonóra honésto génere, parénte Joánne Girláni Tridíni in Monte Ferráto orta est. Puérula moniálibus sancti Benedícti instituénda trádita, singulári in Deum pietáte et caritáte in próximum, quantæ futúra esset sanctitátis íllico osténdit. Nihil enim pueríle cógitans, grávia ac sublímia meditabátur ánimo, quæ sursum sunt únice sapiébat. Domum revocáta a patre, atque ad núptias illécta, non modo in propósito virginitátis Deo consecrándæ perseverávit, sed duas étiam soróres germánas, Maríam et Scholásticam, felíciter in eámdem senténtiam addúxit, eo quidem éxitu, ut religiónis, et virtútum sócias habére merúerit. Tandem non sine manifésto Dei nutu ad Perménse Ordinis beatíssimæ Vírginis de Monte Carmélo monastérium convolávit, ubi reguláris disciplínæ studiosíssima, exsúltans ad curréndas perfectiónis vias, ipsas provectióres in virtútum palæstra brevi superávit. Cumque paucos post annos de moderatríce eligénda agerétur in eam moniálium suffrágia conspirárunt.
%
%\begin{tightlist}{}
%\item \textred{℟.} Pulchra fácie, sed púlchrior fide, beáta es Archángela Virgo; réspuens mundum, lætáberis cum Angelis: * Intercéde pro ómnibus nobis.
%\item \textred{℣.} Spécie tua et pulchritúdine tua inténde, próspere procéde, et regna. Intercéde.
%\end{tightlist}
%
%\subsubsection*{\center\textred{Lectio v.}}
%\lettrine{E}o in múnere, facis ad instar super candelábrum pósitæ, exímia Archángelæ sánctitas luculéntius enítuit: cum ipsa perfectíssime, ac longe præstántius quam céteræ, in se exprímeret illud Apóstoli: múlier innúpta et virgo cógitat quæ Dómini sunt, ut sit sancta córpore et spíritu. Quare paulo post munus eídem concréditum fuit novi regéndi ascetérii ab insígni pietáte príncipum Gonzága Mántuæ excitáti. Huc digréssa pusíllum gregem ita excóluit, ut brevi in amplam beátæ Vírginis famíliam créverit. Soróribus tanta præíbat humilitáte, suavitáte morum et patiéntia, ut eam non solum Matrem, sed quasi Angelum e cœlo demíssum veneraréntur. Disciplínæ reguláris et pietátis erga sacrum beátæ Vírginis Scapuláre observantíssima, alúmnas edocébat véstium exteriórum mundítiem, sine interióri cordis puritáte, esse spernéndam; satagebátque ut quemádmodum títulus monastério ínerat sanctæ Maríæ de Paradíso, ita ipsa et soróres, vitam degéntes in terris, conversatiónem habérent in cœlis.
%
%\begin{tightlist}{}
%\item \textred{℟.} Spécie tua et pulchritúdine tua: * Inténde, próspere procéde, et regna. 
%\item \textred{℣.} Propter veritátem, et mansuetúdinem, et justítiam. Inténde.
%\end{tightlist}
%
%\subsubsection*{\center\textred{Lectio vj.}}
%\lettrine{I}n fídei mystériis assídue recoléndis non semel in éxtasim eláta visa est. Miris étiam prodígiis sanctitátem sponsæ suæ Deus, ut fertur, non raro osténdit, eámque ætáte floréntem, sed cœlo matúram ad præmium vocávit. Assíduis quippe labóribus, jejúniis, et aspérrima córporis castigatióne attríta, cum nondum triénnium Mántuæ explevísset, gravíssimo corrépta morbo, suprémum sibi imminére diem ubi intelléxit, colléctas soróres sapientíssimis mónitis cohortáta, ut in cœpto perfectiónis trámite pérgerent alácriter, Ecclésiæ sacraméntis muníta, óculis in Christi crucifíxi effígiem inténtis, et consuéta sibi verba ingéminans: Jesus amor meus; placidíssime ánimam efflávit octávo Kaléndas Februárii anno salútis millésimo quadringentésimo nonagésimo quarto. Públicum ecclesiásticum cultum ab immemorábili témpore beátæ Archángelæ delátum a Sacra Rítuum Congregatióne recógnitum, Summus Póntifex Pius Nonus approbávit, et offícium in ejus honórem recitári cleménter indúlsit.
%
%\begin{tightlist}{}
%\item \textred{℟.} Induit me Dóminus vestiménto salútis: et induménto lætítiæ circúmdedit me, et tamquam sponsam: * Decorávit me coróna.
%\item \textred{℣.} Quasi sponsum decorátum coróna, et tamquam sponsam ornátam. Decorávit.
%\item Glória Patri. Decorávit.
%\end{tightlist}
%
%\subsubsection*{\center \textsc{IN TERTIO NOCTURNO.}}
%\begin{tightlist}{}
%\item \textred{Lect. de Hom. in Evang.} Símile erit regnum cœlórum, \textred{de Communi.} [111]
%\item \textred{Nisi legi debeat \textbf{Lectio ix}. de aliqua Homilia, legatur:}
%\end{tightlist}
%\subsubsection*{\center Pro sanctis Vedásto et Amándo Epíscopis et Confessóribus.}
%\begin{center}\textred{\textbf{Lectio ix}.}\end{center}
%\lettrine{V}edástus témpore Clodovéi Francórum regis, quem ipse, ántequam ad Episcopátum assumerétur, in fide instrúxerat, a sancto Remígio Epíscopus consecrátus, quadragínta annis Atrebaténsi Ecclésiæ, magna vitæ sanctitáte, et miraculórum glória clarus præfuit. Ad cujus quoque mérita declaránda e cœlo colúmna ædibus ægrotántis illucéscens illíus óbitum præsignávit. Amándus vero Epíscopus illústri génere in Aquitánia natus, monásticæ vitæ disciplínam puer ampléxus, paréntis rogátu ab ea dimovéri nullo modo pótuit. Cumque annis quíndecim mira abstinéntia in monastério vixísset, peregrinatiónis deínde causa Romam proféctus est. At cœlésti admónitu in Gálliam revérsus, ibi multis in locis fidem prædicávit. Complúres hábuit discípulos sanctitátis laude præstántes: innumerábiles ab impietáte ad Christiánæ religiónis cultum perdúxit. Dénique in tota vita sanctíssime versátus, fideíque causa multis labóribus pie perfúnctus, Heraclío Imperatóre, quiévit in Dómino, anno sexcentésimo vigésimo.
%
%\begin{tightlist}{}
%\item \textred{In Laud. fit comm. Octavæ Purificationis B. Mariæ Virginis, Aña.} Senex púerum portábat, puer autem senem regébat, quem Virgo péperit, et post partum Virgo permánsit: ipsum, quem génuit, adorávit.
%\item \textred{℣.} Elégit eam. [176]
%\end{tightlist}
%\subsubsection*{\center \textred{Oratio.}}
%\lettrine{O}mnípotens sempitérne Deus, majestátem tuam súpplices exorámus: ut sicut Unigénitus Fílius tuus hodiérna die, cum nostræ carnis substántia in templo est præsentátus, ita nos fácias purificátis tibi méntibus præsentári. Per eúmdem Dóminum.
%
%\begin{tightlist}{}
%\item \textred{Item fit comm. Octavæ S. Andreæ Corsini, Aña.} Fecit quod rectum erat in óculis Dómini, et diréxit eum in ómnibus, et adjúvit eum Deus, et divulgátum est nomen ejus.
%\item \textred{℣.} Justus germinábit. [74]
%\end{tightlist}
%\subsubsection*{\center \textred{Oratio.}}
%\lettrine{D}eus, qui in Ecclésia tua nova semper instáuras exémpla virtútum, da pópulo tuo beáti Andréæ Confessóris tui atque Pontíficis ita sequi vestígia, ut assequátur et præmia. Per Dóminum nostrum.
%
%\begin{tightlist}{}
%\item \textred{Et Comm. SS. Vedasti, et Amandi, Aña.} Sint lumbi vestri præcíncti, et lucérnæ ardéntes in mánibus vestris: et vos símiles homínibus exspectántibus dóminum suum, quando revertátur a núptiis.
%\item \textred{℣.} Mirábilis Deus. [104]
%\end{tightlist}
%\subsubsection*{\center \textred{Oratio.}}
%\lettrine{D}eus, qui nos sanctórum tuórum Confessórum paritérque Pontíficum Vedásti et Amándi confessiónibus gloriósis circúmdas, et prótegis: da nobis; et eórum imitatióne profícere, et intercessióne gaudére. Per Dóminum. 
%
%\begin{tightlist}{}
%\item \textred{Vesp. a cap. de S. Romualde Abbatis. 860. Comm. præc. et Octavarum.}
%\end{tightlist}

%\begin{center}\textcolor{red}{\rule{\linewidth}{0.5pt}} \end{center}

\subsection*{\center \textsc{DIE IX. FEBRUARII.}}
\fancyhead[C]{\textsc{9 Feb - In Octava Purificationis B.M.V.}}
\subsubsection*{\center \textred{In Octava}}
\subsection*{\center \textsc{Purificationis\\ B. Mariæ Virginis.}}
\subsubsection*{\center \textred{Duplex minus.}}
\begin{tightlist}{}
\item \textred{A capitulo omnia ut in die festo,} 828. \textred{sed si Octava celebretur post Officium non duplex, in j. Vesp. prima Aña. Laudum, et Psalmi de Officio præc.}
\item \textred{In j. Vesp. fit Comm. S. Cyrilii, etc.} 865. 
\end{tightlist}

\subsubsection*{\center \textsc{IN PRIMO NOCTURNO.}}
\begin{tightlist}{}
\item {\small \textred{Bened.} Alma Virgo vírginum intercédat pro nobis ad Dóminum.}
\end{tightlist}
\subsubsection*{\center De Cánticis Canticórum.}
\begin{center}\textred{\textbf{Lectio j}. \hfill Cap. 5.}\end{center}
\lettrineThree{A}djúro vos fíliæ Jerúsalem, si invenéritis diléctum meum, ut nuntiétis ei, quia amóre lángueo. Qualis est diléctus tuus ex dilécto, o pulchérrima mulíerum? qualis est diléctus tuus ex dilécto, quia sic adjurásti nos? Diléctus meus cándidus, et rubicúndus, eléctus ex míllibus. Caput ejus aurum óptimum, comæ ejus sicut élatæ palmárum, nigræ quasi corvus. Oculi ejus sicut colúmbæ super rívulos aquárum, quæ lacte sunt lotæ, et résident juxta fluénta pleníssima.

\begin{tightlist}{}
\item \textred{℟.} Adórna thálamum tuum Sion, et súscipe Regem Christum: \textred{*} Quem Virgo concépit, Virgo péperit, Virgo post partum, quem génuit adorávit. \textred{℣.} Accípiens Símeon púerum in mánibus, grátias agens benedíxit Dóminum. \textred{Q}uem Virgo.
\item {\small \textred{Bened.} Nos cum prole pia benedícat Virgo María.}
\end{tightlist}

\subsubsection*{\center\textred{Lectio ij. \hfill Cap. 6.}}
\lettrine{D}iléctus meus descéndit in hortum suum ad aréolam arómatum, ut pascátur in hortis, et lília cólligat. Ego dilécto meo, et diléctus meus mihi, qui páscitur inter lília. Pulchra es amíca mea, suávis, et decóra sicut Jerúsalem, terríbilis ut castrórum ácies ordináta. Avérte óculos tuos a me, quia ipsi me avoláre fecérunt. Capílli tui sicut grex caprárum, quæ apparuérunt de Gálaad. Dentes tui sicut grex óvium, quæ ascendérunt de lavácro, omnes geméllis fœ́tibus, et stérilis non est in eis.

\begin{tightlist}{}
\item \textred{℟.} Postquam impléti sunt dies purgatiónis Maríæ secúndum legem Móysi, tulérunt Jesum in Jerúsalem, ut sísterent eum Dómino: \textred{*} Sicut scriptum est in lege Dómini: Quia omne masculínum adapériens vulvam, sanctum Dómino vocábitur. \textred{℣.} Obtulérunt pro eo Dómino par túrturum, aut duos pullos columbárum. \textred{S}icut.
\item {\small \textred{Bened.} Sancta Dei Génitrix sit nobis auxiliátrix.}
\end{tightlist}

\subsubsection*{\center\textred{Lectio iij.}}
\lettrine{U}na est colúmba mea, perfécta mea: una est matris suæ, elécta genitríci suæ. Vidérunt eam fíliæ, et beatíssimam prædicavérunt: regínæ et concubínæ, et laudavérunt eam. Quæ est ista, quæ progréditur, quasi auróra consúrgens, pulchra ut luna, elécta ut sol, terríbilis ut castrórum ácies ordináta? Descéndi in hortum nucum, ut vidérem poma convállium, et inspícerem si floruísset vínea, et germinássent mala púnica. Nescívi: ánima mea conturbávit me, propter quadrígas Amínadab. Revértere, revértere, Sunamítis: revértere, revértere, ut intueámur te.

\begin{tightlist}{}
\item \textred{℟.} Obtulérunt pro eo Dómino par túrturum, aut duos pullos columbárum, \textred{*} Sicut scriptum est in lege Dómini. \textred{℣.} Postquam impléti sunt dies purgatiónis Maríæ secúndum legem Móysi, tulérunt Jesum in Jerúsalem, ut sísterent eum Dómino. \textred{S}icut. Glória Patri. \textred{S}icut.
\end{tightlist}

\subsubsection*{\center \textsc{IN SECUNDO NOCTURNO.}}
\begin{tightlist}{}
\item {\small \textred{Bened.} Précibus suæ Matris benedícat nos Fílius Dei Patris.}
\end{tightlist}
\subsubsection*{\center Sermo sancti Bernárdi Abbátis.}
\begin{center}\textred{Serm. I. de Purificat.}\\
\textred{\textbf{Lectio iv}.}\end{center}
\lettrineThree{H}odie templi Dóminum, in templum Dómini Virgo mater indúcit. Joseph quoque sistit Dómino non suum, sed ejus Fílium diléctum, in quo ei bene complácuit. Agnóscit justus, quem exspectábat. Anna quoque vídua confitétur. Ab his quátuor primo hodiérna procéssio celebráta est, quæ póstmodum exsultatióne univérsæ terræ in omni loco, et ab omni Gente celebrarétur. Nec sane mirum si tunc parva fuit, quandóquidem parvus erat qui suscipiebátur. Nullum ibi locum peccátor hábuit, omnes justi, omnes sancti, omnes fuére perfécti. Sed numquid hos tantum salvábis Dómine, cum multiplicáveris misericórdiam tuam Deus?

\begin{tightlist}{}
\item \textred{℟.} Símeon justus, et timorátus exspectábat consolatiónem Israel: \textred{*} Et Spíritus sanctus erat in eo. \textred{℣.} Respónsum accéperat Símeon a Spíritu sancto, non visúrum se mortem, nisi prius vidéret Christum Dómini. \textred{E}t.
\item {\small \textred{Bened.} Jesus Maríæ Fílius sit nobis clemens et propítius.}
\end{tightlist}

\subsubsection*{\center\textred{Lectio v.}}
\lettrine{I}n secúnda jam processióne prǽeunt turbæ, turbæ sequúntur, nec Virgo portat, sed aséllus. Néminem ítaque dedignátur, nec ipsos, qui computruérunt tamquam juménta in stércore suo. Non dedignátur, inquam, sed si Apostólica vestiménta non desint, si doctrína eórum, si morum justítia, si obediéntia, si cáritas coopériat multitúdinem peccatórum, ex hoc jam processiónis suæ glória non censébit indígnos. Magis vero et ipsam, quæ tam paucis colláta vidétur, nobis quoque ab ipso inveníre est reservátum. Quidni reserváret pósteris, quam et prærogávit antíquis? David rex et Prophéta exsultávit ut vidéret diem istum, et vidit, et gavísus est. Alióquin si non víderat, unde psallébat: Suscépimus, inquit, Deus, misericórdiam tuam in médio templi tui?

\begin{tightlist}{}
\item \textred{℟.} Respónsum accéperat Símeon a Spíritu sancto, non visúrum se mortem, nisi prius vidéret Christum Dómini. Et benedíxit Deum et dixit: \textred{*} Nunc dimíttis servum tuum in pace, quia vidérunt óculi mei salutáre tuum Dómine. \textred{℣.} Cum indúcerent púerum Jesum paréntes ejus, ut fácerent secúndum consuetúdinem legis pro eo: accépit eum Símeon in ulnas suas, et dixit. \textred{N}unc.
\item {\small \textred{Bened.} Fílius Vírginis Maríæ det nobis gáudia vitæ.}
\end{tightlist}

\subsubsection*{\center\textred{Lectio vi.}}
\lettrine{S}uscépit hanc misericórdiam Dómini David, suscépit Símeon, suscépimus et nos, et quicúmque sunt præordináti ad vitam: síquidem Christus heri, et hódie, et in ætérnum. In médio enim templi misericórdia est: non in ángulo, aut diversório, quia non est accéptio personárum apud Deum. In commúni pósita est, offértur ómnibus, et nemo illíus expers, nisi qui rénuit. Derivántur aquæ tuæ foras, Dómine Deus, nihilóminus tamen fons tuus tibi próprius est, et non bibit aliénus ex eo. Qui tuus est, non vidébit mortem donec víderit Christum Dómini, ut secúrus dimittátur in pace. Quidni dimittátur in pace: qui Christum Dómini habet in péctore? Ipse enim est pax nostra, qui per fidem hábitat in córdibus nostris.

\begin{tightlist}{}
\item \textred{℟.} Cum indúcerent púerum Jesum paréntes ejus, accépit eum Símeon in ulnas suas, et benedíxit Deum, et dixit: \textred{*} Nunc dimíttis Dómine servum tuum in pace. \textred{℣.} Suscípiens Símeon púerum in mánibus, exclamávit, et dixit. \textred{N}unc. Glória Patri. \textred{N}unc.
\end{tightlist}

\subsubsection*{\center \textsc{IN TERTIO NOCTURNO.}}
\begin{tightlist}{}
\item {\small \textred{Bened.} Fons Evangélii répleat nos dógmate cæli.}
\end{tightlist}
\subsubsection*{\center Léctio sancti Evangélii secúndum Lucam.}
\textred{\textbf{Lectio vij}. \hfill Cap. 2.}\\
\lettrine{I}n illo témpore: Postquam impléti sunt dies purgatiónis Maríæ secúndum legem Móysi, tulérunt Jesum in Jerúsalem, ut sísterent eum Dómino, sicut scriptum est in lege Dómini. Et réliqua.

\begin{center}Homília venerábilis Bedæ Presbýteri.\\
\textred{Ex Comment. in Luc.}\end{center}
\lettrineThree{E}t venit in spíritu in templum. Signíficat eum eádem spíritus grátia, qua olim ventúrum præcognóverat, étiam nunc veniéntem, et jam jamque a se vidéndum cognovísse Salvatórem. Et cum indúcerent púerum Jesum paréntes ejus. Magna quidem Dómini poténtia, sed non minor claret humílitas, ut qui cœlo terráque non cápitur, grandǽvi hóminis gestétur totus in ulnis. 

\begin{tightlist}{}
\item \textred{℟.} Suscípiens Jesum in ulnas suas Símeon exclamávit, et dixit: \textred{*} Tu es vere lumen ad illuminatiónem Géntium, et glóriam plebis tuæ Israel. \textred{℣.} Símeon in mánibus infirmitátem accépit: sed majestátem intus agnóvit, et dixit. \textred{T}u es.
\item {\small \textred{Bened.} Ab hoste malígno deféndat nos Dei Génitrix Virgo.}
\end{tightlist}

\subsubsection*{\center\textred{Lectio viij.}}
\lettrine{S}ed et trópice accépit Símeon Christum, veteránus infántem, ut dóceat nos exúere véterem hóminem qui corrúmpitur cum áctibus suis, et renovátos spíritu mentis nostræ indúere eum, qui secúndum Deum creátus est in justítia, sanctitáte et veritáte, hoc est, deponéntes mendácium, loqui veritátem, et cétera, quæ ad novi hóminis statum pértinent, ore, corde et ópere perfícere.

\begin{tightlist}{}
\item \textred{℟.} Senex púerum portábat, puer autem senem regébat: \textred{*} Quem Virgo concépit, Virgo péperit, Virgo post partum, quem génuit, adorávit. \textred{℣.} Accípiens Símeon púerum in mánibus, grátias agens, benedíxit Dóminum. \textred{Q}uem Virgo.
\item {\small \textred{Bened.} In omni tribulatióne et angústia succúrrat nobis pia Virgo María.}
\end{tightlist}

\subsubsection*{\center\textred{Lectio ix.}}
\lettrine{A}ccipit sénior justus secúndum legem, et timorátus púerum Jesum in ulnas suas, ut signíficet justítiam óperum, quæ ex lege erat (quis enim nésciat, ópera per manus et bráchia solére figurári?) húmili quidem, sed salutári fídei Evangélicæ grátia mutándam. Accipit sénior infántem Christum, ut insínuet hoc sǽculum, quasi sénio jam, et longǽva ætáte deféssum, ad innocéntiam, et (ut ita díxerim) infántiam Christiánæ conversatiónis reditúrum, et sicut áquilæ juventútem illíus esse renovándam. 

\begin{tightlist}{}
\item \textred{℟.} Christi Virgo dilectíssima, virtútum operátrix, opem fer míseris: \textred{*} Súbveni Dómina clamántibus ad te júgiter. \textred{℣.} Quóniam peccatórum mole prémimur, et non est, qui ádjuvet. \textred{S}úbveni. Glória Patri. \textred{S}úbveni.
\item Te Deum laudámus. 6.
\item \textred{Sacerdot. ℣.} Accípiens Símeon púerum in mánibus. \textred{℟.} Grátias agens, benedíxit Dóminum, (allelúja). 
\end{tightlist}

\begin{tightlist}{}
\item \textred{In Laud. Comm. S. Cyrilli, Aña.} Euge serve bone et fidélis, quia super pauca fuísti fidélis, supra multa te constítuam, dicit Dóminus.
\item \textred{℣.} Justus germinábit sicut lílium. \textred{℟.} Et florébit in ætérnum ante Dóminum, (allelúja).
\item Orémus. \hfill \textred{Oratio.}
\end{tightlist}
\lettrine{S}upérni lúminis Pater, qui mentem beáti Cyrílli Confessóris tui atque Pontíficis, ut Genitrícis Fílii tui semper Vírginis Maríæ honórem ab hæreticórum nébulis vindicáret, veræ fídei et sapiéntiæ luce illustráre dignátus es: ipso interveniénte, concéde; ut errántium corda ad veritátis tuæ rédeant unitátem, et in tua simus voluntáte concórdes. Per eúmdem Christum.

\begin{tightlist}{}
\item \textred{Et Comm. S. Apolloniæ, Aña.} Símile est regnum cœlórum sagénæ missæ in mare, et ex omni génere píscium congregánti: quam cum impléta esset, educéntes, et secus littus sedéntes, elegérunt bonos in vasa sua, malos autem foras misérunt.
\item \textred{℣.} Elégit eam Deus, et præelégit eam. \textred{℟.} Et habitáre eam facit in tabernáculo suo.
\item Orémus. \hfill \textred{Oratio.}
\end{tightlist}
\lettrine{D}eus, qui nos ánnua beátæ Apollóniæ Vírginis et Mártyris tuæ solemnitáte lætíficas: da; ut quam venerámur offício, étiam piæ conversatiónis sequámur exémplo. Per Christum.

\begin{tightlist}{}
\item \textred{Et fit commem. Oct. S. Andreæ. Aña.} Fecit quod rectum erat in óculis Dómini, et diréxit eum in ómnibus, et adjúvit eum Deus, et divulgátum est nomen ejus.
\item \textred{℣.} Justus germinábit sicut lílium. \textred{℟.} Et florébit in ætérnum ante Dóminum.
\item Orémus. \hfill \textred{Oratio.}
\end{tightlist}
\lettrine{D}eus, qui in Ecclésia tua nova semper instáuras exémpla virtútum, da pópulo tuo beáti Andréæ Confessóris tui atque Pontíficis ita sequi vestígia, ut assequátur et prǽmia. Per Dóminum.

\begin{tightlist}{}
\item \textred{Deinde, in Quadragesima, Commemoratio Feriæ.}
\end{tightlist}

\begin{tightlist}{}
\item \textred{In Vesp. comm. S. Scholasticæ, Aña.} Accínxit fortitúdine lumbos suos, et roborávit bráchium suum: ideóque lucérna ejus non exstinguétur in sempitérnum.
\item \textred{℣.} Diffúsa est grátia in lábiis tuis. \textred{℟.} Proptérea benedíxit te Deus in ætérnum, (allelúja).
\item Orémus. \hfill \textred{Oratio.}
\end{tightlist}
\lettrine{D}eus, qui ánimam beátæ Vírginis tuæ Scholásticæ ad ostendéndam innocéntiæ viam in colúmbæ spécie cælum penetráre fecísti: da nobis ejus méritis et précibus ita innocénter vívere; ut ad ætérna mereámur gáudia perveníre. Per Dóminum.

\begin{tightlist}{}
\item \textred{Et Comm. S. Cyrilli, Aña.} Iste est, qui ante Deum magnas virtútes operátus est, et omnis terra doctrína ejus repléta est: ipse intercédat pro peccátis ómnium populórum.
\item \textred{℣.} Amávit eum Dóminus et ornávit eum. \textred{℟.} Stolam glóriæ índuit eum, (allelúja).
\item \textred{Oratio} Supérni lúminis, \textred{ut supra.}
\item \textred{Et Comm. S. Appoloniæ, Aña.} Quinque prudéntes Vírgines accepérunt óleum in vasis suis cum lampádibus. Média autem nocte clamor factus est: Ecce sponsus venit, exíte óbviam ei.
\item \textred{℣.} Diffúsa est grátia in lábiis tuis. \textred{℟.} Proptérea benedíxit te Deus in ætérnum.
\item \textred{Oratio} Deus, qui nos ánnua, \textred{ut supra.}
\item \textred{Et com. Oct. S. Andreæ, Aña.} Hic est Fratrum amátor, et pópuli: hic est qui multum orat pro pópulo, et univérsa civitáte.
\item \textred{℣.} Ora pro nobis sancte Andréa. \textred{℟.} Ut digni efficiámur promissiónibus Christi.
\item \textred{Oratio} Deus, qui in Ecclésia, \textred{ut supra.}
\item \textred{Postea, in Quadragesima, Commemoratio Feriæ.}
\end{tightlist}


\begin{center}\textcolor{red}{\rule{\linewidth}{0.5pt}} \end{center}

\subsection*{\center \textsc{DIE XI. FEBRUARII.}}
\subsubsection*{\center \textred{Dies Octava}}
\subsection*{\center \textsc{S. Andreæ Corsini.}}
\subsubsection*{\center \textred{Duplex minus.}}
\fancyhead[C]{\textsc{11 Feb - Dies Octava S. Andreæ Corsini}}

\begin{tightlist}{}
\item \textred{Omnia ut in Officio ejusdem menstruo, infra, sed si non præcedat Festum duplex, in j. Vesp. prima Aña Laudum} Diléxit Andréam.
\item \textred{Psalmi de Officio præc.}
\item \textred{Post capit. ℟.} Spíritu magno vidit última, consolátus est lugéntes, osténdit futúra, et abscóndita. \textred{*} Quæsívit bona genti suæ, et fecit pacem. \textred{℣.} Adjúvit páuperem de inópia, et pósuit sicut oves famílias. \textred{Q}uæsívit. Glória Patri. \textred{Q}uæsívit.
\item \textred{In j. Nocturno, Lectiones de Scriptura occur.}
\item \textred{In ij. et iij. Noct. Lect. de Offic. menstr.}
\item \textred{Ad Vesp. Añæ Laudum, Ps. de Dominica, et loco ultimi, Ps. 131} Meménto Dómine David. [76]
\end{tightlist}

\begin{center}\textcolor{red}{\rule{\linewidth}{0.5pt}} \end{center}

%\subsection*{\center \textsc{DIE XXVI. FEBRUARII.}}
%%\subsubsection*{\center \textred{}}
%\subsection*{\center \textsc{S. Margaritæ de Cortona.}}
%\subsubsection*{\center \textred{Semiduplex.}}
%\fancyhead[C]{\textsc{26 Feb - S. Margaritæ de Cortona}}
%
%\begin{tightlist}{}
%\item \textred{Omnia de Comm. unius Matronæ} [127] \textred{præter hic signata.}
%\end{tightlist}
%
%\subsubsection*{\center \textred{Oratio.}}
%\lettrine{D}eus, qui fámulam tuam Margarítam de perditiónis via ad salútis trámitem misericórditer deduxísti: eádem nobis miseratióne concéde; ut, quam prius errántem sectári non erubúimus, mox pœniténtem ímpigre sequi gloriémur. Per Dóminum.
%
%\begin{tightlist}{}
%\item \textred{In primo Nocturno, Lectiones de Script. occurrente.}
%\item \textred{In Quadrag.} Mulíerem fortem, \textred{de Communi.} [128]
%\end{tightlist}
%
%\subsubsection*{\center \textsc{IN SECUNDO NOCTURNO.}}
%\subsubsection*{\center \textred{Lectio iv.}}
%
%\lettrineThree{M}argaríta, a loco dormitiónis, Cortonénsis appelláta, Laviáni in Túscia ortum hábuit. Primis adolescéntiæ suæ annis mundi voluptátibus capta in Montis Politiáni civitáte vanam, et lúbricam vitam duxit: sed cum amásium ab hóstibus fœde transfóssum indício canis in fóvea sub strue lignórum tumulátum fortúito reperísset, íllico facta est magnus Dómini super eam: quæ magno culpárum suárum mœróre tacta, éxiit foras, et flevit amáre. Itaque Laviánum revérsa crine detónso, neglécto cápite, pulláque veste contécta, erróribus suis, mundíque illécebris núntium misit; inque ǽdibus Deo sacris fune ad collum alligáto, humi procúmbens, ab ómnibus, quos ántea móribus suis palam offénderat, véniam exorávit. Mox Cortónam profécta in cínere, et cilício a se læsam Dei majestátem placáre stúduit; donec post triennále virtútum experiméntum a Frátribus Minóribus spirituális vitæ dúcibus Tértii Ordinis hábitum impetrávit. Uberes exínde lácrymæ ei familiáres fuérunt, atque ima suspíria, tanta ánimi contritióne ducta, ut diu elínguis consísteret. Léctulum nuda humus, cervícal lapis aut lignum porréxit, atque ita noctes insómnes in cœléstium meditatióne tráhere consuévit, nullum ámplius pravum desidérium perpéssa, dum bonus spíritus prómptior infírmam carnem ad subeúndos labóres erigébat.\\
%
%\begin{tightlist}{}
%\item \textred{℟.} Audi fília, et vide, et inclína aurem tuam: et oblivíscere pópulum tuum, et domum patris tui. Et concupíscet Rex decórem tuum: * Quóniam ipse est Dóminus Deus tuus.
%\textred{℣.} Induit te Dóminus vestiménto salútis: et induménto lætítiæ circúmdedit te. \textred{Q}uóniam.
%\end{tightlist}
%
%\subsubsection*{\center \textred{Lectio v.}}
%\lettrine{A} dǽmone insídiis funestísque conátibus lacessíta, múlier fortis hostem ex verbis detéctum semel atque íterum invícta répulit. Ad eludéndum vanæ glóriæ lenocínium, quo a malo spíritu petebátur, prætéritos mores suos per vicos, et platéas alta voce accusáre non déstitit, omni supplício se ream inclámans: nec nisi a confessário detérrita in speciósam fáciem olim impúri amóris causam sævíre abstínuit, ægre ferens suam formam longa carnis maceratióne non aboléri. Quibus, aliísque magnæ pœniténtiæ arguméntis suórum críminum labe expiáta, atque ita de se triumphátrix, ut sensus plane omnes a mundi illécebris custodíret, digna facta est, quæ sæpe Dómini consuetúdine fruerétur. Ejúsdem quoque Christi, et Vírginis Matris dolórum, quod ipsa ardénter expetíerat, párticeps facta, cunctis sénsibus destitúta, et vere mórtua intérdum visa est. Ad eam proínde véluti ad perfectiónis magístram ex díssitis étiam regiónibus plúrimi conveniébant: ipsa vero cœlésti, quo erat perfúsa lúmine, córdium secréta, consciéntias hóminum, imo et peccáta in remótis licet pártibus Deum offendéntium cum dolóre, et lácrymis détegens, summáque in Deum, et próximum caritáte fervens ingéntem animárum fructum operáta est. Ægris ad se veniéntibus salútem, obséssis a dǽmone liberatiónem impetrávit. Púerum defúnctum, lugénte matre, ad vitam redúxit. Imminéntes bellórum tumúltus assíduis oratiónibus sedávit. Dénique summæ pietátis opéribus vivos, et mórtuos sibi deméruit.
%
%\begin{tightlist}{}
%\item \textred{℟.} Spécie tua, et pulchritúdine tua: * Inténde, próspere procéde, et regna.
%\textred{℣.} Propter veritátem, et mansuetúdinem, et justítiam. \textred{I}nténde.
%\end{tightlist}
%
%\subsubsection*{\center \textred{Lectio vj.}}
%\lettrine{T}ot sanctis opéribus occupáta, de rigóre, quo assídue corpus suum exercébat, nihil remísit, neque a stúdio cœléstia meditándi se avélli passa est, in utróque vitæ génere plane admiránda, utrámque sorórem Magdalénam, et Martham réferens. Tandem pro se Dóminum orans, ut ex hac valle lacrymárum sursum in cœléstem pátriam revocarétur, exaudíta est orátio ejus, die atque hora dormitiónis ei patefáctis. Méritis ítaque, et labóribus plena, ac cœléstibus donis cumuláta cœpit córporis víribus destítui, perque dies decem, et septem nullo cibo, sed divínis tantum collóquiis refécta est: tum, sanctíssimis Ecclésiæ Sacraméntis rite suscéptis, vultu hílari, atque óculis in cœlum convérsis octávo Kaléndas Mártias anno ætátis quinquagésimo, suæ conversiónis vigésimo tértio, humánæ vero salútis millésimo ducentésimo nonagésimo séptimo, felix migrávit ad Sponsum. Corpus in hanc usque diem végetum, incorrúptum illǽsum, et suáviter olens summa religióne cólitur in ecclésia Fratrum Minórum, quæ jam ab eádem Margaríta appellátur. Miráculis contínuo flóruit; quibus permóti Románi Pontífices ad augéndum ejus cultum plúrima liberáliter indulsérunt. Benedíctus vero Décimus tértius in festo Pentecóstes die sexta décima Maii anni millésimi septingentésimi vigésimi octávi solémnem ejus Canonizatiónem religiosíssime celebrávit.
%
%\begin{tightlist}{}
%\item \textred{℟.} In sanctitáte fúlgida, et plena felicitáte complévit dies suos sanctíssima Margaríta; * Hacque die clara suscépit cœléstia gáudia.
%\textred{℣.} Pulchra fácie, sed púlchrior fide felicíssima Margaríta. \textred{H}acque. Glória Patri. \textred{H}acque.
%\end{tightlist}
%
%\begin{tightlist}{}
%\item \textred{In tertio Nocturno, Lectiones de Homiilia in Evangelium} Símile est regnum cœlórum thesáuro, \textred{de Communi unius Virginis.} [116]
%\end{tightlist}
%
%\subsubsection*{\center \textsc{IN LAUDIBUS.}}
%\begin{tightlist}{}
%\item \textred{Ad Bened. Aña.} Ut cognóvit, resípuit Seráphici Ordinis Magdaléna, dimíssa sunt ei peccáta multa, quia diléxit multum.
%\end{tightlist}
%
%\subsubsection*{\center \textsc{IN II. VESPERIS.}}
%\begin{tightlist}{}
%\item \textred{Ad Magnificat, vel pro comm. Aña.} Ego dilécto meo, et ad me convérsio ejus: invéni quem díligit ánima mea, ténui eum, nec dimíttam.
%\end{tightlist}
%
%\begin{center}\textcolor{red}{\rule{\linewidth}{0.5pt}} \end{center}
%
%\subsection*{\center \textsc{DIE V. MARTII.}}
%%\subsubsection*{\center \textred{}}
%\subsection*{\center \textsc{S. Casimiri.}}
%\subsubsection*{\center \textred{Confessoris.\\Semiduplex.}}
%\fancyhead[C]{\textsc{5 Mar. - S. Casimiri}}
%
%\begin{tightlist}{}
%\item \textred{Omnia de Comm. unius Conf. non Pont.} [85] \textred{præter ea, quæ hic habentur.}
%\end{tightlist}
%
%\subsubsection*{\center \textred{Oratio.}}
%\lettrine{D}eus, qui inter regáles delícias, et mundi illécebras, sanctum Casimírum virtúte constántiæ roborásti: quǽsumus; ut ejus intercessióne fidéles tui terréna despíciant, et ad cœléstia semper adspírent. Per Dóminum.
%
%\begin{tightlist}{}
%\item \textred{In primo Nocturno, si venerit in Quadrag. Lect.} Justus si morte, \textred{de Communi,} [87] \textred{alias de Scriptura occurrente.}
%\end{tightlist}
%
%\subsubsection*{\center \textsc{IN SECUNDO NOCTURNO.}}
%\subsubsection*{\center \textred{Lectio iv.}}
%\lettrineThree{C}asimírus, patre Casimíro, matre Elisabétha Austríaca, Polóniæ régibus ortus, a puerítia sub óptimis magístris, pietáte et bonis ártibus instrúctus, juveníles artus áspero domábat cilício, et assíduis extenuábat jejúniis. Régii spreta lecti mollítie, dura cubábat humo, et clam intempésta nocte, præ fóribus templórum pronus in terra divínam exorábat cleméntiam. In Christi contemplánda passióne assíduus, Missárum solémniis ádeo erécta in Deum mente solébat adésse, ut extra se rapi viderétur.
%
%\begin{tightlist}{}
%\item \textred{℟.} Desidérium ánimæ ejus tribuísti ei Dómine. * Et voluntáte labiórum ejus non fraudásti eum.
%\textred{℣.} Posuísti Dómine super caput ejus corónam de lápide pretióso. \textred{E}t.
%\end{tightlist}
%
%\subsubsection*{\center \textred{Lectio v.}}
%\lettrine{C}athólicam promovére fidem summópere stúduit, et Ruthenórum schisma abolére: quaprópter Casimírum patrem indúxit, ut legem ferret, ne schismátici nova templa contrúerent, nec vétera collabéntia restaurárent. Erga páuperes et calamitátibus oppréssos benéficus et miséricors, patris et defensóris egenórum nomen obtínuit. Virginitátem, quem ab incunábulis servávit illǽsam, sub extrémo vitæ término fórtiter asséruit, dum gravi pressus infirmitáte, mori pótius, quam castitátis jactúram ex medicórum consílio subíre, constánter decrévit.
%
%\begin{tightlist}{}
%\item \textred{℟.} Amávit eum Dóminus, et ornávit eum: stolam glóriæ índuit eum: * Et ad portas paradísi coronávit eum.
%\textred{℣.} Induit eum Dóminum lorícam fídei, et ornávit eum. \textred{E}t.
%\end{tightlist}
%
%\subsubsection*{\center \textred{Lectio vj.}}
%\lettrine{C}onsummátus in brevi, virtútibus et méritis plenus, prænuntiáto mortis die, inter sacerdótum, et religiosórum choros, spíritum Deo réddidit, anno ætátis vigésimo quinto. Corpus Vilnam delátum, multis claret miráculis. Etenim prætérquam quod puélla defúncta vitam, cæci visum, claudi gressum, et várii infírmi sanitátem ad ejus sepúlcrum recuperárunt: Lithuánis exíguo número ad potentíssimi hostis insperátam irruptiónem trepidántibus in áere appárens, insígnem tríbuit victóriam. Quibus permótus Leo Décimus eúmdem Sanctórum catálogo adscrípsit.
%
%\begin{tightlist}{}
%\item \textred{℟.} Iste sanctus digne in memóriam vértitur hóminum: qui ad gáudium tránsiit Angelórum: quóniam in hac peregrinatióne solo córpore constitútus, cogitatióne, et aviditáte, * In illa ætérna pátria conversátus est.
%\textred{℣.} Vínculis carnis absolútus, taléntum sibi créditum Dómino suo duplicátum reportávit. \textred{I}n. Glória Patri. \textred{I}n.
%\end{tightlist}
%
%\begin{tightlist}{}
%\item \textred{In tertio Nocturno, Lectiones de Homilia in Evang.} Sint lumbi vestri, \textred{de Comm. plurimorum Conf.} [102]
%\item \textred{Vesp. seq. Ps. et comm. præced. Aña.} Hic vir despíciens mundum, et terréna, triúmphans, divítias cœlo cóndidit ore, manu.
%\item \textred{℣.} Amávit eum Dóminus, et ornávit eum. \textred{℟.} Stolam glóriæ índuit eum.
%\item Orémus. \hfill \textred{Oratio.}
%\end{tightlist}
%\lettrine{D}eus, qui inter regáles delícias, et mundi illécebras, sanctum Casimírum virtúte constántiæ roborásti: quǽsumus; ut ejus intercessióne fidéles tui terréna despíciant, et ad cœléstia semper adspírent. Per Dóminum.
%
%\begin{center}\textcolor{red}{\rule{\linewidth}{0.5pt}} \end{center}

%\subsection*{\center \textsc{DIE XVI. FEBRUARII.}\\
%Commemorationis Sanctorum,}
%\subsubsection*{\center \textred{quorum Corpora, vel Reliquiæ asservantur in Ecclesiis Ordinis Nostri.}\\
%\textred{Duplex minus.}}
%
%\begin{tightlist}{}
%\item \textred{Omnia de Communi plurimorum Martyrum [42] exceptis hic signatis.}
%\end{tightlist}
%\fancyhead[C]{\textsc{16 Feb - Sanctorum Reliquorum}}
%
%\subsubsection*{\center \textred{Oratio.}}
%\lettrine{D}eus, qui hanc sacrosánctam Ecclésiam tot voluísti Sanctórum decorári Relíquiis: concéde nobis fámulis; ut quorum memóriam venerámur in terris, eórum consórtio perfruámur in cœlis. Per Dóminum.
%
%\begin{tightlist}{}
%\item \textred{In primo Nocturno, Lectiones de Script. occurr.}
%\item \textred{In Quadragesima vero} Justórum ánimæ, \textred{de Communi.} [46]
%\end{tightlist}
%
%\subsubsection*{\center \textsc{IN SECUNDO NOCTURNO.}}
%\subsubsection*{\center Ex Tractátu sancti Joánnis Damascéni de Fide Orthodóxa.}
%\begin{center}\textred{Lib. 16. cap. 4.}\\
%\textred{\textbf{Lectio iv}.}\end{center}
%
%\lettrineThree{C}hristus Dóminus Sanctórum Relíquias velut salutíferos fontes prǽbuit, ex quibus plúrima ad nos benefícia manant, suavissimúmque unguéntum prófluit. Nec vero quisquam huic sermóni fidem détrahat; nam si aqua in desérto ex áspera, et sólida rupe, atque ex ásini maxílla ad sedándam Sampsónis sitim, Deo ita volénte prosíliit; erítne, cur cuíquam incredíbile videátur, ex Mártyrum relíquiis suáve unguéntum scaturíre? Mínime certe iis quidem, quibus Dei poténtia, honórque, quo Sanctos suos áfficit, perspéctus, et explorátus est. Atque illud quidem véteri lege comparátum erat, ut quisquis mórtuum contigísset, immúndus censerétur; verum hi in mortuórum número mínime habéndi sunt.
%
%\begin{tightlist}{}
%\item \textred{℟.} Sancti tui, Dómine, mirábile consecúti sunt iter, serviéntes præcéptis tuis, ut inveniréntur illǽsi in aquis válidis: * Terra appáruit árida, et in mari rubro via sine impediménto. \textred{℣.} Victrícem manum tuam laudavérunt páriter, et decantavérunt Dómine, nomen sanctum tuum. \textred{T}erra.
%\end{tightlist}
%
%\subsubsection*{\center\textred{Lectio v.}}
%\lettrine{E}x quo enim ípsamet vita, ac vitæ Auctor inter mórtuos adscríptus est, eos, qui in spe resurrectiónis, fidéque erga eum, diem extrémum clausérunt, mórtuos haudquáquam appellámus. Qui enim mórtuum corpus mirácula édere queat? Per ea dǽmones expellúntur, morbi profligántur, ægróti sanántur, cæci prospíciunt, leprósi mundántur, tentatiónes, ac mœróres discutiúntur: ac dénique omne donum óptimum a Patre lúminum, eórum ópera ad eos, qui fide mínime dúbia póstulant, descéndit. Quid non labóris suscípias, ut patrónum quemdam nanciscáris, qui te mortáli regi ófferat, ac tuo nómine ad eum verba fáciat? An non ígitur hi honorándi sunt, qui totíus humáni géneris patrónos se profiténtur; ac Deo nostra causa súpplicant?
%
%\begin{tightlist}{}
%\item \textred{℟.} Vérbera carníficum non timuérunt sancti Dei, moriéntes pro Christi nómine, * Ut herédes fíerent in domo Dómini. \textred{℣.} Immánia enim pro Christo in suis corpóribus pertulérunt torménta. \textred{U}t.
%\end{tightlist}
%
%\subsubsection*{\center\textred{Lectio vi.}}
%\lettrine{H}onorándi certe, et quidem ita, ut in eórum nómine Deo templa exstruámus, dona offerámus, eórum memóriam colámus, atque in ea spirituáliter oblectémur; quo nimírum lætítia nostra iis, a quibus invitámur, cóngruat; ne, dum eos cólere et demeréri studémus, offendámus pótius, et irritémus. Quibus enim rebus Deus cólitur, iísdem quoque servi ipsíus oblectántur; quibus autem Deus offénditur, iísdem étiam ipsíus mílites offendúntur. Quocírca in psalmis, et hymnis, et cánticis spirituálibus, et compunctióne, et eórum, qui in egestáte versántur, commiseratióne, (quibus obséquiis Deus potíssimum conciliátur) Sanctos colámus, státuas ipsis ac visíbiles imágines erigámus, imo ipsi virtútibus eórum imitándis hoc consequámur, ut vivæ eórum státuæ atque imágines simus. 
%
%\begin{tightlist}{}
%\item \textred{℟.} O constántia Mártyrum laudábilis! o cáritas inexstinguíbilis! o patiéntia invincíbilis! quæ licet inter pressúras persequéntium visa sit despicábilis, * Inveniétur in laudem, glóriam, et honórem in témpore retributiónis. \textred{℣.} Nobis ergo pétimus, piis subvéniant méritis, honorificáti a Patre, qui est in cœlis. \textred{I}nveniétur. Glória Patri. \textred{I}nveniétur.
%\end{tightlist}
%
%\begin{tightlist}{}
%\item \textred{In tertio Noct. Lect. de Homil. in Evang.} Descéndens Jesus, \textred{de Comm.} [53]
%%\item \textred{Vesp. a cap. seq. comm. præced.}
%\end{tightlist}
%
%\begin{center}\textcolor{red}{\rule{\linewidth}{0.5pt}} \end{center}
\clearpage
\pagenumbering{arabic}
\renewcommand{\thepage}{[\arabic{page}]*}

\section*{\center \textsc{\textred{Officia Frequentius Persolvenda}}}
%\begin{center}\textcolor{red}{\rule{\linewidth}{0.5pt}} \end{center}
\subsection*{\center \textsc{\textred{Officia Hebdomadales}}}
\subsubsection*{\center \textred{Dicenda singulis hebomadibus}}
\begin{center}\textcolor{red}{\rule{\linewidth}{0.5pt}} \end{center}

\subsection*{\center \textsc{Officium Sanctissimi Sacramenti.}}
\subsubsection*{\center \textred{Semiduplex.}}
\begin{center}\textsc{recitandum semel qualibet hebdomada.}\end{center}
\subsubsection*{\center \textsc{AD VESPERAS.}}
\fancyhead[C]{\textsc{Officia Hebdomadales - Feria V - SS. Sacramenti}}
\subsubsection*{\textred{Capitulum. \hfill 1 Cor. 11}}
\lettrineThree{D}ominus Jesus, in qua nocte tradebátur accépit panem, et grátias agens fregit, et dixit: Accípite, et manducáte: hoc est Corpus meum, quod pro vobis tradétur.

\subsubsection*{\center\textred{Hymnus.}}
\lettrine{P}ange lingua gloriósi\\ 
Córporis mystérium,\\ 
Sanguinísque pretiósi,\\ 
Quem in mundi prétium\\ 
Fructus ventris generósi\\ 
Rex effúdit Géntium.\\ 
\\
Nobis datus, nobis natus\\ 
Ex intácta Vírgine,\\ 
Et in mundo conversátus,\\ 
Sparso verbi sémine,\\ 
Sui moras incolátus\\ 
Miro clausit órdine.\\ 
\\
In suprémæ nocte cœnæ\\ 
Recúmbens cum frátribus,\\ 
Observáta lege plene\\ 
Cibis in legálibus,\\ 
Cibum turbæ duodénæ\\ 
Se dat suis mánibus.\\ 
\\
Verbum caro, panem verum\\ 
Verbo carnem éfficit:\\ 
Fitque sanguis Christi merum;\\ 
Et si sensus déficit,\\ 
Ad firmándum cor sincérum\\ 
Sola fides súfficit.\\ 
\\
Tantum ergo Sacraméntum\\ 
Venerémur cérnui:\\ 
Et antíquum documéntum\\ 
Novo cedat rítui:\\ 
Prǽstet fides suppleméntum\\ 
Sénsuum deféctui.\\ 
\\
Genitóri, Genitóque\\ 
Laus, et jubilátio,\\ 
Salus, honor, virtus quoque\\ 
Sit et benedíctio,\\ 
Procedénti ab utróque\\ 
Compar sit laudátio.\\
Amen.

\begin{tightlist}{}
\item \textred{℣.} Panem de cœlo præstitísti eis. \textred{℟.} Omne delectaméntum in se habéntem.
\item \textred{Ad Magnificat, Aña.} O quam suávis est Dómine Spíritus tuus, qui ut dulcédinem tuam in fílios demonstráres, pane suavíssimo de cœlo prǽstito, esuriéntes reples bonis, fastidiósos dívites dimíttens inánes!
\end{tightlist}
\subsubsection*{\center \textred{Oratio.}}
\lettrine{D}eus, qui nobis sub Sacraménto mirábili passiónis tuæ memóriam reliquísti: tríbue, quǽsumus, ita nos córporis, et sánguinis tui sacra mystéria venerári, ut redemptiónis tuæ fructum in nobis júgiter sentiámus. Qui vivis.

\begin{tightlist}{}
\item \textred{Ad Completorium in fine Hymni dicitur ℣.} Glória tibi Dómine, qui natus es de Vírgine, \textred{et similiter ad Horas.}
\item \textred{Ad Nunc dimittis, Aña.} Allelúja, Panis: \textred{ut in Psalt. Post Septuagesima dicitur eadem Aña. omisso} Allelúja.
\end{tightlist}

\subsubsection*{\center \textsc{AD MATUTINUM.}}
\begin{tightlist}{}
\item \textred{Invitatorium.} Christum Regem adorémus dominántem Géntibus: * Qui se manducántibus dat spíritus pinguédinem.
\item \textred{Ps.} Veníte, exsultémus. 2.
\end{tightlist}

\subsubsection*{\center\textred{Hymnus.}}
\lettrine{S}acris solémniis juncta sint gáudia,\\ Et ex præcórdiis sonent præcónia;\\ Recédant vétera, nova sint ómnia,\\ Corda, voces, et ópera.\\
\\
Noctis recólitur cœna novíssima,\\ Qua Christus créditur agnum, et ázyma\\ Dedísse frátribus, juxta legítima\\ Priscis indúlta pátribus.\\
\\
Post agnum týpicum, explétis épulis,\\ Corpus Domínicum datum discípulis,\\ Sic totum ómnibus, quod totum síngulis,\\ Ejus fatémur mánibus.\\
\\
Dedit fragílibus córporis férculum,\\ Dedit et trístibus sánguinis póculum,\\ Dicens: Accípite quod trado vásculum,\\ Omnes ex eo bíbite.\\
\\
Sic sacrifícium istud instítuit,\\ Cujus offícium commítti vóluit\\ Solis Presbýteris, quibus sic cóngruit,\\ Ut sumant, et dent céteris.\\
\\
Panis Angélicus fit panis hóminum:\\ Dat panis cǽlicus figúris términum:\\ O res mirábilis! mandúcat Dóminum\\ Pauper servus, et húmilis.\\
\\
Te trina Déitas, únaque, póscimus,\\ Sic nos tu vísita, sicut te cólimus,\\ Per tuas sémitas duc nos quo téndimus,\\ Ad lucem, quam inhábitas.\\
Amen.

\subsubsection*{\center \textsc{IN PRIMO NOCTURNO.}}
\textred{Aña.} Fructum salutíferum.
\subsubsection*{\center Psalmus j.}
\lettrineThree{B}eátus vir, qui non ábiit in consílio impiórum, et in via peccatórum non stetit, * et in cáthedra pestiléntiæ non sedit :

Sed in lege Dómini volúntas ejus, * et in lege ejus meditábitur die ac nocte.

Et erit tamquam lignum, quod plantátum est secus decúrsus aquárum, * quod fructum suum dabit in témpore suo:

Et fólium ejus non défluet: * et ómnia quæcúmque fáciet, prosperabúntur.

Non sic ímpii, non sic: * sed tamquam pulvis, quem prójicit ventus a fácie terræ.

Ideo non resúrgent ímpii in judício: * neque peccatóres in concílio justórum.

Quóniam novit Dóminus viam justórum: * et iter impiórum períbit.

Glória Patri.

\begin{tightlist}{}
\item \textred{Aña.} Fructum salutíferum gustándum dedit Dóminus mortis suæ témpore.
\item \textred{Aña.} A fructu fruménti.
\end{tightlist}

\subsubsection*{\center Psalmus iv.}
\lettrine{C}um invocárem exaudívit me Deus justítiæ meæ : * in tribulatióne dilatásti mihi.

Miserére mei, * et exáudi oratiónem meam.

Fílii hóminum, úsquequo gravi corde? * ut quid dilígitis vanitátem, et quǽritis mendácium?

Et scitóte quóniam mirificávit Dóminus sanctum suum: * Dóminus exáudiet me, cum clamávero ad eum.

Irascímini, et nolíte peccáre: * quæ dícitis in córdibus vestris, in cubílibus vestris compungímini.

Sacrificáte sacrifícium justítiæ, et speráte in Dómino. * Multi dicunt : Quis osténdit nobis bona?

Signátum est super nos lumen vultus tui Dómine: * dedísti lætítiam in corde meo.

A fructu fruménti, vini, et ólei sui * multiplicáti sunt.

In pace in idípsum * dórmiam, et requiéscam.

Quóniam tu Dómine singuláriter in spe * constituísti me.

Glória Patri.

\begin{tightlist}{}
\item \textred{Aña.} A fructu fruménti, et vini multiplicáti fidéles, in pace Christi requiéscunt.
\item \textred{Aña.} Communióne cálicis.
\end{tightlist}
\subsubsection*{\center Psalmus xv.}
\lettrine{C}onsérva me Dómine, quóniam sperávi in te. * Dixi Dómino : Deus meus es tu, quóniam bonórum meórum non eges.

Sanctis, qui sunt in terra ejus: * mirificávit omnes voluntátes meas in eis.

Multiplicátæ sunt infirmitátes eórum : * póstea acceleravérunt.

Non congregábo conventícula eórum de sanguínibus: * nec memor ero nóminum eórum per lábia mea.

Dóminus pars hereditátis meæ, et cálicis mei : * tu es, qui restítues hereditátem meam mihi.

Funes cecidérunt mihi in præcláris: * étenim heréditas mea præclára est mihi.

Benedícam Dóminum, qui tríbuit mihi intelléctum : * ínsuper et usque ad noctem increpuérunt me renes mei.

Providébam Dóminum in conspéctu meo semper : * quóniam a dextris est mihi, ne commóvear.

Propter hoc lætátum est cor meum, et exsultávit lingua mea : * ínsuper et caro mea requiéscet in spe.

Quóniam non derelínques ánimam meam in inférno : * nec dabis sanctum tuum vidére corruptiónem.

Notas mihi fecísti vias vitæ, adimplébis me lætítia cum vultu tuo : * delectatiónes in déxtera tua usque in finem.

Gloria Patri.

\begin{tightlist}{}
\item \textred{Aña.} Communióne cálicis, quo Deus ipse súmitur, non vitulórum sánguine, congregávit nos Dóminus.
\item \textred{℣.} Panem cœli dedit eis. \textred{℟.} Panem Angelórum manducávit homo.
\item \textred{Lectiones de Scriptura occurrente; si desint, quæ sequuntur.}
\end{tightlist}

\subsubsection*{\center De Epístola prima beáti Pauli Apóstoli ad Corínthios.}
\textred{\textbf{Lectio j}. \hfill {Cap. 11, 20.}}
\lettrineThree{C}onveniéntibus ergo vobis in unum, jam non est Domínicam cœnam manducáre. Unusquísque enim suam cœnam præsúmit ad manducándum. Et álius quidem ésurit: álius autem ébrius est. Numquid domos non habétis ad manducándum, et bibéndum? aut Ecclésiam Dei contémnitis, et confúnditis eos, qui non habent? Quid dicam vobis? Laudo vos? in hoc non laudo.

\begin{tightlist}{}
\item \textred{℟.} Immolábit hœdum multitúdo filiórum Israel ad vésperam Paschæ: \textred{*} Et edent carnes, et ázymos panes. \textred{℣.} Pascha nostrum immolátus est Christus: ítaque epulémur in ázymis sinceritátis, et veritátis. \textred{E}t edent.
\end{tightlist}

\subsubsection*{\center\textred{Lectio ij.}}
\lettrine{E}go enim accépi a Dómino quod et trádidi vobis, quóniam Dóminus Jesus in qua nocte tradebátur, accépit panem, et grátias agens fregit, et dixit: Accípite, et manducáte: hoc est corpus meum, quod pro vobis tradétur: hoc fácite in meam commemoratiónem. Simíliter et cálicem, postquam cœnávit, dicens: Hic calix novum testaméntum est in meo sánguine: hoc fácite quotiescúmque bibétis, in meam commemoratiónem. Quotiescúmque enim manducábitis panem hunc, et cálicem bibétis: mortem Dómini annuntiábitis donec véniat.

\begin{tightlist}{}
\item \textred{℟.} Comedétis carnes, et saturabímini pánibus. \textred{*} Iste est panis, quem Dóminus dedit vobis ad vescéndum. \textred{℣.} Non Móyses dedit vobis panem de cœlo, sed Pater meus dat vobis panem de cœlo verum. \textred{I}ste.
\end{tightlist}

\subsubsection*{\center\textred{Lectio iij.}}
\lettrine{I}taque quicúmque manducáverit panem hunc, vel bíberit cálicem Dómini indígne: reus erit córporis, et sánguinis Dómini. Probet autem seípsum homo: et sic de pane illo edat, et de cálice bibat. Qui enim mandúcat, et bibit indígne, judícium sibi mandúcat, et bibit; non dijúdicans corpus Dómini. Ideo inter vos multi infírmi et imbecílles, et dórmiunt multi. Quod si nosmetípsos dijudicarémus, non útique judicarémur. Dum judicámur autem, a Dómino corrípimur, ut non cum hoc mundo damnémur.

\begin{tightlist}{}
\item \textred{℟.} Respéxit Elías ad caput suum subcinerícium panem, qui surgens comédit, et bibit, \textred{*} Et ambulávit in fortitúdine cibi illíus usque ad montem Dei. \textred{℣.} Si quis manducáverit ex hoc pane, vivet in ætérnum. \textred{E}t ambulávit. Glória Patri. \textred{E}t ambulávit.
\end{tightlist}

\subsubsection*{\center \textsc{IN SECUNDO NOCTURNO}}
\textred{Aña.} Memor sit Dóminus.
\subsubsection*{\center Psalmus xix.}
\lettrineThree{E}xáudiat te Dóminus in die tribulatiónis: * prótegat te nomen Dei Jacob.

Mittat tibi auxílium de sancto: * et de Sion tueátur te.

Memor sit omnis sacrifícii tui: * et holocáustum tuum pingue fiat.

Tríbuat tibi secúndum cor tuum: * et omne consílium tuum confírmet.

Lætábimur in salutári tuo: * et in nómine Dei nostri magnificábimur.

Impleat Dóminus omnes petitiónes tuas: * nunc cognóvi quóniam salvum fecit Dóminus Christum suum.

Exáudiet illum de cælo sancto suo: * in potentátibus salus déxteræ ejus.

Hi in cúrribus, et hi in equis: * nos autem in nómine Dómini Dei nostri invocábimus.

Ipsi obligáti sunt, et cecidérunt: * nos autem surréximus et erécti sumus.

Dómine salvum fac regem: * et exáudi nos in die, qua invocavérimus te.

Glória Patri.

\begin{tightlist}{}
\item \textred{Aña} Memor sit Dóminus omnis sacrifícii nostri, et holocáustum nostrum pingue fiat.
\item \textred{Aña.} Parátur nobis.
\end{tightlist}

\subsubsection*{\center Psalmus xxij.}
\lettrine{D}óminus regit me, et nihil mihi déerit: * in loco páscuæ ibi me collocávit.

Super aquam refectiónis educávit me: * ánimam meam convértit.

Dedúxit me super sémitas justítiæ : * propter nomen suum.

Nam, et si ambulávero in médio umbræ mortis, non timébo mala: * quóniam tu mecum es.

Virga tua, et báculus tuus: * ipsa me consoláta sunt.

Parásti in conspéctu meo mensam, * advérsus eos, qui tríbulant me.

Impinguásti in óleo caput meum: * et calix meus inébrians quam præclárus est!

Et misericórdia tua subsequétur me * ómnibus diébus vitæ meæ.

Et ut inhábitem in domo Dómini, * in longitúdinem diérum.

Glória Patri.

\begin{tightlist}{}
\item \textred{Aña.} Parátur nobis mensa Dómini, advérsus omnes qui tríbulant nos.
\item \textred{Aña.} In voce exsultatiónis.
\end{tightlist}

\subsubsection*{\center Psalmus xlj.}
\lettrine{Q}uemádmodum desíderat cervus ad fontes aquárum: * ita desíderat ánima mea ad te, Deus.

Sitívit ánima mea ad Deum fortem vivum: * quando véniam, et apparébo ante fáciem Dei?

Fuérunt mihi lácrimæ meæ panes die ac nocte: * dum dícitur mihi quotídie: Ubi est Deus tuus?

Hæc recordátus sum, et effúdi in me ánimam meam: * quóniam transíbo in locum tabernáculi admirábilis, usque ad domum Dei.

In voce exsultatiónis, et confessiónis: * sonus epulántis.

Quare tristis es, ánima mea? * et quare contúrbas me?

Spera in Deo, quóniam adhuc confitébor illi: * salutáre vultus mei, et Deus meus.

Ad meípsum ánima mea conturbáta est : * proptérea memor ero tui de terra Jordánis, et Hermóniim a monte módico.

Abýssus abýssum ínvocat, * in voce cataractárum tuárum.

Omnia excélsa tua, et fluctus tui * super me transiérunt.

In die mandávit Dóminus misericórdiam suam : * et nocte cánticum ejus.

Apud me orátio Deo vitæ meæ, * dicam Deo : Suscéptor meus es.

Quare oblítus es mei? * et quare contristátus incédo, dum afflígit me inimícus?

Dum confringúntur ossa mea, * exprobravérunt mihi qui tríbulant me inimíci mei.

Dum dicunt mihi per síngulos dies: Ubi est Deus tuus? * quare tristis es, ánima mea? et quare contúrbas me?

Spera in Deo, quóniam adhuc confitébor illi : * salutáre vultus mei, et Deus meus.

Glória Patri.

\begin{tightlist}{}
\item \textred{Aña} In voce exsultatiónis résonent epulántes in mensa Dómini.
\item \textred{℣.} Cibávit eos ex ádipe fruménti. \textred{℟.} Et de petra melle saturávit eos.
\end{tightlist}

\subsubsection*{\center Sermo sancti Cyrílli Epíscopi Jerosolymitáni.}
\textred{\textbf{Lectio iv}. \hfill Catechesi myst. 4.}
\lettrineThree{I}psa beáti Pauli doctrína abúnde suffícere vidétur, ut certam vobis de divínis mystériis fidem fáciat: quibus digni rédditi, concorpórei, ut ita dicam, et consanguínei Christi facti estis. Ipse enim modo clamábat, quod in nocte, qua tradebátur Dóminus noster Jesus Christus, accípiens panem et grátias agens, fregit, et dedit discípulis suis, dicens: Accípite et manducáte, hoc est corpus meum. Et accípiens cálicem, et grátias agens, dixit: Accípite, et bíbite: hic est sanguis meus. Cum ígitur ipse de pane pronuntiáverit, ac díxerit: Hoc est corpus meum: quis audébit deínceps ambígere; et cum idem ipse tam asseveránter díxerit: Hic est sanguis meus: quis unquam dubitáverit, ut dicat non esse ejus sánguinem?

\begin{tightlist}{}
\item \textred{℟.} Panis, quem ego dabo, caro mea est pro mundi vita. Litigábant ergo Judǽi ad ínvicem, dicéntes: \textred{*} Quómodo potest hic nobis carnem suam dare ad manducándum? \textred{℣.} Locútus est pópulus contra Dóminum: Anima nostra náuseat super cibo isto levíssimo. \textred{Q}uómodo.
\end{tightlist}

\subsubsection*{\center \textred{Lectio v.}}
\lettrine{A}quam olim in vinum convértit in Cana Galilǽæ: quod habet quamdam cum sánguine propinquitátem: et eum parum dignum existimábimus, cui credámus quod vinum in sánguinem transmutárit? Ad eas núptias, quibus córpora copulántur, vocátus, præter opiniónem ómnium, hoc fecit miráculum: et non multo magis sic eum corpus et sánguinem suum fruénda nobis donásse, persuásum fírmiter habébimus; ut ea cum omni certitúdine tamquam corpus ipsíus, et sánguinem sumámus? nam in spécie panis dat nobis corpus, et in spécie vini dat nobis sánguinem: ut cum súmpseris, gustes corpus, et sánguinem Christi, factus ejúsdem córporis et sánguinis párticeps. Sic enim effícimur Christíferi, hoc est, Christum in corpóribus nostris feréntes, cum corpus ejus et sánguinem in membra nostra recípimus; sic secúndum beátum Petrum, divínæ natúræ consórtes réddimur. 

\begin{tightlist}{}
\item \textred{℟.} Cœnántibus eis, accépit Jesus panem, et benedíxit, ac fregit, dedítque discípulis suis, et ait: \textred{*} Accípite et comédite: hoc est corpus meum. \textred{℣.} Dixérunt viri tabernáculi mei: Quis det de cárnibus ejus, ut saturémur? \textred{A}ccípite.
\end{tightlist}

\subsubsection*{\center \textred{Lectio vj.}}
\lettrine{O}lim cum Judǽis Christus dísserens, nisi manducavéritis, inquit, carnem meam, et bibéritis meum sánguinem, non habébitis vitam in vobis. Cum autem illi, quæ dicta fúerant, non spiritáliter accepíssent, offénsi abiérunt retro: putábant enim, quod eos ad manducándas carnes hortarétur. Erant et in véteri testaménto panes propositiónis; verum illi cum fúerint véteris testaménti, finem jam accepérunt. In novo vero testaménto, panis est cœléstis et calix salutáris: qui et ánimam et corpus sanctíficant. Quámobrem non sic hæc atténdas velim, tamquam sint nudus et simplex panis, nudum et simplex vinum: corpus enim sunt et sanguis Christi. Nam etiámsi sensus illud tibi renúntiat, fides tamen te confírmet. Ne júdices rem ex gustu, sed te citra ullam dubitatiónem fides certum reddat, quod sis dignus factus qui córporis et sánguinis Christi párticeps fíeres.

\begin{tightlist}{}
\item \textred{℟.} Accépit Jesus cálicem, postquam cœnávit, dicens: Hic calix novum testaméntum est in meo sánguine: \textred{*} Hoc fácite in meam commemoratiónem. \textred{℣.} Memória memor ero, et tabéscet in me ánima mea. \textred{H}oc fácite. Glória Patri. \textred{H}oc fácite.
\end{tightlist}

\subsubsection*{\center \textsc{IN TERTIO NOCTURNO.}}
\begin{tightlist}{}
\item \textred{Aña.} Introíbo.
\end{tightlist}
\subsubsection*{\center Psalmus xlij.}
\lettrineThree{J}údica me, Deus, et discérne causam meam de gente non sancta, * ab hómine iníquo, et dolóso érue me.

Quia tu es Deus fortitúdo mea: * quare me repulísti? et quare tristis incédo, dum afflígit me inimícus?

Emítte lucem tuam, et veritátem tuam: * ipsa me deduxérunt, et adduxérunt in montem sanctum tuum, et in tabernácula tua.

Et introíbo ad altáre Dei: * ad Deum, qui lætíficat juventútem meam.

Confitébor tibi in cíthara Deus Deus meus: * Quare tristis es, ánima mea? et quare contúrbas me?

Spera in Deo, quóniam adhuc confitébor illi: * salutáre vultus mei, et Deus meus.

Glória Patri.

\begin{tightlist}{}
\item \textred{Aña.} Introíbo ad altáre Dei: sumam Christum, qui rénovat juventútem meam.
\item \textred{Aña.} Cibávit nos.
\end{tightlist}

\subsubsection*{\center Psalmus lxxx.}
\lettrine{E}xsultáte Deo, adjutóri nostro: * jubiláte Deo Jacob.

Súmite psalmum, et date týmpanum: * psaltérium jucúndum cum cíthara.

Buccináte in Neoménia tuba, * in insígni die solemnitátis vestræ.

Quia præcéptum in Israël est: * et judícium Deo Jacob.

Testimónium in Joseph pósuit illud, cum exíret de terra Ægýpti: * linguam, quam non nóverat, audívit.

Divértit ab onéribus dorsum ejus: * manus ejus in cóphino serviérunt.

In tribulatióne invocásti me, et liberávi te: exaudívi te in abscóndito tempestátis: * probávi te apud aquam contradictiónis.

Audi, pópulus meus, et contestábor te: * Israël, si audíeris me, non erit in te deus recens, neque adorábis deum aliénum.

Ego enim sum Dóminus Deus tuus, qui edúxi te de terra Ægýpti: * diláta os tuum, et implébo illud.

Et non audívit pópulus meus vocem meam: * et Israël non inténdit mihi.

Et dimísi eos secúndum desidéria cordis eórum: * ibunt in adinventiónibus suis.

Si pópulus meus audísset me: * Israël si in viis meis ambulásset:

Pro níhilo fórsitan inimícos eórum humiliássem: * et super tribulántes eos misíssem manum meam.

Inimíci Dómini mentíti sunt ei: * et erit tempus eórum in sǽcula.

Et cibávit eos ex ádipe fruménti: * et de petra, melle saturávit eos.

Glória Patri.

\begin{tightlist}{}
\item \textred{Aña.} Cibávit nos Dóminus ex ádipe fruménti: et de petra, melle saturávit nos.
\item \textred{Aña.} Ex altári tuo.
\end{tightlist}

\subsubsection*{\center Psalmus lxxxiij.}
\lettrine{Q}uam dilécta tabernácula tua, Dómine virtútum: * concupíscit, et déficit ánima mea in átria Dómini.

Cor meum, et caro mea: * exsultavérunt in Deum vivum.

Etenim passer invénit sibi domum: * et turtur nidum sibi, ubi ponat pullos suos.

Altária tua, Dómine virtútum: * Rex meus, et Deus meus.

Beáti, qui hábitant in domo tua, Dómine: * in sǽcula sæculórum laudábunt te.

Beátus vir, cujus est auxílium abs te: * ascensiónes in corde suo dispósuit, in valle lacrimárum in loco, quem pósuit.

Etenim benedictiónem dabit legislátor, ibunt de virtúte in virtútem: * vidébitur Deus deórum in Sion.

Dómine, Deus virtútum, exáudi oratiónem meam: * áuribus pércipe, Deus Jacob.

Protéctor noster, áspice, Deus: * et réspice in fáciem Christi tui.

Quia mélior est dies una in átriis tuis: * super míllia.

Elégi abjéctus esse in domo Dei mei: * magis quam habitáre in tabernáculis peccatórum.

Quia misericórdiam et veritátem díligit Deus: * grátiam et glóriam dabit Dóminus.

Non privábit bonis eos qui ámbulant in innocéntia: * Dómine virtútum, beátus homo qui sperat in te.

Glória Patri.

\begin{tightlist}{}
\item \textred{Aña.} Ex altári tuo Dómine, Christum súmimus: in quem cor, et caro nostra exsúltant.
\item \textred{℣.} Edúcas panem de terra. \textred{℟.} Et vinum lætíficet cor hóminis.
\end{tightlist}

\subsubsection*{\center Léctio sancti Evangélii secúndum Joánnem.}
\begin{center} \textred{\textbf{Lectio vij}. \hfill Cap. 6.} \end{center}
\lettrine{I}n illo témpore: Dixit Jesus turbis Judæórum: Caro mea vere est cibus, et sanguis meus vere est potus. Et réliqua.

\begin{center}Homilía sancti Cyrílli Epíscopi Alexandríni.\\
\textred{Lib. 4. in Joan. cap. 17.} \end{center}
\lettrineThree{Q}ui mandúcat, inquit, carnem meam et bibit sánguinem meum, in me manet, et ego in illo. Sícuti enim si quis liquefáctæ ceræ áliam ceram infúderit, álteram cum áltera per totum commísceat necésse est; sic qui carnem et sánguinem Dómini récipit, cum ipso ita conjúngitur ut Christus in ipso, et ipse in Christo inveniátur. Símile quodámmodo apud Matthǽum compéries: símile est, inquit, regnum cœlórum ferménto, quod accéptum abscóndit múlier in farínæ satis tribus. Sicut parum, ut Paulus ait, ferménti totam massam ferméntat; sic párvula benedíctio totum hóminem in seípsam áttrahit, et sua grátia replet: et hoc modo in nobis Christus manet, et nos in Christo.

\begin{tightlist}{}
\item \textred{℟.} Melchísedech vero rex Salem, próferens panem et vinum, erat enim sacérdos Dei altíssimi, \textred{*} Benedíxit Abram, et ait. \textred{℣.} Benedíctus Abram Deo excélso, qui creávit cœlum, et terram. \textred{B}enedíxit.
\end{tightlist}

\subsubsection*{\center \textred{Lectio viij.}}
\lettrine{N}os vero, si vitam ætérnam cónsequi vólumus, si largitórem immortalitátis habére in nobis desiderámus, ad recipiéndam benedictiónem libénter concurrámus: caveamúsque ne loco láquei damnósam religiónem diábolus nobis præténdat. Recte (inquit) dicis: scriptum tamen esse non ignorámus, judícium sibi comédere atque bíbere illum, qui de pane cómedit, et de cálice bibit indígne. Ego ígitur probo meípsum, et indígnum invénio. Quando ígitur, quicúmque tu es, qui ista dicis, dignus eris? quando Christo teípsum ófferes? Nam si peccándo indígnus es, et peccáre non désinis (quis enim delícta intélligit? secúndum Psalmístam) expers omníno eris vivíficæ hujus sanctificatiónis.

\begin{tightlist}{}
\item \textred{℟.} Calix benedictiónis, cui benedícimus, nonne communicátio sánguinis Christi est? \textred{*} Et panis, quem frángimus, nonne participátio córporis Christi est? \textred{℣.} Quóniam unus panis, et unum corpus multi sumus omnes, qui de uno pane, et de uno cálice participámus. \textred{E}t panis.
\end{tightlist}

\subsubsection*{\center \textred{Lectio ix.}}
\lettrine{Q}uare pias quæso cogitatiónes sucípias, studióse sanctéque vivas, et benedictiónem partícipes: quæ (mihi crede) non mortem solum, verum étiam morbos omnes depéllit. Sedat enim, cum in nobis máneat Christus, sæviéntem membrórum nostrórum legem, pietátem corróborat, perturbatiónes ánimi exstínguit, ægrótos curat, collísos redíntegrat: et sicut Pastor bonus, qui ánimam suam pro óvibus pósuit, ab omni nos érigit casu.

\begin{tightlist}{}
\item \textred{℟.} Unus panis, et unum corpus multi sumus \textred{*} Omnes, qui de uno pane, et de uno cálice participámus. 
\item \textred{℣.} Parásti in dulcédine tua páuperi Deus, qui habitáre facis unánimes in domo. \textred{O}mnes. Glória Patri. \textred{O}mnes.
\item Te Deum laudámus. 6.
\item \textred{Sacerdotalis ℣.} Pluit illis manna ad manducándum. \textred{℟.} Et panem cœli dedit eis.
\end{tightlist}

\subsubsection*{\center \textsc{AD LAUDES.}}
\begin{center}\textred{et per Horas, Aña.}\end{center}
\lettrine{S}apiéntia ædificávit sibi domum, míscuit vinum, et pósuit mensam, allelúja.

\begin{tightlist}{}
\item \textred{Ps. de Dominica.} 30.
\item \textred{Aña.} Angelórum esca nutrivísti pópulum tuum, et panem de cœlo præstitísti eis, allelúja.
\item \textred{Aña.} Pinguis est panis Christi, et præbébit delícias régibus, allelúja.
\item \textred{Aña.} Sacerdótes sancti incénsum, et panes ófferunt Deo, allelúja.
\item \textred{Aña.} Vincénti dabo manna abscónditum, et nomen novum, allelúja.
\end{tightlist}

\subsubsection*{\textred{Capitulum. \hfill 1 Cor. 11.}}
\lettrineThree{D}óminus Jesus in qua nocte tradebátur, accépit panem, et grátias agens fregit, et dixit: Accípite, et manducáte, hoc est corpus meum quod pro vobis tradétur.

\subsubsection*{\center \textred{Hymnus.}}
\begin{center}\textred{Temp. Pasch. nihil additur, neque mutatur.}\end{center}
\lettrine{V}erbum supérnum pródiens,\\
Nec Patris linquens déxteram.\\
Ad opus suum éxiens\\
Venit ad vitæ vésperam.\\
\\
In mortem a discípulo\\
Suis tradéndus ǽmulis,\\
Prius in vitæ férculo\\
Se trádidit discípulis.\\
\\
Quibus sub bina spécie\\
Carnem dedit et sánguinem:\\
Ut dúplicis substántiæ\\
Totum cibáret hóminem.\\
\\
Se nascens dedit sócium,\\
Convéscens in edúlium,\\
Se móriens in prétium,\\
Se regnans dat in prǽmium.\\
\\
O salutáris Hóstia:\\
Quæ cœli pandis óstium,\\
Bella premunt hostília,\\
Da robur, fer auxílium.\\
\\
Uni, trinóque Dómino\\
Sit sempitérna glória:\\
Qui vitam sine término\\
Nobis donet in pátria.\\
Amen.

\begin{tightlist}{}
\item \textred{℣.} Pósuit fines tuos pacem. \textred{℟.} Et ádipe fruménti sátiat te.
\item \textred{Ad Benedictus, Aña.} Ego sum panis vivus, qui de cœlo descéndi. Si quis manducáverit ex hoc pane, vivet in ætérnum, allelúja. 
\end{tightlist}

\subsubsection*{\center \textred{Oratio.}}
\lettrine{D}eus, qui nobis sub Sacraménto mirábili, passiónis tuæ memóriam reliquísti: tríbue, quǽsumus, ita nos Córporis et Sánguinis tui sacra mystéria venerári, ut redemptiónis tuæ fructum in nobis júgiter sentiámus. Qui vivis.

\begin{tightlist}{}
\item \textred{Ad Horas Psalmi de Dominica, ad Primam tamen ut in Festis.}
\end{tightlist}

\subsubsection*{\center \textsc{AD PRIMAM.}}
\noindent \textred{℣.} Qui corpus tuum nobis dedísti in cibum, et sánguinem in potum.

\subsubsection*{\center \textsc{AD TERTIAM.}}
\subsubsection*{\textred{Capitulum. \hfill 1 Cor. 11.}}
\lettrine{D}óminus Jesus in qua nocte tradebátur, accépit panem, et grátias agens fregit, et dixit: Accípite, et manducáte, hoc est corpus meum quod pro vobis tradétur.

\begin{tightlist}{}
%\item \textred{Capitulum.} Dóminus Jesus, \textred{ut supra.}
\item \textred{℟.} Panem cœli \textred{*} Dedit eis. \textred{P}anem. \textred{℣.} Panem Angelórum manducávit homo. \textred{D}edit. Glória Patri. \textred{P}anem.
\item \textred{℣.} Cibávit eos ex ádipe fruménti. \textred{℟.} Et de petra, melle saturávit eos.
\end{tightlist}

\subsubsection*{\center \textsc{AD SEXTAM.}}
\subsubsection*{\textred{Capitulum. \hfill 1 Cor. 11.}}
\lettrine{Q}uotiescúmque manducábitis panem hunc, et cálicem bibétis, mortem Dómini annuntiábitis, donec véniat. 

\begin{tightlist}{}
\item \textred{℟.} Cibávit eos \textred{*} Ex ádipe fruménti. \textred{C}ibávit. \textred{℣.} Et de petra, melle saturávit eos. \textred{E}x. Glória Patri. \textred{C}ibávit.
\item \textred{℣.} Edúcas panem de terra. \textred{℟.} Et vinum lætíficet cor hóminis.
\end{tightlist}

\subsubsection*{\center \textsc{AD NONAM.}}
\subsubsection*{\textred{Capitulum. \hfill 1 Cor. 11.}}
\lettrine{Q}uicúmque manducáverit panem hunc, vel bíberit cálicem Dómini indígne: reus erit córporis et sánguinis Dómini.

\begin{tightlist}{}
\item \textred{℟.} Edúcas panem \textred{*} De terra. \textred{E}dúcas. \textred{℣.} Et vinum lætíficet cor hóminis. \textred{D}e. Glória Patri. \textred{E}dúcas.
\item \textred{℣.} Pósuit fines tuos pacem. \textred{℟.} Et ádipe fruménti sátiat te.
\end{tightlist}

\subsubsection*{\center \textsc{AD VESPERAS.}}
\begin{tightlist}{}
\item \textred{Prima Aña Laud.} Sapiéntia ædificávit sibi domum, míscuit vinum, et pósuit mensam, allelúja.
\item \textred{Ps. 109.} Dixit Dóminus. [56]. 
\item \textred{Ps. 110.} Confitébor. [57]. 
\item \textred{Ps. 115.} Crédidi. [57]. 
\item \textred{Ps. 127.} Beati omnes, qui. 163. 
\item \textred{Ps. 147.} Lauda, Jerúsalem. [59].
%\item \textred{Capit. Hymn. et ℣. de j. Vesperis.}
\end{tightlist}

\subsubsection*{\textred{Capitulum. \hfill 1 Cor. 11}}
\lettrineThree{D}ominus Jesus, in qua nocte tradebátur accépit panem, et grátias agens fregit, et dixit: Accípite, et manducáte: hoc est Corpus meum, quod pro vobis tradétur.

\subsubsection*{\center\textred{Hymnus.}}
\lettrine{P}ange lingua gloriósi\\ 
Córporis mystérium,\\ 
Sanguinísque pretiósi,\\ 
Quem in mundi prétium\\ 
Fructus ventris generósi\\ 
Rex effúdit Géntium.\\ 
\\
Nobis datus, nobis natus\\ 
Ex intácta Vírgine,\\ 
Et in mundo conversátus,\\ 
Sparso verbi sémine,\\ 
Sui moras incolátus\\ 
Miro clausit órdine.\\ 
\\
In suprémæ nocte cœnæ\\ 
Recúmbens cum frátribus,\\ 
Observáta lege plene\\ 
Cibis in legálibus,\\ 
Cibum turbæ duodénæ\\ 
Se dat suis mánibus.\\ 
\\
Verbum caro, panem verum\\ 
Verbo carnem éfficit:\\ 
Fitque sanguis Christi merum;\\ 
Et si sensus déficit,\\ 
Ad firmándum cor sincérum\\ 
Sola fides súfficit.\\ 
\\
Tantum ergo Sacraméntum\\ 
Venerémur cérnui:\\ 
Et antíquum documéntum\\ 
Novo cedat rítui:\\ 
Prǽstet fides suppleméntum\\ 
Sénsuum deféctui.\\ 
\\
Genitóri, Genitóque\\ 
Laus, et jubilátio,\\ 
Salus, honor, virtus quoque\\ 
Sit et benedíctio,\\ 
Procedénti ab utróque\\ 
Compar sit laudátio.\\
Amen.

\begin{tightlist}{}
\item \textred{℣.} Panem de cœlo præstitísti eis. \textred{℟.} Omne delectaméntum in se habéntem.
\item \textred{Ad Magnificat, Aña.} O sacrum convívium, in quo Christus súmitur: recólitur memória passiónis ejus: mens implétur grátia: et futúræ glóriæ nobis pignus datur, allelúja.
\item \textred{Ad Nunc dimittis, Aña.} Allelúja: Panis quem ego dédero, allelúja: caro mea est pro mundi vita, allelúja, allelúja.
\item \textred{A Septuagesima autem usque ad Pascha:} Panis, quem ego dédero, caro mea est pro mundi vita.
\end{tightlist}

\subsubsection*{\center \textsc{A SEPTUAGESIMA.}}
\noindent \textred{Omnia ut supra; sed omittitur ubique} Allelúja.

\subsubsection*{\center \textsc{TEMPORE PASCHALI.}}
\begin{tightlist}{}
\item \textred{Omnia ut supra; sed tres Psalmi cujuslibet Nocturni dicuntur sub prima sola Antiphona, et in fine Añarum, et Responsoriorum, additur} Allelúja.
\item \textred{Item in fine omnium Hymnorum (exceptis Hymnis Vesp. Matut. et Laud.) dicitur ℣.} Quǽsumus Auctor. \textred{et ℣.} Glória tibi Dómine, qui natus es de Vírgine.
\end{tightlist}
\begin{center}\textcolor{red}{\rule{\linewidth}{0.5pt}} \end{center}

\subsection*{\center \textsc{Officium Commem.\\ Beatissimæ V. Mariæ\\ in Sabbato}}
\subsubsection*{\center \textred{Semiduplex.}}
\begin{center}\textsc{recitandum semel qualibet hebdomada.}\end{center}
\subsubsection*{\center \textsc{AD VESPERAS.}}
%\begin{center}\textred{in quibusdam Festis B. Mariæ Virginis.}\end{center}
\begin{tightlist}{}
\item \textred{Omnia de commune B.M.V.} [161] \textred{præter sequentia.}
%\item \textred{In quibusdam Festis B.M.V., in j. Vesperis, Añæ., et Ps. ex j. Vesperis communis B.M.V.}
%\item \textred{Si duplex fieri debet in quibusdam Festis B.M.V. dicitur:}
%\item \textred{℟.} Christi Virgo dilectíssima, virtútum operátrix, opem fer míseris: * Súbveni Dómine clamántibus ad te júgiter. \textred{℣.} Quóniam peccatórum mole prémimur, et non est qui ádjuvet. \textred{S}úbveni. Glória Patri. \textred{S}úbveni.
\end{tightlist}

\fancyhead[C]{\textsc{Officia Hebdomadales - Sabbato - B.V.M.}}

\subsubsection*{\center \textred{Oratio.}}
\lettrine{M}ajestátem tuam omnípotens Deus, súpplices deprecámur; ut qui exsultántibus ánimis beatíssimæ, sempérque Vírginis Maríæ commemoratiónem celebráre gaudémus, ejus gloriósis précibus et méritis, ab instántibus malis, et a morte perpétua liberémur. Per Dóminum.

\begin{tightlist}{}
\item \textred{In fine Hymni Completorii et Horarum dicitur ℣.} Glória tibi Dómine, qui natus es.
\item \textred{Ad Nunc dimittis, Aña.} Sub tuum præsídium.
\end{tightlist}

\begin{center}\textcolor{red}{\rule{\linewidth}{0.5pt}} \end{center}

\subsection*{\center \textred{\textsc{IN ADVENTU.}}}
\begin{tightlist}{}
\item \textred{Omnia ut in die Annuntiationis B.V.M., 25. Martii, ut infra.}
\item \textred{Lect. j. Noct. erunt de Scriptura occurrenti si adsint, alioquin etiam ipsæ accipiantur de eodem Festo Annuntiationis.}
\item \textred{Quum autem non sit dies vacans, quo recitari posset, ideo hic omittitur hoc Officium.}
\end{tightlist}
\fancyhead[C]{\textsc{Officia Hebd. - Sabbato B.V.M. - In Adventu}}
%\begin{center}\textred{\textit{De festo Annuntiationis B. Mariæ. V.}}\end{center}
\subsubsection*{\center \textsc{AD VESPERAS.}}
%\begin{tightlist}{}
%\item \textred{Aña.} Hæc est Regína Vírginum, \textred{cum reliquis Añis. et Psalmis ut in Comm. B.V.M.} [161]
%\end{tightlist}

\subsubsection*{\textred{Capitulum. \hfill Isaiæ 7.}}
\lettrineThree{E}cce Virgo concípiet, et páriet fílium, et vocábitur nomen ejus Emmánuel. Butýrum, et mel cómedet, ut sciat reprobáre malum, et elígere bonum.

\begin{tightlist}{}
%\item \textred{Si fuerit Duplex,}
%\item \textred{℟.} Gaude, María Virgo, cunctas hǽreses sola interemísti, quæ Gabriélis Archángeli dictis credidísti: * Dum Virgo Deum et hóminem genuísti, et post partum Virgo invioláta permansísti. \textred{℣.} Gabriélem Archángelum scimus divínitus te esse affátum: úterum tuum de Spíritu sancto crédimus fecundátum. Erubéscat Judǽus infélix, qui dicit Christum ex Joseph sémine esse natum. \textred{D}um. Glória Patri. \textred{D}um.
\item \textred{Hymn.} Ave, maris stella. [164]
\item \textred{℣.} Ave María, grátia plena. \textred{℟.} Dóminus tecum.
\item \textred{Ad Magnif. Aña.} Spíritus sanctus supervéniet in te, María, et virtus Altíssimi obumbrábit tibi.
\end{tightlist}

\subsubsection*{\center \textred{Oratio.}}
\lettrine{D}eus, qui de beátæ Maríæ Vírginis útero, Verbum tuum, Angelo nuntiánte, carnem suscípere voluísti: præsta supplícibus tuis; ut qui vere eam Genitrícem Dei crédimus, ejus apud te intercessiónibus adjuvémur. Per eúmdem Dóminum.

\begin{tightlist}{}
\item \textred{℣.} Glória tibi Dómine, Qui natus es de Vírgine, \textred{dicitur in fine Hymnorum Complet. et Horarum.}
\item \textred{Ad Nunc dimittis, Aña.} Ecce ancílla Dómini: fiat mihi secúndum verbum tuum.
\end{tightlist}

\subsubsection*{\center \textsc{AD MATUTINUM.}}
\begin{tightlist}{}
\item \textred{Invitat.} Ave María, grátia plena: * Dóminus tecum.
\item \textred{Ps.} Veníte, exsultémus. 2.
\item \textred{Hymn.} Quem terra. [165]
\end{tightlist}

\subsubsection*{\center \textsc{IN PRIMO NOCTURNO.}}
\begin{tightlist}{}
\item \textred{Aña.} Ingréssus Angelus ad Maríam dixit: Ave María, grátia plena, Dóminus tecum.
\item \textred{Ps.} Dómine, Dóminus noster, \textred{cum reliquis de Comm. B.V.M.} [165]
\item \textred{Aña.} Benedícta tu in muliéribus, et benedíctus fructus ventris tui.
\item \textred{Aña.} Angelus Dómini nuntiávit Maríæ, et concépit de Spíritu sancto.
\item \textred{℣.} Sancta Dei Génitrix, Virgo semper María. \textred{℟.} Intercéde pro nobis ad Dóminum, Deum nostrum.
\item \textred{Lectiones j. Noct. sumendi sunt de Scriptura occurente, vel si non habeantur, de Festo Annuntiationis, ut infra:}
\item {\small \textred{Bened.} Alma Virgo vírginum intercédat pro nobis ad Dóminum.}
\end{tightlist}

\subsubsection*{\center De Isaía Prophéta}
\begin{center}\textred{\textbf{Lectio j}. \hfill Cap. 7.}\end{center}
\lettrineThree{E}t adjécit Dóminus loqui ad Achaz, dicens:
Pete tibi signum a Dómino Deo tuo in profúndum inférni, sive in excélsum supra.
Et dixit Achaz: Non petam, et non tentábo Dóminum.
Et dixit: Audíte ergo domus David: Numquid parum vobis est moléstos esse homínibus, quia molésti estis et Deo meo?
Propter hoc dabit Dóminus ipse vobis signum. Ecce Virgo concípiet, et páriet fílium, et vocábitur nomen ejus Emmánuel.
Butýrum et mel cómedet, ut sciat reprobáre malum, et elígere bonum. Hæc dicit Dóminus.

\begin{tightlist}{}
\item \textred{℟.} Missus est Angelus Gábriel ad Maríam Vírginem desponsátam Joseph, núntians ei verbum: et expavéscit Virgo de lúmine. Ne tímeas María; invenísti enim grátiam apud Deum: \textred{*} Ecce concípies, et páries, et Fílius Altíssimi vocábitur.  \textred{℣.} Ave María, grátia plena: Dóminus tecum. \textred{E}cce.
\item {\small \textred{Bened.} Nos cum prole pia benedícat Virgo María.}
\end{tightlist}

\subsubsection*{\textred{Lectio ij. \hfill Cap. 11.}}
\lettrine{E}t egrediétur virga de radíce Jesse, et flos de radíce ejus ascéndet.
Et requiéscet super eum Spíritus Dómini, spíritus sapiéntiæ, et intelléctus, spíritus consílii et fortitúdinis, spíritus sciéntiæ, et pietátis,
et replébit eum spíritus timóris Dómini. Non secúndum visiónem oculórum judicábit, neque secúndum audítum áurium árguet,
sed judicábit in justítia páuperes, et árguet in æquitáte pro mansuétis terræ, et percútiet terram virga oris sui, et spíritu labiórum suórum interfíciet ímpium.
Et erit justítia cíngulum lumbórum ejus, et fides cinctórium renum ejus.
Hæc dicit Dóminus.

\begin{tightlist}{}
\item \textred{℟.} Ingréssus Angelus ad Maríam dixit: \textred{*} Ave María, grátia plena: Dóminus tecum. \textred{℣.} Benedícta tu in muliéribus, et benedíctus fructus ventris tui. \textred{A}ve.
\item {\small \textred{Bened.} Sancta Dei Génitrix sit nobis auxiliátrix.}
\end{tightlist}

\subsubsection*{\textred{Lectio iij. \hfill Cap. 35.}}
\lettrine{L}ætábitur desérta et ínvia, et exsultábit solitúdo, et florébit quasi lílium.
Gérminans germinábit, et exsultábit lætabúnda et laudans: glória Líbani data est ei: decor Carméli, et Saron, ipsi vidébunt glóriam Dómini, et decórem Dei nostri.
Confortáte manus dissolútas, et génua debília roboráte.
Dícite pusillánimis: Confortámini, et nolíte timére: ecce Deus vester ultiónem addúcet retributiónis: Deus ipse véniet, et salvábit vos.
Tunc aperiéntur óculi cæcórum, et aures surdórum patébunt.
Tunc sáliet sicut cervus claudus, et apérta erit lingua mutórum: quia scissæ sunt in desérto aquæ, et torréntes in solitúdine.
Et quæ erat árida, erit in stagnum, et sítiens in fontes aquárum. Hæc dicit Dóminus.

\begin{tightlist}{}
\item \textred{℟.} María cum audísset, turbáta est in sermóne Angeli: \textred{*} Et cogitábat qualis esset ista salutátio. \textred{℣.} Quómodo fiet istud, quóniam virum non cognósco? \textred{E}t. Glória Patri. \textred{E}t.
\end{tightlist}

\subsubsection*{\center \textsc{IN SECUNDO NOCTURNO.}}
\begin{tightlist}{}
\item \textred{Aña.} Ne tímeas María, invenísti enim grátiam apud Deum: ecce concípies in útero, et páries fílium.
%\item \textred{Ps.} Eructávit cor. [169]
\item \textred{Aña.} Dabit illi Dóminus Deus sedem David patris ejus: et regnábit in ætérnum.
\item \textred{Aña.} Super sólium David, et super regnum ejus sedébit in ætérnum.
\item \textred{℣.} Post partum, Virgo, invioláta permansísti. \textred{℟.} Dei Génitrix, intercéde pro nobis.
\item {\small \textred{Bened.} Précibus suæ Matris benedícat nos Fílius Dei Patris.}
\end{tightlist}

\subsubsection*{\center Sermo sancti Leónis Papæ.}
\begin{center}\textred{Serm. 2. de Nat. Dom.\\
\textbf{Lectio iv}.}\end{center}
\lettrineThree{D}eus omnípotens et clemens, cujus natúra bónitas, cujus volúntas poténtia, cujus opus misericórdia est; statim ut nos diabólica malígnitas venéno suæ mortificávit invídiæ, prædestináta renovándis mortálibus suæ pietátis remédia, inter ipsa mundi primórdia præsignávit, denúntians serpénti futúrum semen mulíeris, quod nóxii cápitis elatiónem sua virtúte contéreret, Christum scílicet in carne ventúrum: Deum hominémque signans, qui natus ex Vírgine, violatórem humánæ propáginis, incorrúpta nativitáte damnáret.

\begin{tightlist}{}
\item \textred{℟.} Ait Angelus Maríæ: Ne tímeas, \textred{*} Invenísti enim grátiam apud Deum. \textred{℣.} Ecce concípies in útero, et páries fílium, et vocábis nomen ejus Jesum. \textred{I}nvenísti.
\item {\small \textred{Bened.} Jesus Maríæ Fílius sit nobis clemens et propítius.}
\end{tightlist}

\subsubsection*{\center \textred{Lectio v.}}
\lettrine{N}am quia gloriabátur diábolus, hóminem sua fraude decéptum divínis caruísse munéribus, et immortalitátis dote nudátum, dura mortis subiísse senténtiam, seque in malis suis quoddam de prævaricatóris consórtio invenísse solátium: Deum quoque, justæ severitátis exigénte ratióne, erga hóminem, quem in tanto honóre condíderat, antíquam mutásse senténtiam: opus fuit, dilectíssimi, secréti dispensatióne consílii, ut incommutábilis Deus (cujus volúntas non potest sua benignitáte privári) primam pietátis suæ dispositiónem sacraménto occultióre compléret: et homo diabólicæ iniquitátis versútia actus in culpam, contra Dei propósitum non períret.

\begin{tightlist}{}
\item \textred{℟.} Ecce concípies in útero, et páries fílium, \textred{*} Et vocábis nomen ejus Jesum. \textred{℣.} Hic erit magnus, et Fílius Altíssimi vocábitur, et dabit illi Dóminus Deus sedem David patris ejus. \textred{E}t.
\item {\small \textred{Bened.} Fílius Vírginis Maríæ det nobis gáudia vitæ.}
\end{tightlist}

\subsubsection*{\center \textred{Lectio vj.}}
\lettrine{A}dveniéntibus ergo tempóribus, dilectíssimi, quæ redemptióni hóminum fúerant præstitúta, ingréditur hæc ínfima Jesus Christus Dóminus noster, de cœlésti sede descéndens, et a patérna glória non recédens, novo órdine, nova nativitáte generátus: novo órdine, quia invisíbilis in suis, visíbilis factus est in nostris, incomprehensíbilis vóluit comprehéndi: ante témpora manens, esse cœpit ex témpore: universitátis Dóminus, servílem formam obumbráta majestátis suæ dignitáte suscépit; impassíbilis Deus, non dedignátus est homo esse passíbilis; et immortális, mortis légibus subjacére.

\begin{tightlist}{}
\item \textred{℟.} Dixit María ad Angelum. Quómodo fiet istud, quóniam virum non cognósco? Et respóndens Angelus dixit ei: \textred{*} Spíritus sanctus supervéniet in te, et virtus Altíssimi obumbrábit tibi. \textred{℣.} Ideóque et quod nascétur ex te Sanctum, vocábitur Fílius Dei. \textred{S}píritus. Glória Patri. \textred{S}píritus.
\end{tightlist}

\subsubsection*{\center \textsc{IN TERTIO NOCTURNO.}}
\begin{tightlist}{}
\item \textred{Aña.} Spíritus Sanctus in te descéndet María, ne tímeas habébis in útero Fílium Dei.
%\item \textred{Ps.} Cantáte. [172]
\item \textred{Aña.} Beáta es María, quæ credidísti: quóniam perficiéntur ea, quæ dicta sunt tibi a Dómino.
\item \textred{Aña.} Beátam me dicent omnes generatiónes: quia respéxit Dóminus humilitátem ancíllæ suæ.
\item \textred{℣.} Speciósa facta es et suávis. \textred{℟.} In delíciis tuis, sancta Dei Génitrix.
\item {\small \textred{Bened.} Fons Evangélii répleat nos dógmate cæli.}
\end{tightlist}

\subsubsection*{\center Léctio sancti Evangélii secúndum Lucam.}
\begin{center}\textred{Lecio vij. \hfill Cap. 1.}\end{center}
\lettrine{I}n illo témpore: Missus est Angelus Gábriel a Deo in civitátem Galilǽæ, cui nomen Názareth, ad Vírginem desponsátam viro, cui nomen erat Joseph, de domo David, et nomen Vírginis María. Et réliqua.

\begin{center}Homilía sancti Ambrósii Epíscopi.\\
\textred{Lib. 2. in Lucam.}\end{center}
\lettrineThree{L}atent quidem divína mystéria: nec fácile, juxta prophéticum dictum, quisquam hóminum potest scire consílium Dei: sed tamen ex céteris factis, atque præcéptis Dómini Salvatóris póssumus intellígere, et hoc propensióris fuísse consílii, quod ea potíssimum elécta est, ut Dóminum páreret, quæ erat desponsáta viro. Cur autem non ántequam desponsarétur, impléta est? Fortásse ne dicerétur, quod concéperat ex adultério.

\begin{tightlist}{}
\item \textred{℟.} Dabit illi Dóminus Deus sedem David patris ejus: \textred{*} Et regnábit in domo Jacob in ætérnum. \textred{℣.} Et regni ejus non erit finis. \textred{E}t.
\item {\small \textred{Bened.} Ab hoste malígno deféndat nos Dei Génitrix Virgo.}
\item \textred{Quando de hac Homilia debent legi tres Lectiones, sequens Lectio dividatur in octavam, et nonam.}
\end{tightlist}

\subsubsection*{\center \textred{Lectio viij.}}
\lettrine{E}t ingréssus ad eam Angelus. Disce vírginem móribus, disce vírginem verecúndia, disce oráculo, disce mystério. Trepidáre vírginum est, et ad omnes viri ingréssus pavére, omnes viri affátus veréri. Discant mulíeres propósitum pudóris imitári. Sola in penetrálibus, quam nemo virórum víderit, solus Angelus repérerit; sola sine cómite, sola sine teste, ne quo degénere depravarétur affátu, ab Angelo salutátur.

¶ Tanti namque mandáti mystérium non hóminis fuit, sed Angeli ore proméndum. Hódie primum audítur: Spíritus sanctus supervéniet in te. Et audítur, et créditur. Dénique: Ecce, inquit, ancílla Dómini, contíngat mihi secúndum verbum tuum. Vide humilitátem, vide devotiónem! Ancíllam se dicit Dómini, quæ mater elígitur: nec repentíno exaltáta promísso est.

\begin{tightlist}{}
\item \textred{℟.} Dixit autem María ad Angelum: Ecce ancílla Dómini: \textred{*} Fiat mihi secúndum verbum tuum. \textred{℣.} Spíritus Sanctus supervéniet in te, et virtus Altíssimi obumbrábit tibi. \textred{F}iat.
\item \textred{Si ix Lectio legenda sit de alia Homilia, Bened.} Verba Redemptóris pia sint suffrágia nobis.
\item \textred{Alias, Bened.} In omni tribulatióne et angústia succúrrat nobis pia Virgo María.
\item \textred{℟. ix.} Gaude María Virgo, cunctas hǽreses sola interemísti, quæ Gabriélis Archángeli dictis credidísti: \textred{*} Dum Virgo Deum et hóminem genuísti, et post partum Virgo invioláta permansísti. \textred{℣.} Gabriélem Archángelum scimus divínitus te esse affátum: úterum tuum de Spíritu sancto crédimus fecundátum. Erubéscat Judǽus infélix, qui dicit Christum ex Joseph sémine esse natum. \textred{D}um. Glória Patri. \textred{D}um.
\item Te Deum laudámus. 6.
\item \textred{Sacerd. ℣.} Ora pro nobis, sancta Dei Génitrix. \textred{℟.} Ut digni efficiámur promissiónibus Christi.
\end{tightlist}

\subsubsection*{\center \textsc{AD LAUDES.}}
\begin{center}\textred{et per Horas, Aña.}\end{center}
\lettrine{M}issus est Angelus Gábriel a Deo ad Maríam Vírginem desponsátam Joseph.

\begin{tightlist}{}
\item \textred{Ps. de Dominica.} 30.
\item \textred{Aña.} Ave María, grátia plena: Dóminus tecum: Benedícta tu in muliéribus.
\item \textred{Aña.} Ne tímeas María, invenísti enim grátiam apud Deum: ecce concípies in útero, et páries fílium.
\item \textred{Aña.} Dabit illi Dóminus Deus sedem David patris ejus, et regnábit in ætérnum.
\item \textred{Aña.} Ecce ancílla Dómini, fiat mihi secúndum verbum tuum.
\end{tightlist}

\subsubsection*{\textred{Capitulum. \hfill Isa. 7.}}
\lettrineThree{E}cce Virgo concípiet, et páriet fílium, et vocábitur nomen ejus Emmánuel. Butýrum et mel cómedet, ut sciat reprobáre malum, et elígere bonum.

\begin{tightlist}{}
\item \textred{Hymn.} O gloriósa Dómina, \textred{de communi.} [176]
\item \textred{℣.} Benedícta tu in muliéribus. \textred{℟.} Et benedíctus fructus ventris tui.
\item \textred{Ad Bened. Aña.} Quómodo fiet istud, Angele Dei, quóniam virum non cognósco? Audi María Virgo: Spíritus Sanctus supervéniet in te, et virtus Altíssimi obumbrábit tibi.
\end{tightlist}

\subsubsection*{\center \textred{Oratio.}}
\lettrine{D}eus, qui de beátæ Maríæ Vírginis útero, Verbum tuum, Angelo nuntiánte, carnem suscípere voluísti: præsta supplícibus tuis; ut qui vere eam Genitrícem Dei crédimus, ejus apud te intercessiónibus adjuvémur. Per eúmdem Dóminum.

\begin{tightlist}{}
\item \textred{Ad Horas Añæ de Laudibus, Psalmi de Dominica, ad Primam tamen ut in Festis.}
\end{tightlist}

\subsubsection*{\center \textsc{AD PRIMAM.}}
\noindent \textred{℣.} Qui de Vírgine dignátus es nasci.

\subsubsection*{\center \textsc{AD TERTIAM.}}
\subsubsection*{\textred{Capitulum. \hfill Isa. 7.}}
\lettrine{E}cce Virgo concípiet, et páriet fílium, et vocábitur nomen ejus Emmánuel. Butýrum et mel cómedet, ut sciat reprobáre malum, et elígere bonum.

\begin{tightlist}{}
\item \textred{℟.} Sancta Dei Génitrix, \textred{*} Virgo semper María. \textred{S}ancta. \textred{℣.} Intercéde pro nobis ad Dóminum Deum nostrum. \textred{V}irgo. Glória Patri. \textred{S}ancta.
\item \textred{℣.} Post partum Virgo invioláta permansísti. \textred{℟.} Dei Génitrix intercéde pro nobis.
\end{tightlist}

\subsubsection*{\center \textsc{AD SEXTAM.}}
\subsubsection*{\textred{Capitulum. \hfill Isa. 11.}}
\lettrine{E}grediétur virga de radíce Jesse, et flos de radíce ejus ascéndet. Et requiéscet super eum Spíritus Dómini.

\begin{tightlist}{}
\item \textred{℟.} Post partum, Virgo \textred{*} Invioláta permansísti. \textred{P}ost. \textred{℣.} Dei Génitrix intercéde pro nobis. \textred{I}nvioláta. Glória Patri. \textred{P}ost.
\item \textred{℣.} Speciósa facta es, et suávis. \textred{℟.} In delíciis tuis, sancta Dei Génitrix.
\end{tightlist}

\subsubsection*{\center \textsc{AD NONAM.}}
\subsubsection*{\textred{Capitulum. \hfill Isa. 45.}}
\lettrine{R}oráte cœli désuper, et nubes pluant justum: aperiátur terra, et gérminet Salvatórem: et justítia oriátur simul: ego Dóminus creávi eum.

\begin{tightlist}{}
\item \textred{℟.} Speciósa facta es, \textred{*} Et suávis. \textred{S}peciósa. \textred{℣.} In delíciis tuis, sancta Dei Génitrix. \textred{E}t. Glória Patri. \textred{S}peciósa.
\item \textred{℣.} Elégit eam Deus, et præelégit eam. \textred{℟.} Et habitáre eam facit in tabernáculo suo.
\end{tightlist}

\subsubsection*{\center \textsc{IN II. VESPERAS.}}
\begin{tightlist}{}
\item \textred{Añæ. Laudum.}
\item \textred{Ps.} Dixit Dóminus, \textred{cum reliquis de Comm. B.V.M.} [177]
%\item \textred{Capit. Hymn. et ℣. ut in j. Vesp.}
\end{tightlist}

\subsubsection*{\textred{Capitulum. \hfill Isaiæ 7.}}
\lettrineThree{E}cce Virgo concípiet et páriet fílium, et vocábitur nomen ejus Emmánuel. Butyrum, et mel cómedet, ut sciat reprobáre malum, et elígere bonum.

\begin{tightlist}{}
%\item \textred{Si fuerit Duplex,}
%\item \textred{℟.} Gaude, María Virgo, cunctas hǽreses sola interemísti, quæ Gabriélis Archángeli dictis credidísti: * Dum Virgo Deum et hóminem genuísti, et post partum Virgo invioláta permansísti. \textred{℣.} Gabriélem Archángelum scimus divínitus te esse affátum: úterum tuum de Spíritu sancto crédimus fecundátum. Erubéscat Judǽus infélix, qui dicit Christum ex Joseph sémine esse natum. \textred{D}um. Glória Patri. \textred{D}um.
\item \textred{Hymn.} Ave maris stella. [180]
\item \textred{℣.} Ave María grátia plena. \textred{℟.} Dóminus tecum.

\item \textred{Ad Magnif. Aña.} O Virgo vírginum, quómodo fiet istud, quia nec primam símilem visa es, nec habére sequéntem? Fíliæ Jerúsalem quid me admirámini? divínum est mystérium hoc, quod cérnitis.
\end{tightlist}

\subsubsection*{\center \textred{Oratio.}}
\lettrine{D}eus, qui de beátæ Maríæ Vírginis útero, Verbum tuum, Angelo nuntiánte, carnem suscípere voluísti: præsta supplícibus tuis; ut qui vere eam Genitrícem Dei crédimus, ejus apud te intercessiónibus adjuvémur. Per eúmdem Dóminum.

\begin{center}\textcolor{red}{\rule{\linewidth}{0.5pt}} \end{center}

\subsection*{\center \textred{\textsc{A SEPTUAGESIMA}}}
\begin{center}\textred{usque ad Pascha.}\end{center}

\subsubsection*{\center \textsc{AD VESPERAS.}}
\begin{center}\textred{Si fuerit Duplex,}\end{center}
\begin{tightlist}{}
\item \textred{Aña.} O vos omnes, qui transítis per viam, atténdite et vidéte, si est dolor sicut dolor meus.
\item \textred{Psalmi de Officio præc.}
\end{tightlist}
%\fancyhead[C]{\textsc{Sabbato - B.V.M. in Septuagesima}}
\fancyhead[C]{\textsc{Officia Hebd. - Sabbato B.V.M. - Septuagesima}}
\subsubsection*{\textred{Capitulum. \hfill Thren. 1.}}
\lettrineThree{V}ide, Dómine, quóniam tríbulor, conturbátus est venter meus: subvérsum est cor meum in memetípsa, quóniam amaritúdine plena sum.

\begin{tightlist}{}
\item \textred{Si fuerit Duplex,}
\item \textred{℟.} Vúlnera Christi moriéntis erant vúlnera matris doléntis; dolóres sævi fuérunt tortóres in ánima matris, \textred{*} Mater erat laniáta morte cari pígnoris.
\item \textred{℣.} Mente Mater erat percússa cúspide teli, quo membra Christi servi fodérunt iníqui. \textred{M}ater. Glória Patri. \textred{M}ater.
\end{tightlist}

\subsubsection*{\center \textred{Hymnus.}}
\lettrine{L}acrymæ Matris prófluunt,\\
Cordis erúmpunt gémitus,\\
Dum crucis orbis Cónditor\\
Suspénsus est patíbulo.\\
\\
Dum tot tuétur vúlnera,\\
Rivos fluéntis sánguinis,\\
Fossúmque pectus láncea,\\
Suo crucem fert péctore.\\
\\
Dolóris ensem séntiunt\\
Piæ paréntis víscera,\\
Cum Jesus dire móritur,\\
Pro reis factus hóstia.\\
\\
Illi tantum compátitur,\\
Quod secum fere móritur,\\
Nullus dolor amárior,\\
Nam nulla proles cárior.\\
\\
O María, mater pia,\\
Spes única fidélium,\\
Tua prece mellíflua,\\
Natum redde propítium.\\
\\
Quo mortis suæ mérito\\
In hujus vitæ stádio\\
Nos servet, nosque líberet\\
A gehénnæ incéndio.\\
\\
Glória tibi, Dómine,\\
Qui natus es de Vírgine,\\
Cum Patre et sancto Spíritu,\\
In sempitérna sǽcula.\\
Amen.

\begin{tightlist}{}
\item \textred{Sic terminantur omnes Hymni.}
\item \textred{℣.} Sancta Dei Génitrix, Virgo dulcis atque pia. \textred{℟.} Christum morti tráditum pro nobis semper exóra.
\item \textred{Ad Magnificat, Aña.} Jesu nate de Vírgine sine viríli sémine, qui mortem crucis páteris, suffóssus plagis ásperis: amóre matris ánxiæ nostræ parcas misériæ; ac ejus piis précibus cœli fruámur gáudiis.
\end{tightlist}

\subsubsection*{\center \textred{Oratio.}}
\lettrine{I}ntervéniat pro nobis, quǽsumus, Dómine Jesu Christe, apud tuam cleméntiam, nunc et in hora mortis nostræ, piíssima Virgo María mater tua: cujus sacratíssimam ánimam in hora passiónis tuæ dolóris gládius pertransívit, et in gloriósa resurrectióne tua ingens gáudium lætificávit. Qui vivis et regnas.

\begin{tightlist}{}
\item \textred{Ad Nunc dimittis, Aña.} Sub tuum præsídium.
\end{tightlist}

\subsubsection*{\center \textsc{AD MATUTINUM.}}
\begin{tightlist}{}
\item \textred{Invitatorium.} Jesum pro nobis in cruce mórtuum, a Vírgine Matre lamentátum, * Veníte adorémus.
\item \textred{Ps.} Venite, exsultémus. 2.
\end{tightlist}

\subsubsection*{\center \textred{Hymnus.}}
\lettrine{M}entes pungant fidélium\\
Dolóres et suspíria\\
Matris Christi, dum víscera\\
Pœna torquéntur áspera.\\
\\
Dum stat crucis sub stípite,\\
Videns pendéntem Fílium,\\
Præ spíritus angústia\\
Est liquefácta pénitus.\\
\\
Exors viri quæ cóncipit,\\
Expers dolóris péperit:\\
Mœrens dolóres pértulit,\\
Quos páriens non sénserat.\\
\\
Hoc ille justus Símeon\\
In spíritu prædíxerat,\\
Ampléxibus dum párvulum\\
Fovens piis adstríngeret.\\
\\
Auxit dolóres Fílii\\
Matris dolor, ut ánxiæ\\
Matris dolor dolóribus\\
Sic ingravéscit Fílii.\\
\\
Qui Maríam dilígitis,\\
Vos ejus et suspíria\\
Et gémitus recólite,\\
Et dira Christi vúlnera.\\
\\
Glória tibi, Dómine,\\
Qui natus es de Vírgine,\\
Cum Patre et Sancto Spíritu,\\
In sempitérna sǽcula.\\
Amen.

\subsubsection*{\center \textsc{IN PRIMO NOCTURNO.}}
\begin{tightlist}{}
\item \textred{Aña.} Surgam et circuíbo civitátem, per vicos et platéas quæram quem díligit ánima mea.
\item \textred{Ps.} Dómine, Dóminus, \textred{cum reliquis de Comm. BMV.} [165]
\item \textred{Aña.} Quæsívi illum, et non invéni, vocávi, et non respóndit mihi.
\item \textred{Aña.} Puer non compáret, et ego quo ibo?
\item \textred{℣.} Sancta Dei Génitrix, Virgo dulcis atque pia. \textred{℟.} Christum morti tráditum, pro nobis semper exóra.
\item \textred{Ante Quadragesimam Lectiones primi Nocturni de Scriptura occurrente, sed in Quadragesima sequentes.}
\item {\small \textred{Bened.} Alma Virgo vírginum intercédat pro nobis ad Dóminum.}
\end{tightlist}

\subsubsection*{\center De Jeremía Prophéta.}
\begin{center}\textred{\textbf{Lectio j}.\\ Thren. 1. V. 20.}\end{center}
\lettrineThree{V}ide Dómine quóniam tríbulor, conturbátus est venter meus: subvérsum est cor meum in memetípsa, quóniam amaritúdine plena sum. Foris intérficit gládius, et domi mors símilis est.
Audiérunt quia ingemísco ego, et non est qui consolétur me. Omnes inimíci mei audiérunt malum meum, lætáti sunt, quóniam tu fecísti: Adduxísti diem consolatiónis, et fient símiles mei. Hæc dicit Dóminus Deus: Convertímini ad me, et salvi éritis.

\begin{tightlist}{}
\item \textred{℟.} Adspiciébat dulcem natum Virgo María in cruce pendéntem, et morte turpíssima moriéntem \textred{*} Tantóque dolóre et tristítia vexabátur in mente, ut non possit explicári sermóne. \textred{℣.} Erat enim Jesus adspéctu dulcis, collóquio suávis, et omni conversatióne benigníssimus. \textred{T}antóque.
\item {\small \textred{Bened.} Nos cum prole pia benedícat Virgo María.}
\end{tightlist}

\subsubsection*{\center \textred{Lectio ij. \hfill Cap. 3. V. 48.}}
\lettrine{D}ivisiónes aquárum dedúxit óculus meus, in contritióne fíliæ pópuli mei.
Oculus meus afflíctus est, nec tácuit, eo quod non esset réquies,
donec respíceret et vidéret Dóminus de cælis.
Oculus meus deprædátus est ánimam meam in cunctis filiábus urbis meæ.
Venatióne cepérunt me quasi avem inimíci mei gratis.
Lapsa est in lacum vita mea, et posuérunt lápidem super me.
Inundavérunt aquæ super caput meum: dixi: Périi.
Invocávi nomen tuum Dómine de lacu novíssimo.
Vocem meam audísti: ne avértas aurem tuam a singúltu meo, et clamóribus.
Appropinquásti in die, quando invocávi te: dixísti: ne tímeas. Hæc dicit Dóminus.

\begin{tightlist}{}
\item \textred{℟.} Manábat sanguis ejus ex quátuor pártibus rigántibus undis, ligno mánibus pedibúsque confíxis: de vultu illíus pulchritúdo efflúxerat omnis, \textred{*} Et qui erat speciósus forma præ fíliis hóminum, videbátur ómnium indecórus. \textred{℣.} Non erat spécies ei, neque decor; vídimus eum, et non erat adspéctus. \textred{E}t.
\item {\small \textred{Bened.} Sancta Dei Génitrix sit nobis auxiliátrix.}
\end{tightlist}

\subsubsection*{\center \textred{Lectio iij.}}
\lettrine{J}udicásti Dómine causam ánimæ meæ, redémptor vitæ meæ.
Vidísti Dómine iniquitátem illórum advérsum me: júdica judícium meum.
Vidísti omnem furórem, univérsas cogitatiónes eórum advérsum me.
Audísti oppróbrium eórum Dómine, omnes cogitatiónes eórum advérsum me,
lábia insurgéntium mihi; et meditatiónes eórum advérsum me tota die.
Sessiónem eórum, et resurrectiónem eórum vide: ego sum psalmus eórum.
Reddes eis vicem, Dómine, juxta ópera mánuum suárum.
Dabis eis scutum cordis labórem tuum.
Persequéris in furóre, et cónteres eos sub cælis, Dómine. Hæc dicit Dóminus.

\begin{tightlist}{}
\item \textred{℟.} Iste erat máximus dolor Maríæ, quia vidébat se relínqui ab eo quem genúerat, nec supérerat álius, quia illi erat únicus: \textred{*} Vox triste sonábat foris vulnus denúntians mentis. \textred{℣.} Amor únice, fili dulcíssime, quis dabit mihi, ut pro te móriar, fili mi? \textred{V}ox. Glória Patri. \textred{V}ox.
\end{tightlist}

\subsubsection*{\center \textsc{IN SECUNDO NOCTURNO.}}
\begin{tightlist}{}
\item \textred{Aña.} Adjúro vos, fíliæ Jerúsalem, si invenéritis diléctum meum, ut nuntiétis ei quia amóre lángueo.
%\item \textred{Ps.} Eructávit. [169]
\item \textred{Aña.} Qualis est diléctus tuus ex dilécto, o pulchérrima mulíerum, quia sic adjurásti nos?
\item \textred{Aña.} Diléctus meus cándidus et rubicúndus, eléctus ex míllibus.
\item \textred{℣.} Virgo mater Christi juxta Crucem adstitísti. \textred{℟.} Præparáta mori propter amórem Fílii.
\item {\small \textred{Bened.} Précibus suæ Matris benedícat nos Fílius Dei Patris.}
\end{tightlist}

\subsubsection*{\center Sermo sancti Bernárdi Abbátis.}
\begin{center}\textred{De Lament. B. Mariæ.\\
\textbf{Lectio iv}.}\end{center}
\lettrineThree{J}uxta crucem Christi stabat emórtua María mater, quæ ipsum ex Spíritu sancto concépit. Vox illi non erat, quia dolóre attríta jacens pallébat. Quasi mórtua vivens, vivébat móriens; moriebátur vivens, nec mori póterat, quæ vivens mórtua erat. In illíus ánima dolor sæve sæviébat: optábat mori magis quam vívere post mortem Christi, quæ male vivens mórtua erat. Ibi stabat dolens, sævo dolóre confécta. O verum elóquium justi Simeónis, quem prædíxit, gládium sentiébat dolóris.

\begin{tightlist}{}
\item \textred{℟.} Nunc solvis Virgo cum usúra, quod in partu non habuísti a natúra: \textred{*} Dolórem pariéndo Fílium non sensísti, quem míllies replicátum Fílio moriénte passa fuísti. \textred{℣.} Juxta Crucem Christi stabat velut emórtua mater, quæ ipsum ex Spíritu sancto concépit. \textred{D}olórem.
\item {\small \textred{Bened.} Jesus Maríæ Fílius sit nobis clemens et propítius.}
\end{tightlist}

\subsubsection*{\center \textred{Lectio v.}}
\lettrine{S}tabat juxta crucem María, íntuens vultu benígno Christum pendéntem in patíbulo, pedúmque summitátibus innítens, manus levábat in altum, ampléctens rubricátam crucem, ac in óscula ejus ruens ea parte, qua unda pretiosíssimi sánguinis defluébat: sursum manus, nisu quo póterat, extendébat, únicum suum amplécti desíderans, nec valébat. Sperat enim amor multa, quæ nunquam, vel raro fíeri possunt: impátiens síquidem amor credit quod sibi débeant cédere univérsa. Volébat amplécti Christum in alto pendéntem, sed manus frustra proténsæ, in se complósæ, compléxæ redíbant.

\begin{tightlist}{}
\item \textred{℟.} Vox illi non erat, quia dolóre attríta intus conflictabátur, quasi mórtua vivens, vivébat móriens, \textred{*} Nec mori póterat, quæ vivens mórtua erat. \textred{℣.} Optábat mori magis quam vívere post mortem Christi, in cujus ánima dolor sæve sæviébat. \textred{N}ec.
\item {\small \textred{Bened.} Fílius Vírginis Maríæ det nobis gáudia vitæ.}
\end{tightlist}

\subsubsection*{\center \textred{Lectio vj.}}
\lettrine{S}ed cum de cruce corpus ejus fuísset totáliter depósitum, præ dolóris veheméntia, et amóris immensitáte, quasi exánimis facta fuit. Stabat ad caput exstíncti fílii mœstíssima mater, et ejus regálem fáciem mortis obfuscátam pallóribus, magna rigábat affluéntia lacrymárum. Adspiciébat illud reverendíssimum caput corónæ spíneæ diris acúleis perforátum, manus illas et pedes sanctos clavis férreis crudelíssime perforátos, latúsque perfóssum láncea, cum céteris membris lacerátis. Verúmtamen rectum erat amóris et mœróris cóntinens modum. Non desperábat, sed pie et juste dolébat, sperans tamen fírmiter, fortitérque tenens ipsum, secúndum promíssum suum, tértia die resúrgere.

\begin{tightlist}{}
\item \textred{℟.} O verum elóquium justi Simeónis, quem prædíxit gládium sentiébat dolóris; cum juxta crucem staret, \textred{*} Intuens vultu benígno Christum pendéntem in ligno. \textred{℣.} Ibi stabat dolens sævo dolóre confécta. \textred{I}ntuens. Glória Patri. \textred{I}ntuens.
\end{tightlist}

\subsubsection*{\center \textsc{IN TERTIO NOCTURNO.}}
\begin{tightlist}{}
\item \textred{Aña.} Jesu Fili mi, Fili mi Jesu, quis mihi tríbuat, ut pro te móriar?
%\item \textred{Ps.} Cantáte. [172]
\item \textred{Aña.} Amor únice, Fili dulcíssime, noli me derelínquere post te, trahe me ad te ipsum, ut ipsa móriar tecum.
\item \textred{Aña.} In tuo me súscipe patíbulo, ut qui una carne vívimus, una morte pereámus.
\item \textred{℣.} Virgo tuum Fílium in cruce clavátum. \textred{℟.} Redde tuis précibus míseris placátum. 
\item {\small \textred{Bened.} Fons Evangélii répleat nos dógmate cæli.}
\end{tightlist}

\subsubsection*{\center Lectio sancti Evangélii secúndum Lucam.}
\begin{center}\textred{\textbf{Lectio vij}. \hfill Cap. 2. ℣. 33.}\end{center}
\lettrine{I}n illo témpore: Erat Pater Jesu et Mater mirántes super his, quæ dicebántur de illo. Et réliqua.

\begin{center}Homilía sancti Ambrósii Epíscopi.\\
\textred{Libro 2. in Luc.}\end{center}
\lettrineThree{V}ides úberem in omnes grátiam Dómini generatióne diffúsam, et prophetíam incrédulis negátam esse, non justis. Ecce et Símeon prophétat in ruínam et resurrectiónem plurimórum venísse Dóminum Jesum Christum, ut justórum iniquorúmque mérita discérnat; et pro nostrórum qualitáte factórum, judex verus et justus, aut supplícia decérnat, aut prǽmia.

\begin{tightlist}{}
\item \textred{℟.} Quis æstimáre potest quanta tristítia, mœróre et dolóre vulnerátum fúerit cor matérnum, super únici Fílii passióne? \textred{*} Quæ tam caste concépit, tam lætánter péperit, tam dúlciter áluit. \textred{℣.} Vidit eum ligátum, et sólvere non pótuit: vidit corpus vulnerátum, et vúlnera ligáre non pótuit. \textred{Q}uæ.
\item {\small \textred{Bened.} Ab hoste malígno deféndat nos Dei Génitrix Virgo.}
\item \textred{Extra Quadrag. et Vigiliam \textbf{Lectio viij}. dividatur in duas,}
\item {\small \textred{et pro ix. Bened.} In omni tribulatióne et angústia succúrrat nobis pia Virgo María.}
\end{tightlist}

\subsubsection*{\center \textred{Lectio viij.}}
\lettrine{E}t tuam, inquit, ipsíus ánimam pertransíbit gládius. Nec líttera, nec história docet, ex hac vita Maríam corporális necis passióne migrásse. Non enim ánima, sed corpus materiáli gládio transverberátur.

¶ Et ídeo prudéntiam Maríæ haud ignáram mystérii cœléstis osténdit: vivum enim verbum Dei et válidum, et acútius omni gládio acutíssimo, pénetrans usque ad divisiónem ánimæ et spíritus, ártuum quoque et medullárum, cogitatiónes cordis, et secréta scrutátur animórum, quia nuda et apérta sunt ómnia Dei Fílio, quem consciéntiæ secréta non fallunt.

\begin{tightlist}{}
\item \textred{℟.} Fóderis, o dulcis Jesu, diris clavis extra; sed aspérius intra, ex compassióne Vírginis matris, tecum ómnia participántis. \textred{*} O mirábilis censúræ condítio! peccávit iníquus, et vápulat justus. \textred{℣.} Delínquit homo reus, et occíditur pro eo homo Deus. \textred{O}.
\item \textred{In Quadrag. vel Vigilia, ix. Lect. et comm. de ea.}
\item {\small \textred{Bened.} Verba Redemptóris pia sint suffrágia nobis.}
\item \textred{℟. ix.} O dulcíssime Jesu, o benigníssima María Virgo, quid commisístis, quod tantis cruciátibus, tot et tantis afflictiónibus addícti estis? \textred{*} Nos sumus qui peccávimus, qui injúste égimus; et tu innocentíssima Virgo, cum único nato pro nobis dolétis et affligímini. \textred{℣.} O peccátor, vide quómodo patiúntur, et flent pro te Jesus et María. \textred{N}os. Glória Patri. \textred{N}os.
\item Te Deum laudámus. 6.
\item \textred{Sacerdotalis ℣.} Tuam ipsíus ánimam dolóris gládius pertransíbit. \textred{℟.} Ut reveléntur ex multis córdibus cogitatiónes.
\end{tightlist}

\subsubsection*{\center \textsc{AD LAUDES.}}
\begin{center}\textred{et per Horas, Aña.}\end{center}
\lettrine{S}tabant juxta crucem Jesu mater ejus, et soror matris ejus, María Cléophæ, et María Magdaléne.

\begin{tightlist}{}
\item \textred{Ps. de Dominica.} 30.
\item \textred{Aña.} Cum vidísset ergo Jesus matrem et discípulum stantem quem diligébat, dixit matri suæ: Múlier ecce fílius tuus.
\item \textred{Aña.} Benedíxit illis Símeon, et dixit ad Maríam matrem ejus: Ecce pósitus est hic in ruínam, et in resurrectiónem multórum in Israel, et in signum, cui contradicétur.
\item \textred{Aña.} Et tuam ipsíus ánimam pertransíbit gládius, ut reveléntur ex multis córdibus cogitatiónes.
\item \textred{Aña.} O vos omnes qui transítis per viam, atténdite et vidéte, si est dolor sicut dolor meus.
\end{tightlist}

\subsubsection*{\textred{Capitulum. \hfill Thren. 1.}}
\lettrineThree{V}ide, Dómine, quóniam tríbulor, conturbátus est venter meus, subvérsum est cor meum in memetípsa, quóniam amaritúdine plena sum.

\subsubsection*{\center \textred{Hymnus.}}
\lettrine{C}ondoléntes mœróribus\\
Maríæ matris Vírginis,\\
Jungámus laudes flétibus\\
Ad laudem sui nóminis.\\
\\
Pensémus matris sérvuli\\
Spinas cordis acérrimas,\\
Dum vidit corpus Fílii\\
Pœnas pati tetérrimas.\\
\\
Pensémus quid perpénderet;\\
Adstans cruci puérpera,\\
Dum testaméntum cónderet\\
Christus in vitæ véspera.\\
\\
Matrem piam discípulo,\\
Patri comméndat spíritum:\\
Acéti sumpto póculo,\\
Carnis exsólvit débitum.\\
\\
Cuncti Maríæ sérvuli,\\
Hujus osóres sǽculi:\\
Jesum occísum cólite,\\
Et cum María plángite.\\
\\
Glória tibi Dómine,\\
Qui natus es de Vírgine,\\
Cum Patre et sancto Spíritu,\\
In sempitérna sǽcula.\\
Amen.

\begin{tightlist}{}
\item \textred{℣.} Sit nobis vita, mors Fílii Dei Patris. \textred{℟.} Prece devóta suæ piíssimæ matris.
\item \textred{Ad Benedictus, Aña.} De cruce depósitum videns corpus Christi mœsta mater lácrymans, atque vultu tristi, dixit: O dulcíssime Fili, quid fecísti, quod has pœnas ásperas et mortem subiísti?
\end{tightlist}

\subsubsection*{\center \textred{Oratio.}}
\lettrine{I}ntervéniat pro nobis, quǽsumus, Dómine Jesu Christe, apud tuam cleméntiam, nunc et in hora mortis nostræ, piíssima Virgo María mater tua: cujus sacratíssimam ánimam in hora passiónis tuæ dolóris gládius pertransívit, et in gloriósa resurrectióne tua ingens gáudium lætificávit. Qui vivis.

\subsubsection*{\center \textsc{AD PRIMAM.}}
\begin{tightlist}{}
\item \textred{℣.} Qui de Vírgine dignátus es nasci.
\end{tightlist}

\subsubsection*{\center \textsc{AD TERTIAM.}}
\subsubsection*{\textred{Capitulum. \hfill Thren. 1.}}
\lettrine{V}ide, Dómine, quóniam tríbulor, conturbátus est venter meus, subvérsum est cor meum in memetípsa, quóniam amaritúdine plena sum.

\begin{tightlist}{}
%\item \textred{Capitulum.} Vide Dómine, \textred{ut supra.}
\item \textred{℟.} Sancta Dei Génitrix, \textred{*} Virgo dulcis atque pia. \textred{S}ancta. \textred{℣.} Christum morti tráditum pro nobis semper exóra. \textred{V}irgo. Glória Patri. \textred{S}ancta.
\item \textred{℣.} Virgo mater Christi, juxta Crucem adstitísti. \textred{℟.} Præparáta mori propter amórem Fílii.
\end{tightlist}

\subsubsection*{\center \textsc{AD SEXTAM.}}
\subsubsection*{\textred{Capitulum. \hfill Ruth. 1.}}
\lettrine{N}e vocétis me Noémi, id est pulchram: sed vocáte me Mara, id est, amáram; quia amaritúdine valde me replévit Omnípotens.

\begin{tightlist}{}
\item \textred{℟.} Virgo mater Christi, \textred{*} Juxta Crucem adstitísti. \textred{V}irgo. \textred{℣.} Præparáta mori propter amórem Fílii. \textred{J}uxta. Glória Patri. \textred{V}irgo.
\item \textred{℣.} Virgo tuum Fílium in cruce clavátum. \textred{℟.} Redde tuis précibus, míseris placátum.
\end{tightlist}

\subsubsection*{\center \textsc{AD NONAM.}}
\subsubsection*{\textred{Capitulum. \hfill Thren. 1.}}
\lettrine{D}educ, quasi torréntem lácrymas per diem et noctem, neque táceat pupílla óculi tui.

\begin{tightlist}{}
\item \textred{℟.} Virgo tuum Fílium \textred{*} In cruce clavátum. Virgo. \textred{℣.} Redde tuis précibus, míseris placátum. \textred{I}n. Glória Patri. \textred{V}irgo.
\item \textred{℣.} Sit nobis vita mors Fílii Dei Patris. \textred{℟.} Prece devóta suæ piíssimæ matris.
\item \textred{Ad secundas Vesp. sola prima Aña Laudum, Ps. de Comm. B.V.M. [177]}
\item \textred{Capitul. Hymn. et ℣. de primis Vesperis.}
\item \textred{Ad Magnificat, Aña.} Ante thronum Trinitátis miserórum miseráta, pia mater pietátis, sis pro nobis advocáta: et tuo apud Fílium intervéntu véniam peccatórum servis tuis óbtine.
\item \textred{Quando dicitur pro commemoratione, ℣.} Sancta Dei Génitrix, Virgo dulcis atque pia. \textred{℟.} Christum morti tráditum pro nobis semper exóra.
\end{tightlist}

\begin{center}\textcolor{red}{\rule{\linewidth}{0.5pt}} \end{center}

\subsubsection*{\center \textsc{TEMP. PASCHALI.}}
\begin{tightlist}{}
\item \textred{Ut in Offic. per annum, sed ante ult. ℣. Hymnorum (Hymno Vesp. excepto) additur ℣.} Quǽsumus auctor ómnium, \textred{vel ℣.} Tu esto nostrum gáudium, \textred{pro Temporis exigentia, et ad Matut. tres Psalmi cujuslibet Noct. dicuntur sub j. sola Aña, et in fine ubique additur} Allelúja, \textred{ut in aliis Officiis.} 
\item \textred{Ad Nunc dimittis, Aña.} Allelúja: Sancta Dei Génitrix.
\item \textred{Lectiones j. Nocturni de Scriptura occurr. sed cum non habentur, legantur ut per annum,} Ego Sapiéntia hábito. [167]
\end{tightlist}

\subsubsection*{\center \textsc{IN SECUNDO NOCTURNO.}}
\begin{tightlist}{}
\item {\small \textred{Bened.} Précibus suæ Matris benedícat nos Fílius Dei Patris.}
\end{tightlist}
\subsubsection*{\center Sermo sancti Augustíni Epíscopi.}
\begin{center}\textred{Serm. 83. In. Append.\\
\textbf{Lectio iv}.}\end{center}
\lettrineThree{L}oquámur áliquid, fratres in láudibus sacratíssimæ Vírginis Maríæ. Sed quid nos tantílli, quid actióne pusílli, quid in ejus láudibus referémus: cum etsi ómnium nostrum membra verteréntur in linguas, eam laudáre suffíceret nullus? Altior cœlo est, de qua lóquimur: Abýsso profúndior, cui laudes dícere conámur. Hæc est enim quæ sola méruit mater et sponsa vocári. Hæc primæ matris damna resólvit. Hæc hómini pérdito redemptiónem addúxit.

\begin{tightlist}{}
\item \textred{℟.} Sicut cedrus. [171]
\item {\small \textred{Bened.} Jesus Maríæ Fílius sit nobis clemens et propítius.}
\end{tightlist}

\subsubsection*{\center \textred{Lectio v.}}
\lettrine{M}ater enim géneris nostri pœnam íntulit mundo: Génitrix Dómini nostri salútem édidit mundo: auctrix peccáti Heva, auctrix mériti María: Heva occidéndo óbfuit: María vivificándo prófuit: illa percússit, ista sanávit. Hæc enim mirábili atque inæstimábili modo ómnium rerum et suum péperit Salvatórem. Quæ hæc Virgo tam sancta, ad quam veníre dignarétur Spíritus sanctus? Quæ tam speciósa, ut eam Deus sibi elígeret Sponsam? Quæ tam casta, ut esse possit Virgo post partum? Hæc est Dei templum, fons ille signátus, et porta in domo Dei clausa. Ad hanc namque, ut dixi, Spíritus sanctus descéndit: hanc virtus Altíssimi obumbrávit.

\begin{tightlist}{}
\item \textred{℟.} Quæ est ista. [171]
\item {\small \textred{Bened.} Fílius Vírginis Maríæ det nobis gáudia vitæ.}
\end{tightlist}

\subsubsection*{\center \textred{Lectio vj.}}
\lettrine{H}æc est immaculáta cóitu, fœcúnda partu, Virgo lactans, Angelórum et hóminum cibum nútriens. Mérito namque beáta singulári a nobis præcónio extóllitur, quæ singuláre commércium mundo prǽbuit. Dénique tantum se ad cœli fastígia sublevávit, ut Verbum in princípio apud Deum de supérna cœli arce suscíperet. O felix María et omni laude digníssima! o Virgo Dei Génitrix gloriósa! o sublímis puérpera! cujus viscéribus auctor cœli terrǽque commíttitur.

\begin{tightlist}{}
\item \textred{℟.} Ornátam monílibus. [172]
\end{tightlist}

\subsubsection*{\center \textsc{IN TERTIO NOCTURNO.}}
\begin{tightlist}{}
\item {\small \textred{Bened.} Fons Evangélii répleat nos dógmate cæli.}
\end{tightlist}
\subsubsection*{\center Léctio sancti Evangélii secúndum Joánnem.}
\begin{center}\textred{\textbf{Lectio vij}. \hfill Cap. 19.}\end{center}
\lettrine{I}n illo témpore: Stabant juxta crucem Jesu mater ejus, et soror matris ejus María Cléophæ, et María Magdaléne. Et réliqua.

\begin{center}Homilía sancti Ambrósii Epíscopi.\\
\textred{Exhortat. ad Virginit.}\end{center}
\lettrineThree{H}ic Christus in vestris pectóribus alta máneat defíxus radíce, ut et fructus vester exúberet, et liquóre grátiæ spiritális, corpórei vapóris incéndia temperéntur. Hic est, qui venit in nube levi, sicut dixit Prophéta: Ecce Dóminus sedet super nubem levem, et véniet in Ægýptum. Signíficans quod in Ægýptum, id est, in afflictiónem istam mundi hujus, veníret per Vírginem. Nubem ítaque Maríam dixit, quia carnem gerébat: levem, quia Virgo erat, nullis onéribus graváta conjúgii. Ipsa est virga gérminans florem, quia pura, et ad Dóminum líbero corde dirécta virgínitas, quæ nullis in hoc sǽculo curárum anfráctibus reflectátur. Ideo illam Dóminus de cruce discípulo suo dilectíssimo sibi trádidit sancto Joánni, qui dixit patri et matri suæ, non novi vos. Dénique vocátus a Christo, relíquit patrem, Verbum secútus.

\begin{tightlist}{}
\item \textred{℟.} Beátam me dicent. [174]
\item {\small \textred{Bened.} Ab hoste malígno deféndat nos Dei Génitrix Virgo.}
\end{tightlist}

\subsubsection*{\center \textred{Lectio viij.}}
\lettrine{H}uic tráditur Virgo, qui suos nésciat. Huic tráditur Virgo, qui de péctore Christi háuriat sapiéntiam. Huic tráditur Virgo, qui fratres suos non cognóvit, et fílios suos non scivit. Ideóque benedícit eum lex: date Levi veros ejus, date Levi sortes ejus. Unde et matrem Dómini ipse suscépit; scriptum enim habémus; quia ex illa hora suscépit eam discípulus in sua. Quid est, in sua, cum relíquerit patrem et Christum secútus sit? aut quómodo in sua, cum ipsi Apóstoli díxerint: ecce nos relíquimus ómnia et secúti sumus te? Quæ sua habébat Joánnes, qui mundána et sæculária non habébat, quia de mundo non erat? quæ ergo habébat sua, nisi ea quæ a Christo accéperat? bonus verbi sapientiǽque posséssor, bonus recéptor grátiæ.

\begin{tightlist}{}
\item \textred{℟.} Felix namque. [175]
\item {\small \textred{Bened.} In omni tribulatióne et angústia succúrrat nobis pia Virgo María.}
\end{tightlist}

\subsubsection*{\center \textred{Lectio ix.}}
\lettrine{A}udíte quæ Apóstoli a Christo accéperint. Accípite, inquit Spíritum sanctum; quorum remiséritis peccáta, remíssa erunt; et quorum detinuéritis, deténta. Neque enim mater Dómini Jesu, nisi ad possessórem grátiæ demigráret, ubi Christus habébat habitáculum. Date ergo et vos fílii mei, vero illi Levi veros ejus. Estóte nubes, sed leves. Eritis profécto, si ábluat virgínitas ónera conditiónis, et ténebras, hujus limósæ carnis illúminet. Ideo dicit illa, nigra sum et decóra, fíliæ Jerúsalem. Nigra per carnem, decóra per Virginitátem.

\begin{tightlist}{}
\item \textred{℟.} Christi Virgo. [175]
\end{tightlist}

\begin{center}\textcolor{red}{\rule{\linewidth}{0.5pt}} \end{center}

\subsection*{\center \textsc{Officium Gloriosissimi Nominis beatæ Virginis Mariæ}}
\subsubsection*{\center \textred{Duplex Majus S.C.}}
\begin{center}\textred{In Festo Septembr. duplex majus S.C. quater in anno duplex minus P.C. juxta Rubr. 16.}\end{center}
\fancyhead[C]{\textsc{Quater in Anno - Officium de Nomine B.V.M.}}

%\begin{center}\textcolor{red}{\rule{\linewidth}{0.5pt}} \end{center}

\subsection*{\center \textsc{AD VESPERAS.}}
\begin{center}\textred{Antiphona.}\end{center}
\lettrine{A} solis ortu usque ad occásum, laudábile nomen Dómini, et Maríæ Matris ejus.

\begin{tightlist}{}
\item \textred{Psalmi de j. Vesperis de commune B.M.V.} [161]
\item \textred{Aña.} Missus est Angelus Gábriel ad Vírginem desponsátam Joseph, et nomen Vírginis María.
\item \textred{Aña.} Ave María, grátia plena: Dóminus tecum: Benedícta tu in muliéribus.
\item \textred{Aña.} Permáneat, et magnificétur nomen tuum usque in sempitérnum, o María.
\item \textred{Aña.} Lauda Jerúsalem Dóminum, lauda ínclytum nomen Maríæ.
\item \textred{Quando hoc Officium recitatur sub ritu duplici minori primæ classis, dicitur in primis Vesperis Aña.} Mirábiles elatiónes maris, sed mirabílior in Nómine Maríæ Dóminus. \textred{prima Aña de Laudibus, Psalmi de Officio præcedente.}
\end{tightlist}

\subsubsection*{\textred{Capitulum. \hfill Judith 13.}}
\lettrineThree{B}enedícta es tu fília a Dómino Deo excélso, præ ómnibus muliéribus super terram: quia hódie nomen tuum ita magnificávit, ut non recédat laus tua de ore hóminum, qui mémores fúerint virtútis Dómini in ætérnum.

\begin{tightlist}{}
\item \textred{℟.} Felix namque es sacra Virgo María, et omni laude digníssima: \textred{*} Quia ex te ortus est Sol justítiæ, Christus Deus noster.
\textred{℣.} Ora pro pópulo, intérveni pro clero, intercéde pro devóto femíneo sexu: séntiant omnes tuum juvámen, quicúmque célebrant tui Nóminis commemoratiónem. \textred{Q}uia. Glória Patri. \textred{Q}uia.
\item \textred{Hymnus.} Ave, maris stella. [164]
\item \textred{℣.} Sit nomen Dóminæ benedíctum. \textred{℟.} Ex hoc nunc, et usque in sǽculum, allelúja.
\item \textred{Ad Magnificat, Aña.} Maríæ Nomen cunctas illústrat Ecclésias, cui fecit magna qui potens est, et sanctum Nomen ejus.
\end{tightlist}

\subsubsection*{\center \textred{Oratio.}}
\lettrine{D}eus, qui gloriósam Matrem tuam nominári Maríam voluísti: concéde quǽsumus; ut, qui dulce ejúsdem Maríæ Nomen implórant, perpétuum séntiant tuæ benedictiónis efféctum. Qui vivis et regnas.

\begin{tightlist}{}
\item \textred{Ad Complet. et alias Horas, in fine Hymn. dicitur:} Glória tibi Dómine, Qui natus es.
\item \textred{Ad Nunc dimittis, Aña.} Allelúja: Sancta Dei Génitrix.
\item \textred{Sed post Septuagesimam, Aña.} Sub tuum præsídium.
\end{tightlist}

\subsubsection*{\center \textsc{AD MATUTINUM.}}
\begin{tightlist}{}
\item \textred{Invitatorium.} Ave María, grátia plena: * Dóminus tecum.
\item \textred{Ps.} Veníte, exsultémus. 2.
\item \textred{Hymn.} Quem terra. [165]
\end{tightlist}

\subsubsection*{\center \textsc{IN PRIMO NOCTURNO.}}
\begin{tightlist}{}
\item \textred{Aña.} In univérsa terra admirábile est Nomen tuum, María.
\item \textred{Ps.} Dómine, Dóminus noster, \textred{cum reliquis de Comm. B.V.M.} [165]
\item \textred{Aña.} In omnem terram diffúsa est laus Nóminis Maríæ.
\item \textred{Aña.} Super mária et flúmina fundávit Maríam, et dedit illi Nomen virtútum Rex glóriæ.
\item \textred{℣.} Magnificáte Maríam vírginem mecum. \textred{℟.} Et exaltémus Nomen ejus in ætérnum, allelúja.
\item {\small \textred{Bened.} Alma Virgo vírginum intercédat pro nobis ad Dóminum.}
\end{tightlist}

\subsubsection*{\center De libro Exodi.}
\noindent \textred{\textbf{Lectio j}. \hfill Cap. 15.}
\lettrineThree{S}umpsit María Prophetíssa, soror Aáron, týmpanum in manu sua, egressǽque sunt omnes mulíeres post eam cum týmpanis et choris, quibus præcinébat, dicens: Cantémus Dómino; glorióse enim magnificátus est, equum et ascensórem ejus dejécit in mare.

\begin{tightlist}{}
\item \textred{℟.} Laudémus Nomen sacrum atque insígne Maríæ: \textred{*} Mundo enim péperit Virgo salútis opem. \textred{℣.} Læténtur cœli de tali Nómine et orbis. \textred{M}undo.
\item {\small \textred{Bened.} Nos cum prole pia benedícat Virgo María.}
\end{tightlist}

\subsubsection*{\center \textred{Lectio ij.}}
\lettrine{T}ulit autem Móyses Israel de mari Rubro, et egréssi sunt in desértum Sur: ambulaverúntque tribus diébus per solitúdinem, et non inveniébant aquam. Et venérunt in Mara, nec póterant bíbere aquas de Mara, eo quod essent amáræ: unde et cóngruum loco nomen impósuit, vocans illum Mara, id est, amaritúdinem.

\begin{tightlist}{}
\item \textred{℟.} Benedíctum sit Nomen Maríæ, \textred{*} Cujus manus tornátiles plenæ hyacínthis. \textred{℣.} Implorántibus Maríæ Nomen, María præsto adest. \textred{C}ujus.
\item {\small \textred{Bened.} Sancta Dei Génitrix sit nobis auxiliátrix.}
\end{tightlist}

\subsubsection*{\center \textred{Lectio iij.}}
\lettrine{E}t murmurávit pópulus contra Móysen, dicens: Quid bibémus? At ille clamávit ad Dóminum, qui osténdit ei lignum: quod cum misísset in aquas, in dulcédinem versæ sunt. Ibi constítuit ei præcépta atque judícia, et ibi tentávit eum, dicens: Si audíeris vocem Dómini Dei tui, et quod rectum est coram eo féceris, et obedíeris mandátis ejus, custodierísque ómnia præcépta illíus, cunctum languórem, quem pósui in Ægýpto, non indúcam super te: ego enim Dóminus sanátor tuus.

\begin{tightlist}{}
\item \textred{℟.} Salutémus Maríam, pro nobis Fílium exorántem, \textred{*} Et illi dicámus: Ave María, grátia plena, Dóminus tecum. \textred{℣.} Hæc est Gedeónis vellus, quæ levíssima oratiónis compressióne, larga éxhibet gratiárum stillicídia. \textred{E}t. Glória Patri. \textred{E}t.
\end{tightlist}

\subsubsection*{\center \textsc{IN SECUNDO NOCTURNO.}}
\begin{tightlist}{}
\item \textred{Aña.} Speciósum Maríæ Nomen deprecabúntur omnes dívites plebis, cujus decórem concupívit Rex.
%\item \textred{Ps.} Eructávit. [169]
\item \textred{Aña.} Refúgium est in tribulatiónibus Maríæ Nomen; ómnibus illud invocántibus.
\item \textred{Aña.} María cívitas Dei sancta, gloriósa dicta sunt de Nómine sancto tuo.
\item \textred{℣.} Repleátur os meum laude Maríæ. \textred{℟.} Ut cantem glóriam Nómini sancto ejus, allelúja.
\item {\small \textred{Bened.} Précibus suæ Matris benedícat nos Fílius Dei Patris.}
\end{tightlist}

\subsubsection*{\center De Sermóne sancti Epiphánii Epíscopi.}
\begin{center}\textred{Orat. de Laud. B.V.M.\\
\textbf{Lectio iv}.}\end{center}
\lettrineThree{B}eáta Anna grávida effécta, cœlum, et thronum Cherúbicum péperit sanctam puéllam Maríam: illa enim reperítur esse cœlum, templum, et thronus; quia Maríam interpretári solémus Dóminam, atque étiam spem. Péperit enim Dóminum, qui est spes totíus mundi, nempe Christum. Rursum, nomen illud, María, interpretátur, myrrha maris; Myrrham vero dicet áliquis (quod ego étiam ajo) de immortalitáte, eo quod paritúra esset gemmam immortálem in mari, hoc est in mundo. Mare autem dicet univérsum mundum, cui Virgo serenitátem, et tranquillitátem cóntulit, dum portum péperit Christum.

\begin{tightlist}{}
\item \textred{℟.} Benedícta tu a Dómino Deo tuo in omni tabernáculo Jacob: \textred{*} Quóniam in omni Gente, quæ audíerit nomen tuum, magnificábitur super te Deus. \textred{℣.} Tu glória Jerúsalem, tu lætítia Israel, tu María honorificéntia pópuli nostri. \textred{Q}uóniam.
\item {\small \textred{Bened.} Jesus Maríæ Fílius sit nobis clemens et propítius.}
\end{tightlist}

\subsubsection*{\center \textred{Lectio v.}}
\lettrine{R}ursum ítaque præcláræ puéllæ Maríæ beátum nomen interpretátur Illustráta, útpote quæ illustráta sit a Fílio Dei, et cunctos usque ad términos terræ credéntes Trinitáti illumináverit. Sancta enim puélla Virgo María est Sponsa Trinitátis, et thesáurus dispensatiónis plane arcánus, cui Gábriel inquit: Ave grátia plena, Dóminus tecum. Subjúnxit Gábriel: Et Pater misit arrhabónem de cœlo Spíritum sanctum, præparávit Vírginem unigénito Fílio, cœlésti sponso, quam Pater diléxit, Fílius inhabitávit, Spíritus sanctus percupívit.

\begin{tightlist}{}
\item \textred{℟.} Mélius est nomen bonum, quam divítiæ multæ: \textred{*} Turris fortíssima Maríæ nomen. \textred{℣.} Propugnáculum est ómnibus, qui ad illud confúgiunt. \textred{T}urris.
\item {\small \textred{Bened.} Fílius Vírginis Maríæ det nobis gáudia vitæ.}
\end{tightlist}

\subsubsection*{\center \textred{Lectio vj.}}
\lettrine{I}psa est enim sponsa, et thálamus: et ex ea Sponsus procédit Christus virgíneum induméntum juxta suávem Davídis Prophétæ cantiónem: In sole pósuit tabernáculum suum, et ipse tamquam sponsus procédens de thálamo suo. O sponsi desidérium summæ puritátis! própriam ancíllam propensióre afféctu constítuit sponsam et matrem.

\begin{tightlist}{}
\item \textred{℟.} Nomen sanctum meum notum fáciam in médio pópuli mei: \textred{*} Et scient, quia ego María. \textred{℣.} Delíciæ meæ, esse cum fíliis hóminum. \textred{E}t. Glória Patri. \textred{E}t.
\end{tightlist}

\subsubsection*{\center \textsc{IN TERTIO NOCTURNO.}}
\begin{tightlist}{}
\item \textred{Aña.} Nuntiáte inter Gentes glóriam Maríæ, in ómnibus pópulis mirabília nóminis ejus.
%\item \textred{Ps.} Cantáte. [172]
\item \textred{Aña.} Annuntiavérunt cœli nomen Maríæ: et vidérunt omnes pópuli glóriam ejus.
\item \textred{Aña.} Notum fecit Dóminus sacrum Maríæ nomen, et in conspéctu Géntium revelávit virtútem nóminis ejus.
\item \textred{℣.} Ave María grátia plena. \textred{℟.} Dóminus tecum, allelúja.
\item {\small \textred{Bened.} Fons Evangélii répleat nos dógmate cæli.}
\end{tightlist}

\subsubsection*{\center Léctio sancti Evangélii secúndum Lucam.}
\textred{\textbf{Lectio vij}. \hfill Cap. 1.}\\
\lettrine{I}n illo témpore: Missus est Angelus Gábriel a Deo in civitátem Galilǽæ, cui nomen Názareth, ad Vírginem desponsátam viro, cui nomen erat Joseph, de domo David, et nomen Vírginis María. Et réliqua.

\begin{center}Homilía sancti Bernárdi Abbátis.\\
\textred{Hom. super Missus est.}\end{center}
\lettrineThree{E}t nomen, inquit, Vírginis María. Loquámur pauca et super hoc nómine, quod interpretátum, maris stella dícitur, et matri Vírgini valde conveniénter aptátur. Ipsa namque aptíssime síderi comparátur: quia sicut sine sui corruptióne sidus suum emíttit rádium, sic absque sui læsióne Virgo partúriit Fílium. Nec síderi rádius suam mínuit claritátem, nec Vírgini Fílius suam integritátem. Ipsa est ígitur nóbilis illa stella ex Jacob orta, cujus rádius univérsum orbem illúminat, cujus splendor et præfúlget in supérnis, et ínferos pénetrat, terras étiam perlústrans et calefáciens magis mentes, quam córpora: fovet virtútes, éxcoquit vítia. Ipsa, inquam, est præclára et exímia stella, super hoc mare magnum et spatiósum necessário subleváta, micans méritis, illústrans exémplis.

\begin{tightlist}{}
\item \textred{℟.} Claríssimum Maríæ Nomen ipsa stella est præfúlgens: \textred{*} Peccatóres illúminat, et dírigit in portum salútis. \textred{℣.} Consolátur tristes, míseros levat, et jacéntes érigit. \textred{P}eccatóres.
\item {\small \textred{Bened.} Ab hoste malígno deféndat nos Dei Génitrix Virgo.}
\end{tightlist}

\subsubsection*{\center \textred{Lectio viij.}}
\lettrine{O} quisquis te intélligis in hujus sǽculi proflúvio magis inter procéllas et temptestátes fluctuáre, quam per terram ambuláre, ne avértas óculos a fulgóre hujus síderis, si non vis óbrui procéllis. Si insúrgant venti tentatiónum, si incúrras scópulos tribulatiónum, réspice stellam, voca Maríam. Si jactáris supérbiæ undis, si ambitiónis, si detractiónis, si æmulatiónis, réspice stellam, voca Maríam. Si iracúndia, aut avarítia, aut carnis illécebra navículam concússerit mentis, réspice ad Maríam. Si críminum immanitáte turbátus, consciéntiæ fœditáte confúsus, júdicis horróre pertérritus, bárathro incípias absorbéri tristítiæ, desperatiónis abýsso, cógita Maríam.

\begin{tightlist}{}
\item \textred{℟.} Omnia flúmina intrant in mare, et ipsum non redúndat. \textred{*} Omnia gratiárum charísmata in Maríam, et ipsa non intuméscit. \textred{℣.} Dixit illi Angelus: Ave grátia plena. At ipsa respóndit: Ecce ancílla Dómini. \textred{O}mnia.
\item {\small \textred{Bened.} In omni tribulatióne et angústia succúrrat nobis pia Virgo María.}
\item {\small \textred{Si ix. Lectio fuerit de Hom. Evang., erit}}
\item {\small \textred{Bened.} Verba Redemptóris pia sint suffrágia nobis.}
%\item \textred{Si ix. Lectio fuerit de Hom. Evang. Benedictio.} Verba Redemptóris pia sint suffrágia nobis.
\end{tightlist}

\subsubsection*{\center \textred{Lectio ix.}}
\lettrine{I}n perículis, in angústiis, in rebus dúbiis Maríam cógita, Maríam ínvoca. Non recédat ab ore, non recédat a corde: et ut ímpetres ejus oratiónis suffrágium, non déseras conversatiónis exémplum. Ipsam sequens, non dévias: ipsam rogans, non despéras: ipsam cógitans, non erras: ipsa tenénte, non córruis: ipsa protegénte, non métuis: ipsa duce, non fatigáris: ipsa propítia, pérvenis: et sic in temetípso experíris, quam mérito dictum sit: Et nomen Vírginis María.

\begin{tightlist}{}
\item \textred{℟.} Felix namque es sacra Virgo María, et omni laude digníssima: \textred{*} Quia ex te ortus est Sol justítiæ, Christus Deus noster. \textred{℣.} Ora pro pópulo, intérveni pro clero, intercéde pro devóto femíneo sexu: séntiant omnes tuum juvámen, quicúmque célebrant tui Nóminis commemoratiónem. \textred{Q}uia. Glória Patri. \textred{Q}uia.
\item Te Deum laudámus. 6.
\item \textred{Sacerdotalis ℣.} Sit nomen Dóminæ benedíctum. \textred{℟.} Ex hoc nunc et usque in sǽculum, allelúja.
\end{tightlist}

\subsubsection*{\center \textsc{AD LAUDES.}}
\begin{center}\textred{et per Horas, Aña.}\end{center}
\lettrine{M}irábiles elatiónes maris, sed mirabílior in Nómine Maríæ Dóminus.

\begin{tightlist}{}
\item \textred{Ps. de Dominica.} 30.
\item \textred{Aña.} Laudáte in exsultatióne Maríæ Nomen, quóniam suáve est.
\item \textred{Aña.} Lábia mea laudábunt te, o María, quia dulce est Nomen tuum super mel ori meo.
\item \textred{Aña.} Spíritus et ánimæ justórum, benedícant Maríam in Nómine sancto ejus.
\item \textred{Aña.} Beneplácitum est Dómino in Nómine Maríæ, in quo exaltávit mansuétos in salútem.
\end{tightlist}

\subsubsection*{\textred{Capitulum. \hfill Judith 13.}}
\lettrineThree{B}enedícta es tu fília a Dómino Deo excélso, præ ómnibus muliéribus super terram: quia hódie nomen tuum ita magnificávit, ut non recédat laus tua de ore hóminum, qui mémores fúerint virtútis Dómini in ætérnum.

\begin{tightlist}{}
\item \textred{Hymn.} O gloriósa Dómina. [176]
\item \textred{℣.} Lætábor et exsultábo in te, María. \textred{℟.} Psalmum dicam Nómini tuo, allelúja.
\item \textred{Ad Benedictus, Aña.} Benedíctus Deus, qui Maríæ Vírgini lucis Nomen donávit, ut illumináret eos, qui in ténebris sedent: et ut stella maris errántes dirígeret in portum pacis.
\end{tightlist}

\subsubsection*{\center \textred{Oratio.}}
\lettrine{D}eus, qui gloriósam Matrem tuam nominári Maríam voluísti: concéde quǽsumus; ut, qui dulce ejúsdem Maríæ Nomen implórant, perpétuum séntiant tuæ benedictiónis efféctum. Qui vivis et regnas.

\subsubsection*{\center \textsc{AD PRIMAM.}}
\begin{tightlist}{}
\item \textred{Ad Horas Antiphonæ de Laudibus, Psalmi de Dominica, ad Primam tamen ut in Festis.}
\item \textred{℣.} Qui de Vírgine dignátus es nasci.
\end{tightlist}

\subsubsection*{\center \textsc{AD TERTIAM.}}
\subsubsection*{\textred{Capitulum. \hfill Judith 13.}}
\lettrine{B}enedícta es tu fília a Dómino Deo excélso, præ ómnibus muliéribus super terram: quia hódie nomen tuum ita magnificávit, ut non recédat laus tua de ore hóminum, qui mémores fúerint virtútis Dómini in ætérnum.

\begin{tightlist}{}
%\item \textred{Capit.} Benedícta es tu, \textred{ut supra ad Laudes.}
\item \textred{℟.} Magnificáte Maríam Vírginem mecum, \textred{*} Allelúja, allelúja. \textred{M}agnificáte. \textred{℣.} Et exaltémus Nomen ejus in ætérnum. \textred{A}llelúja, allelúja. Glória Patri. \textred{M}agnificáte.
\item \textred{℣.} Repleátur os meum laude Maríæ. \textred{℟.} Ut cantem glóriam nómini sancto ejus, allelúja.
\end{tightlist}

\subsubsection*{\center \textsc{AD SEXTAM.}}
\subsubsection*{\textred{Capitulum. \hfill Judith 13.}}
\lettrine{B}enedícta tu a Deo tuo in omni tabernáculo Jacob; quóniam in omni Gente, quæ audíerit nomen tuum, magnificábitur super te Deus Israel.

\begin{tightlist}{}
\item \textred{℟.} Repleátur os meum laude Maríæ, \textred{*} Allelúja, alleluja. \textred{R}epleátur. \textred{℣.} Ut cantem glóriam nómini sancto ejus. \textred{A}llelúja, allelúja. Glória Patri. \textred{R}epleátur.
\item \textred{℣.} Ave María, grátia plena. \textred{℟.} Dóminus tecum, allelúja.
\end{tightlist}

\subsubsection*{\center \textsc{AD NONAM.}}
\subsubsection*{\textred{Capitulum. \hfill Ezech. 16.}}
\lettrine{E}gréssum est nomen tuum in Gentes propter spéciem tuam: quia perfécta eras in decóre meo, quem posúeram super te: dicit Dóminus Deus.

\begin{tightlist}{}
\item \textred{℟.} Ave María, grátia plena, \textred{*} Alleluja, alleluja. \textred{A}ve María. \textred{℣.} Dóminus tecum. \textred{A}llelúja, allelúja. Glória Patri. \textred{A}ve María.
\item \textred{℣.} Lætábor et exsultábo in te, María. \textred{℟.} Psalmum dicam Nómini tuo, alleluja.
\end{tightlist}

\subsubsection*{\center \textsc{AD II VESPERAS.}}
\begin{tightlist}{}
\item \textred{Añæ Laudum}
\item \textred{Psalmi de Comm. B.M.V.} [177]
%\item \textred{Capitulum Hymnus, et ℣. ut in j Vesperis.}
\end{tightlist}

\subsubsection*{\textred{Capitulum. \hfill Judith 13.}}
\lettrineThree{B}enedícta es tu fília a Dómino Deo excélso, præ ómnibus muliéribus super terram: quia hódie nomen tuum ita magnificávit, ut non recédat laus tua de ore hóminum, qui mémores fúerint virtútis Dómini in ætérnum.

\begin{tightlist}{}
%\item \textred{℟.} Felix namque es sacra Virgo María, et omni laude digníssima: * Quia ex te ortus est Sol justítiæ, Christus Deus noster.
%\textred{℣.} Ora pro pópulo, intérveni pro clero, intercéde pro devóto femíneo sexu: séntiant omnes tuum juvámen, quicúmque célebrant tui Nóminis commemoratiónem. \textred{Q}uia. Glória Patri. \textred{Q}uia.
\item \textred{Hymnus.} Ave maris stella. [180]
\item \textred{℣.} Sit nomen Dóminæ benedíctum. \textred{℟.} Ex hoc nunc, et usque in sǽculum, allelúja.
\item \textred{Ad Magnificat, Aña.} María Virgo ómnium óculis gratiósa, nominátum est nomen glóriæ ejus usque in extrémum terræ.
\end{tightlist}

%\subsubsection*{\center \textred{Oratio.}}
%\lettrine{D}eus, qui gloriósam Matrem tuam nominári Maríam voluísti: concéde quǽsumus; ut, qui dulce ejúsdem Maríæ Nomen implórant, perpétuum séntiant tuæ benedictiónis efféctum. Qui vivis et regnas.

\begin{center}\textcolor{red}{\rule{\linewidth}{0.5pt}} \end{center}
\clearpage
\subsection*{\center \textred{\textsc{Officia Mensualia}}}
\subsubsection*{\center \textred{Dicenda singulis mensibus}}

\subsection*{\center \textsc{Officium S.P.N. Eliæ}}
\subsubsection*{\center \textred{Semiduplex.}}
\begin{center}\textred{Recitandum semel in mense, juxta Rubr. 17.}\end{center}
\fancyhead[C]{\textsc{Officia Mensualia - Officium S.P.N. Eliæ}}

%\begin{center}\textcolor{red}{\rule{\linewidth}{0.5pt}} \end{center}
\begin{tightlist}{}
\item \textred{Temp. Pasch. tres Psalmi cujuslibet Nocturni dicuntur sub una sola Aña; et in fine omnium Añarum, et Respons. Versuum additur} allelúja, \textred{et ultimus versus Hymni ad Mat. mutatur.}
\item \textred{Reliqua omnia ut infra.}
\end{tightlist}

\subsubsection*{\center \textsc{AD I VESPERAS}}
\subsubsection*{\textred{Capitulum. \hfill Eccli. 48.}}
\lettrineThree{S}urréxit Elías Prophéta quasi ignis, et verbum ipsíus quasi fácula ardébat. Qui indúxit in pópulos famem, et irritántes illum invídia sua pauci facti sunt.

\subsubsection*{\center \textred{Hymnus.}}
\lettrine{N}unc juvat celsi properáre ad alta\\
Mente Carméli juga: et ore pleno\\
Vatis Elíæ méritos reférre\\
Laudis honóres.\\
\\
Hic sacri est ductor, cólumen, decúsque\\
Ordinis nostri, superísque gratus:\\
O Pater salve, plácidus tuósque\\
Excipe alúmnos.\\
\\
Dum fames terras prémeret, minístro\\
Uteris corvo; víduæ farínam\\
Súfficis solers; oleóque manat\\
Hydria pingui.\\
\\
Tu, Deo fusis précibus, redúctum\\
Fínibus mortis púerum parénti\\
Illius reddis víduæ, senéctæ\\
Dulce levámen.\\
\\
Claudis et nubes, aperísque verbo,\\
Imperas post tres properánter annos,\\
Décidant cœlo, sitiénsque tellus\\
Háuriat imbres.\\
\\
Sit Patri summo, genitǽque Proli,\\
Et tibi compar utriúsque sancte\\
Spíritus, laus, impérium, potéstas\\
Tempus in omne.\\
Amen.

\begin{tightlist}{}
\item \textred{℣.} Verbo Dómini contínuit cœlum. \textred{℟.} Et dejécit de cœlo ignem ter.
\item \textred{Ad Magnificat, Aña.} Ecce ego mittam vobis Elíam Prophétam, ántequam véniat dies Dómini magnus, et horríbilis. Et convértet cor patrum ad fílios, et cor filiórum ad patres eórum.
\end{tightlist}

\subsubsection*{\center \textred{Oratio.}}
\lettrine{P}ræsta, quǽsumus omnípotens Deus, ut sicut beátum Elíam Prophétam tuum, et Patrem nostrum ante commúnem mortem curru ígneo aéreum elevásti ad cœlum: ita nos fácias, eo interveniénte, dum vívimus, a terrénis semper ad cœléstia spíritu sublevári; et cum eo in resurrectióne justórum páriter gaudére. Per Dóminum.

\subsubsection*{\center \textsc{AD MATUTINUM}}
\begin{tightlist}{}
\item \textred{Invitatorium.} Regem Prophetárum Dóminum, * Veníte adorémus.
\item \textred{Ps.} Venite, exsultémus. 2.
\end{tightlist}

\subsubsection*{\center \textred{Hymnus.}}
\lettrine{T}e magne rerum Cónditor,\\
Mens nostra gliscit láudibus\\
In hoc Thesbíte máximo,\\
Quem díligis extóllere.\\
\\
Hic namque sacri nóminis\\
Tui zelátor próvocat\\
Vates Baal nequíssimos,\\
Victósque jure intérficit.\\
\\
Illo precánte, víctimas\\
Absúmit ignis cœ́litus:\\
Hi perstrepéntes ácriter,\\
Sunt ómnibus ludíbrio\\
\\
Tum Jezabélis ímpiæ\\
Vitans furórem nóxium,\\
Sub junípero dórmiens \\
Adésse cernit Angelum.\\
\\
Panem sibi qui próferens,\\
Lymphámque, jussit pérgere\\
Cibo reféctum strénue\\
Ad montis Horeb vérticem.\\
\\
Hoc in cibo jejúnium\\
Tulit quadragenárium:\\
Hac vi dapes sunt prǽditæ,\\
Quas dextra Dei pórrigit.\\
\\
Omnis tibi sit glória\\
Verbum, Pater, Paráclite,\\
Inseparáta Trínitas,\\
Quæ condidísti ómnia.\\
Amen.

\subsubsection*{\center \textsc{IN PRIMO NOCTURNO.}}
\begin{tightlist}{}
\item \textred{Aña.} Factum est autem, ut ægrotáret fílius mulíeris matris famílias, et erat languor fortíssimus.
\item \textred{Ps.} Beátus vir, \textred{cum reliquis, de Communi unius Conf. Pont.} [64]
\item \textred{Aña.} Dixit múlier ad Elíam: Quid mihi et tibi, vir Dei? Ingréssus es ad me, ut rememoraréntur iniquitátes meæ, et interfíceres fílium meum?
\item \textred{Aña.} Ait Elías ad víduam: Da mihi fílium tuum. Et clamávit ad Dóminum; et revíxit puer.
\item \textred{℣.} Verbo Dómini contínuit cœlum. \textred{℟.} Et dejécit de cœlo ignem ter.
\item \textred{Lectiones de Scriptura occurrente.}
\item \textred{℟. j.} Recéde hinc, et vade contra Oriéntem, dixit Dóminus ad Elíam. \textred{*} Qui ábiit, et fecit juxta verbum Dómini. \textred{℣.} Abscóndere in Torrénte Carith, et ibi de torrénte bibes: corvísque præcépi, ut pascant te ibi. \textred{Q}ui.
\item \textred{℟. ij.} Abiit múlier, et fecit juxta verbum Prophétæ: \textred{*} Et comédit ipse, et illa, et domus ejus. \textred{℣.} Ex illa die hydria farínæ non defécit, et lécythus ólei non est imminútus, juxta verbum Elíæ. \textred{E}t.
\item \textred{℟. iij.} Revérsa est ánima púeri intra eum, et revíxit, et trádidit matri suæ, et ait illi: \textred{*} En vivit fílius tuus. \textred{℣.} Tulit Elías púerum, et depósuit eum in inferiórem domum, dicens. \textred{E}n. Glória Patri. \textred{E}n.
\end{tightlist}

\subsubsection*{\center \textsc{IN SECUNDO NOCTURNO.}}
\begin{tightlist}{}
\item \textred{Aña.} Factum est verbum Dómini ad Elíam, dicens: Vade, et osténde te Achab, ut dem plúviam super fáciem terræ.
%\item \textred{Ps.} Cum invocárem. [67]
\item \textred{Aña.} Ait Elías ad Achab: Ascénde, cómede et bibe: quia sonus multæ plúviæ est.
\item \textred{Aña.} Manus Dómini facta est super Elíam, accinctísque lumbis currébat ante Achab, donec veníret in Jézræl.
\item \textred{℣.} Tu ungis reges ad pœniténtiam. \textred{℟.} Et Prophétas facis successóres post te.
\end{tightlist}

\subsubsection*{\center Sermo sancti Joánnis Chrysóstomi.}
\begin{center}\textred{Hom. 2. ad pop. Antioch.\\
\textbf{Lectio iv}.}\end{center}
\lettrineThree{I}n ómnibus ígitur philosophémur, et nihil nos inópia lǽdere póterit: quin et máxime adjuvábit, et clarióres reddet, et ómnibus divítibus copiosióres. Quid enim, dic mihi, Elía paupérius? Sed proptérea omnes dívites vincébat, quóniam pauper hic erat; ipsam vero paupertátem ex mentis opuléntia elégit. Quóniam enim omnem pecuniárum cópiam ánimi sui magnitúdine esse putávit inferiórem, nec philosophía sua dignam, ídeo tantam paupertátem amplexátus est.

\begin{tightlist}{}
\item \textred{℟.} Ait Abdías ad Elíam: Cum recéssero a te, Spíritus Dómini asportábit te in locum, quem ego ignóro, et tu dicis: \textred{*} Dic dómino tuo: Adest Elías. \textred{℣.} Dixit Elías: Vivit Dóminus exercítuum, ante cujus vultum sto, quia hódie apparébo ei. \textred{D}ic.
\end{tightlist}

\subsubsection*{\center \textred{Lectio v.}}
\lettrine{Q}uod si magna præséntia putásset, non solam melótem possedísset: sed ita vitæ vanitátem condemnávit, et tamquam lutum projéctum, aurum omne despéxit, ut præter amíctum illum, nihil ámplius possidéret. Idcírco rex páuperis indigébat, et verbis nihil plus quam melótem habéntis inhiábat, qui tantum habébat auri, ádeo splendídior trábea erat melóte, et regálibus aulis justi spelúnca. Proptérea et in cœlum ascéndens, nihil áliud quam melótem discípulo relíquit. Cum hac (inquit) cum diábolo pugnávi, et tu contra illum armátus hanc cápies.

\begin{tightlist}{}
\item \textred{℟.} Cum venísset Elías ad portam civitátis, appáruit ei múlier vídua cólligens ligna, et vocávit eam, dixítque ei: \textred{*} Da mihi páululum aquæ in vase, ut bibam. \textred{℣.} Cumque illa pérgeret ut afférret, clamávit post tergum ejus, dicens: Affer mihi, óbsecro, et buccéllam panis in manu tua. \textred{D}a.
\end{tightlist}

\subsubsection*{\center \textred{Lectio vj.}}
\lettrine{V}alidum enim telum inópia, et inexpugnábile habitáculum, et turris inconcússa. Tamquam máximam hereditátem Eliséus melótem suscépit: étenim vere máxima fuit heréditas, et omni auro pretiósior. Et erat posthac duplex Elías ille: et erat sursum Elías, et deórsum Elías. Novi quod justum illum beátum putátis, et vellétis quisque esse ut ille. Quod ígitur si vobis demonstrávero quid áliud, quod illo multo majus omnes sacris mystériis imbúti recípimus? Elías nempe melótem discípulo relíquit: Fílius autem Dei ascéndens, suam nobis carnem dimísit: sed Elías quidem exútus, Christus autem et nobis relíquit, et ipsam habens ascéndit.

\begin{tightlist}{}
\item \textred{℟.} Expándit se, atque mensus est Elías super púerum tribus vícibus, et clamávit ad Dóminum: \textred{*} Et exaudívit Dóminus vocem Elíæ.
\item \textred{℣.} Dómine Deus meus, revertátur óbsecro ánima púeri hujus in víscera ejus. \textred{E}t. Glória Patri. \textred{E}t.
\end{tightlist}

\subsubsection*{\center \textsc{IN TERTIO NOCTURNO}}
\begin{tightlist}{}
\item \textred{Aña.} Dixit Elías prophétis Baal: Elígite vobis bovem unum, et invocáte nómina deórum vestrórum. Qui clamábant, et non erat qui respondéret eis.
%\item \textred{Ps.} Dómine, quis. [70]
\item \textred{Aña.} Cum jam tempus esset, ut offerrétur holocáustum, accédens Elías orávit ad Dóminum.
\item \textred{Aña.} Oránte Elía cécidit ignis Dómini, et vorávit holocáustum, et ligna, et lápides, púlverem quoque, et aquam quæ erat in aquædúctu lambens.
\item \textred{℣.} Beáti sunt qui te vidérunt.
\item \textred{℟.} Et in amicítia tua decoráti sunt.
\end{tightlist}

\subsubsection*{\center Léctio sancti Evangélii secúndum Lucam.}
\noindent \textred{\textbf{Lectio vij}. \hfill Cap. 9.}\\
\lettrine{I}n illo témpore: Assúmpsit Jesus Petrum, et Jacóbum, et Joánnem, et ascéndit in montem ut oráret, et facta est dum oráret, spécies vultus ejus áltera. Et réliqua. 

\begin{center}Homilía venerábilis Bedæ Presbyteri.\end{center}
\lettrineThree{E}rant autem Móyses et Elías visi in majestáte, et dicébant excéssum ejus, quem completúrus erat in Jerúsalem. Móyses ergo et Elías, qui locúti sunt cum Dómino in monte, et passiónem ejus ac resurrectiónem dicébant, orácula legis et prophetíæ desígnant, quæ in Dómino compléta, et nunc doctis quibúsque patent, et eléctis ómnibus in futúro maniféstius patébunt.

\begin{tightlist}{}
\item \textred{℟.} Ait Elías núntiis regis samaríæ: Numquid non est Deus in Israel, ut eátis ad Beélzebub? Quam ob rem hæc dicit Dóminus: \textred{*} De léctulo, super quem ascendísti, non descéndes, sed morte moriéris. \textred{℣.} Revérsi núntii ad Ochozíam dixérunt: Vir occúrrit nobis, dicens: Ite, et revertímini ad regem, et dicétis ei. \textred{D}e.
\end{tightlist}

\subsubsection*{\center \textred{Lectio viij.}}
\lettrine{Q}ui bene visi in majestáte dicúntur, quia tunc apértius videbúntur, quando veritátis omnis divinórum eloquiórum non solum sensus, sed et sermo fúerit prolátus: quamvis cum Móyse et Elía omnes, qui cum Dómino regnatúri sunt, recte intélligi possint. In Móyse quidem, qui mórtuus et sepúltus est, hi qui in judício resuscitándi a morte; in Elía vero, qui necdum mortis débitum solvit, hi qui in advéntu Júdicis, vivi inveniéndi sunt in carne. Qui utríque páriter in uno eodémque moménto rapti in núbibus óbviam Dómino in áera, mox perácto judício ad vitam inducéntur ætérnam.

\begin{tightlist}{}
\item \textred{℟.} Ascéndit quinquagenárius, et quinquagínta qui erant sub eo, ad Elíam sedéntem in vértice montis. \textred{*} Et descéndit ignis de cœlo, et devorávit eos. \textred{℣.} Dixit Elías: Si homo Dei sum, descéndat ignis de cœlo, et dévoret te et quinquagínta tuos. \textred{E}t.
\end{tightlist}

\subsubsection*{\center \textred{Lectio ix.}}
\lettrine{Q}uibus apte cónvenit, quod Móyses et Elías in majestáte fuísse visi perhibéntur: per excelléntiam namque majestátis, insígne múneris, quo sunt coronándi, monstrátur. Bene cónvenit et hoc, quod excéssum ejus, quem completúrus erat in Jerúsalem, dixísse memorántur: quia nimírum única laudis matéria fidélibus suis sit pássio Redemptóris, tanto suæ reparatiónis memóriam devóto semper in péctore versant, devóta confessióne testántur.

\begin{tightlist}{}
\item \textred{℟.} Factum est, dum tólleret Dóminus Elíam per túrbinem in cœlum, \textred{*} Eliséus clamábat, dicens: Pater mi, pater mi, currus Israel, et auríga ejus. \textred{℣.} Cumque pérgerent, et incedéntes sermocinaréntur: ecce currus ígneus et equi ígnei divisérunt utrúmque, et ascéndit Elías per túrbinem in cœlum. \textred{E}liséus. Glória Patri. \textred{E}liséus.
\item Te Deum laudámus. 6.
\item \textred{Sacerd. ℣.} Surréxit Elías prophéta quasi ignis. \textred{℟.} Et verbum ipsíus quasi fácula ardébat.
\end{tightlist}

\subsubsection*{\center \textsc{AD LAUDES}}
\begin{center}\textred{et per Horas, Aña.}\end{center}
\lettrine{E}lías, dum zelat zelum legis, recéptus est in cœlum.

\begin{tightlist}
\item \textred{Ps. de Dominica.} 30.
\item \textred{Aña.} Factum est autem, cum leváre vellet Dóminus Elíam per túrbinem in cœlum, ibant Elías et Eliséus de Gálgalis.
\item \textred{Aña.} Tulit Elías pállium suum, et percússit aquas, et divísæ sunt, et transiérunt ambo per siccum.
\item \textred{Aña.} Cum Elías et Eliséus pérgerent, et incedéntes sermocinaréntur, ecce currus ígneus et equi ígnei divisérunt utrúmque.
\item \textred{Aña.} Ascéndit Elías per túrbinem in cœlum, Eliséus autem vidébat, et clamábat: Pater mi, pater mi, currus Israel et auríga ejus.
\end{tightlist}

\subsubsection*{\textred{Capitulum. \hfill Eccli. 48.}}
\lettrineThree{S}urréxit Elías Prophéta quasi ignis, et verbum ipsíus quasi fácula ardébat. Qui indúxit in pópulos famem: et irritántes illum invídia sua pauci facti sunt.

\subsubsection*{\center \textred{Hymnus.}}
\lettrine{P}ergámus sócii, tóllere cánticis\\
Thesbítæ, exímii Vatis, in ǽthera\\
Res gestas, hílares: nam decet hoc die\\
Nostri nos mémores ducis.\\
\\
Audívit Dómini verba recóndita\\
Auræ sub ténuis múrmure síbili,\\
Jussus sævítiam Jézabel ímprobæ,\\
Atque Achab scélera ínsequi.\\
\\
Succínctus túnica pellícea latus,\\
Carmelíque hábitans antra viréntia,\\
Prædícit mínitans trístia fúnera\\
Regi Ochozíæ ímpio.\\
\\
Hujus bis précibus ignis ab ǽthere\\
Absórpsit véniens, flámmaque mílites:\\
Findit dein próprio flúmina pállio,\\
Et sicco pede prǽterit.\\
\\
Nos nunquam auxílium destítuat, Pater,\\
Optatúsque tuus perpétuo favor:\\
Nobis quin pótius múnera Spíritus\\
Augéri prece póstula.\\
\\
Laus Patri ingénito, summaque glória,\\
Sit semper génito laus quoque Fílio,\\
Necnon par tibi sit, ínclyte Spíritus,\\
Simplex únaque Trínitas.\\
Amen.

\begin{tightlist}
\item \textred{℣.} Elías in túrbine tectus est. 
\item \textred{℟.} Et in Eliséo complétus est spíritus ejus.
\item \textred{Ad Benedictus, Aña.} Elías homo erat símilis nobis passíbilis, et oratióne orávit, ut non plúeret super terram; et non pluit annos tres et menses sex. Et rursum orávit: et cœlum dedit plúviam, et terra dedit fructum suum.
\end{tightlist}

\subsubsection*{\center \textred{Oratio.}}
\lettrine{P}ræsta, quǽsumus omnípotens Deus, ut sicut beátum Elíam Prophétam tuum, et Patrem nostrum ante commúnem mortem curru ígneo aéreum elevásti ad cœlum: ita nos fácias, eo interveniénte, dum vívimus, a terrénis semper ad cœléstia spíritu sublevári; et cum eo in resurrectióne justórum páriter gaudére. Per Dóminum.

\subsubsection*{\center \textsc{AD TERTIAM.}}
\subsubsection*{\textred{Capitulum. \hfill Eccli. 48.}}
\lettrine{S}urréxit Elías Prophéta quasi ignis, et verbum ipsíus quasi fácula ardébat. Qui indúxit in pópulos famem: et irritántes illum invídia sua pauci facti sunt.

\begin{tightlist}{}
%\item \textred{Capit.} Surréxit Elías, \textred{ut supra ad Laudes.}
\item \textred{℟.} Verbo Dómini \textred{*} Contínuit cœlum. \textred{V}erbo. \textred{℣.} Et dejécit de cœlo ignem ter. \textred{C}ontínuit. Glória Patri. \textred{V}erbo.
\item \textred{℣.} Tu ungis reges ad pœniténtiam. \textred{℟.} Et Prophétas facis successóres post te.
\end{tightlist}

\subsubsection*{\center \textsc{AD SEXTAM.}}
\subsubsection*{\textred{Capitulum \hfill Eccli. 48.}}
\lettrine{V}erbo Dómini contínuit cœlum, et dejécit de cœlo ignem ter: sic amplificátus est Elías in mirabílibus suis.

\begin{tightlist}
\item \textred{℟.} Tu ungis reges \textred{*} Ad pœniténtiam. \textred{T}u ungis. \textred{℣.} Et Prophétas facis successóres post te. \textred{A}d. Glória Patri. \textred{T}u.
\item \textred{℣.} Beáti sunt qui te vidérunt. 
\item \textred{℟.} Et in amicítia tua decoráti sunt.
\end{tightlist}

\subsubsection*{\center \textsc{AD NONAM.}}
\subsubsection*{\textred{Capitulum \hfill Eccli. 48.}}
\lettrine{T}u sustulísti mórtuum ab ínferis de sorte mortis in verbo Dómini Dei. Qui scriptus es in judíciis témporum leníre iracúndiam Dómini.

\begin{tightlist}
\item \textred{℟.} Beáti sunt, \textred{*} Qui te vidérunt. \textred{B}eáti. \textred{℣.} Et in amicítia tua decoráti sunt. \textred{Q}ui. Glória Patri. \textred{B}eáti.
\item \textred{℣.} Elías in túrbine tectus est. 
\item \textred{℟.} Et in Eliséo complétus est spíritus ejus.
\end{tightlist}

\subsubsection*{\center \textsc{AD II. VESPERAS.}}
\begin{tightlist}
\item \textred{Nisi sequatur Festum duplex, dicitur prima Aña. Laudum, sub qua sola dicuntur Psalmi de Dominica; et loco ult.} Lauda Jerúsalem. 216.
\end{tightlist}

\subsubsection*{\textred{Capitulum. \hfill Eccli. 48.}}
\lettrineThree{S}urréxit Elías Prophéta quasi ignis, et verbum ipsíus quasi fácula ardébat. Qui indúxit in pópulos famem, et irritántes illum invídia sua pauci facti sunt.

\subsubsection*{\center \textred{Hymnus.}}
\lettrine{A}udiat miras óriens cadénsque\\
Sol tuas laudes, mare, terra, et aer,\\
Ordinis nostri cólumen, tuǽque\\
Glória gentis.\\
\\
O Deo multum (Pater alme) care,\\
O potens signis, meritísque felix,\\
Quem Deus gestis adhibére suévit\\
Grándibus olim.\\
\\
Te triumpháli super alta curru\\
Igneis vectum céleres quadrígis\\
Angeli attóllunt, nímium corúsca\\
Luce micántem\\
\\
Unico ætérni Patris adfuísti\\
Fílio testis, sócius fidélis\\
Mosis, in summo pósitus Thabóris\\
Vértice montis.\\
\\
Nos tui præsens clypeo favóris,\\
Ac tua magna bonitáte fretos,\\
Súpplices dextra fámulos benígna\\
Prótege semper.\\
\\
Sit Patri summo, genitǽque Proli,\\
Et tibi compar utriúsque sancte\\
Spíritus, laus, impérium, potéstas\\
Tempus in omne.\\
Amen.

\begin{tightlist}
\item \textred{℣.} Verbo Dómini contínuit cœlum. \textred{℟.} Et dejécit de cœlo ignem ter.
\item \textred{Ad Magnificat, vel pro commemor. Aña.} Tulit Elías pállium suum, et invólvit illud, et percússit aquas Jordánis, quæ divísæ sunt in utrámque partem, et transiérunt ipse, et Eliséus per siccum: qui incedéntes, ecce currus ígneus, et equi ígnei divisérunt utrúmque: et ascéndit Elías per túrbinem in cœlum: Eliséus autem non vidit eum ámplius.
%\item \textred{Temp. Pasch. tres Psalmi cujuslibet Nocturni dicuntur sub una sola Aña; et in fine omnium Añarum, et Respons. Versuum additur} allelúja, \textred{et ultimus versus Hymni ad Mat. mutatur.}
%\item \textred{Cetera omnia ut supra.}
\end{tightlist}

\begin{center}\textcolor{red}{\rule{\linewidth}{0.5pt}} \end{center}

\subsection*{\center \textsc{Officium S. Andreæ Corsini}}
\subsubsection*{\center \textred{Ep. et Conf. Ordinis Nostri.\\ Semiduplex}}
\begin{center}\textred{Recitandum in toto Ordine nostro bis in Mense per Annum, et semel pariter in Mense per Quadragesimam sub ritu Semiduplici, juxta Rubr. 17.}\end{center}
\fancyhead[C]{\textsc{Officia Mensualia - Officium S. Andreæ Corsini}}

\begin{tightlist}{}
\item \textred{Tempore Paschali tres Psalmi cujuslibet Nocturni dicuntur sub una sola Aña, et in fine omnium Añarum etc. additur} Allelúja. \textred{Cetera omnia ut infra.}
\end{tightlist}

\subsubsection*{\center \textsc{AD I VESPERAS}}
\subsubsection*{\textred{Capitulum. \hfill Eccli. 50.}}
\lettrineThree{E}cce Sacérdos magnus, qui in vita sua suffúlsit domum: et in diébus suis corroborávit templum. Qui præváluit amplificáre civitátem, qui adéptus est glóriam in conversatióne Gentis.

\subsubsection*{\center \textred{Hymnus.}}
\lettrine{A}ndréam grádibus sídera lúcidis\\
Scandéntem, sócii, et Prǽsule cérnitis\\
Stipátum gémino: ferte citi lyras,\\
Plectrúmque, et nova cántica.\\
\\
Magnis ex átavis éditus áureæ\\
Sortis dum séquitur blanda perícula;\\
Ibat forte puer, qua via flórido\\
Fallit nóxia trámite.\\
\\
Inclámat génitrix: ille pedem loco\\
Festínat rétrahens; áltaque per juga\\
Carméli evólitat, certus inánia\\
Mundi línquere gáudia.\\
\\
Ast illum cúpide nóbilis ínfula\\
Quærit: sed látitans défugit ínfulam\\
Frustra: nam púero vox data trímulo\\
Noléntem ad sólium rapit.\\
\\
Ergo divítias fúndere páuperi\\
Et rectum sérere, et prava revéllere;\\
Cœlóque, ac pópulo, non sibi vívere;\\
Hæc laus Pontíficis fuit.\\
\\
Pastor summe, óvium qui regis ágmina,\\
Pastórum páriter Rector, et Arbiter\\
Christe, ad perpétuæ duc sata Pátriæ\\
Una cum dúcibus gregem.\\
Amen.

\begin{tightlist}{}
\item \textred{℣.} Ora pro nobis sancte Andréa. \textred{℟.} Ut digni efficiámur promissiónibus Christi.
\item \textred{Ad Magnificat, Aña.} Natus est firmaméntum gentis, rector fratrum stabiliméntum pópuli.
\end{tightlist}

\subsubsection*{\center \textred{Oratio.}}
\lettrine{D}eus, qui in Ecclésia tua nova semper instáuras exémpla virtútum: da pópulo tuo; beáti Andréæ Confessóris tui, atque Pontíficis ita sequi vestígia, ut assequátur et prǽmia. Per Dóminum.

\subsubsection*{\center \textsc{AD MATUTINUM}}
\begin{tightlist}{}
\item \textred{Invitatorium.} Dóminum misericórdiæ, qui prædestinávit Andréam confórmem fíeri agno immaculáto: * Veníte adorémus.
\item \textred{Ps.} Veníte, exsultémus. 2.
\end{tightlist}

\subsubsection*{\center \textred{Hymnus.}}
\lettrine{V}irgo, Carméli decor et volúptas,\\
Monte subjéctos speculáta campos,\\
Ut lupum vidit; míserans corúsco\\
Térruit arco.\\
\\
Tunc feræ (mirum!) cecidére vultus,\\
Ora tunc agni subiére mitis,\\
tunc datum cursu per amœna ferri\\
Cúlmina montis.\\
\\
Vírginis cordi pius hæret ardor,\\
Vírginis formam sedet æmulári\\
Móribus puris: alit illa sacros\\
Altius ignes.\\
\\
Mox jubet certam pópulis salútem\\
Ferre commíssis: mora nulla, late\\
Italas Præsul rádians per oras\\
Lucet, et urit.\\
\\
Núminis plenus rétegit futúra,\\
Urbibus rixas fugat, et mináces\\
Artubus morbos; epulásque turbis\\
Dévocat axe.\\
\\
Laus Patri, et Nato; tibi laus et alme\\
Spíritus semper; duce quo triúmphat\\
Victor Andréas, Mariána vivens\\
Signa secútus.\\
Amen.

\subsubsection*{\center \textsc{IN PRIMO NOCTURNO.}}
\begin{tightlist}{}
\item \textred{Aña.} Stérilis virtútem in conceptiónem accépit, et unigénitum offerébat.
\item \textred{Ps.} Beátus vir, \textred{cum reliquis de Communi unius Confess. Pontif.} [64]
\item \textred{Aña.} Oculus Dómini respéxit eum in bono, et eréxit eum ab humilitáte sua.
\item \textred{Aña.} Miráti sunt in illo multi, et laudavérunt Deum.
\item \textred{℣.} Amávit eum Dóminus, et ornávit eum. \textred{℟.} Stolam glóriæ índuit eum.
\item \textred{Lect. j. Noct. de Script. occurr. quæ si desint, legantur sequentes.}
\end{tightlist}

\subsubsection*{\center De libro Job.}
\begin{center}\textred{\textbf{Lectio j}. \hfill Cap. 29. V. 7.}\end{center}
\lettrineThree{Q}uando procedébam ad portam civitátis, et in platéa parábant cáthedram mihi: auris aúdiens beatificábat me, et óculus videns testimónium reddébat mihi. Eo quod liberássem páuperem vociferántem, et pupíllum, cui non esset adjútor. Benedíctio peritúri super me veniébat, et cor víduæ consolátus sum. Justítia indútus sum: et vestívi me, sicut vestiménto, et diadémate, judício meo. Oculus fui cæco, et pes claudo. Pater eram páuperum, et causam, quam nesciébam, diligentíssime investigábam. Dicebámque: in nídulo meo móriar et sicut palma multiplicábo dies. Radix mea apérta est secus aquas, et ros morábitur in messióne mea. Glória mea semper innovábitur.

\begin{tightlist}{}
\item \textred{℟.} Concépit stérilis et péperit fílium quia magnificábit Dóminus misericórdiam suam cum illa. \textred{*} Et congratulabántur ei. \textred{℣.} Cœpítque Spíritus Dómini esse cum eo. \textred{E}t.
\end{tightlist}

\subsubsection*{\center\textred{Lectio ij. \hfill Cap. 31. V. 1.}}
\lettrine{P}epigi fœdus cum óculis meis, ut ne cogitárem quidem de vírgine. Quam enim partem habéret in me Deus désuper, et hereditátem Omnípotens de excélsis? Numquid non perdítio est iníquo, et alienátio operántibus injustítiam? Si ambulávi in vanitáte, et festinávit in dolo pes meus: Si declinávit gressus meus de via, et si secútum est óculos meos cor meum, et si mánibus meis adhǽsit mácula: seram, et álius cómedat. Si decéptum est cor meum super mulíere, et si ad óstium amíci mei insidiátus sum. Hoc enim nefas est, et iníquitas máxima. Ignis est usque ad perditiónem dévorans, et ómnia eradícans genímina. Semper quasi tuméntes super me fluctus tímui Deum, et pondus ejus ferre non pótui. 

\begin{tightlist}{}
\item \textred{℟.} Humiliátus est póstea, eo quod exaltátum fuísset cor ejus. \textred{*} Et opéruit cilício carnem suam et jejunávit. \textred{℣.} Ascéndit ut immoláret Dómino hóstiam solémnem et votum suum. \textred{E}t.
\end{tightlist}

\subsubsection*{\center \textred{Lectio iij.}}
\lettrine{S}i putávi autum robur meum, et obrízo dixi fidúcia mea. Si gavísus sum ad ruínam ejus, qui me óderat: et exsultávi, quod invenísset eum malum. Non enim dedi ad peccándum guttur meum, ut expéterem maledícens ánimam ejus. Si abscóndi quasi homo peccátum meum, et celávi in sino meo iniquitátem meam. Si expávi ad multitúdinem nímiam, et despéctio propinquórum térruit me: et non magis tácui, nec egréssus sum óstium. Si advérsum me terra mea clamat, et cum ipsa sulci ejus deflent. Si fructus ejus comédi absque pecúnia, et ánimam agricolárum ejus afflíxi: pro fruménto oriátur mihi tríbulus, et pro hórdeo spina. 

\begin{tightlist}{}
\item \textred{℟.} Invénit hóminem, qui erat cæcus, et ait: sanat te Dóminus Jesus Christus: \textred{*} Et conféstim vidit. 
\item \textred{℣.} Est aqua lustratiónis aspérsus. \textred{E}t. Glória Patri. \textred{E}t.
\end{tightlist}

\subsubsection*{\center \textsc{IN SECUNDO NOCTURNO.}}
\begin{tightlist}{}
\item \textred{Aña.} Fecit quod plácuit Deo, et fórtiter ivit in via, quam mandávit Prophéta magnus in conspéctu Dei.
%\item \textred{Ps.} Cum invocárem. [67]
\item \textred{Aña.} Elégit Dóminus servum suum páscere hereditátem suam, et in intelléctibus mánuum suárum dedúxit eos.
\item \textred{Aña.} Dedit eis petitiónem ipsórum, misit saturitátem in ánimas eórum.
\item \textred{℣.} Justum dedúxit Dóminus per vias rectas. \textred{℟.} Et osténdit illi regnum Dei.
\end{tightlist}

\subsubsection*{\center Sermo sancti Gregórii Nysséni.}
\begin{center}\textred{In fun. M. Milet.\\
\textbf{Lectio iv}.}\end{center}
\lettrineThree{C}um primum bene moráta, ac modésta Ecclésia beátum Andréam vidit, vidit fáciem ad imáginem Dei vere formátam: vidit dilectiónem fontis modo scaturiéntem: vidit grátiam lábiis circumfúsam, ánimi demissiónis summum gradum: post quem ámplius quidquam cogitári non potest: vidit mansuetúdinem atque cleméntiam, qualis in Davíde fuit: qualis in Salomóne, intelligéntiam atque prudéntiam: qualis in Móyse, bonitátem: qualis in Samuéle, perfectiónem: qualis in Josépho, continéntiam pudicitiámque: qualis in Daniéle, sapiéntiam: quemádmodum magnus Elías, zelo fídei prǽditum: sícuti sublímis Joánnes, integritáte córporis ornátum: sícuti Paulus, inexsuperábili prǽditum dilectióne. Vidit tot bonórum circa unam ánimam concúrsum: amóre beáto vulneráta est, casto bonóque Sponsum suum amóre atque benevoléntia prosecúta diléxit.

\begin{tightlist}{}
\item \textred{℟.} In médio fratrum rector illórum loquebátur de regno Dei, et eos qui cura indigébant, sanábat: \textred{*} Omnibus mirántibus. \textred{℣.} Ecce quidam hydrópicus erat ante illum; ipse vero sanávit eum. \textred{O}mnibus.
\end{tightlist}

\subsubsection*{\center \textred{Lectio v.}}
\lettrine{S}ed priúsquam cupiditátem expléret, ántequam desidérium recreáret atque sedáret, amóris vi fervens, tentatiónibus athlétam ad certámina vocántibus, sola relícta est. Atque ille quidem in certamínibus pro pietáte suscéptis desudábat: hæc vero durábat in castitáte matrimónium consérvans. Non ablátus est a nobis Sponsus, in médio nostrum stat, étiam si nos non videámus. In ádytis, ac penetrálibus Sacérdos est. In interióribus veli, quo præcúrsor pro nobis ingréssus est Christus, relíquit carnis teguméntum. Non ámplius signo, et umbræ cœléstium servit; sed in ipsam rerum imáginem intuétur: non ámplius per spéculum, atque per transénam, et ænígma; sed ipse fácie cum fácie colláta intercédit apud Deum. 

\begin{tightlist}{}
\item \textred{℟.} Non seípsum clarificávit, ut Póntifex fíeret \textred{*} sed vocátur a Deo tamquam Aáron. \textred{℣.} Quasi sol refúlgens, sic ille effúlsit in templo Dei. \textred{S}ed.
\end{tightlist}

\subsubsection*{\center \textred{Lectio vj.}}
\lettrine{I}ntercédit autem pro nobis, et pópuli errátis: depósuit túnicas pellíceas, neque enim tálibus túnicis opus habent, qui in Paradíso degunt: sed habet induménta, quæ puritáte vitæ suæ contéxuit, iísque sese ornávit. Honoráta, ac pretiósa coram Dómino talis viri mors est; immo vero non mors, sed ruptúra membrórum est. Dirupísti enim, inquit, víncula mea. Dimíssus est Símeon, liberátus est a vínculis córporis. Láqueus contrítus est, et avícula avolávit. Ingréssus est terram promissiónis: in monte cum Deo philosophátur. Solvit ánimæ calceaméntum, ut pura planta mentis terram sanctam, ubi conspícitur Deus, conscénderet.

\begin{tightlist}{}
\item \textred{℟.} Cum turba multa esset, nec habérent quod manducárent, accípiens septem panes fregit, et dabat: \textred{*} Et manducavérunt, et saturáti sunt. \textred{℣.} Non sunt fraudáti a desidério suo. \textred{E}t. Glória Patri. \textred{E}t.
\end{tightlist}

\subsubsection*{\center \textsc{IN TERTIO NOCTURNO.}}
\begin{tightlist}{}
\item \textred{Aña.} Minúitur crescens mirabíliter in consummatióne, respléndens glorióse.
%\item \textred{Ps.} Dómine quis. [70]
\item \textred{Aña.} Dóminus dedit illi sedem glóriæ, et propter illum dejécit omnem poténtiam inimicórum.
\item \textred{Aña.} Certe in géntibus multis non erat símilis illi diléctus Deo suo.
\item \textred{℣.} Justus ut palma florébit in domo Dómini. \textred{℟.} Sicut cedrus Líbani multiplicábitur.
\end{tightlist}

\subsubsection*{\center Léctio sancti Evangélii secúndum Matthǽum.}
\begin{center}\textred{\textbf{Lectio vij}.} \hfill \textred{Cap. 25.}\end{center}
\lettrine{I}n illo témpore: Dixit Jesus discípulis suis: Homo quidam péregre proficíscens, vocávit servos suos, et trádidit illis bona sua. Et réliqua.

\begin{center}Homilía sancti Joánnis Chrysóstomi.\\
\textred{Hom. 79. ad med.}\end{center}
\lettrineThree{T}alénta hic pro eo, quod unusquísque fácere potest accípimus, sive auctoritáte protégere, sive pecúniis juváre, sive doctrína admonére, sive ália quápiam re próximis prodésse queas. Nemo secum dicat: Cum unum taléntum hábeam, nihil possum effícere: potest profécto ex ea una re sola comprobári. Non es paupérior illa vídua; non es Petro, atque Joánne rustícior: qui, quamvis rudes simul atque illiteráti fuérunt, quóniam magno stúdio commúnem utilitátem compléxi fúerant, cœlórum Príncipes facti sunt.

\begin{tightlist}{}
\item \textred{℟.} In illa hora multos curávit a languóribus et plagis: \textred{*} Quia virtus de illo exíbat et sanábat. \textred{℣.} In vita sua fecit monstra, et in morte mirabília operátus est. \textred{Q}uia.
\end{tightlist}

\subsubsection*{\center \textred{Lectio viij.}}
\lettrine{Q}uippe nulla res Deo grátior est, quam ut univérsam vitam ad commúne cómmodum cónferas. Idcírco ratióne, atque oratióne nos Deus decorávit, mentem, atque ingénium concéssit, manus, pedes, vires córporis dedit, ut his ómnibus, et nos ipsos, et próximos tutémur. Non enim ad agéndas grátias solúmmodo confert orátio, verum étiam ad docéndum, et ad monéndum perútilis est: qua in re, si ea diligénter utémur, Dóminum; sin vero in contráriis, diábolum imitábimur.

\begin{tightlist}{}
\item \textred{℟.} Repúlsi sunt inimíci præ timóre ejus: \textred{*} Et dirécta est salus in manu ejus. \textred{℣.} Cógnitus est in verbis suis fidélis. \textred{E}t.
\end{tightlist}

\subsubsection*{\center \textred{Lectio ix.}}
\lettrine{N}am et Petrus quidem, quando Christum confitebátur, quasi ea dicens, quæ Patris sunt, beatificabátur: quando autem crucem horrébat, atque repellébat, quasi ea, quæ diáboli sunt sápiens, ácriter increpabátur. Si ergo ignoratióne crucem recusándo reprehénsus est, quam nos sponte peccántes véniam impetrábimus? Ea ergo loquámur, quæ Christi verba esse, conféstim cognoscúntur.

\begin{tightlist}{}
\item \textred{℟.} Spíritu magno vidit última, consolátus est lugéntes, osténdit futúra, et abscóndita. \textred{*} Quæsívit bona genti suæ, et fecit pacem. \textred{℣.} Adjúvit páuperem de inópia, et pósuit sicut oves famílias. \textred{Q}uæsívit. Glória Patri. \textred{Q}uæsívit.
\item Te Deum laudámus. 6.
\item \textred{Sacerdotalis ℣.} Elégit dum Dóminus sacerdótem sibi. \textred{℟.} Ad sacrificándum ei hóstiam laudis.
\end{tightlist}

\subsubsection*{\center \textsc{AD LAUDES}}
\begin{center} \textred{et per Horas, Aña.}\end{center}
\lettrine{D}iléxit Andréam Dóminus in odórem suavitátis.

\begin{tightlist}{}
\item \textred{Ps. de Dominica.} 30.
\item \textred{Aña.} Jucunditátem et exsultatiónem thesaurizávit super illum, et nómine ætérno hereditávit illum.
\item \textred{Aña.} In médio pópuli sui exaltátus est, et in multitúdine sanctórum habet laudem.
\item \textred{Aña.} Perféctus invéntus est, et erit illi glória ætérna.
\item \textred{Aña.} Implétus est benedictiónibus, et vidéntes illum laudábunt.
\end{tightlist}

\subsubsection*{\textred{Capitulum. \hfill Eccli. 50.}}
\lettrineThree{E}cce Sacérdos magnus, qui in vita sua suffúlsit domum: et in diébus suis corroborávit templum. Qui præváluit amplificáre civitátem, qui adéptus est glóriam in conversatióne Gentis.

\subsubsection*{\center \textred{Hymnus.}}
\begin{center}\textred{In quo Temp. Pasch. nihil additur, neque mutatur.}\end{center}
\lettrine{Q}ua nocte mundi Cónditor\\
Tenéllus infans náscitur,\\
María amánti Prǽsuli\\
Mortem propínquam núntiat.\\
\\
Qua stella nocte ad Párvuli\\
Tres ducit antra Príncipes,\\
María, quæ stella est maris,\\
Ad astra ducit Prǽsulem.\\
\\
Artus relícti léctulum,\\
Seu pulchra lília, aut rosæ,\\
Odóre perflant; ac salus\\
Odor fit unus lánguidis.\\
\\
Quin luce mira prǽnitent,\\
Qua pulsa luce nox fugit:\\
Hæc nempe sunt vestígia\\
Nitentióris spíritus.\\
\\
O qui per alta, spíritus,\\
Spátia vagáris ǽtheris,\\
Torus Deo jam fúlgurans,\\
Ne sperne res mortálium.\\
\\
Te fons, Redémptor, lúminis,\\
Te campi odórum lílium,\\
Trahe nos odóre, póscimus,\\
Ad lumen almum glóriæ.\\
Amen.

\begin{tightlist}{}
\item \textred{℣.} Justus germinábit sicut lílium. \textred{℟.} Et florébit in ætérnum ante Dóminum.
\item \textred{Ad Benedictus, Aña.} Fecit quod rectum erat in óculis Dómini, et diréxit eum in ómnibus, et adjúvit eum Deus, et divulgátum est nomen ejus.
\end{tightlist}

\subsubsection*{\center \textred{Oratio.}}
\lettrine{D}eus, qui in Ecclésia tua nova semper instáuras exémpla virtútum, da pópulo tuo beáti Andréæ Confessóris tui atque Pontíficis ita sequi vestígia, ut assequátur et prǽmia. Per Dóminum.

\begin{tightlist}{}
\item \textred{Ad Horas Antiphonæ de Laudibus, Psalmi de Dominica, ad Primam tamen ut in Festis. Capit, et ℟℟. de Communi.} [75]
\item \textred{Ad ij. Vesp. Aña. prima Laudum, Psalmi de Dominica, et loco ultimi, Ps. 131.} Meménto Dómine David. [76].
%\item \textred{Capit. Hymn. et ℣. ut in j. Vesperis.}
\end{tightlist}

\subsubsection*{\textred{Capitulum. \hfill Eccli. 50.}}
\lettrineThree{E}cce Sacérdos magnus, qui in vita sua suffúlsit domum: et in diébus suis corroborávit templum. Qui præváluit amplificáre civitátem, qui adéptus est glóriam in conversatióne Gentis.

\subsubsection*{\center \textred{Hymnus.}}
\lettrine{A}ndréam grádibus sídera lúcidis\\
Scandéntem, sócii, et Prǽsule cérnitis\\
Stipátum gémino: ferte citi lyras,\\
Plectrúmque, et nova cántica.\\
\\
Magnis ex átavis éditus áureæ\\
Sortis dum séquitur blanda perícula;\\
Ibat forte puer, qua via flórido\\
Fallit nóxia trámite.\\
\\
Inclámat génitrix: ille pedem loco\\
Festínat rétrahens; áltaque per juga\\
Carméli evólitat, certus inánia\\
Mundi línquere gáudia.\\
\\
Ast illum cúpide nóbilis ínfula\\
Quærit: sed látitans défugit ínfulam\\
Frustra: nam púero vox data trímulo\\
Noléntem ad sólium rapit.\\
\\
Ergo divítias fúndere páuperi\\
Et rectum sérere, et prava revéllere;\\
Cœlóque, ac pópulo, non sibi vívere;\\
Hæc laus Pontíficis fuit.\\
\\
Pastor summe, óvium qui regis ágmina,\\
Pastórum páriter Rector, et Arbiter\\
Christe, ad perpétuæ duc sata Pátriæ\\
Una cum dúcibus gregem.\\
Amen.

\begin{tightlist}{}
\item \textred{℣.} Ora pro nobis sancte Andréa. \textred{℟.} Ut digni efficiámur promissiónibus Christi.
\item \textred{Ad Magnificat, Aña.} Hic est Fratrum amátor, et pópuli: hic est qui multum orat pro pópulo, et univérsa civitáte.
%\item \textred{Tempore Paschali tres Psalmi cujuslibet Nocturni dicuntur sub una sola Aña, et in fine omnium Añarum etc. additur} Allelúja. \textred{Cetera omnia ut supra.}
\end{tightlist}

\begin{center}\textcolor{red}{\rule{\linewidth}{0.5pt}} \end{center}

\section*{\center \textsc{Officium Sanctæ Teresiæ}}
\subsubsection*{\center \textred{Virginis Ordinis Nostri.\\ Semiduplex}}
\begin{center}\textred{Semel in mense sub Ritu semid. juxta Rubr. 17.}\end{center}
\fancyhead[C]{\textsc{Officia Mensualia - Officium S. Teresiæ}}

%\begin{center}\textcolor{red}{\rule{\linewidth}{0.5pt}} \end{center}

\subsubsection*{\center \textsc{AD I VESPERAS}}
\subsubsection*{\textred{Capitulum. \hfill Sap. 7.}}
\lettrineThree{O}ptávi, et datus est mihi sensus; et invocávi, et venit in me spíritus sapiéntiæ; et præpósui illam regnis, et sédibus, et divítias nihil esse duxi in comparatióne illíus.

\subsubsection*{\center \textred{Hymnus.}}
\begin{center}\textred{In quo Temp. Pasch. nihil additur, neque mutatur.}\end{center}
\lettrine{R}egis supérni núntia\\
Domum patérnam déseris,\\
Terris Terésa bárbaris\\
Christum datúra, aut sánguinem.\\
\\
Sed te manet suávior\\
Mors, pœna poscit dúlcior:\\
Divíni amóris cúspide\\
In vulnus icta cóncides.\\
\\
O caritátis víctima!\\
Tu corda nostra cóncrema,\\
Tibíque Gentes créditas\\
Avérni ab igne líbera.\\
\\
Sit laus Patri cum Fílio,\\
Et Spíritu Paráclito,\\
Tibíque sancta Trínitas\\
Nunc, et per omne sǽculum.\\
Amen.

\begin{tightlist}{}
\item \textred{℣.} Oblivíscere pópulum tuum et domum patris tui. \textred{℟.} Et concupíscet Rex decórem tuum.
\item \textred{Ad Magnificat, Aña.} Quæsívi in sponsam mihi eam assúmere. Doctrix enim est disciplínæ Dei, et eléctrix óperum illíus.
\end{tightlist}

\subsubsection*{\center \textred{Oratio.}}
\lettrine{E}xáudi nos Deus salutáris noster, ut sicut de beátæ Terésiæ Vírgnis tuæ commemoratióne gaudémus: ita cœléstis ejus doctrínæ pábulo nutriámur, et piæ devotiónis erudiámur afféctu. Per Dóminum.

\subsubsection*{\center \textsc{AD MATUTINUM.}}
\begin{tightlist}{}
\item \textred{Invitatorium.} Regem Vírginum Dóminum, * Veníte adorémus.
\item \textred{Ps.} Veníte, exsultémus. 2.
\item \textred{Hymnus ut in j. Vesperis.}
\end{tightlist}

\subsubsection*{\center \textsc{IN PRIMO NOCTURNO.}}
\begin{tightlist}{}
\item \textred{Aña.} Quæ est ista, quæ ascéndit per desértum sicut vírgula fumi ex aromátibus myrrhæ et thuris?
\item \textred{Ps.} Dómine Dóminus noster, \textred{cum reliquis de Communi B.M.V.} [165]
\item \textred{Aña.} En diléctus meus lóquitur mihi: Surge, própera amíca mea, colúmba mea, et veni.
\item \textred{Aña.} Revértere, revértere Sunamítis: revértere, revértere, ut intueámur te.
\item \textred{℣.} Fiat cor meum immaculátum in justificatiónibus tuis. \textred{℟.} Ut non confúndar.
\item \textred{Lect. de Script. occurr. sed si desunt, legantur sequentes.}
\end{tightlist}

\subsubsection*{\center De libro Sapiéntiæ.}
\textred{\textbf{Lectio j}.} \hfill \textred{Cap. 7.}\\
\lettrineThree{O}ptávi, et datus est mihi sensus: et invocávi, et venit in me spíritus sapiéntiæ: et præpósui illam regnis, et sédibus, et divítias nihil esse duxi in comparatióne illíus. Nec comparávi illi lápidem pretiósum: quóniam omne aurum in comparatióne illíus aréna est exígua, et tamquam lutum æstimábitur argéntum in conspéctu illíus. Super salútem et spéciem diléxi illam, et propósui pro luce habére illam: quóniam enexstinguíbile est lumen illíus. Venérunt autem mihi ómnia bona páriter cum illa, et innumerábilis honéstas per manus illíus, et lætáta sum in ómnibus: quóniam antecedébat me ista sapiéntia, et ignorábam, quóniam horum ómnium mater est. Quam sine fictióne dídici, et sine invídia commúnico, et honestátem illíus non abscóndo. Infinítus enim thesáurus est homínibus: quo qui usi sunt, partícipes facti sunt amicítiæ Dei, propter disciplínæ dona commendáti.

\begin{tightlist}{}
\item \textred{℟.} Habitábit in solitúdine judícium, et justítia in Carmélo sedébit, et erit opus justítiæ pax, \textred{*} Et cultus justítiæ siléntium, et secúritas usque in sempitérnum. \textred{℣.} Et sedébit pópulus meus in pulchritúdine pacis, et in tabernáculis fidúciæ. \textred{E}t.
\end{tightlist}

\subsubsection*{\center \textred{Lectio ij.}}
\lettrine{O}mnium enim ártifex dócuit me sapiéntia. Est enim in illa spíritus intelligéntiæ, sanctus, únicus, múltiplex, subtílis, disértus, móbilis, incoinquinátus, certus, suávis, amans bonum, acútus, quem nihil vetat, benefáciens, humánus, benígnus, stábilis, certus, secúrus, omnem habens virtútem, ómnia prospíciens, et qui cápiat omnes spíritus: intelligíbilis, mundus, subtílis. Omnibus enim mobílibus mobílior est sapiéntia: attíngit autem ubíque propter suam mundítiam. Vapor est enim virtútis Dei, et emanátio quædam est claritátis omnipoténtis Dei sincéra: et ídeo nihil inquinátum in eam incúrrit. Candor est enim lucis ætérnæ, et spéculum sine mácula Dei majestátis, et imágo bonitátis illíus. Et cum sit una, ómnia potest, et in se pérmanens ómnia ínnovat, et per natiónes in ánimas sanctas se transfert, amícos Dei, et Prophétas constítuit.

\begin{tightlist}{}
\item \textred{℟.} Diffúsa est grátia in lábiis tuis: \textred{*} Proptérea benedíxit te Deus in ætérnum. \textred{℣.} Spécia tua, et pulchritúdine tua inténde, próspere procéde, et regna. \textred{P}roptérea.
\end{tightlist}

\subsubsection*{\center \textred{Lectio iij.}}
\lettrine{N}eminem enim díligit Deus, nisi eum, qui cum sapiéntia inhábitat. Est enim hæc speciósior sole, et super omnem dispositiónem stellárum, luci comparáta invenítur prior. Illi enim succédit nox, sapiéntiam autem non vincit malítia. Attíngit ergo a fine usque ad finem fórtiter, et dispónit ómnia suáviter. Hanc amávi et exquisívi a juventúte mea, et quæsívi sponsam mihi eam assúmere, et amátrix facta sum formæ illíus. Generositátem illíus gloríficat, contubérnium habens Dei: sed et ómnium Dóminus diléxit illam. Doctrix enim est disciplínæ Dei, et eléctrix óperum illíus. Et si divítiæ appetúntur in vita, quid sapiéntia locuplétius, quæ operátur ómnia? Si autem sensus operátur: quis horum, quæ sunt, magis, quam illa est ártifex? Et si justítiam quis díligit: labóres hujus magnas habent virtútes: sobrietátem enim, et prudéntiam docet, et justítiam, et virtútem, quibus utílius nihil est in vita homínibus.

\begin{tightlist}{}
\item \textred{℟.} Audi fília, et vide, et inclína aurem tuam; \textred{*} Quia deínceps ut vera sponsa meum zelábis honórem. \textred{℣.} Sponsábo te mihi in justítia, et judício, et in misericórdia, et miseratiónibus. \textred{Q}uia. Glória Patri. \textred{Q}uia.
\end{tightlist}

\subsubsection*{\center \textsc{IN SECUNDO NOCTURNO.}}
\begin{tightlist}{}
\item \textred{Aña.} Vidi Dóminum, qui dedit corónam decóris in cápite meo.
%\item \textred{Ps.} Eructávit. [169]
\item \textred{Aña.} Induit me Génitrix Dómini vestiméntis salútis, et datum est mihi, ut coopériar býssino splendénti, et cándido.
\item \textred{Aña.} Ego autem in Dómino gaudébo, et exsultábo in Deo Jesu meo.
\item \textred{℣.} Sitívit in te ánima mea. 
\item \textred{℟.} Quam multiplíciter tibi caro mea.
\end{tightlist}

\subsubsection*{\center Sermo sancti Máximi Epíscopi.}
\begin{center}\textred{Hom. in nat. S. Agnetis.\\
\textbf{Lectio iv}.}\end{center}
\lettrineThree{C}um in toto mundo virgíneus flos Maríæ immarcescíbiles corónas innéctat, et sceptrígeram pudóris aulam immaculáto consérvet afféctu, eoúsque intégritas perseverávit ad palmam, ut in puéllis trophǽum sanctitátis arríperet, et per vestígia matris Vírginis ad cœléstem thálamum perveníret. O quam imitábile virgínibus, Virgo, tui amóris tradidísti exémplum! O quam sanctum pectóribus dandum virgíneis tribuísti respónsum, contemnéndo mundi divítias, exsecrándo carnis concupiscéntias, et solam Christi pulchritúdinem adamándo! Accédite ad puéllam puéllæ, et díscite, qualis in péctore suo circa amórem Christi flammas accéndit. Díscite, vírgines, Christi amórem in ea fervéntem, et omnes mundi divítias velut stércora recusántem.

\begin{tightlist}{}
\item \textred{℟.} Dixit mihi Dóminus meus: Ne tímeas fília: \textred{*} Non rápiet te quisquam de manu mea. \textred{℣.} Neque mors, neque vita, neque creatúra ália póterit me separáre a caritáte Dei. \textred{N}on.
\end{tightlist}

\subsubsection*{\center \textred{Lectio v.}}
\lettrine{I}deo enim sanctórum per lítteras tradúntur exémpla, ut unusquísque pro qualitáte sexus modum sui propósiti téneat, et sicut per vírginem vírgines, ita per cónjugem cónjuges víncere sǽculum pro Christi amóre stúdeant, ac repugnatióne ánimi, mundi dúlcia recusántes, et ejus non recusántes austéra, ad delectatiónes perpétuas, et ad gáudia ætérna pertíngant. Hæc dicéntes, non vituperámus cónjugum casta commércia, sed vírginum perseverántiam adornámus. Non enim virgínitas malis rebus superponénda est, sed est rebus óptimis præferénda: quæ enim mundis múndior, et sanctis est sánctior, non índiget malárum rerum vituperatióne laudári.

\begin{tightlist}{}
\item \textred{℟.} Gaudens gaudébo in Dómino, et exsultábit ánima mea in Deo meo, quia índuit me vestiméntis salútis; \textred{*} Et induménto justítiæ circúmdedit me. \textred{℣.} Quasi sponsam decorátam coróna, et quasi sponsam ornátam monílibus suis. \textred{E}t.
\end{tightlist}

\subsubsection*{\center \textred{Lectio vj.}}
\lettrine{V}eni jam, Virgo, ad thálamum, quem quæsísti, ut sponsam te virgínitas immaculáta cómprobet Christi. Felix catérva vírginum, quibus stúdium est tuæ sequi confessiónis vestígia. Habébunt enim tecum corónam in cœlo, quæ tecum superavérunt inimícum in mundo. Nam sicut tibi vestígia Maríæ prosequénti participátio ætérnæ remuneratiónis appáruit: ita eis te imitántibus ætérna crédimus gáudia non negánda. Itaque o spléndida Christo Dei Fílio, et ómnibus Angelis, atque Archángelis grata! ut nostri meminísse dignéris, quibus póssumus précibus exorámus, ut nobis ille tríbuat indulgéntiam, qui tibi tuórum ómnium labórum trádidit palmam. 

\begin{tightlist}{}
\item \textred{℟.} Pulchra fácie, sed púlchrior fide beáta es Terésia virgo, réspuens mundum, lætáberis cum Angelis: \textred{*} Intercéde pro ómnibus nobis. \textred{℣.} Gaudeámus, et exsultémus, et demus glóriam ei, quia venérunt núptiæ Agni. \textred{I}ntercéde. Glória Patri. \textred{I}ntercéde.
\end{tightlist}

\subsubsection*{\center \textsc{IN TERTIO NOCTURNO.}}
\begin{tightlist}{}
\item \textred{Aña.} Diléctus meus mihi, et ego illi, qui páscitur inter lília.
%\item \textred{Ps.} Cantáte. [172]
\item \textred{Aña.} Dómine prævenísti eam in benedictiónibus dulcédinis: posuísti in cápite ejus corónam de lápide pretióso.
\item \textred{Aña.} Os suum apéruit sapiéntiæ, et lex cleméntiæ in lingua ejus. 
\item \textred{℣.} Deus præcínxit me virtúte. \textred{℟.} Et pósuit immaculátam viam meam.
\end{tightlist}

\subsubsection*{\center Lectio sancti Evangélii secúndum Matthǽum.}
\begin{center}\textred{\textbf{Lectio vij}. \hfill Cap. 11.}\end{center}
\lettrine{I}n illo témpore: Respóndens Jesus dixit: Confíteor tibi Pater, Dómine cœli et terræ, quia abscondísti hæc a sapiéntibus, et prudéntibus, et revelásti ea párvulis. Et réliqua.

\begin{center}Homilía sancti Augustíni Epíscopi.\\
\textred{Serm. 9. de Verbis Dom.}\end{center}
\lettrineThree{N}on immérito ille vas electiónis cum ingénti lætítia dixit: Non sunt condígnæ passiónes hujus témporis ad futúram glóriam, quæ revelábitur in nobis. Ecce unde illud jugum suáve est, et sárcina levis, et si angústa est paucis eligéntibus, fácilis tamen ómnibus diligéntibus. Dixit Psalmísta: Propter verba labiórum tuórum ego custodívi vias duras. Sed quæ dura sunt laborántibus, éadem ipsa mitéscunt amántibus.

\begin{tightlist}{}
\item \textred{℟.} Habitáre fecit eam Deus stérilem in domo: \textred{*} Matrem filiórum lætántem. \textred{℣.} Cibávit eam pane vitæ, et intelléctus, et aqua sapiéntiæ salutáris potávit illam. \textred{M}atrem.
\end{tightlist}

\subsubsection*{\center \textred{Lectio viij.}}
\lettrine{P}ropter quod ita divínæ pietátis dispensatióne actum est, ut intérior homo, qui renovátur de die in diem, non adhuc sub lege pósitus, sed jam sub grátia exonerátus sárcinis innumerabílium observatiónum, quod erat revéra grave jugum, sed duræ cervíci conveniénter impósitum, facilitáte símplicis fídei, et bonæ spei, et sanctæ caritátis, quidquid molestiárum exterióri hómini forínsecus intulísset ille princeps, qui missus est foras, interióri gáudio fíeret leve. Nihil enim tam fácile bonæ voluntáti est, quam ipsa sibi: et hæc súfficit Deo.

\begin{tightlist}{}
\item \textred{℟.} Tu glória Jerúsalem, tu lætítia Israel, tu honorificéntia pópuli nostri, quia fecísti viríliter, et confortátum est cor tuum: ídeo et manus Dómini confortábit te: \textred{*} Et eris benedícta in ætérnum. \textred{℣.} Surrexérunt fílii tui, et. beatíssimam prædicavérunt. \textred{E}t.
\end{tightlist}

\subsubsection*{\center \textred{Lectio ix.}}
\lettrine{Q}uantúmlibet ergo sǽviat iste mundus, veríssime Angeli, nato in carne Dómino clamavérunt: Glória in excélsis Deo, et in terra pax homínibus bonæ voluntátis; quia, qui natus erat, suáve jugum est, et sárcina levis. Et sicut dicit Apóstolus: Fidélis Deus, qui nos non sinit tentári supra id, quod póssumus ferre, sed facit cum tentatióne étiam éxitum, ut possímus sustinére.

\begin{tightlist}{}
\item \textred{℟.} Super salútem et spéciem diléxi sapiéntiam, et propósui pro luce habére illam: \textred{*} Quóniam inexstinguíbile est lumen illíus. \textred{℣.} Lætáta sum in ómnibus, quóniam antecedébat me ista sapiéntia. \textred{Q}uóniam. Glória Patri. \textred{Q}uóniam.
\item Te Deum laudámus. 6.
\item \textred{Sacerdotalis ℣.} Veni, Sponsa Christi, áccipe corónam. \textred{℟.} Quam tibi Dóminus præparávit in ætérnum.
\end{tightlist}

\subsubsection*{\center \textsc{AD LAUDES.}}
\begin{center}\textred{et per Horas, Aña.}\end{center}
\lettrine{H}æc est virgo sápiens, et una de número prudéntum. 

\begin{tightlist}{}
\item \textred{Ps. de Dominica.} 30.
\item \textred{Aña.} Date ei de fructu mánuum suárum, et laudent eam in portis ópera ejus.
\item \textred{Aña.} Sitívit in te Dómine ánima sponsæ tuæ, et lætáta est in his, quæ dicta sunt illi.
\item \textred{Aña.} Trahe me post te, in odórem currémus unguentórum tuórum: óleum effúsum nomen tuum.
\item \textred{Aña.} Lætáre Terésia in eo, qui te fecit: quóniam benedíxit tibi Dóminus, et ádipe fruménti sátiat te.
\end{tightlist}

\subsubsection*{\center \textred{Capitulum. \hfill Sap. 7.}}
\lettrineThree{O}ptávi, et datus est mihi sensus; et invocávi, et venit in me spíritus sapiéntiæ, et præpósui illam regnis et sédibus, et divítias nihil esse duxi in comparatióne illíus.

\subsubsection*{\center \textred{Hymnus.}}
\begin{center}\textred{In quo Temp. Pasch. nihil additur, neque mutatur.}\end{center}
\lettrine{F}elix dies, qua cándidæ\\
Instar colúmbæ, cœ́litum\\
Ad sacra templa spíritus\\
Se tránstulit Terésiæ;\\
\\
Sponsíque voces áudiit:\\
Veni soror de vértice\\
Carméli, ad Agni núptias;\\
Veni ad corónam glóriæ.\\
\\
Te sponse Jesu Vírginum\\
Beáti adórent órdines,\\
Et nuptiáli cántico\\
Laudent per omne sǽculum.\\
Amen.

\begin{tightlist}{}
\item \textred{℣.} Adjuvábit eam Deus vultu suo. \textred{℟.} Deus in médio ejus non commovébitur.
\item \textred{Ad Benedictus, Aña.} Dedit ei Dóminus sapiéntiam et prudéntiam multam nimis, et latitúdinem cordis, quasi arénam, quæ est in líttore maris.
\end{tightlist}

\subsubsection*{\center \textred{Oratio.}}
\lettrine{E}xáudi nos Deus salutáris noster, ut sicut de beátæ Terésiæ Vírginis tuæ commemoratióne gaudémus: ita cœléstis ejus doctrínæ pábulo nutriámur, et piæ devotiónis erudiámur afféctu. Per Dóminum nostrum.

\begin{tightlist}{}
\item \textred{Ad Horas Antiphonæ de Laudibus, Psalmi de Dominica, ad Primam tamen ut in Festis.}
\end{tightlist}

\subsubsection*{\center \textsc{AD TERTIAM.}}
\subsubsection*{\center \textred{Capitulum. \hfill Sap. 7.}}
\lettrine{O}ptávi, et datus est mihi sensus; et invocávi, et venit in me spíritus sapiéntiæ, et præpósui illam regnis et sédibus, et divítias nihil esse duxi in comparatióne illíus.

\begin{tightlist}{}
%\item \textred{Capit. Optávi, ut supra.}
\item \textred{℟.} Tibi dixit cor meum, \textred{*} Exquisívit te fácies mea. \textred{T}ibi dixit. \textred{℣.} Fáciem tuam Dómine requíram. \textred{E}xquisívit. Glória Patri. \textred{T}ibi.
\item \textred{℣.} Anima mea exsultábit in Dómino. \textred{℟.} Et delectábitur super salutári suo.
\end{tightlist}

\subsubsection*{\center \textsc{AD SEXTAM.}}
\subsubsection*{\textred{Capitulum. \hfill Sap. 7.}}
\lettrine{S}uper salútem et spéciem diléxi illam, et propósui pro luce habére illam, quóniam inexstinguíbile est lumen illíus.

\begin{tightlist}{}
\item \textred{℟.} Anima mea exsultábit \textred{*} In Dómino. \textred{A}nima mea. \textred{℣.} Et delectábitur super salutári suo. \textred{I}n. Glória Patri. \textred{A}nima.
\item \textred{℣.} Apud me orátio Deo vitæ meæ. \textred{℟.} Dicam Deo: suscéptor meus es.
\end{tightlist}

\subsubsection*{\center \textsc{AD NONAM.}}
\subsubsection*{\textred{Capitulum. \hfill Sap. 7.}}
\lettrine{V}enérunt autem mihi ómnia bona páriter cum illa, et innumerábilis honéstas per manus illíus, et lætáta sum in ómnibus: quóniam antecedébat me ista sapiéntia.

\begin{tightlist}{}
\item \textred{℟.} Apud me orátio \textred{*} Deo vitæ meæ. \textred{A}pud me. \textred{℣.} Dicam Deo: suscéptor meus es. \textred{D}eo. Glória Patri. \textred{A}pud.
\item \textred{℣.} Passer invénit sibi domum. \textred{℟.} Et turtur nidum sibi ubi ponat pullos suos.
\end{tightlist}

\subsubsection*{\center \textsc{AD VESPERAS.}}
\begin{tightlist}{}
\item \textred{Aña j. Laud. Psalmi de Comm. B.V.M. [177]}
%\item \textred{Capit. Hymn. et ℣. de j. Vesperis.}
\end{tightlist}

\subsubsection*{\textred{Capitulum. \hfill Sap. 7.}}
\lettrineThree{O}ptávi, et datus est mihi sensus; et invocávi, et venit in me spíritus sapiéntiæ; et præpósui illam regnis, et sédibus, et divítias nihil esse duxi in comparatióne illíus.

\subsubsection*{\center \textred{Hymnus.}}
\begin{center}\textred{In quo Temp. Pasch. nihil additur, neque mutatur.}\end{center}
\lettrine{R}egis supérni núntia\\
Domum patérnam déseris,\\
Terris Terésa bárbaris\\
Christum datúra, aut sánguinem.\\
\\
Sed te manet suávior\\
Mors, pœna poscit dúlcior:\\
Divíni amóris cúspide\\
In vulnus icta cóncides.\\
\\
O caritátis víctima!\\
Tu corda nostra cóncrema,\\
Tibíque Gentes créditas\\
Avérni ab igne líbera.\\
\\
Sit laus Patri cum Fílio,\\
Et Spíritu Paráclito,\\
Tibíque sancta Trínitas\\
Nunc, et per omne sǽculum.\\
Amen.

\begin{tightlist}{}
\item \textred{℣.} Oblivíscere pópulum tuum et domum patris tui. \textred{℟.} Et concupíscet Rex decórem tuum.
\item \textred{Ad Magnificat, Aña.} Sapiéntiam ejus enarrábunt gentes, et laudem ejus enuntiábit Ecclésia.
\item \textred{Tempore Paschali tres Psalmi cujuslibet Nocturni dicuntur sub una sola Aña; et in fine Antiphonar. Responsor. et ℣. additur} allelúja, \textred{cetera omnia ut supra.}
\end{tightlist}

\begin{center}\textcolor{red}{\rule{\linewidth}{0.5pt}} \end{center}

\subsection*{\center \textsc{Officium Sanctæ Mariæ Magdalenæ de Pazzis.}}
\subsubsection*{\center \textred{Virginis Ordinis Nostri.\\ Semiduplex}}
\begin{center}\textred{Semel in mense sub Ritu semid. juxta Rubr. 17.}\end{center}
\fancyhead[C]{\textsc{Officia Mensualia - Officium S. Mariæ Magdalenæ}}

%\begin{center}\textcolor{red}{\rule{\linewidth}{0.5pt}} \end{center}

\subsubsection*{\textred{Capitulum. \hfill Cant. 3.}}
\lettrineThree{I}n léctulo meo per noctes quæsívi quem díligit ánima mea: quæsívi illum, et non invéni. Surgam, et circuíbo civitátem: per vicos et platéas quæram quem díligit ánima mea.

\subsubsection*{\center \textred{Hymnus.}}
\lettrine{P}ange Carméli specióse vertex\\
Vírginis sanctæ memoránda gesta,\\
Inclytum cujus decus univérsus\\
Cóncinit orbis.\\
\\
Cándido velans scápulas amíctu,\\
Síngulos Christi pátitur dolóres:\\
Pérvigil cunctos súperat furéntis\\
Dǽmonis astus.\\
\\
Inter urtícas, rigidásque spinas\\
Innocens puro rútilat cruóre:\\
Cándida sacris mánibus Maríæ\\
Veste beátur.\\
\\
Cordis humáni pénetrat secréta,\\
Orbis evéntus mémorat futúros,\\
Imperat morbis, fébribus medétur,\\
Ulcera sanat.\\
\\
Próprium nunquam plácitum secúta,\\
In nono vitæ moritúra lustro,\\
Spíritum jussu cóhibens patérno,\\
Vincit agónem.\\
\\
In lacu quinos Erebi per annos\\
Sústinet larvas, móritur; nec unquam\\
Gnara quid fœdet sacra virginális\\
Lília cordis.\\
\\
Horret in tumba júvenem protérvum:\\
Cláruit vivens, moriénsque signis:\\
Integrum miros hódie profúndit\\
Corpus odóres.\\
\\
Laus Deo summo, Genitóque Patris,\\
Flámini sancto sit honor perénnis:\\
Et Deum trinum venerémur uno\\
Semper honóre.\\
Amen.

\begin{tightlist}{}
\item \textred{℣.} Ora pro nobis sancta María Magdaléna. \textred{℟.} Ut digni efficiámur promissiónibus Christi.
\item \textred{Ad Magnificat, Aña.} Aquæ multæ non potuérunt exstínguere caritátem ejus, nec flúmina óbruent illam; quia fortis est ut mors diléctio.
\end{tightlist}

\subsubsection*{\center \textred{Oratio.}}
\lettrine{D}eus virginitátis amátor, qui beátam Maríam Magdalénam vírginem, tuo amóre succénsam, cœléstibus donis decorásti; da, ut quam devóta commemoratióne venerámur, puritáte, et caritáte imitémur. Per Dóminum nostrum.

\subsubsection*{\center \textsc{AD MATUTINUM.}}
\begin{tightlist}{}
\item \textred{Invitat.} Jesum sponsum Vírginum, quem Magdaléna diléxit multum, * Veníte adorémus.
\item \textred{Ps.} Veníte, exsultémus. 2.
\end{tightlist}

\subsubsection*{\center \textred{Hymnus.}}
\begin{center}\textred{in quo Temp. Pach. nihil additur neque mutatur.}\end{center}
\lettrine{H}æc est Etrúsca Mágdala,\\
Magístra veri spíritus:\\
Partem prehéndit óptimam,\\
Non auferéndam sǽculis.\\
\\
Cum Rege sacras cœ́litum\\
Vovens decénnis núptias,\\
Exspéctat indelébili\\
Sponsum coléndum lámpade.\\
\\
Manúque dotem própria\\
Crucem, flagélla, vúlnera,\\
Cor, ánnulum, spinas, dapes\\
Dat Sponsus Eucharísticas.\\
\\
Scriptóris augústi manu\\
In corde descríptum suo,\\
Verbum caro factum est, legit,\\
Verbúmque grande intélligit.\\
\\
Mirábili sub schémate\\
Ignis, colúmnæ, flúminis,\\
Linguæ, colúmbæ, fláminis,\\
Sanctum recépit Spíritum.\\
\\
Amóris æstro pércita,\\
Compúlsat æra túrrium,\\
Ut totus orbis cónvolet,\\
Deúmque amórem díligat.\\
\\
Te Sponse Jesu Vírginum\\
Beáti adórent Ordines,\\
Et nuptiáli cántico\\
Laudent per omne sǽculum.\\
Amen.

\subsubsection*{\center \textsc{IN PRIMO NOCTURNO.}}
\begin{tightlist}{}
\item \textred{Aña.} Ex ore puéllæ perfecísti laudem Dómine, quia novit quam admirábile est nomen tuum.
\item \textred{Ps.} Dómine, Dóminus noster, \textred{cum reliquis de Comm. B.V.M.} [165]
\item \textred{Aña.} Dóminus prǽstitit sapiéntiam párvulæ desiderábilem super aurum et lápidem pretiósum multum.
\item \textred{Aña.} Accépit benedictiónem a Dómino, quia fuit ínnocens, et mundo corde.
\item \textred{℣.} Diffúsa est grátia in lábiis tuis. \textred{℟.} Proptérea benedíxit te Deus in ætérnum.
\item \textred{Lectiones de Scriptura occurr. si desint, legantur sequentes.}
\end{tightlist}

\subsubsection*{\center De Cánticis Canticórum.}
\textred{\textbf{Lectio j}. \hfill Cant. 3.}
\lettrineThree{I}n léctulo meo per noctes quæsívi quem díligit ánima mea: quæsívi illum, et non invéni. Surgam, et circuíbo civitátem: per vicos et platéas quæram quem díligit ánima mea: quæsívi illum, et non invéni. Invenérunt me vígiles, qui custódiunt civitátem: Num quem díligit ánima mea vidístis? Páululum cum pertransíssem eos, invéni quem díligit ánima mea: ténui eum, nec dimíttam, donec introdúcam illum in domum matris meæ, et in cubículum genitrícis meæ.

\begin{tightlist}{}
\item \textred{℟.} Flos campi diléctus meus Jesus, cujus adspéctu lætántur Angeli: \textred{*} Cujus pulchritúdine pascúntur beáti, cujus odóre revivíscunt mórtui. \textred{℣.} Sicut malus inter ligna silvárum, sic diléctus meus inter fílias. \textred{C}ujus.
\end{tightlist}

\subsubsection*{\center \textred{Lectio ij. \hfill Cant. 8.}}
\lettrine{Q}uis mihi det te fratrem meum sugéntem úbera matris meæ, ut invéniam te foris, et deósculer te, et jam me nemo despíciat? Apprehéndam te, et ducam in domum matris meæ: ibi me docébis, et dabo tibi póculum ex vino condíto, et mustum malórum granatórum meórum. Læva ejus sub cápite meo, et déxtera illíus amplexábitur me. Adjúro vos, fíliæ Jerúsalem, ne suscitétis, neque evigiláre faciátis diléctam, donec ipsa velit.

\begin{tightlist}{}
\item \textred{℟.} Sponsa Jesu ardet caritáte, fervet desidério, \textred{*} Amóris æstu rapta quo se vertat nescit. \textred{℣.} Fulcíte me flóribus, stipáte me malis, quia amóre lángueo. \textred{A}móris.
\end{tightlist}

\subsubsection*{\center \textred{Lectio iij.}}
\lettrine{Q}uæ est ista, quæ ascéndit de desérto delíciis áffluens inníxa super diléctum suum? Sub árbore malo suscitávi te: ibi corrúpta est mater tua, ibi violáta est génitrix tua. Pone me ut signáculum super cor tuum, ut signáculum super bráchium tuum: quia fortis est ut mors diléctio, dura sicut inférnus æmulátio: lámpades ejus lámpades ignis, atque flammárum. Aquæ multæ non potuérunt exstínguere caritátem, nec flúmina óbruent illam. Si déderit homo omnem substántiam domus suæ pro dilectióne, quasi nihil despíciet eam.

\begin{tightlist}{}
\item \textred{℟.} O felícem Magdalénæ ánimam caritátis ádeo succénsam ardóribus, \textred{*} Ut præ amóre languens defíceret. \textred{℣.} Déficit cor meum, et caro mea, Deus cordis mei, et pars mea Deus in ætérnum. \textred{U}t. Glória Patri. \textred{U}t.
\end{tightlist}

\subsubsection*{\center \textsc{IN SECUNDO NOCTURNO.}}
\begin{tightlist}{}
\item \textred{Aña.} Magdaléna oblíta est pópulum suum, et domum patris sui, et concupívit Rex decórem illíus.
%\item \textred{Ps.} Eructávit. [169]
\item \textred{Aña.} Deus, adjútor in tribulatiónibus, fuit cum Magdaléna: ídeo non tímuit, cum transferrétur in lacum leónum.
\item \textred{Aña.} Magdalénæ fundaménta in Monte Carmélo, et diléxit eam Dóminus super ómnia tabernácula Jacob.
\item \textred{℣.} Spécie tua et pulchritúdine tua. \textred{℟.} Inténde, próspere procéde, et regna.
\end{tightlist}

\subsubsection*{\center Sermo sancti Ambrósii Epíscopi.}
\begin{center}\textred{\textbf{Lectio iv}. \hfill Lib. 1. de Virg.}\end{center}
\lettrineThree{E}go licet laudatiónem non suscéperim Virginitátis, sed expressiónem, ad rem tamen pertinére árbitror, ut quæ sit ei pátria, et quis auctor appáreat. Ac prius ubi sit pátria definiámus. Si enim ibi est pátria, ubi genitále domicílium: in cœlo profécto est pátria castitátis: ítaque hic ádvena, ibi íncola est. Quid autem est cástitas virginális, nisi expers contagiónis intégritas? atque ejus auctórem, quem póssumus existimáre, nisi immaculátum Dei Fílium, cujus caro non vidit corruptiónem, divínitas non est expérta contagiónem?

\begin{tightlist}{}
\item \textred{℟.} Cum adhuc júnior essem quæsívi sapiéntiam in oratióne mea: \textred{*} Animam meam diréxi ad illam, et excépi eam. \textred{℣.} Spíritus tuus bonus dedúxit me in terram rectam, et docuísti me facere voluntátem tuam. \textred{A}nimam.
\end{tightlist}

\subsubsection*{\center \textred{Lectio v.}}
\lettrine{V}idéte ígitur quanta virginitátis mérita sint. Christus ante Vírginem, Christus ex Vírgine: a Patre quidem natus ante sǽcula, sed ex Vírgine natus ob sǽcula. Illud natúræ suæ, hoc nostræ utilitátis est. Illud erat semper, hoc esse vóluit. Spectáte áliud virginitátis méritum, Christus virgo vírginis Sponsus, et si dici potest, Christus virgíneæ castitátis: virgínitas enim Christi, non virginitátis est Christus. Virgo est ergo quæ nupsit, Virgo, quæ nos suo útero portávit, Virgo, quæ génuit, Virgo, quæ próprio lacte nutrívit. Sed jam a matre descendámus ad fílias.

\begin{tightlist}{}
\item \textred{℟.} Certámen forte dedit illi ut vínceret, et sciret quóniam ómnium poténtior est sapiéntia, descendítque cum illa in fóveam, \textred{*} Et non derelíquit eam. \textred{℣.} Secúndum multitúdinem dolórum meórum consolatiónes tuæ lætificavérunt ánimam meam. \textred{E}t.
\end{tightlist}

\subsubsection*{\center \textred{Lectio vj.}}
\lettrine{D}e virgínibus, ínquit Apóstolus, præcéptum Dómini non hábeo. Si Doctor géntium non hábuit, habére quis pótuit? Et præcéptum quidem non hábuit, sed hábuit exémplum. Non enim imperári potest virgínitas; sed optári. Nam quæ super nos sunt, in voto magis, quam in magistério sunt. Sed volo vos, inquit, sine sollicitúdine esse: nam qui sine uxóre est, sollícitus est quæ Dómini sunt, quómodo pláceat Deo: et virgo cógitat, quæ Dómini sunt, ut sit sancta córpore et spíritu. Non ego quidem dissuádeo matrimónium, sed virginitátis attéxo benefícium.

\begin{tightlist}{}
\item \textred{℟.} Absit mihi gloriári nisi in Cruce Dómini nostri Jesu Christi, \textred{*} Per quem mihi mundus crucifíxus est, et ego mundo. \textred{℣.} Vivo ego, jam non ego, vivit vero in me Christus. \textred{P}er. Glória Patri. \textred{P}er.
\end{tightlist}

\subsubsection*{\center \textsc{IN TERTIO NOCTURNO.}}
\begin{tightlist}{}
\item \textred{Aña.} Introívit in átria sancta Religiónis, et Dómino hóstiam própriæ voluntátis óbtulit.
%\item \textred{Ps.} Cantáte. [172]
\item \textred{Aña.} De manu dǽmonum liberávit Magdalénam Dóminus, et vidérunt omnes pópuli glóriam sanctificatiónis ejus.
\item \textred{Aña.} Mons Carméli exsúltet in conspéctu Dómini, qui mirabília fecit in Magdaléna.
\item \textred{℣.} Adjuvábit eam Deus vultu suo. \textred{℟.} Deus in médio ejus, non commovébitur.
\end{tightlist}

\subsubsection*{\center Lectio sancti Evangélii secúndum Matthǽum.}
\begin{center}\textred{\textbf{Lectio vij}. \hfill Cap. 11.}\end{center}
\lettrine{I}n illo témpore: Respóndens Jesus dixit: Confíteor tibi Pater, Dómine cœli et terræ, quia abscondísti hæc a sapiéntibus et prudéntibus, et revelásti ea párvulis. Et réliqua.

\begin{center}Homilía sancti Augustíni Epíscopi.\\
\textred{Serm. 9. de Verb. Dom.}\end{center}
\lettrineThree{O}mnia sæva, et immánia prorsus facília et prope nulla éfficit amor. Quanto ergo cértius ac facílius ad veram beatitúdinem cáritas facit, quod ad misériam, quantum pótuit, cupíditas fecit! quam fácile tolerátur quǽlibet advérsitas temporális, ut ætérna pœna vitétur, ætérna réquies comparétur. Non immérito ille vas electiónis cum ingénti lætítia dixit: Non sunt condígnæ passiónes hujus témporis ad superventúram glóriam, quæ revelábitur in nobis. Ecce unde illud jugum suáve est, et sárcina levis, et si angústa est paucis eligéntibus, fácilis tamen ómnibus diligéntibus. Dixit Psalmísta: Propter verba labiórum tuórum ego custodívi vias duras.

\begin{tightlist}{}
\item \textred{℟.} Diléxi decórem domus tuæ, et locum habitatiónis glóriæ tuæ, \textred{*} Elégi abjécta esse in domo tua, magis quam habitáre in tabernáculis peccatórum. \textred{℣.} Quam dilécta tabernácula tua Dómine virtútum, concupíscit, et déficit ánima mea in átria sanctæ Religiónis. \textred{E}légi.
\end{tightlist}

\subsubsection*{\center \textred{Lectio viij.}}
Sed quæ dura sunt laborántibus, éadem ipsa mitéscunt amántibus. Propter quod ita divínæ pietátis dispensatióne actum est, ut intérior homo, qui renovátur de die in diem, non adhuc sub lege pósitus, sed jam sub grátia exonerátus sárcinis innumerabílium observatiónum, quod erat revéra grave jugum, sed duræ cervíci conveniénter impósitum, facilitáte símplicis fídei, et bonæ spei et sanctæ caritátis, quidquid molestiárum exterióri hómini forínsecus intulísset ille princeps, qui missus est foras, interióri gáudio fíeret leve.

\begin{tightlist}{}
\item \textred{℟.} Fascículus myrrhæ diléctus meus mihi, inter úbera mea commorábitur: \textred{*} Pósui eum ut signáculum, super cor meum, et super bráchium meum. \textred{℣.} Adjúro vos fíliæ Jerúsalem, si invenéritis diléctum meum, ut nuntiétis ei, quia amóre lángueo. \textred{P}ósui.
\end{tightlist}

\subsubsection*{\center \textred{Lectio ix.}}
Nihil enim tam fácile bonæ voluntáti est quam ipsa sibi, et hæc súfficit Deo. Quantúmlibet ergo sǽviat iste mundus, veríssime Angeli, nato in carne Dómino, clamavérunt: Glória in excélsis Deo, et in terra pax homínibus bonæ voluntátis. Quia, qui natus erat, suáve jugum est, et sárcina levis. Et sicut dicit Apóstolus: Fidélis Deus, qui nos non sinit tentári supra id, quod póssumus ferre, sed facit cum tentatióne étiam éxitum, ut possímus sustinére.

\begin{tightlist}{}
\item \textred{℟.} Egredímini fíliæ Sion, et vidéte Sponsam in diadémate, quo coronávit eam Sponsus suus: \textred{*} In die desponsatiónis illíus, et lætítiæ cordis ejus. \textred{℣.} Collocávit super caput sponsæ suæ corónam spíneam, quæ pósita fuit super caput suum. \textred{I}n. Glória Patri. \textred{I}n.
\item Te Deum laudámus. 6.
\item \textred{Sacerdotalis ℣.} Ora pro nobis sancta María Magdaléna. \textred{℟.} Ut digni efficiámur promissiónibus Christi, allelúja.
\end{tightlist}

\subsubsection*{\center \textsc{AD LAUDES.}}
\begin{center}\textred{et per Horas, Aña.}\end{center}
\lettrine{M}irábilis in Magdaléna Dóminus, quam e vulnéribus suis decórem índuit, et advérsus dǽmonum insídias fortitúdine præcínxit.

\begin{tightlist}{}
\item \textred{Ps. de Dominica.} 30.
\item \textred{Aña.} Relíctis sǽculi pompis, introívit in conspéctu Dei Magdaléna, quam ánnulo suo subarrhávit Dóminus.
\item \textred{Aña.} In terra desérta tentatiónum appáruit Magdaléna, ut vidéret virtútem et glóriam tuam.
\item \textred{Aña.} Patrem, et Fílium cum sancto Spíritu puérula benedicébat, dum sanctæ Trinitátis fidem párvulos docébat.
\item \textred{Aña.} In týmpano et choro, in chordis, et órgano omnis spíritus laudet Dóminum, cujus in Magdaléna nomen exaltátum est.
\end{tightlist}

\subsubsection*{\textred{Capitulum. \hfill Cant. 3.}}
\lettrineThree{I}n léctulo meo per noctes quæsívi quem díligit ánima mea: quæsívi illum, et non invéni. Surgam, et circuíbo civitátem: per vicos et platéas quæram quem díligit ánima mea.

\subsubsection*{\center \textred{Hymnus.}}
\lettrine{C}orporis nexus óbitu rescíndens,\\
Plena virtútum, méritis onústa,\\
Scandit ætérnas, comitánte Sponso,\\
Mágdala sedes.\\
\\
Abdito ventris látitans recéssu,\\
Non gravat matris grémium, nec unquam\\
Críminis pressit própriam molésto\\
Póndere mentem.\\
\\
Obtegit nudos, docet imperítos,\\
Cárceres visit, juvat univérsos:\\
Plúrimas fusis ánimas ab orco\\
Flétibus arcet.\\
\\
Déserens mundum, crucifíxa Christo,\\
Corpus afflíxit précibus, caténis\\
Férreis, flagris, fame, felle, spinis,\\
Frígore, flamma.\\
\\
Cum suis Horas cécinit Patrónis,\\
Exstases, raptus pátitur frequénter,\\
Osculo Sponsi recreáta, Jesum\\
Portat in ulnis.\\
\\
Nunc opem nobis miseráta confer,\\
Qui tuos læti cánimus triúmphos;\\
Erudi mentem, refluátque amóris\\
Flamma perénnis.\\
\\
Laus Deo summo, Genitóque Patris,\\
Flámini sancto sit honor perénnis:\\
Et Deum Trinum venerémur uno\\
Semper honóre.\\
Amen.

\begin{tightlist}{}
\item \textred{℣.} Elégit eam Deus, et præelégit eam. \textred{℟.} Et habitáre facit in tabernáculo suo.
\item \textred{Ad Benedictus, Aña.} Magdaléna Virgo óptimam partem elégit, quæ non auferétur ab ea: puríssima in vita, adhuc post mortem ab omni corruptióne manet immúnis.
\end{tightlist}

\subsubsection*{\center \textred{Oratio.}}
\lettrine{D}eus Virginitátis amátor, qui beátam Maríam Magdalénam Vírginem, tuo amóre succénsam, cœléstibus donis decorásti: da, ut quam devóta commemoratióne venerámur, puritáte et caritáte imitémur. Per Dóminum nostrum.

\begin{tightlist}{}
\item \textred{Ad Horas Antiphonæ de Laudibus, Psalmi de Dominica, ad Primam tamen ut in Festis.}
\end{tightlist}

\subsubsection*{\center \textsc{AD TERTIAM.}}
\subsubsection*{\textred{Capitulum. \hfill Cant. 3.}}
\lettrine{I}n léctulo meo per noctes quæsívi quem díligit ánima mea: quæsívi illum, et non invéni. Surgam, et circuíbo civitátem: per vicos et platéas quæram quem díligit ánima mea.

\begin{tightlist}{}
%\item \textred{Capit. In léctulo meo.}
\item \textred{℟.} Diffúsa est grátia \textred{*} In lábiis tuis. \textred{D}iffúsa. \textred{℣.} Proptérea benedíxit te Deus in ætérnum. \textred{I}n. Glória Patri. \textred{D}iffúsa.
\item \textred{℣.} Spécie tua et pulchritúdine tua. \textred{℟.} Inténde, próspere procéde, et regna.
\end{tightlist}

\subsubsection*{\center \textsc{AD SEXTAM.}}
\subsubsection*{\textred{Capitulum. \hfill Cant. 3.}}
\lettrine{I}nvenérunt me vígiles, qui custódiunt civitátem: num quem díligit ánima mea vidístis? Páululum cum pertransíssem eos, invéni quem díligit ánima mea: ténui eum nec dimíttam.

\begin{tightlist}{}
\item \textred{℟.} Spécie tua \textred{*} Et pulchritúdine tua. \textred{S}pécie. \textred{℣.} Inténde, próspere procéde, et regna. \textred{E}t. Glória Patri. \textred{S}pécie.
\item \textred{℣.} Adjuvábit eam Deus vultu suo. \textred{℟.} Deus in médio ejus, non commovébitur.
\end{tightlist}

\subsubsection*{\center \textsc{AD NONAM.}}
\subsubsection*{\textred{Capitulum. \hfill Cant. 3.}}
\lettrine{A}djúro vos, fíliæ Jerúsalem, si invenéritis diléctum meum, ut nuntiétis ei, quia amóre lángueo.

\begin{tightlist}{}
\item \textred{℟.} Adjuvábit eam \textred{*} Deus vultu suo. \textred{A}djuvábit. \textred{℣.} Deus in médio ejus, non commovébitur. \textred{D}eus. Glória Patri. \textred{A}djuvábit.
\item \textred{℣.} Elégit eam Deus, et præelégit eam. \textred{℟.} Et habitáre eam facit in tabernáculo suo.
\end{tightlist}

\subsubsection*{\center \textsc{AD II. VESPERAS.}}
\begin{tightlist}{}
\item \textred{Aña j. Laud. Psalmi de Comm. B.V.M. [177]}
%\item \textred{Capit. Hymn. et ℣. ut in j. Vesperis.}
\end{tightlist}

\subsubsection*{\textred{Capitulum. \hfill Philipp. 2.}}
\lettrineThree{I}n léctulo meo per noctes quæsívi quem díligit ánima mea: quæsívi illum, et non invéni. Surgam, et circuíbo civitátem: per vicos et platéas quæram quem díligit ánima mea.

\subsubsection*{\center \textred{Hymnus.}}
\lettrine{P}ange Carméli specióse vertex\\
Vírginis sanctæ memoránda gesta,\\
Inclytum cujus decus univérsus\\
Cóncinit orbis.\\
\\
Cándido velans scápulas amíctu,\\
Síngulos Christi pátitur dolóres:\\
Pérvigil cunctos súperat furéntis\\
Dǽmonis astus.\\
\\
Inter urtícas, rigidásque spinas\\
Innocens puro rútilat cruóre:\\
Cándida sacris mánibus Maríæ\\
Veste beátur.\\
\\
Cordis humáni pénetrat secréta,\\
Orbis evéntus mémorat futúros,\\
Imperat morbis, fébribus medétur,\\
Ulcera sanat.\\
\\
Próprium nunquam plácitum secúta,\\
In nono vitæ moritúra lustro,\\
Spíritum jussu cóhibens patérno,\\
Vincit agónem.\\
\\
In lacu quinos Erebi per annos\\
Sústinet larvas, móritur; nec unquam\\
Gnara quid fœdet sacra virginális\\
Lília cordis.\\
\\
Horret in tumba júvenem protérvum:\\
Cláruit vivens, moriénsque signis:\\
Integrum miros hódie profúndit\\
Corpus odóres.\\
\\
Laus Deo summo, Genitóque Patris,\\
Flámini sancto sit honor perénnis:\\
Et Deum trinum venerémur uno\\
Semper honóre.\\
Amen.

\begin{tightlist}{}
\item \textred{℣.} Ora pro nobis sancta María Magdaléna. \textred{℟.} Ut digni efficiámur promissiónibus Christi.
\item \textred{Ad Magnificat, Aña.} Veni Sponsa de Líbano, veni, coronáberis de cápite Amana, de vértice Sanir, et Hermon, de cubílibus leónum, et móntibus pardórum.
\item \textred{Completorium de Dominica.}
\item \textred{Tempore Paschali tres Psalmi cujuslibet Nocturni dicuntur sub una sola Antiphona, et in fine Añarum etc. additur} Allelúja. \textred{Cetera omnia ut supra.}
\end{tightlist}
\begin{center}\textcolor{red}{\rule{\linewidth}{0.5pt}} \end{center}
\clearpage
%\begin{center}\textcolor{red}{\rule{\linewidth}{0.5pt}} \end{center}
\section*{\center \textred{\textsc{OFFICIA VOTIVA PER ANNUM.}}}
\subsubsection*{\center \textred{\textsc{Dicenda ad libitum}}}

\lettrineThree{O}\textred{fficia hæc Votiva per annum ritu Semiduplici pro singulis Hebdomadæ Feriis a Summo Pontifice Leone XIII. per Decretum \textsc{Urbis et Orbis} Sacrorum Rituum Congregationis die 5 Julii 1883 generali Indulto concessa, et ad Nostrum Ordinem extensa Decr. 20 Martii 1884, persolvi possunt, juxta idem Indultum, loco Officiorum Ferialium occurrentium, quocumque anni tempore; exceptis Feria IV. Cinerum, Feriis totius temporis Passionis, ac Sacri Adventus a die 17 ad 24 Decembris inclusive.}

\begin{center}\textcolor{red}{\rule{\linewidth}{0.5pt}} \end{center}

\subsection*{\center \textred{\textsc{PRO FERIA II.}}}
\subsection*{\center \textsc{\textred{OFFICIUM VOTIVUM}\\
De Sanctis Angelis.}}
\subsubsection*{\center \textred{Semiduplex.}}
\fancyhead[C]{\textsc{Officia Votiva - Feria II - De Angelis.}}


\subsubsection*{\center \textsc{AD VESPERAS.}}
\subsubsection*{\textred{Capitulum. \hfill Exod. 23.}}
\lettrineThree{E}cce ego mittam Angelum meum, qui præcédat te, et custódiat in via, et introdúcat in locum quem parávi. Obsérva eum, et audi vocem ejus.

\subsubsection*{\center \textred{Hymnus.}}
\lettrine{T}ibi, Christe, splendor Patris,\\
Vita, virtus córdium,\\
In conspéctu Angelórum,\\
Votis, voce psállimus:\\
Alternántes concrepándo\\
Melos damus vócibus.\\
\\
Collaudámus venerántes\\
Omnes cæli mílites;\\
Sed præcípue primátem\\
Cæléstis exércitus\\
Michaélem, in virtúte,\\
Conteréntem zábulum.\\
\\
Quo custóde, procul pelle\\
Rex Christe piíssime,\\
Omne nefas inimíci,\\
Mundo corde, et córpore,\\
Paradíso redde tuo\\
Nos sola cleméntia.\\
\\
Glóriam Patri melódis\\
Personémus vócibus,\\
Glóriam Christo canámus,\\
Glóriam Paráclito;\\
Qui Trinus, et unus Deus\\
Exstat ante sǽcula.\\
Amen.

\begin{tightlist}{}
\item \textred{℣.} In conspéctu Angelórum psallam tibi, Deus meus. \textred{℟.} Adorábo ad templum sanctum tuum, et confitébor nómini tuo.
\item \textred{Ad Magnif. Aña.} Omnes sunt administratórii spíritus, in ministérium missi propter eos, qui hereditátem cápiunt salútis.
\end{tightlist}

\subsubsection*{\center \textred{Oratio.}}
\lettrine{D}eus, qui miro órdine Angelórum ministéria, hominúmque dispénsas: concéde propítius; ut a quibus tibi ministrántibus in cœlo semper assístitur, ab his in terra vita nostra muniátur. Per Dóminum.

\subsubsection*{\center \textsc{AD MATUTINUM.}}
\fancyhead[C]{\textsc{Officia Votiva - Feria II - De Angelis.}}
\begin{tightlist}{}
\item \textred{Invit.} Regem Angelórum Dóminum, * Veníte, adorémus.
\item \textred{Ps.} Veníte, exsultémus. 2.
\item \textred{Hymnus.} Tibi Christe, \textred{ut supra.}
\end{tightlist}

\subsubsection*{\center \textsc{IN PRIMO NOCTURNO.}}
\fancyhead[C]{\textsc{Officia Votiva - Feria II - De Angelis.}}
\begin{tightlist}{}
\item \textred{Aña.} Multa magnália de Michaéle Archángelo, qui fortis in prælio fecit victóriam.
\item \textred{Ps. 8.} Dómine, Dóminus. [27] 
\item \textred{Aña.} Gábriel Archángelus locútus est Maríæ dicens: Ecce concípies in útero et páries Fílium et vocábis nomen ejus Jesum.
\item \textred{Ps. 10.} In Dómino confído. [29]
\item \textred{Aña.} Benedícite Deum cœli, dixit Angelus Ráphael, et coram ómnibus vivéntibus confitémini ei, quia fecit vobíscum misericórdiam suam.
\item \textred{Ps. 14.} Dómine, quis. [30]
\item \textred{℣.} In conspéctu Angelórum psallam tibi, Deus meus. \textred{℟.} Adorábo ad templum sanctum tuum, et confitébor nómini tuo.
\item \textred{Lectiones de Scriptura occurr. sed si desint, legantur sequentes.}
\end{tightlist}

\subsubsection*{\center De libro Exodi.}
\begin{center}\textred{\textbf{Lectio j}. \hfill Cap. 23.}\end{center}
\lettrineThree{E}cce ego mittam Angelum meum, qui præcédat te, et custódiat in via, et introdúcat in locum quem parávi.
Obsérva eum, et audi vocem ejus, nec contemnéndum putes: quia non dimíttet cum peccáveris, et est nomen meum in illo.
Quod si audíeris vocem ejus, et féceris ómnia quæ loquor, inimícus ero inimícis tuis, et afflígam affligéntes te. Præcedétque te Angelus meus.

\begin{tightlist}{}
\item \textred{℟.} Venit Míchael Archángelus cum multitúdine Angelórum, cui trádidit Deus ánimas sanctórum, \textred{*} Ut perdúcat eas in paradísum exsultatiónis. \textred{℣.} Emítte, Dómine, Spíritum sanctum tuum de cœlis, Spíritum sapiéntiæ et intelléctus. \textred{U}t.
\end{tightlist}

\subsubsection*{\center Ex Zacharía Prophéta.}
\begin{center}\textred{\textbf{Lectio ij}. \hfill Cap. 1.}\end{center}
\lettrine{F}actum est verbum Dómini ad Zacharíam, fílium Barachíæ fílii Addo, prophétam, dicens: Vidi per noctem, et ecce vir ascéndens super equum rufum, et ipse stabat inter myrtéta, quæ erant in profúndo: et post eum equi rufi, várii, et albi.
Et dixi: Quid sunt isti, Dómine mi? Et dixit ad me Angelus, qui loquebátur in me: Ego osténdam tibi quid sint hæc.
Et respóndit vir, qui stabat inter myrtéta, et dixit: Isti sunt, quos misit Dóminus ut perámbulent terram.
Et respondérunt Angelo Dómini, qui stabat inter myrtéta, et dixérunt: Perambulávimus terram, et ecce omnis terra habitátur, et quiéscit. Hæc dicit Dóminus.

\begin{tightlist}{}
\item \textred{℟.} Respóndit Angelus Dómini, et dixit: Dómine exercítuum, \textred{*} Usquequo tu non miseréberis Jerúsalem et úrbium Juda, quibus irátus es? \textred{℣.} Iste enim septuagésimus annus est. \textred{U}squequo.
\end{tightlist}

\subsubsection*{\center \textred{Lectio iij. \hfill Cap. 2.}}
\lettrine{E}t levávi óculos meos, et vidi: et ecce vir, et in manu ejus funículus mensórum.
Et dixi: Quo tu vadis? Et dixit ad me: Ut métiar Jerúsalem, et vídeam quanta sit latitúdo ejus, et quanta longitúdo ejus.
Et ecce Angelus, qui loquebátur in me, egrediebátur, et Angelus álius egrediebátur in occúrsum ejus.
Et dixit ad eum: Curre, lóquere ad púerum istum, dicens: Absque muro habitábitur Jerúsalem præ multitúdine hóminum, et jumentórum in médio ejus.
Et ego ero ei, ait Dóminus, murus ignis in circúitu: et in glória ero in médio ejus. Hæc dicit Dóminus.

\begin{tightlist}{}
\item \textred{℟.} In conspéctu Géntium nolíte timére: vos enim in córdibus vestris adoráte et timéte Dóminum: \textred{*} Angelus enim ejus vobíscum est. \textred{℣.} Stetit Angelus juxta aram templi, habens thuríbulum áureum in manu sua. \textred{A}ngelus. Glória Patri. \textred{A}ngelus.
\end{tightlist}

\subsubsection*{\center \textsc{IN SECUNDO NOCTURNO.}}
%\fancyhead[C]{\textsc{Officia Votiva - Feria II - De Angelis.}}
\begin{tightlist}{}
\item \textred{Aña.} Ascéndit fumus arómatum in conspéctu Dómini de manu Angeli.
\item \textred{Ps. 18.} Cœli enárrant. [166]
\item \textred{Aña.} Immíttet Angelus Dómini in circúitu timéntium eum, et erípiet eos.
\item \textred{Ps. 23.} Dómini est terra. [167]
\item \textred{Aña.} Laudémus Dóminum, quem laudant Angeli, quem Chérubim et Séraphim, Sanctus, Sanctus, Sanctus, proclámant.
\item \textred{Ps. 33.} Benedícam Dóminum. [7]
\item \textred{℣.} Stétit Angelus juxta aram templi. \textred{℟.} Habens thuríbulum áureum in manu sua.
\end{tightlist}

\subsubsection*{\center De Sermóne sancti Gregórii Papæ.}
\begin{center}\textred{Ex. Homilia 34. in Evang.\\
\textbf{Lectio iv}.}\end{center}
\lettrineThree{Q}uia Angelórum nómina interpretándo perstrínximus, súperest, ut ipsa officiórum vocábula bréviter exsequámur. Virtútes étenim vocántur illi nimírum spíritus, per quos signa et mirácula frequéntius fiunt. Potestátes étiam vocántur hi, qui hoc poténtius céteris in suo órdine percepérunt, ut eórum ditióni virtútes advérsæ subjéctæ sint, quorum potestáte refrænántur, nec corda hóminum tantum tentáre præváleant, quantum volunt.

\begin{tightlist}{}
\item \textred{℟.} Omnes Angeli stabant in circúitu throni, et seniórum, et quátuor animálium: \textred{*} Et cecidérunt in conspéctu throni in fácies suas, et adoravérunt Deum. \textred{℣.} Adoráte Dóminum omnes Angeli ejus. \textred{E}t.
\end{tightlist}

\subsubsection*{\center \textred{Lectio v.}}
\lettrine{P}rincipátus étiam vocántur, qui ipsis quoque bonis Angelórum spirítibus præsunt, qui subjéctis áliis dum quæquæ sunt agénda dispónunt, eis ad explénda divína ministéria principántur. Dominatiónes autem vocántur, qui étiam potestátes principátuum dissimilitúdine alta transcéndunt. Nam principári est inter réliquos priórem exsístere: dominári vero est étiam subjéctos quosque possidére. Throni quoque illa ágmina sunt vocáta, quibus ad exercénda judícia semper Deus omnípotens prǽsidet. Unde et per Psalmístam dícitur: Sedes super thronum, qui júdicas æquitátem.

\begin{tightlist}{}
\item \textred{℟.} Stetit Angelus juxta aram templi, habens thuríbulum áureum in manu sua, et data sunt ei incénsa multa: \textred{*} Et ascéndit fumus arómatum de manu Angeli in conspéctu Dómini. \textred{℣.} In conspéctu Angelórum psallam tibi: adorábo ad templum sanctum tuum, et confitébor nómini tuo, Dómine. \textred{E}t.
\end{tightlist}

\subsubsection*{\center \textred{Lectio vj.}}
\lettrine{C}hérubim quoque plenitúdo sciéntiæ dícitur; et sublimióra illa ágmina idcírco Chérubim vocáta sunt, quia tanto perfectióri sciéntia plena sunt, quanto claritátem Dei vicínius contemplántur. Séraphim étiam vocántur illa spirítuum sanctórum ágmina, quæ ex singulári propinquitáte Conditóris sui incomparábili ardent amóre. Séraphim namque ardéntes vel incendéntes vocántur; quorum profécto flamma amor est: quia quo subtílius claritátem divinitátis ejus adspíciunt, eo valídius in ejus amóre flamméscunt.

\begin{tightlist}{}
\item \textred{℟.} In conspéctu Angelórum psallam tibi, et adorábo ad templum sanctum tuum: \textred{*} Et confitébor nómini tuo, Dómine. \textred{℣.} Super misericórdia tua et veritáte tua: quóniam magnificásti super nos nomen sanctum tuum. \textred{E}t. Glória Patri. \textred{E}t confitébor.
\end{tightlist}

\subsubsection*{\center \textsc{IN TERTIO NOCTURNO.}}
\begin{tightlist}{}
\item \textred{Aña.} Misit Dóminus Angelum suum, qui percússit omnem virum robústum, et bellatórem, et príncipem exércitus regis Assyriórum.
\item \textred{Ps. 95.} Cantáte Dómino. [172]
\item \textred{Aña.} Adoráte Dóminum omnes Angeli ejus: audívit, et lætáta est Sion.
\item \textred{Ps. 96.} Dóminus regnávit. [173]
\item \textred{Aña.} Benedícite Dómino omnes Angeli ejus, poténtes virtúte, faciéntes verbum illíus, ad audiéndam vocem sermónum ejus.
\item \textred{Ps. 102} Bénedic, ánima. 278.
\item \textred{℣.} Ascendit fumus arómatum in conspéctu Dómini. \textred{℟.} De manu Angeli.
\end{tightlist}

\subsubsection*{\center Léctio sancti Evangélii secúndum Joánnem.}
\begin{center}\textred{\textbf{Lectio vij}. \hfill Cap. 1.}\end{center}
\lettrine{I}n illo témpore: Vidit Jesus Nathánael veniéntem ad se, et dicit de eo: Ecce vere Israelíta, in quo dolus non est. Et réliqua.

\begin{center}Homilía sancti Bernárdi Abbátis.\\
\textred{In Ps. Qui habitat.\\
Sermo 11.}\end{center}
\lettrineThree{Q}uæ vero sunt sanctórum Angelórum viæ? Profécto illæ, quas Unigénitus enarrávit, dicens: Vidébitis Angelos ascendéntes et descendéntes super Fílium hóminis. Ascénsio ígitur et descénsio viæ illórum; ascénsio propter se, descénsio, vel pótius condescénsio propter nos. Sic beáti illi spíritus ascéndunt per contemplatiónem Dei, descéndunt per compassiónem tui, ut custódiant te in ómnibus viis tuis. Ascéndunt ad vultum ejus, descéndunt ad nutum ejus: quóniam Angelis suis mandávit de te. Nec tamen vel descendéndo visióne glóriæ fraudántur, quia semper vident fáciem Patris.

\begin{tightlist}{}
\item \textred{℟.} Angelus Dómini descéndit cum Azaría et sóciis ejus in fornácem, et excússit flammam ignis de fornáce: \textred{*} Et non tétigit eos omníno ignis, neque contristávit. \textred{℣.} Benedíctus Deus eórum, qui misit Angelum suum, et éruit servos suos, qui credidérunt in eum. \textred{E}t.
\end{tightlist}

\subsubsection*{\center \textred{Lectio viij.}}
\lettrine{C}um ascéndunt ad contemplatiónem, inquírunt veritátem, de qua et desiderándo satiántur, et satiándo desíderant. Cum vero descéndunt, fáciunt nobíscum misericórdiam, ut custódiant nos in ómnibus viis nostris. Administratórii enim spíritus sunt, missi in ministérium propter nos. Plane minístri nostri, non dómini nostri. Et in hoc Unigéniti formam imitántur, qui non venit ministrári, sed ministráre, qui stetit inter discípulos tamquam qui minístrat. Fructus angelicárum viárum, quod ad ipsos spectat, sua ipsórum beatitúdo, et obediéntia est caritátis: quod autem ad nos pértinet, inde quidem obténtus divínæ grátiæ, hinc vero custódia viæ nostræ: síquidem Angelis suis mandávit de te, ut custódiant te in ómnibus indigéntiis tuis.

\begin{tightlist}{}
\item \textred{℟.} Tu, Dómine, qui misísti Angelum tuum sub Ezechía, rege Juda, et interfecísti de castris Senácherib centum octogínta quinque míllia: \textred{*} Et nunc Dominátor cœlórum mitte Angelum tuum bonum ante nos, in timóre et tremóre magnitúdinis bráchii tui. \textred{℣.} Ut métuant qui cum blasphémia véniunt advérsus sanctum pópulum tuum. \textred{E}t.
\end{tightlist}

\subsubsection*{\center \textred{Lectio ix.}}
\lettrine{I}nterim ergo mandávit Angelis suis Deus, non quidem ut amóveant te a viis tuis, sed in eis ipsis te custódiant, et quasi per vias suas, vias tuas dírigant in vias ipsíus. Quonam modo, inquis? Nempe ut quod Angelus ex sola caritáte púrius agit, ut própria saltem necessitáte compúlsus et admónitus, descéndas, et condescéndas; exhibére scílicet próximo misericórdiam: et rursum cum eódem Angelo levans desidéria tua, tota ánimi cupiditáte ascéndere stúdeas ad summa et perpétuam Veritátem. Hinc enim monémur leváre corda nostra cum mánibus; hinc audímus quotídie: Sursum corda. Hinc quoque negligéntes argúimur, et dícitur nobis: Fílii hóminum, úsquequo gravi corde? ut quid dilígitis vanitátem et quǽritis mendácium? Exonerátum enim cor et leve levátur magis, ut quærat et díligat veritátem.

\begin{tightlist}{}
\item \textred{℟.} Te sanctum Dóminum in excélsis laudant omnes Angeli, dicéntes: \textred{*} Te decet laus, et honor Dómine. \textred{℣.} Chérubim quoque, ac Séraphim, Sanctus proclámant, et omnis cœlicus ordo, dicéntes. \textred{T}e. Glória Patri. \textred{T}e.
\item Te Deum laudámus. 6.
\item \textred{Sacerd. ℣.} Sancti Angeli, Custódes nostri, deféndite nos in prǽlio. \textred{℟.} Ut non pereámus in treméndo judício.
\end{tightlist}

\subsubsection*{\center \textsc{AD LAUDES.}}
\begin{center}\textred{et per Horas, Aña.}\end{center}
\lettrine{A}ngeli Dómini Dóminum benedícite in ætérnum.

\begin{tightlist}{}
\item \textred{Ps. de Dominica.} 30.
\item \textred{Aña.} Angelis suis Deus mandávit de te, ut custódiant te in ómnibus viis tuis.
\item \textred{Aña.} Angeli eórum semper vident fáciem Patris mei, qui est in cœlis.
\item \textred{Aña.} Laudáte Deum, omnes Angeli ejus, laudáte eum, omnes virtútes ejus.
\item \textred{Aña.} Angeli, Archángeli, Throni et Dominatiónes, Principátus et Potestátes, Virtútes cœlórum, laudáte Dóminum de cœlis, allelúja.
\end{tightlist}

\subsubsection*{\textred{Capitulum. \hfill Exod. 23.}}
\lettrineThree{E}cce ego mittam Angelum meum, qui præcédat te, et custódiat in via, et introdúcat in locum quem parávi. Obsérva eum, et audi vocem ejus.

\subsubsection*{\center \textred{Hymnus.}}
\lettrine{C}hriste, sanctórum decus Angelórum,\\
Rector humáni géneris, et auctor,\\
Nobis ætérnum tríbue benígnus\\
Scándere cœlum.\\
\\
Angelum pacis Míchael ad istam\\
Cœ́litus mitti rogitámus aulam:\\
Nobis ut crebro veniénte, crescant\\
Próspera cuncta.\\
\\
Angelus fortis Gábriel, ut hostem\\
Pellat antíquum, vólitet ab alto,\\
Sǽpius templum véniens ad istud\\
Vísere nostrum.\\
\\
Angelum nobis médicum salútis\\
Mite de cœlis Ráphael, ut omnes\\
Sanet ægrótos, paritérque nostros\\
Dírigat actus.\\
\\
Hinc Dei nostri Génitrix María,\\
Totus et nobis chorus Angelórum\\
Semper assístat, simul et beáta\\
Cóncio tota.\\
\\
Præstet hoc nobis Déitas beáta\\
Patris, ac Nati, paritérque sancti\\
Spíritus cujus réboat in omni\\
Glória mundo.\\
Amen.

\begin{tightlist}{}
\item \textred{℣.} Adoráte Deum. \textred{℟.} Omnes Angeli ejus.
\item \textred{Ad Benedictus, Aña.} Angelus, qui loquebátur in me, et suscitávit me, quasi virum qui suscitátur a somno suo.
\end{tightlist}

\subsubsection*{\center \textred{Oratio.}}
\lettrine{D}eus, qui miro órdine Angelórum ministéria, hominúmque dispénsas: concéde propítius; ut a quibus tibi ministrántibus in cœlo semper assístitur, ab his in terra vita nostra muniátur. Per Dóminum.

\begin{tightlist}{}
\item \textred{Ad Horas Añæ de Laudibus, Psalmi de Dominica, ad Primam tamen ut in Festis.}
\end{tightlist}

\subsubsection*{\center \textsc{AD TERTIAM.}}
\subsubsection*{\textred{Capitulum. \hfill Exod. 23.}}
\lettrine{E}cce ego mittam Angelum meum, qui præcédat te, et custódiat in via, et introdúcat in locum quem parávi. Obsérva eum, et audi vocem ejus.

\begin{tightlist}{}
\item \textred{℟.} In conspéctu Angelórum \textred{*} Psallam tibi, Deus meus. \textred{I}n. \textred{℣.} Adorábo ad templum sanctum tuum, et confitébor nómini tuo. \textred{P}sallam. Glória Patri. \textred{I}n.
\item \textred{℣.} Stetit Angelus juxta aram templi. \textred{℟.} Habens thuríbulum áureum in manu sua.
\end{tightlist}

\subsubsection*{\center \textsc{AD SEXTAM.}}
\subsubsection*{\textred{Capitulum. \hfill Apoc. 5.}}
\lettrine{A}udívi vocem Angelórum multórum in circúitu throni, et animálium, et seniórum: et erat númerus eórum míllia míllium, voce magna dicéntium: Salus Deo nostro.

\begin{tightlist}{}
\item \textred{℟.} Stetit Angelus \textred{*} Juxta aram templi. \textred{S}tetit. \textred{℣.} Habens thuríbulum áureum in manu sua. \textred{J}uxta. Glória Patri. \textred{S}tetit.
\item \textred{℣.} Ascéndit fumus arómatum in conspéctu Dómini. \textred{℟.} De manu Angeli.
\end{tightlist}

\subsubsection*{\center \textsc{AD NONAM.}}
\subsubsection*{\textred{Capitulum. \hfill Apoc. 12.}}
\lettrine{F}actum est prǽlium magnum in cœlo: Míchael et Angeli ejus præliabántur cum dracóne, et draco pugnábat, et ángeli ejus: et non prævaluérunt, neque locus invéntus est eórum ámplius in cœlo.

\begin{tightlist}{}
\item \textred{℟.} Ascéndit fumus arómatum \textred{*} In conspéctu Dómini. \textred{A}scéndit. \textred{℣.} De manu Angeli. \textred{I}n. Glória Patri. \textred{A}scéndit.
\item \textred{℣.} Adoráte Deum. \textred{℟.} Omnes Angeli ejus.
\end{tightlist}

\subsubsection*{\center \textsc{AD II. VESPERAS.}}
\begin{tightlist}{}
\item \textred{Aña.} Angeli Dómini, \textred{prima de Laudibus, Ps. de Dominica et loco ultimi, Ps. 137.} Confitébor tibi, Dómine. 200.
%\item \textred{Capit. Hymnus et ℣. ut in j. Vesperis.}
\end{tightlist}

\subsubsection*{\textred{Capitulum. \hfill Exod. 23.}}
\lettrineThree{E}cce ego mittam Angelum meum, qui præcédat te, et custódiat in via, et introdúcat in locum quem parávi. Obsérva eum, et audi vocem ejus.

\subsubsection*{\center \textred{Hymnus.}}
\lettrine{T}ibi, Christe, splendor Patris,\\
Vita, virtus córdium,\\
In conspéctu Angelórum\\
Votis, voce psállimus:\\
Alternantes concrepándo\\
Melos damus vócibus.\\
\\
Collaudámus venerántes\\
Omnes Cæli mílites;\\
Sed præcípue Primátem\\
Cæléstis exércitus\\
Michaélem, in virtúte\\
Conteréntem zábulum.\\
\\
Quo custóde, procul pelle\\
Rex, Christe, piíssime,\\
Omne nefas inimíci,\\
Mundo corde et córpore,\\
Paradíso redde tuo\\
Nos sola cleméntia.\\
\\
Glóriam Patri melódis\\
Personémus vócibus,\\
Glóriam Christo canámus,\\
Glóriam Paráclito:\\
Qui Trinus et unus Deus\\
Exstat ante sǽcula.\\
Amen.

\begin{tightlist}{}
\item \textred{℣.} In conspéctu Angelórum psallam tibi, Deus meus. \textred{℟.} Adorábo ad templum sanctum tuum, et confitébor nómini tuo.
\item \textred{Ad Magnif. Aña.} Sancti Angeli, qui Deo in cœlis semper assístitis, deféndite nos in prǽlio, ut non pereámus in treméndo judício.
\end{tightlist}

\begin{center}\textcolor{red}{\rule{\linewidth}{0.5pt}} \end{center}

\subsection*{\center \textred{\textsc{PRO FERIA III.}}}
\subsection*{\center \textsc{\textred{OFFICIUM VOTIVUM}\\ De Sanctis Apostolis.}}
\subsubsection*{\center \textred{Semiduplex}}

\begin{tightlist}{}
\item \textred{Omnia de Communi Apostolorum [3] juxta diversitatem temporis, præter sequentia.}
\end{tightlist}

\subsubsection*{\center \textred{Oratio.}}
\lettrine{D}eus, qui nos per beátos Apóstolos tuos ad agnitiónem tui nóminis veníre tribuísti: da nobis eórum glóriam sempitérnam et proficiéndo celebráre, et celebrándo profícere. Per Dóminum.

\subsubsection*{\center \textsc{IN PRIMO NOCTURNO.}}
\fancyhead[C]{\textsc{Officia Votiva - Feria III - De Apostolis.}}
\begin{tightlist}{}
\item \textred{Lectiones de Scriptura occurrente. In Quadragesima vero:} Sic nos exístimet homo, \textred{de Communi.} [9]
\end{tightlist}

\subsubsection*{\center \textsc{IN SECUNDO NOCTURNO.}}

\subsubsection*{\center De Sermóne sancti Augustíni Epíscopi.}
\begin{center}\textred{Serm. 43. De Sanctis.\\
\textbf{Lectio iv}.}\end{center}
\lettrineThree{I}sti sunt lux mundi: quia per ipsos lumen fídei, et veræ sciéntæ primum Dóminus huic mundo trádidit, et ab errórum ac peccatórum ténebris gentes et pópulos éruit. Isti sunt sal terræ: quia per eos terrígenæ condiméntum sapóris vitæ ætérnæ percepérunt, ut restríngerent carnis lascíviam, et a putrédine peccatórum et vitiórum vérmibus servaréntur illǽsi. Isti sunt lápides pretiósi, quos in fundaménto cœléstis ædifícii pósitos Joánnes in Apocalýpsi sua descríbit: quia prædicátio eórum Ecclésiæ fundaménta locávit. Unde et Paulus ait: Vos estis cives Sanctórum et doméstici Dei, superædificáti super fundaméntum Apostolórum et Prophetárum.

\begin{tightlist}{}
\item \textred{℟.} Vidi conjúnctos viros habéntes spléndidas vestes: et Angelus Dómini locútus est ad me dicens: \textred{*} Isti sunt viri sancti, facti amíci Dei.
\textred{℣.} Vidi Angelum Dei fortem volántem per médium cœli, voce magna clamántem et dicéntem. \textred{I}sti.
\end{tightlist}

\subsubsection*{\center \textred{Lectio v.}}
\lettrine{H}oc ergo considerántes, fratres caríssimi, quæ tanti duces pópuli nos docuérunt, studeámus factis implére, quæ illi jussérunt. Discámus eórum exémplo mundi divítias contémnere, præséntis sǽculi voluptátem non amáre, regnum cœléste desideráre, Christo nihil præpónere, sed ejus mandátis in ómnibus obedíre, paupertátem præséntium rerum dilígere, virtútum divítias habére, sapiéntiæ thesáurum appétere, spiritáles delícias quǽrere, nulli invidére, sed omnes hómines dilígere, amícos in Deo et inimícos propter Deum; quia hæc est vera et sola diléctio.

\begin{tightlist}{}
\item \textred{℟.} Isti sunt triumphatóres, et amíci Dei, qui contemnéntes jussa príncipum, meruérunt prǽmia ætérna: \textred{*} Modo coronántur, et accípiunt palmam.
\textred{℣.} Isti sunt, qui venérunt ex magna tribulatióne, et lavérunt stolas suas in sánguine Agni. \textred{M}odo
\end{tightlist}

\subsubsection*{\center \textred{Lectio vj.}}
\lettrine{I}sti ergo Príncipes nostri, in amóre Dei perfectíssimi, et in proximórum dilectióne impléti, hinc potuérunt mundi ímpetum víncere, et cruéntum sǽculum domáre, quia nihil amavérunt in omni re, nisi Dei voluntátem. Ita et nos, fratres, diligámus in ómnibus rebus Dei fácere voluntátem, et creatórem nostrum amémus in semetípso, creatúram vero in conditóre suo, et sic habébimus ordinatíssimam caritátem: quia Deus cáritas est, et qui díligit hanc dilectiónem, díligit Deum. Et si ita dilexérimus, díligit nos ipse Deus, díligunt nos et sancti Apóstoli júdices nostri, et exórant pro nobis, ut in universáli judício Christi cum ipsis in perpétuum coronémur.

\begin{tightlist}{}
\item \textred{℟.} Fuérunt sine queréla ante Dóminum, et ab ínvicem non sunt separáti: \textred{*} Cálicem Dómini bibérunt, et amíci Dei facti sunt.
\item \textred{℣.} Tradidérunt córpora sua propter Deum ad supplícia: ídeo coronántur, et accípiunt palmam. \textred{C}álicem.
\end{tightlist}

\subsubsection*{\center \textsc{IN TERTIO NOCTURNO.}}
\subsubsection*{\center Léctio sancti Evangélii secúndum Matthǽum.}
\begin{center}\textred{\textbf{Lectio vij}. \hfill Cap. 19.}\end{center}
\lettrine{I}n illo témpore: Dixit Petrus ad Jesum: Ecce nos relíquimus ómnia, et secúti sumus te: quid ergo erit nobis? Et réliqua.

\begin{center}Homilía sancti Hilárii Epíscopi.\\
\textred{Comm. in Matth. cap. 20.}\end{center}
\lettrineThree{A}póstoli dictis Dómini hæc reddunt, reliquísse se ómnia, et cum ipso esse. Quibus Dóminus, cum séderit in majestátis suæ sede, sessúros super sedes duódecim, ac tótidem tribus Israel judicatúros spopóndit; omnibúsque qui univérsa relíquerint propter nomen ejus, fructum céntupli prǽmii reservátum: multos autem ex novíssimis primos, et ex primis futúros novíssimos. Multa sunt quæ non sinunt nos símplici intelléctu dicta Evangélica suscípere.

\begin{tightlist}{}
\item \textred{℟.} Isti sunt, qui vivéntes in carne, plantavérunt Ecclésiam sánguine suo. Non sunt de terris córpora eórum separáta: \textred{*} Quorum mérita sunt in cælo glorióse cumuláta. \textred{℣.} Sancti per fidem vicérunt regna, operáti sunt justítiam, adépti sunt repromissiónes. \textred{Q}uorum.
\end{tightlist}

\subsubsection*{\center \textred{Lectio viij.}}
\lettrine{I}nterpósitis nonnúllis rebus, quæ ex natúra humáni sensus sibi contráriæ sunt, ratiónem quǽrere cœléstis intelligéntiæ admonémur. Apóstoli dicunt, et sequi se Christum, et se ómnia reliquísse. Quómodo ígitur fiunt tristes, et quómodo métuunt dicéntes, salvum esse néminem posse? Namque et ab áliis fíeri póterat, si quid fecíssent ipsi. Deínde cum fecíssent ipsi; quare metus, vel unde suscéptus est? Additur étiam in responsióne Dómini, hæc apud hómines impossibília, possibília apud Deum. Numquid apud hóminem impossibília erant, quæ et Apóstoli fecísse se gloriántur, et fecísse eos Dóminus agnóscit?

\begin{tightlist}{}
\item \textred{℟.} Isti sunt viri sancti, quos elégit Dóminus in caritáte non ficta, et dedit illis glóriam sempitérnam: \textred{*} Quorum doctrína fulget Ecclésia, ut sole luna. \textred{℣.} In omnem terram exívit sonus eórum, et in fines orbis terræ verba eórum. \textred{Q}uorum.
\end{tightlist}

\subsubsection*{\center \textred{Lectio ix.}}
\lettrine{O}mnis ítaque hic sermo est spiritális. Quid autem Dei virtúti tam possíbile est, quam ut per fidem salvet? ut per aquam regéneret? ut per crucem vincat? ut per Evangélia adóptet? ut per resurrectiónem ex morte vivíficet? Quibus audítis Apóstoli cito credunt, seséque ómnia reliquísse profiténtur. Sed hanc eórum obediéntiam cito Dóminus remunerátur, omnem difficultátem superióris quæstiónis absólvens, cum dicit: Vos, qui secúti estis me, in regeneratióne, judicábitis duódecim tribus Israel. Hæc enim illa regenerátio est, quam Apóstoli sunt secúti, quam lex indulgére non pótuit, quæ eos super duódecim thronos in judicándis duódecim tríbubus Israel, in duódecim Patriarchárum glória copulávit.

\begin{tightlist}{}
\item \textred{℟.} Cives Apostolórum, et doméstici Dei advenérunt hódie, \textred{*} Portántes facem, et illuminántes pátriam, dare pacem Géntibus, et liberáre pópulum Dómini. \textred{℣.} Audíte preces súpplicum, vitæ ætérnæ poscéntes prǽmia, qui fertis in dextris manípulos justítiæ, quique gaudéntes advenístis hódie. \textred{P}ortántes. Glória Patri. \textred{P}ortántes.
\item Te Deum laudámus. 6.
\item \textred{Sacerdotalis ℣.} Dedísti hereditátem. \textred{℟.} Timéntibus nomen tuum Dómine.
\end{tightlist}

\begin{center}\textcolor{red}{\rule{\linewidth}{0.5pt}} \end{center}

\subsection*{\center \textred{PRO FERIA IV.}}
\subsection*{\center \textsc{\textred{OFFICIUM VOTIVUM}\\ De S. Joseph,\\ Sponso B.M.V.}}
\subsubsection*{\center \textred{Confess. et Catholicæ Ecclesiæ Patrono ac Protectore Primario Ordinis Nostri.\\ Semiduplex.}}
\fancyhead[C]{\textsc{Officia Votiva - Feria IV - De S. Joseph.}}

\subsubsection*{\center \textsc{AD VESPERAS.}}
\subsubsection*{\textred{Capitulum. \hfill Gen. 49.}}
\lettrineThree{B}enedictiónes patris tui confortátæ sunt benedictiónibus patrum ejus, donec veníret desidérium cóllium æternórum: fiant in cápite Joseph, et in vértice Nazarǽi inter fratres suos.

\subsubsection*{\center \textred{Hymnus.}}
\lettrine{T}e, Joseph, célebrent ágmina cœ́litum:\\
Te cuncti résonent Christíadum chori,\\
Qui clarus méritis, junctus es ínclytæ\\
Casto fœ́dere Vírgini.\\
\\
Almo cum túmidam gérmine cónjugem\\
Admírans, dúbio tángeris ánxius,\\
Afflátu súperi Fláminis Angelus\\
Concéptum púerum docet.\\
\\
Tu natum Dóminum stringis, ad éxteras\\
Ægypti prófugum tu séqueris plagas;\\
Amíssum Sólymis, quæris, et ínvenis,\\
Miscens gáudia flétibus.\\
\\
Post mortem réliquos mors pia cónsecrat,\\
Palmámque eméritos glória súscipit:\\
Tu vivens, Súperis par, frúeris Deo,\\
Mira sorte beátior.\\
\\
Nobis, summa Trias, parce precántibus,\\
Da Joseph méritis sídera scándere:\\
Ut tandem líceat nos tibi pérpetim\\
Gratum prómere cánticum.\\
Amen.

\begin{tightlist}{}
\item \textred{℣.} Constítuit eum dóminum domus suæ. \textred{℟.} Et príncipem omnis possessiónis suæ.
\item \textred{Ad Magnif. Aña.} Cum esset desponsáta Mater Jesu María Joseph, ántequam convenírent, invénta est in útero habens de Spíritu sancto.
\end{tightlist}

\subsubsection*{\center \textred{Oratio.}}
\lettrine{D}eus, qui ineffábili providéntia beátum Joseph sanctíssimæ Genitrícis tuæ sponsum elígere dignátus es: præsta, quǽsumus; ut quem protectórem venerámur in terris, intercessórem habére mereámur in cœlis: Qui vivis et regnas.

\subsubsection*{\center \textsc{AD MATUTINUM.}}
\begin{tightlist}{}
\item \textred{Invit.} Laudémus Deum nostrum * In veneratióne beáti Joseph, protectóris nostri.
\item \textred{Ps.} Veníte, exsultémus. 2.
\item \textred{Hymnus.} Te, Joseph, \textred{ut supra.}
\end{tightlist}

\subsubsection*{\center \textsc{IN PRIMO NOCTURNO.}}
\begin{tightlist}{}
\item \textred{Aña.} Ascéndit Joseph a Galilǽa de civitáte Názareth in Judǽam, in civitátem David, quæ vocátur Béthlehem, ut profiterétur cum María.
\item \textred{Ps.} Beátus vir, \textred{cum reliquis de Communi unius Conf. Pont.} [64]
\item \textred{Aña.} Venérunt pastóres festinántes: et invenérunt Maríam, et Joseph, et infántem pósitum in præsépio.
\item \textred{Aña.} Ecce Angelus Dómini appáruit in somnis Joseph, dicens: Surge, et áccipe púerum, et matrem ejus, et fuge in Ægýptum.
\item \textred{℣.} Constítuit eum Dóminum domus suæ. \textred{℟.} Et príncipem omnis possessiónis suæ.
\item \textred{Lectiones de Scriptura occurrente; si vero desint, legantur sequentes.}
\end{tightlist}

\subsubsection*{\center De libro Génesis.}
\begin{center}\textred{\textbf{Lectio j}. \hfill Cap. 39.}\end{center}
\lettrineThree{I}gitur Joseph ductus est in Ægýptum, emítque eum Putíphar eunúchus Pharaónis, princeps exércitus, vir Ægýptius, de manu Ismaëlitárum, a quibus perdúctus erat.
Fuítque Dóminus cum eo, et erat vir in cunctis próspere agens: habitavítque in domo dómini sui,
qui óptime nóverat Dóminum esse cum eo, et ómnia, quæ gerébat, ab eo dírigi in manu illíus.
Invenítque Joseph grátiam coram dómino suo, et ministrábat ei: a quo præpósitus ómnibus gubernábat créditam sibi domum, et univérsa quæ ei trádita fúerant:
benedixítque Dóminus dómui Ægýptii propter Joseph, et multiplicávit tam in ǽdibus quam in agris cunctam ejus substántiam:
nec quidquam áliud nóverat, nisi panem quo vescebátur. Erat autem Joseph pulchra fácie, et decórus aspéctu.

\begin{tightlist}{}
\item \textred{℟.} Clamávit pópulus ad regem, aliménta petens: \textred{*} Quibus ille respóndit: Ite ad Joseph. \textred{℣.} Salus nostra in manu tua est: réspice nos tantum, et læti serviémus regi. \textred{Q}uibus.
\end{tightlist}

\subsubsection*{\center \textred{Lectio ij. \hfill Cap. 41.}}
\lettrine{P}lácuit Pharaóni consílium et cunctis minístris ejus:
locutúsque est ad eos: Num inveníre potérimus talem virum, qui spíritu Dei plenus sit?
Dixit ergo ad Joseph: Quia osténdit tibi Deus ómnia quæ locútus es, numquid sapientiórem et consímilem tui inveníre pótero?
Tu eris super domum meam, et ad tui oris impérium cunctus pópulus obédiet: uno tantum regni sólio te præcédam.
Dixítque rursus Phárao ad Joseph: Ecce, constítui te super univérsam terram Ægýpti.
Tulítque ánnulum de manu sua, et dedit eum in manu ejus: vestivítque eum stola býssina, et collo torquem áuream circumpósuit.
Fecítque eum ascéndere super currum suum secúndum, clamánte præcóne, ut omnes coram eo genu flécterent, et præpósitum esse scirent univérsæ terræ Ægýpti.

\begin{tightlist}{}
\item \textred{℟.} Fecit me Deus quasi patrem regis, et dóminum univérsæ domus ejus: \textred{*} Exaltávit me, ut salvos fáceret multos pópulos. \textred{℣.} Veníte ad me, et ego dabo vobis ómnia bona Ægypti, ut comedátis medúllam terræ. \textred{E}xaltávit.
\end{tightlist}

\subsubsection*{\center \textred{Lectio iij.}}
\lettrine{D}ixit quoque rex ad Joseph: Ego sum Phárao: absque tuo império non movébit quisquam manum aut pedem in omni terra Ægýpti.
Vertítque nomen ejus, et vocávit eum, lingua ægyptíaca, Salvatórem mundi. Dedítque illi uxórem Aseneth fíliam Putíphare sacerdótis Heliopóleos. Egréssus est ítaque Joseph ad terram Ægýpti
(trigínta autem annórum erat quando stetit in conspéctu regis Pharaónis), et circuívit omnes regiónes Ægýpti.
Venítque fertílitas septem annórum: et in manípulos redáctæ ségetes congregátæ sunt in hórrea Ægýpti.
Omnis étiam frugum abundántia in síngulis úrbibus cóndita est.
Tantáque fuit abundántia trítici, ut arénæ maris coæquarétur, et cópia mensúram excéderet.

\begin{tightlist}{}
\item \textred{℟.} Jam lætus móriar quia vidi fáciem tuam, et supérstitem te relínquo. Non sum fraudátus adspéctu tuo: \textred{*} Insuper osténdit mihi Dóminus semen tuum. \textred{℣.} Qui pascit me ab adolescéntia mea, benedícat púeris istis, et invocétur super eos nomen meum. \textred{I}nsuper. Glória Patri. \textred{I}nsuper.
\end{tightlist}

\subsubsection*{\center \textsc{IN SECUNDO NOCTURNO.}}
\begin{tightlist}{}
\item \textred{Aña.} Consúrgens Joseph, accépit púerum et matrem ejus nocte, et secéssit in Ægýptum: et erat ibi usque ad óbitum Heródis.
%\item \textred{Ps.} Cum invocárem. [67]
\item \textred{Aña.} Defúncto Heróde, Angelus Dómini appáruit in somnis Joseph in Ægýpto, dicens: Surge, et áccipe púerum et matrem ejus, et vade in terram Israel: defúncti sunt enim qui quærébant ánimam púeri.
\item \textred{Aña.} Accépit Joseph púerum et matrem ejus, et venit in terram Israel.
\item \textred{℣.} De ventre matris meæ tu es protéctor meus. \textred{℟.} In te confirmátus sum ex útero.
\end{tightlist}

\subsubsection*{\center Sermo sancti Bernárdi Abbátis.}
\begin{center}\textred{Homilia 2. super Missus.\\
\textbf{Lectio iv}.}\end{center}
\lettrineThree{D}esponsáta est María Joseph, vel pótius (sicut ponit Evangelísta) viro, cui nomen erat Joseph: virum nóminat, non quia marítus, sed quod homo virtútis erat: vel pótius quia juxta álium Evangelístam, non vir simplíciter, sed vir ejus dictus est, mérito appellátur, quod necessário putátur. Débuit ergo vir ejus appellári, quia necésse fuit et putári: sicut et pater Salvatóris non quidem esse, sed dici méruit, ut putarétur esse, dicénte hoc ipso Evangelísta: Et ipse Jesus erat incípiens quasi annórum trigínta, ut putabátur, fílius Joseph.

\begin{tightlist}{}
\item \textred{℟.} Dedísti mihi protectiónem salútis tuæ, et déxtera tua suscépit me: \textred{*} Protéctor meus, et cornu salútis meæ, et suscéptor meus. \textred{℣.} Ego protéctor tuus sum, et merces tua magna nimis. \textred{P}rotéctor.
\end{tightlist}

\subsubsection*{\center \textred{Lectio v.}}
\lettrine{N}on est dúbium, quin bonus et fidélis homo fúerit iste Joseph, cui Mater desponsáta est Salvatóris. Fidélis, inquam, servus et prudens, quem constítuit Dóminus suæ Matris solátium, suæ carnis nutrítium, solum dénique in terris magni consílii coadjutórem fidíssimum. Huc accédit, quod dícitur fuísse de domo David. Vere enim de domo David. Vere de régia stirpe descéndit vir iste Joseph, nóbilis génere, mente nobílior. Plane fílius David, non degénerans a patre suo David; prorsus, inquam, fílius David, non tantum carne, sed fide, sed sanctitáte, sed devotióne, quem tamquam álterum David Dóminus invénit secúndum cor suum, cui tuto commítteret sacratíssimum atque secretíssimum sui cordis arcánum: cui, tamquam álteri David, incérta et occúlta sapiéntiæ suæ manifestávit, et dedit illi non ignárum esse mystérii, quod nemo príncipum hujus sǽculi agnóvit.

\begin{tightlist}{}
\item \textred{℟.} Státuet fílios suos sub tégmine illíus, et sub ramis ejus morábitur: protegétur sub tégmine illíus a fervóre: \textred{*} Et in glória ejus requiéscet. \textred{℣.} Speráte in eo, omnis congregátio pópuli: effúndite coram illo corda vestra. \textred{E}t.
\end{tightlist}

\subsubsection*{\center \textred{Lectio vj.}}
\lettrine{C}ui dénique datum est, quod multi reges et prophétæ cum vellent vidére, non vidérunt, audíre et non audiérunt, non solum vidére et audíre, sed étiam portáre, dedúcere, amplécti, deosculári, et nutríre et custodíre. Non tantum autem Joseph, sed et María descendísse credénda est de domo David: alióquin non esset desponsáta viro de domo David, si non esset et ipsa de domo David. Ambo ígitur erant de domo David: sed in áltera compléta est véritas, quam jurávit Dóminus David, áltero tamen cónscio et teste adimplétæ promissiónis.

\begin{tightlist}{}
\item \textred{℟.} Si consístant advérsum me castra, non timébit cor meum: \textred{*} Si exsúrgat advérsum me prǽlium, in hoc ego sperábo. \textred{℣.} In te cantátio mea semper, quóniam tu adjútor fortis. \textred{S}i. Glória Patri. \textred{S}i.
\end{tightlist}

\subsubsection*{\center \textsc{IN TERTIO NOCTURNO.}}
\begin{tightlist}{}
\item \textred{Aña.} Audiens Joseph quod Archeláus regnáret in Judǽa pro Heróde, patre suo, tímuit illo ire.
%\item \textred{Ps.} Dómine, quis habitábit. [70]
\item \textred{Aña.} Admónitus in somnis, Joseph secéssit in partes Galilǽæ, et véniens habitávit in civitáte quæ vocátur Názareth; ut adimplerétur quod dictum est per Prophétas: Quóniam Nazarǽus vocábitur.
\item \textred{Aña.} Erat pater Jesu et mater mirántes super his quæ dicebántur de illo; et benedíxit illis Símeon.
\item \textred{℣.} Invocávi Dóminum, Patrem Dómini mei. \textred{℟.} Ut non derelínquat me in die tribulatiónis.
\end{tightlist}

\subsubsection*{\center Lectio sancti Evangélii secúndum Lucam.}
\begin{center}\textred{\textbf{Lectio vij}. \hfill Cap. 3.}\end{center}
\lettrine{I}n illo témpore: Factum est autem, cum baptizarétur omnis pópulus et Jesu baptizáto et oránte, apértum est cœlum. Et réliqua.

\begin{center}Homilía sancti Augustíni Epíscopi.\\
\textred{Lib. 23. contra Faust. cap. 7. 8.}\end{center}
\lettrineThree{S}ic de cœlo dictum est super aquam Jordánis: Hic est Fílius meus diléctus, in quo mihi bene complácui: quemádmodum dictum est et in monte. Neque enim quia et ibi de cœlo vox ipsa sónuit, Fílius Dei ante non fuit: quandóquidem ex útero Vírginis ille accépit formam servi, qui cum in forma Dei esset, non rapínam arbitrátus est, esse se æquálem Deo. Dénique idem Apóstolus Paulus álio loco apertíssime dicit: Cum autem venit plenitúdo témporis, misit Deus Fílium suum factum ex mulíere, factum sub lege, ut eos, qui sub lege erant, redímeret, ut filiórum adoptiónem reciperémus. Ipse ergo est Fílius Dei, qui et Dóminus David secúndum divinitátem, et idem ipse Fílius David ex sémine David secúndum carnem.

\begin{tightlist}{}
\item \textred{℟.} Joseph, fili David, noli timére accípere Maríam, cónjugem tuam; quod enim in ea natum est, de Spíritu sancto est: páriet autem fílium, \textred{*} Et vocábis nomen ejus Jesum. \textred{℣.} Ipse enim salvum fáciet pópulum suum a peccátis eórum. \textred{E}t.
\end{tightlist}

\subsubsection*{\center \textred{Lectio viij.}}
\lettrine{Q}uod si nobis crédere non prodésset, non hoc tam atténte idem Apóstolus Timótheo commendáret dicens: Memor esto Christum Jesum resurrexísse a mórtuis, ex sémine David, secúndum Evangélium meum. Quid ergo jam móveat sancti Evangélii sectatórem, quod sine concúbitu Joseph Christus natus ex Vírgine, fílius tamen David appellátur, cum generatiónum sériem non usque ad Maríam, sed usque ad Joseph Matthǽus Evangelísta perdúcat? Primo quia maríti ejus fúerat propter virílem sexum pótius honoránda persóna: neque enim quia concúbitu non permíxtus, ídeo non marítus, cum ipse Matthǽus narret ab Angelo Maríam cónjugem ipsíus appellátam, qui narrat, quod de Spíritu sancto concéperat.

\begin{tightlist}{}
\item \textred{℟.} Surge, et áccipe púerum et matrem ejus, et fuge in Ægýptum: \textred{*} Et esto ibi, usque dum dicam tibi. \textred{℣.} Ut adimplerétur quod dictum est a Dómino per Prophétam dicéntem: Ex Ægýpto vocávi fílium meum. \textred{E}t.
\end{tightlist}

\subsubsection*{\center \textred{Lectio ix.}}
\lettrine{C}um vero unus idémque narrátor utrúmque dicat, utrúmque comméndet, et virum Maríæ Joseph, et Christi Vírginem Matrem, et Christum ex sémine David, et Joseph in série progeneratórum Christi ex David: quid restat, nisi et Maríam non fuísse extráneam a cognatióne David, et eam Joseph cónjugem non frustra appellátam propter órdinem sexus et animórum confœderatiónem; et Joseph pótius propter dignitátem virílem ab órdine generatiónum illárum non fuísse separándum, ne hoc ipso viderétur ab illa fémina separátus, cui eum conjungébat mentis afféctus?

\begin{tightlist}{}
\item \textred{℟.} Valde honorándus est, et laudábilis nimis beátus Joseph, \textred{*} Quem constítuit Dóminus magni consílii coadjutórem fidíssimum. \textred{℣.} Fidélis servus et prudens, Vírginis Matris Sponsus, et Fílii Dei Nutrítius. \textred{Q}uem. Glória Patri. \textred{Q}uem.
\item Te Deum laudámus. 6.
\item \textred{Sacerdotalis ℣.} Ora pro nobis sancte Joseph. \textred{℟.} Ut digni efficiámur promissiónibus Christi.
\end{tightlist}

\subsubsection*{\center \textsc{AD LAUDES.}}
\begin{center}\textred{er per Horas, Aña.}\end{center}
\lettrine{J}acob autem génuit Joseph, virum Maríæ, de qua natus est Jesus, qui vocátur Christus.

\begin{tightlist}{}
\item \textred{Ps. de Dominica.} 30.
\item \textred{Aña.} Missus est Angelus Gábriel a Deo in civitátem Galilǽæ, cui nomen Názareth, ad Vírginem desponsátam viro, cui nomen erat Joseph.
\item \textred{Aña.} Ascéndit autem Joseph a Galilǽa de civitáte Názareth, in Judǽam, in civitátem David, quæ vocátur Béthlehem.
\item \textred{Aña.} Et venérunt festinántes, et invenérunt Maríam, et Joseph, et infántem pósitum in præsépio.
\item \textred{Aña.} Et ipse Jesus erat incípiens quasi annórum trigínta, ut putabátur fílius Joseph.
\end{tightlist}

\subsubsection*{\textred{Capitulum. \hfill Gen. 49.}}
\lettrineThree{B}enedictiónes patris tui confortátæ sunt benedictiónibus patrum ejus, donec veníret desidérium cóllium æternórum: fiant in cápite Joseph, et in vértice Nazarǽi inter fratres suos.

\subsubsection*{\center \textred{Hymnus.}}
\lettrine{C}œ́litum, Joseph, decus, atque nostræ\\
Certa spes vitæ, columénque mundi,\\
Quas tibi læti cánimus, benígnus\\
Súscipe laudes.\\
\\
Te Sator rerum státuit pudícæ\\
Vírginis sponsum, voluítque Verbi\\
Te patrem dici, dedit et minístrum\\
Esse salútis.\\
\\
Tu Redemptórem stábulo jacéntem,\\
Quem chorus Vatum cécinit futúrum,\\
Adspicis gaudens, humilísque natum\\
Numen adóras.\\
\\
Rex Deus regum, Dominátor orbis,\\
Cujus ad nutum tremit inferórum\\
Turba, cui pronus famulátur æther,\\
Se tibi subdit.\\
\\
Laus sit excélsæ Tríadi perénnis,\\
Quæ tibi præbens súperos honóres,\\
Det tuis nobis méritis beátæ\\
Gáudia vitæ.\\
Amen.

\begin{tightlist}{}
\item \textred{℣.} Dedísti mihi protectiónem salútis tuæ. \textred{℟.} Et déxtera tua suscépit me.
\item \textred{Ad Bened. Aña.} Joseph, fili David, noli timére accípere Maríam cónjugem tuam; quod enim in ea natum est, de Spíritu sancto est.
\end{tightlist}

\subsubsection*{\center \textred{Oratio.}}
\lettrine{D}eus, qui ineffábili providéntia beátum Joseph sanctíssimæ Genitrícis tuæ sponsum elígere dignátus es: præsta, quǽsumus; ut quem protectórem venerámur in terris, intercessórem habére mereámur in cœlis: Qui vivis.

\begin{tightlist}{}
\item \textred{Ad Horas Añæ de Laudibus, Psalmi de Dominica, ad Primam tamen ut in Festis.}
\end{tightlist}

\subsubsection*{\center \textsc{AD TERTIAM.}}
\subsubsection*{\textred{Capitulum. \hfill Gen. 49.}}
\lettrine{B}enedictiónes patris tui confortátæ sunt benedictiónibus patrum ejus, donec veníret desidérium cóllium æternórum: fiant in cápite Joseph, et in vértice Nazarǽi inter fratres suos.

\begin{tightlist}{}
\item \textred{℟.} Constítuit eum \textred{*} Dóminum domus suæ. \textred{C}onstítuit. \textred{℣.} Et príncipem omnis possessiónis suæ. \textred{D}óminum. Glória Patri. \textred{C}onstítuit.
\item \textred{℣.} De ventre matris meæ, tu es protéctor meus. \textred{℟.} In te confirmátus sum ex útero.
\end{tightlist}

\subsubsection*{\center \textsc{AD SEXTAM.}}
\subsubsection*{\textred{Capitulum. \hfill Gen. 49.}}
\lettrine{D}eus patris tui erit adjútor tuus, et Omnípotens benedícet tibi benedictiónibus cœli désuper.

\begin{tightlist}{}
\item \textred{℟.} De ventre matris meæ, \textred{*} Tu es protéctor meus. \textred{D}e. \textred{℣.} In te confirmátus sum ex útero. \textred{T}u. Glória Patri. \textred{D}e.
\item \textred{℣.} Invocávi Dóminum Patrem Dómini mei. \textred{℟.} Ut non derelínquat me in die tribulatiónis.
\end{tightlist}

\subsubsection*{\center \textsc{AD NONAM.}}
\subsubsection*{\textred{Capitulum. \hfill Gen. 49.}}
\lettrine{F}ilius accréscens Joseph, fílius accréscens, et decórus adspéctu: fíliæ discurrérunt super murum.

\begin{tightlist}{}
\item \textred{℟.} Invocávi Dóminum \textred{*} Patrem Dómini mei. \textred{I}nvocávi. \textred{℣.} Ut non derelínquat me in die tribulatiónis. \textred{P}atrem. Glória Patri. \textred{I}nvocávi.
\item \textred{℣.} Dedísti mihi protectiónem salútis tuæ. \textred{℟.} Et déxtera tua suscépit me.
\end{tightlist}

\subsubsection*{\center \textsc{AD II. VESPERAS.}}
\begin{tightlist}{}
\item \textred{Prima Aña Laudum. Ps. de j. Vesp. Comm. Apost. [3]}
%\item \textred{Capit. Hymnus et ℣. ut in j. Vesperis.}
\end{tightlist}

\subsubsection*{\textred{Capitulum. \hfill Gen. 49.}}
\lettrineThree{B}enedictiónes patris tui confortátæ sunt benedictiónibus patrum ejus, donec veníret desidérium cóllium æternórum: fiant in cápite Joseph, et in vértice Nazarǽi inter fratres suos.

\subsubsection*{\center \textred{Hymnus.}}
\lettrine{T}e, Joseph, célebrent ágmina cœ́litum:\\
Te cuncti résonent Christíadum chori,\\
Qui clarus méritis, junctus es ínclytæ\\
Casto fœ́dere Vírgini.\\
\\
Almo cum túmidam gérmine cónjugem\\
Admírans, dúbio tángeris ánxius,\\
Afflátu súperi Fláminis Angelus\\
Concéptum púerum docet.\\
\\
Tu natum Dóminum stringis, ad éxteras\\
Ægypti prófugum tu séqueris plagas;\\
Amíssum Sólymis, quæris, et ínvenis,\\
Miscens gáudia flétibus.\\
\\
Post mortem réliquos mors pia cónsecrat,\\
Palmámque eméritos glória súscipit:\\
Tu vivens, Súperis par, frúeris Deo,\\
Mira sorte beátior.\\
\\
Nobis, summa Trias, parce precántibus,\\
Da Joseph méritis sídera scándere:\\
Ut tandem líceat nos tibi pérpetim\\
Gratum prómere cánticum.\\
Amen.

\begin{tightlist}{}
\item \textred{℣.} Constítuit eum dóminum domus suæ. \textred{℟.} Et príncipem omnis possessiónis suæ.
\item \textred{Ad Magnif. Aña.} Fili, quid fecísti nobis sic? Ecce pater tuus et ego doléntes quærebámus te.
\end{tightlist}

\begin{center}\textcolor{red}{\rule{\linewidth}{0.5pt}} \end{center}

\subsection*{\center \textred{PRO FERIA VI.}}
\subsection*{\center \textsc{\textred{OFFICIUM VOTIVUM}\\ De Passione D.N.J.C.}}
\subsubsection*{\center \textred{Semiduplex.}}
\fancyhead[C]{\textsc{Officia Votiva - Feria VI - De Passione D.N.J.C.}}

\subsubsection*{\center \textsc{AD I. VESPERAS.}}
\subsubsection*{\textred{Capitulum. \hfill Philipp. 2.}}
\lettrineThree{F}ratres: Hoc sentíte in vobis, quod et in Christo Jesu: qui, cum in forma Dei esset, non rapínam arbitrátus est esse se æquálem Deo: sed semetípsum exinanívit, formam servi accípiens, in similitúdinem hóminum factus, et hábitu invéntus ut homo. Humiliávit semetípsum, factus obédiens usque ad mortem, mortem autem crucis.

\subsubsection*{\center \textred{Hymnus.}}
\lettrine{M}œréntes óculi spárgite lácrymas,\\
Et luctu résonent íntima córdium,\\
Illátas réfero Númini ab ímpiis\\
Pœnas, et fera vúlnera.\\
\\
Accíncta heu! gládiis turba satéllitum\\
Arréptum Dóminum fústibus ímpetit;\\
Nunc cædit cólaphis, nunc quatit hórridis\\
Divínum caput íctibus.\\
\\
Haud finis scéleri: tráditur ímprobo\\
Christus carnífici; nec mora, bárbarus\\
In Regem súperum non timet ímpio\\
Ausu, vértere déxteram.\\
\\
Audíte, o pópuli: Numen amábile,\\
Manánte ex húmeris úndique sánguine,\\
Lictóris rábidi sústinet ímpetum,\\
Et vocem premit ínnocens.\\
\\
Quis non illácrymet? Jam nova cóndidit\\
Torménta indómitæ gentis iníquitas;\\
Infígit cérebro proh dolor! éfferis\\
Sertum vépribus ásperum.\\
\\
Heu crimen! tráhitur fúnibus ímprobis\\
Funésti Dóminus supplícii ad locum;\\
Illic oppétiit fúnere, spíritum\\
Patri restítuens suo.\\
\\
Passo pro míseris tam fera vúlnera,\\
In terris résonet débita glória;\\
Sacrúmque assídue Nomen in ǽthera\\
Humánum genus éfferat.\\
Amen.

\begin{tightlist}{}
\item \textred{℣.} Oblátus est quia ipse vóluit. \textred{℟.} Et livóre ejus sanáti sumus.
\item \textred{Ad Magnificat, Aña.} O vos omnes, qui transítis per viam, atténdite, et vidéte, si est dolor sicut dolor meus.
\end{tightlist}

\subsubsection*{\center \textred{Oratio.}}
\lettrine{D}omine Jesu Christe, qui de cœlis ad terram de sinu Patris descendísti et sánguinem tuum pretiósum in remissiónem peccatórum nostrórum fudísti: te humíliter deprecámur; ut in die judícii, ad déxteram tuam audíre mereámur: Veníte benedícti. Qui vivis et regnas.

\begin{tightlist}{}
\item \textred{Ad Complet. et per Horas, in fine Hymnorum dicitur:} ℣. Glória tibi, Dómine, Qui nos salvásti in Cruce.
\item \textred{Completorium de Dominica.}
\item \textred{Ad Nunc dimittis, Aña.} Hoc signum Crucis * erit in cælo, cum Dóminus ad judicándum vénerit.
\end{tightlist}

\subsubsection*{\center \textsc{AD MATUTINUM.}}
\begin{tightlist}{}
\item \textred{Invit.} Christum Regem crucifíxum * Veníte, adorémus.
\item \textred{Ps.} Veníte, exsultémus. 2. 
\end{tightlist}
\subsubsection*{\center\textred{Hymnus.}}
\lettrine{A}dspice, infámi Deus ipse ligno\\
Pendet, effúso mádidus cruóre,\\
Adspice, immíti manus alma clavo\\
Fínditur alte.\\
\\
Hunc, velut pravi scéleris minístrum,\\
Inter indígnos médium latrónes\\
Cerne: crudélis fuit ista genis\\
Dira volúntas.\\
\\
Pallet heu! vultus; caput ecce lassum\\
Fléctitur: clausis óculis Redémptor\\
Spíritum sacro méritis onústum\\
Fundit ab ore.\\
\\
O cor, æs durum súperas, geméndo\\
Ni scelus tergas; tua culpa Christum\\
Stípiti affíxit: tua culpa morti\\
Súbdidit atræ.\\
\\
Sit Deo ætérnum decus omne in ævum\\
Qui pius nostri géneris Redémptor,\\
Críminis labem míseris nocéntem\\
Sánguine tersit.\\
Amen.

\subsubsection*{\center \textsc{IN PRIMO NOCTURNO.}}
\begin{tightlist}{}
\item \textred{Aña.} Adstitérunt reges terræ, et príncipes convenérunt in unum advérsus Dóminum, et advérsus Christum ejus.
\item \textred{Ps. 2.} Quare fremuérunt. 24.
\item \textred{Aña.} Multiplicáti sunt, qui tríbulant me, multi insúrgunt advérsum me.
\item \textred{Ps. 3.} Dómine quid. 25.
\item \textred{Aña.} Divisérunt sibi vestiménta mea, et super vestem meam misérunt sortem.
\item \textred{Ps. 21.} Deus, Deus meus. 222.
\item \textred{℣.} Oblátus est, quia ipse vóluit. \textred{℟.} Et livóre ejus sanáti sumus.
\item \textred{Lectiones de Scriptura occurr. si vero desint, legantur sequentes.}
\end{tightlist}

\subsubsection*{\center De Epístola beáti Pauli Apóstoli ad Romános}
\begin{center}\textred{\textbf{Lectio j}. \hfill Cap. 5. et 6.}\end{center}
\lettrineThree{C}omméndat caritátem suam Deus in nobis: quóniam cum adhuc peccatóres essémus, secúndum tempus, Christus pro nobis mórtuus est: multo ígitur magis nunc justificáti in sánguine ipsíus, salvi érimus ab ira per ipsum. Si enim, cum inimíci essémus, reconciliáti sumus Deo per mortem Fílii ejus, multo magis reconciliáti, salvi érimus in vita ipsíus. Non solum autem, sed et gloriámur in Deo per Dóminum nostrum Jesum Christum, per quem nunc reconciliatiónem accépimus. Si enim, uníus delícto mors regnávit per unum: multo magis abundántiam grátiæ, et donatiónis, et justítiæ accipiéntes, in vita regnábunt per unum Jesum Christum.

\begin{tightlist}{}
\item \textred{℟.} Vínea mea elécta, ego te plantávi, dicit Dóminus, et tu facta es mihi nimis amára; \textred{*} Quia parásti crucem Salvatóri tuo. \textred{℣.} Pópule meus, quid feci tibi, aut in quo contristávi te? respónde mihi. \textred{Q}uia.
\end{tightlist}

\subsubsection*{\center\textred{Lectio ij.}}
\lettrine{I}gitur, sicut per uníus delíctum, in omnes hómines in condemnatiónem, sic et per uníus justítiam, in omnes hómines in justificatiónem vitæ. Sicut enim, per inobediéntiam uníus hóminis, peccatóres constitúti sunt multi: ita et per uníus obeditiónem justi constituéntur multi. Lex autem subintrávit, ut abundáret delíctum: ubi autem abundávit delíctum, superabundávit grátia. Ut sicut regnávit peccátum in mortem, ita et grátia regnet per justítiam in vitam ætérnam, per Jesum Christum Dóminum nostrum.

\begin{tightlist}{}
\item \textred{℟.} Ego edúxi te de Ægýpto, demérso Pharaóne in mare rubrum, et ante te præívi in colúmna nubis. \textred{*} Et tu me tradidísti princípibus sacerdótum, et me duxísti ad prætórium Piláti. \textred{℣.} Pópule meus, quid feci tibi, aut in quo contristávi te? respónde mihi. \textred{E}t.
\end{tightlist}

\subsubsection*{\center\textred{Lectio iij.}}
\lettrine{Q}uid ergo dicémus? Permanébimus in peccáto, ut grátia abúndet? Absit: qui enim mórtui sumus peccáto, quómodo adhuc vivémus in illo? An ignorátis, quia quicúmque baptizáti sumus in Christo Jesu, in morte ipsíus baptizáti sumus? Consepúlti enim sumus cum illo, per Baptísmum, in mortem: ut quómodo Christus surréxit a mórtuis per glóriam Patris; ita et nos in novitáte vitæ ambulémus. Si enim complantáti facti sumus similitúdini mortis ejus; simul et resurrectiónis érimus. Hoc sciéntes, quia vetus homo noster simul crucifíxus est, ut destruátur corpus peccáti, et ultra non serviámus peccáto. Qui enim mórtuus est, justificátus est a peccáto. Si autem mórtui sumus cum Christo, crédimus quia simul étiam vivémus cum Christo.

\begin{tightlist}{}
\item \textred{℟.} Ego propter te flagellávi Ægýptum cum primogénitis suis; \textred{*} Et tu me flagellátum tradidísti, qui tamquam agnus, coram tondénte se, obmútui, et non apérui os meum. \textred{℣.} Pópule meus, quid feci tibi, aut in quo contristávi te? respónde mihi. \textred{E}t.  Glória Patri. \textred{E}t.
\end{tightlist}

\subsubsection*{\center \textsc{IN SECUNDO NOCTURNO.}}
\begin{tightlist}{}
\item \textred{Aña.} Inimíci mei dixérunt mala mihi, quando moriétur, et períbit nomen ejus?
\item \textred{Ps. 40.} Beátus qui intélligit. 122.
\item \textred{Aña.} Aliéni insurrexérunt in me, et fortes quæsiérunt ánimam meam.
\item \textred{Ps. 53.} Deus in nómine. 41.
\item \textred{Aña.} Conculcavérunt me inimíci mei tota die, advérsum me omnes cogitatiónes eórum in malum.
\item \textred{Ps. 55.} Miserére mei. 158.
\item \textred{℣.} Próprio Fílio suo non pepércit Deus. \textred{℟.} Sed pro ómnibus tradidit illum.
\end{tightlist}

\subsubsection*{\center Sermo sancti Bernárdi Abbátis.}
\begin{center}\textred{Feria 4. Hebdom. S. Serm. de Passione Domini. \hfill \textbf{Lectio iv}.}\end{center}
\lettrineThree{I}n passióne domínica, fratres, tria speciáliter cónvenit intuéri: opus, modum, causam. Nam in ópere quidem patiéntia, in modo humílitas, in causa cáritas commendátur. Patiéntia autem singuláris, quod vidélicet: cum supra dorsum ejus fabricárent peccatóres; cum sic extenderétur in ligno, ut dinumeraréntur ómnia ossa ejus; cum fortíssimum illud propugnáculum, quod custódit Israel, úndique forarétur; cum foderéntur manus ejus et pedes, sicut agnus ad occisiónem ductus sit et tamquam ovis coram tondénte non apéruit os suum. Non advérsus Patrem múrmurans, a quo missus fúerat; non advérsus humánum genus, pro quo, quæ non rápuit, exsolvébat; non dénique vel contra pópulum ipsum peculiárem sibi, a quo pro tantis benefíciis tanta mala recipiébat.

\begin{tightlist}{}
\item \textred{℟.} Ego propter te Chananæórum reges percússi, et dedi tibi sceptrum regále; \textred{*} Et tu dedísti cápiti meo spíneam corónam, et percussísti arúndine caput meum. \textred{℣.} Pópule meus, quid feci tibi, aut in quo contristávi te? respónde mihi. \textred{E}t.
\end{tightlist}

\subsubsection*{\center \textred{Lectio v.}}
\lettrine{M}odum ipsum, si diligénter atténdas, non modo mitem, sed et húmilem corde cognósces. Nempe in humilitáte judícium ejus sublátum est, cum nec ad tantas blasphémias, nec ad falsíssima, quæ sibi objiciebántur crímina respondéret. Vídimus, inquit, eum et non erat ei adspéctus, nec speciósum forma præ fíliis hóminum, sed oppróbrium hóminum et tamquam leprósum, novíssimum virórum, plane virum dolórum a Deo percússum et humiliátum; ita ut nulla esset ei spécies, neque decor. O novíssimum et altíssimum! O húmilem et sublímem! O oppróbrium hóminum et glóriam Angelórum! Nemo illo sublímior, neque humílior. Dénique sputis íllitus est, oppróbriis saturátus est, morte turpíssima condemnátus est, cum scelerátis deputátus est. Nihílque merébitur vel ista humílitas, quæ hunc habet modum; imo quæ tam est ultra modum? Sicut est patiéntia singuláris, sic et humílitas admirábilis; útraque sine exémplo.

\begin{tightlist}{}
\item \textred{℟.} Ego duxi te per desértum quadragínta annis, et manna cibávi te; \textred{*} Et tu me cecidísti álapis, et flagéllis. \textred{℣.} Pópule meus, quid feci tibi, aut in quo contristávi te? respónde mihi. \textred{E}t.
\end{tightlist}

\subsubsection*{\center \textred{Lectio vj.}}
\lettrine{U}trámque tamen magnífice causa ipsa comméndat, nimírum cáritas est. Propter nímiam enim caritátem suam, qua diléxit nos Deus, ut servum redímeret, nec Pater Fílio, nec sibi Fílius ipse pepércit. Vere nímiam, quia et hæc mensúram excédit, modum súperat, plane superéminens univérsis. Majórem, inquit, caritátem nemo habet, quam ut ánimam suam ponat quis pro amícis suis. Tu majórem habuísti, Dómine, ponens eam étiam pro inimícis. Cum enim adhuc inimíci essémus, per mortem tuam et tibi reconciliáti sumus et Patri. Quænam ergo ália vidébitur esse, vel fuísse, vel fore huic símilis caritáti? Vix pro justo quis móritur. Tu pro injústis passus es, móriens propter delícta nostra, qui venísti justificáre gratis peccatóres, servos fácere fratres, captívos coherédes, éxsules reges.

\begin{tightlist}{}
\item \textred{℟.} Ego exaltávi te magna virtúte: \textred{*} Et tu me suspendísti in patíbulo crucis, in qua expándi manus meas ad pópulum non credéntem et contradicéntem. \textred{℣.} Pópule meus, quid feci tibi, aut in quo contristávi te? respónde mihi. \textred{E}t. Glória Patri. \textred{E}t.
\end{tightlist}

\subsubsection*{\center \textsc{IN TERTIO NOCTURNO.}}
\begin{tightlist}{}
\item \textred{Aña.} Fílii hóminum dentes eórum arma, et sagíttæ, et lingua eórum gládius acútus. 
\item \textred{Ps. 56.} Miserére mei Deus. 159.
\item \textred{Aña.} Intendérunt arcum rem amáram, ut sagíttent in occúltis immaculátum.
\item \textred{Ps. 63.} Exáudi Deus. 253.
\item \textred{Aña.} Factus sum sicut homo sine adjutório inter mórtuos liber.
\item \textred{Ps. 87.} Dómine Deus salútis. 277.
\item \textred{℣.} Factus est obédiens usque ad mortem. \textred{℟.} Mortem autem crucis.
\end{tightlist}

\subsubsection*{\center Léctio sancti Evangélii secúndum Joánnem.}
\begin{center}\textred{\textbf{Lectio vij}. \hfill Cap. 19.}\end{center}
\lettrine{I}n illo témpore: Sciens Jesus, quia ómnia consummáta sunt, ut consummarétur scriptúra, dixit: Sítio. Et réliqua.

\begin{center}Homilía sancti Cyrílli Epíscopi Alexandríni.\\
\textred{Comment. in Joannem.\\ Lib. 12. cap. 25.}\end{center}
\lettrineThree{C}um truculénter ómnia impietáte Judæórum essent perácta, nec quidquam ad summam crudelitátem deésset, tandem sanctíssima caro Christi naturále quid pátitur. Multis enim, ac váriis dolóribus exsiccáta, siti torquétur. Multum enim possunt dolóres ad commovéndam sitim; innátum enim calórem exagitántes, vel quod in profunditáte húmidum est, consúmunt et ígneis ardóribus doléntis víscera urunt. Sicut ígitur céteras passiónes, sic hanc étiam sponte Christus perpéssus est. Bíbere ígitur petit: illi autem ádeo longe ab omni humanitáte áberant, ut pro juvánte potu atque jucúndo nocéntem et acérbum attúlerint, et humanitátem, qua movéri videbántur, in impietátem convérterint. Nam cum peténti dabant, caritátis hábitum suscípere videbántur; sed non est possíbile mentíri Scriptúram, in qua ex persóna Christi dícitur: Et dedérunt in escam meam fel, et in siti mea potavérunt me acéto.

\begin{tightlist}{}
\item \textred{℟.} Ego te potávi aqua salútis de petra: \textred{*} Et tu me potásti felle, et acéto. \textred{℣.} Pópule meus, quid feci tibi, aut in quo contristávi te? respónde mihi. \textred{E}t.
\end{tightlist}

\subsubsection*{\center\textred{Lectio viij. \hfill Idem. Cap. 36.}}
\lettrine{C}um ergo accepísset Jesus acétum, dixit: Consummátum est: et inclináto cápite trádidit spíritum. Consummátum jam esse Salvátor dicit, quia sæviéndi Judæórum potéstas in ipsum per mortem suam términum hábuit. Quid enim non excogitárunt, aut quid suprémæ inhumanitátis relíctum sibi est? Quis cædéndi, quis contumeliándi modus non adhíbitus? Mérito ígitur consummátum esse, ait. Cumque tempus jam esset, ut étiam spirítibus apud ínferos prædicáret (venit enim, ut tam vivórum quam mortuórum dominarétur), mortem pro nobis suscépit: et hanc natúræ nostræ própriam passiónem sponte súbiit secúndum carnem, quamvis, ut Deus vita naturáliter sit, ut spoliátis ínferis, primítiæ dormiéntium, et primogénitus, (ut Scriptúræ dicunt) mortuórum factus, recúrsum ad vitam natúræ nostræ largirétur.

\begin{tightlist}{}
\item \textred{℟.} Ego ante te apérui mare: \textred{*} Et tu aperuísti láncea latus Salvatóri tuo, qui redémit mundum in sánguine suo. \textred{℣.} Pópule meus, quid feci tibi, aut in quo contristávi te? respónde mihi. \textred{E}t.
\end{tightlist}

\subsubsection*{\center\textred{Lectio ix.}}
\lettrine{I}nclinávit ígitur caput. Accidit enim moriéntibus, vigóre carnis jam consúmpto et ánima, quæ corpus rétinet, abeúnte, inclinátio cápitis. Quámobrem verbis hoc modo usus Evangelísta est. Sic et tradidísse spíritum, a consuetúdine nostra sibi dictum est; multi enim sic exstíngui, et mori dicúntur. Quamvis mihi videátur ídeo non simplíciter mórtuum esse, sed tradidísse spíritum Evangelísta dixísse; quia in manus Dei Patris ánimam suam commendávit, ut ipse dixit: Pater! in manus tuas comméndo spíritum meum. Quod nobis magnæ spei fundaméntum atque oríginem præbet.

\begin{tightlist}{}
\item \textred{℟.} Crucem tuam adorámus, Dómine, sanctámque recólimus passiónem: \textred{*} Ecce enim propter lignum venit gáudium in univérso mundo. \textred{℣.} Deus misereátur nostri, et benedícat nobis, illúminet vultum suum super nos, et misereátur nostri. \textred{E}cce. Glória Patri. \textred{E}cce.
\item Te Deum laudámus. 6.
\item \textred{Sacerd. ℣.} Adorámus te, Christe, et benedícimus tibi. \textred{℟.} Quia per mortem tuam redemísti mundum.
\end{tightlist}

\subsubsection*{\center \textsc{AD LAUDES.}}
\begin{center}\textred{et per Horas, Aña.}\end{center}
\lettrine{F}ui flagellátus tota die, et castigátio mea in matutínis.

\begin{tightlist}{}
\item \textred{Ps. de Dominica.} 30. 
\item \textred{Aña.} Veni ad montem myrrhæ et ad collem thuris; tamquam agnus ad occisiónem ductus obmútui, et non apérui os meum.
\item \textred{Aña.} Fodérunt manus meas, et pedes meos, dinumeravérunt ómnia ossa mea.
\item \textred{Aña.} Consolántem me quæsívi et non invéni: dedérunt in escam meam fel, et in siti mea potavérunt me acéto.
\item \textred{Aña.} Cum accepísset Jesus acétum dixit: Consummátum est, et inclináto cápite, emísit spíritum.
\end{tightlist}

\subsubsection*{\textred{Capitulum. \hfill Philipp. 2.}}
\lettrineThree{F}ratres, Hoc enim sentíte in vobis quod et in Christo Jesu, qui cum in forma Dei esset, non rapínam arbitrátus est esse se æquálem Deo, sed semetípsum exinanívit, formam servi accípiens in similitúdinem hóminum factus, et hábitu invéntus ut homo. Humiliávit semetípsum factus obédiens usque ad mortem, mortem autem crucis.

\subsubsection*{\center \textred{Hymnus.}}
\lettrine{S}ævo dolórum túrbine\\
Jactátur, atris óbrutus\\
Pœnis, acérba sústinens,\\
Redémptor affíxus cruci.\\
\\
Pedes, manúsque horríbili\\
Clavi cruéntant vúlnere,\\
Cor, vultus, artus, péctora\\
Sacro redúndant sánguine.\\
\\
Flet, orat, et clamans obit:\\
Cor Matris ictum cóncidit:\\
Heu, Mater! Heu Fili! dolor\\
Ingráta frangat péctora.\\
\\
Montes, sepúlcra, sáxaque\\
Scindúntur; arva, flúmina,\\
Rupes et æquor cóntremunt,\\
Templíque velum scínditur.\\
\\
Sol, luna, cœlum, sídera\\
Plangunt, et orbis íngemit:\\
O vos viri, vos párvuli,\\
Nuptæ, puéllæ plángite.\\
\\
Adstáre mœréntes cruci,\\
Pedes beátos úngite,\\
Laváte fletu, térgite\\
Comis, et ore lámbite.\\
\\
Tu caritátis víctima,\\
Ut nostra tollas crímina,\\
Nobis salúbri pérficis\\
Adoptiónem sánguine.\\
\\
Nostra ergo pax, et gáudium\\
Sis vita, Jesu, et prǽmium;\\
Sis ductor, et lux in via,\\
Merces, coróna in Pátria.\\
Amen.

\begin{tightlist}{}
\item \textred{℣.} Ipse passus est propter iniquitátes nostras. \textred{℟.} Et attrítus est propter scélera nostra.
\item \textred{Ad Benedictus, Aña.} Ad Jesum autem cum veníssent, ut vidérunt eum jam mórtuum, non fregérunt ejus crura, sed unus mílitum láncea latus ejus apéruit, et contínuo exívit sanguis et aqua.
\end{tightlist}

\subsubsection*{\center\textred{Oratio.}}
\lettrine{D}omine Jesu Christe, qui de cœlis ad terram de sinu Patris descendísti et sánguinem tuum pretiósum in remissiónem peccatórum nostrórum fudísti: te humíliter deprecámur; ut in die judícii, ad déxteram tuam audíre mereámur: Veníte benedícti. Qui vivis et regnas.

\begin{tightlist}{}
\item \textred{Ad Horas Añæ de Laudibus, Psalmi de Dominica, ad Primam tamen ut in Festis.}
\end{tightlist}

\subsubsection*{\center \textsc{AD PRIMAM.}}
\begin{tightlist}{}
\item \textred{℣.} Qui nos salvásti in cruce.
\end{tightlist}

\subsubsection*{\center \textsc{AD TERTIAM.}}
\subsubsection*{\textred{Capitulum. \hfill Philipp. 2.}}
\lettrine{F}ratres, Hoc enim sentíte in vobis quod et in Christo Jesu, qui cum in forma Dei esset, non rapínam arbitrátus est esse se æquálem Deo, sed semetípsum exinanívit, formam servi accípiens in similitúdinem hóminum factus, et hábitu invéntus ut homo. Humiliávit semetípsum factus obédiens usque ad mortem, mortem autem crucis.

\begin{tightlist}{}
\item \textred{℟.} Oblátus est, \textred{*} Quia ipse vóluit. \textred{O}blátus. \textred{℣.} Et livóre ejus sanáti sumus. \textred{Q}uia. Glória Patri. \textred{O}blátus 
\item \textred{℣.} Próprio Fílio suo non pepércit Deus. \textred{℟.} Sed pro nobis ómnibus trádidit illum.
\end{tightlist}

\subsubsection*{\center \textsc{AD SEXTAM.}}
\subsubsection*{\textred{Capitulum \hfill Zach. 13.}}
\lettrine{Q}uid sunt plagæ istæ in médio mánuum tuárum? Et dicet: His plagátus sum in domo eórum, qui diligébant me.

\begin{tightlist}{}
\item \textred{℟.} Próprio Fílio suo \textred{*} Non pepércit Deus. \textred{P}róprio. \textred{℣.} Sed pro nobis ómnibus trádidit illum. \textred{N}on. Glória Patri. \textred{P}róprio.
\item \textred{℣.} Factus est obédiens usque ad mortem. \textred{℟.} Mortem autem crucis.
\end{tightlist}

\subsubsection*{\center \textsc{AD NONAM.}}
\subsubsection*{\textred{Capitulum. \hfill Is. 53.}}
\lettrine{G}eneratiónem ejus quis enarrábit? Quia abscíssus est de terra vivéntium, propter scelus pópuli mei percússi eum.

\begin{tightlist}{}
\item \textred{℟.} Factus est obédiens \textred{*} Usque ad mortem. \textred{F}actus. \textred{℣.} Mortem autem crucis. \textred{U}sque. Glória Patri. \textred{F}actus.
\item \textred{℣.} Ipse passus est propter iniquitátes nostras. \textred{℟.} Et attrítus est propter scélera nostra.
\end{tightlist}

\subsubsection*{\center \textsc{AD II. VESPERAS.}}
\begin{tightlist}{}
\item \textred{Prima Aña Laudum.}
\item \textred{Psalm. 115} Crédidi. 93. 
\item \textred{Ps. 119.} Ad Dóminum. 93. 
\item \textred{Ps. 139.} Eripe me. 236. 
\item \textred{Ps. 140.} Dómine clamávi. 237. 
\item \textred{Ps. 141.} Voce mea. 238.
%\item \textred{Capit, Hymn. et ℣. ut in j. Vesperis.}
\end{tightlist}

\subsubsection*{\textred{Capitulum. \hfill Philipp. 2.}}
\lettrineThree{F}ratres: Hoc sentíte in vobis, quod et in Christo Jesu: qui, cum in forma Dei esset, non rapínam arbitrátus est esse se æquálem Deo: sed semetípsum exinanívit, formam servi accípiens, in similitúdinem hóminum factus, et hábitu invéntus ut homo. Humiliávit semetípsum, factus obédiens usque ad mortem, mortem autem crucis.

\subsubsection*{\center \textred{Hymnus.}}
\lettrine{M}œréntes óculi spárgite lácrymas,\\
Et luctu résonent íntima córdium,\\
Illátas réfero Númini ab ímpiis\\
Pœnas, et fera vúlnera.\\
\\
Accíncta heu! gládiis turba satéllitum\\
Arréptum Dóminum fústibus ímpetit;\\
Nunc cædit cólaphis, nunc quatit hórridis\\
Divínum caput íctibus.\\
\\
Haud finis scéleri: tráditur ímprobo\\
Christus carnífici; nec mora, bárbarus\\
In Regem súperum non timet ímpio\\
Ausu, vértere déxteram.\\
\\
Audíte, o pópuli: Numen amábile,\\
Manánte ex húmeris úndique sánguine,\\
Lictóris rábidi sústinet ímpetum,\\
Et vocem premit ínnocens.\\
\\
Quis non illácrymet? Jam nova cóndidit\\
Torménta indómitæ gentis iníquitas;\\
Infígit cérebro proh dolor! éfferis\\
Sertum vépribus ásperum.\\
\\
Heu crimen! tráhitur fúnibus ímprobis\\
Funésti Dóminus supplícii ad locum;\\
Illic oppétiit fúnere, spíritum\\
Patri restítuens suo.\\
\\
Passo pro míseris tam fera vúlnera,\\
In terris résonet débita glória;\\
Sacrúmque assídue Nomen in ǽthera\\
Humánum genus éfferat.\\
Amen.

\begin{tightlist}{}
\item \textred{℣.} Oblátus est quia ipse vóluit. \textred{℟.} Et livóre ejus sanáti sumus.
\item \textred{Ad Magnif. Aña.} Depónens Joseph corpus Jesu invólvit síndone, et pósuit eum in monuménto excíso. Fuit planctus magnus in illa die; mulíeres autem sedéntes contra sepúlcrum dolébant, sicut doléri solet in morte primogéniti.
\item \textred{Completorium de Dominica.}
\item \textred{Ad Nunc dimittis, Aña.} Hoc signum Crucis * erit in cælo, cum Dóminus ad judicándum vénerit.
\end{tightlist}

\begin{center}\textcolor{red}{\rule{\linewidth}{0.5pt}} \end{center}

\subsection*{\center \textred{ROMÆ\\ PRO FERIA III.}}
\subsection*{\center \textsc{\textred{OFFICIUM VOTIVUM}\\ De SS. Apostolis Petro et Paulo}}
\subsubsection*{\center \textred{Semiduplex.}}
\fancyhead[C]{\textsc{Offic. Vot. - Feria III Romæ - SS. Petro et Paulo}}

\begin{tightlist}{}
\item \textred{Omnia de Comm. Apost.} [3] \textred{pro diversitate temporis, præter sequentia.}
\end{tightlist}

\subsubsection*{\center \textsc{IN UTRISQUE VESPERIS.}}
\begin{tightlist}{}
\item \textred{℣.} Constítues eos príncipes super omnem terram. \textred{℟.} Mémores erunt nóminis tui, Dómine.
\item \textred{Ad Magnificat, Aña.} Petrus Apóstolus, et Paulus Doctor géntium, ipsi nos docuérunt legem tuam, Dómine.
\end{tightlist}

\subsubsection*{\center \textred{Oratio.}}
\lettrine{D}eus, cujus déxtera beátum Petrum ambulántem in flúctibus, ne mergerétur, eréxit, et co-apóstolum ejus Paulum, tértio naufragántem, de profúndo pélagi liberávit: exáudi nos propítius, et concéde; ut ambórum méritis æternitátis glóriam consequámur: Qui vivis.

\subsubsection*{\center \textsc{IN PRIMO NOCTURNO.}}
\begin{tightlist}{}
\item \textred{Lect. de Script. occurr.}
\item \textred{In Quadrag. vero dicuntur sequentes.}
\end{tightlist}

\subsubsection*{\center Incipit Epístola prima beáti Petri Apóstoli.}
\begin{center}\textred{\textbf{Lectio j}. \hfill Cap. 1.}\end{center}
\lettrineThree{P}etrus Apóstolus Jesu Christi, eléctis ádvenis dispersiónis Ponti, Galátiæ, Cappadóciæ, Asiæ, et Bithýniæ,
secúndum præsciéntiam Dei Patris, in sanctificatiónem Spíritus, in obediéntiam, et aspersiónem sánguinis Jesu Christi. Grátia vobis, et pax multiplicétur.
Benedíctus Deus et Pater Dómini nostri Jesu Christi, qui secúndum misericórdiam suam magnam regenerávit nos in spem vivam, per resurrectiónem Jesu Christi ex mórtuis,
in hereditátem incorruptíbilem, et incontaminátam, et immarcescíbilem, conservátam in cælis in vobis,
qui in virtúte Dei custodímini per fidem in salútem, parátam revelári in témpore novíssimo.

\begin{tightlist}{}
\item \textred{℟.} Ecce, ego. [9]
\end{tightlist}

\subsubsection*{\center \textred{Lectio ij.}}
\lettrine{I}n quo exsultábitis, módicum nunc si opórtet contristári in váriis tentatiónibus:
ut probátio vestræ fídei multo pretiósior auro (quod per ignem probátur) inveniátur in laudem, et glóriam, et honórem in revelatióne Jesu Christi:
quem cum non vidéritis, dilígitis: in quem nunc quoque non vidéntes créditis: credéntes autem exsultábitis lætítia inenarrábili, et glorificáta:
reportántes finem fídei vestræ, salútem animárum.

\begin{tightlist}{}
\item \textred{℟.} Tóllite jugum. [10]
\end{tightlist}

\subsubsection*{\center \textred{Lectio iij.}}
\lettrine{D}e qua salúte exquisiérunt, atque scrutáti sunt prophétæ, qui de futúra in vobis grátia prophetavérunt:
scrutántes in quod, vel quale tempus significáret in eis Spíritus Christi: prænúntians eas quæ in Christo sunt passiónes, et posterióres glórias:
quibus revelátum est quia non sibimetípsis, vobis autem ministrábant ea, quæ nunc nuntiáta sunt vobis per eos, qui evangelizavérunt vobis, Spíritu sancto misso de cælo, in quem desíderant ángeli prospícere.

\begin{tightlist}{}
\item \textred{℟.} Dum stetéritis. [10]
\end{tightlist}

\subsubsection*{\center \textsc{IN SECUNDO NOCTURNO.}}
\subsubsection*{\center De Sermóne sancti Leónis Papæ.}
\begin{center}\textred{Ex hom. 82. in Natal. SS. Apost. Petri et Pauli.\\ \textbf{Lectio iv}.}\end{center}
\lettrineThree{N}erónis furor per ómnium vitiórum inflammátus excéssum, in hunc eum usque torréntem suæ præcipitávit insániæ, ut primus nómini Christiáno atrocitátem generális persecutiónis inférret, quasi per sanctórum Apostolórum Petri et Pauli neces grátia Dei posset exstíngui; quibus hoc ipsum erat máximum lucrum, ut contémptus vitæ hujus occíduæ, percéptio fíeret felicitátis ætérnæ. Pretiósa est ergo in conspéctu Dómini mors sanctórum ejus: nec ullo crudelitátis génere déstrui potest sacraménto crucis Christi fundáta relígio. Non minúitur persecutiónibus Ecclésia, sed augétur: et semper domínicus ager ségete ditióri vestítur, dum grana, quæ síngula cadunt, multiplicáta nascúntur. Unde duo ista præclára divíni séminis gérmina in quantam sóbolem germinárint, beatórum míllia Mártyrum protestántur, qui apostolicórum æmuli triumphórum, Urbem nostram purpurátis, et longe latéque rutilántibus pópulis ambiérunt, et quasi ex multárum honóre gemmárum consérto uno diadémate coronárunt.

\begin{tightlist}{}
\item \textred{℟.} Vidi cunjúnctos. [12]
\end{tightlist}

\subsubsection*{\center \textred{Lectio v.}}
\lettrine{D}e quorum præsídio, dilectíssimi, divínitus nobis ad exémplum patiéntiæ et confirmatiónem fídei præparáto, universáliter quidem in ómnium sanctórum commemoratióne lætándum est: sed in horum excelléntia Patrum mérito est excelléntius gloriándum, quos grátia Dei in tantum ápicem inter ómnia Ecclésiæ membra provéxit, ut eos in córpore, cui caput est Christus, quasi géminum constitúeret lumen oculórum. De quorum méritis atque virtútibus, quæ omnem loquéndi súperant facultátem, nihil divérsum, nihil debémus sentíre discrétum: quia illos et eléctio pares, et labor símiles, et finis fecit æquáles.

\begin{tightlist}{}
\item \textred{℟.} Isti sunt triumphatóres. [13]
\end{tightlist}

\subsubsection*{\center \textred{Lectio vj.}}
\lettrine{S}icut autem et nos expérti sumus, et nostri probavére majóres, crédimus atque confídimus, inter omnes labóres istíus vitæ, ad obtinéndam misericórdiam Dei, semper nos speciálium patronórum oratiónibus adjuvándos: ut, quantum própriis peccátis deprímimur, tantum apostólicis méritis erigámur. Per Dóminum nostrum Jesum Christum, cui est cum Patre et sancto Spíritu éadem potéstas, una divínitas, in sǽcula sæculórum. Amen.

\begin{tightlist}{}
\item \textred{℟.} Fuérunt sine queréla. [13]
\end{tightlist}

\subsubsection*{\center \textsc{IN TERTIO NOCTURNO.}}
\subsubsection*{\center Léctio sancti Evangélii secúndum Matthǽum.}
\begin{center}\textred{\textbf{Lectio vij}. \hfill Cap. 19.}\end{center}
\lettrine{I}n illo témpore: Dixit Petrus ad Jesum: Ecce nos relíquimus ómnia, et secúti sumus te: quid ergo erit nobis? Et réliqua.

\begin{center}Homilía sancti Ambrósii Epíscopi.\\
\textred{In Ps. 118. Serm. 8.}\end{center}
\lettrineThree{E}cce nos ómnia derelíquimus, et secúti sumus te. Hoc est, non quæsívimus quæ sǽculi sunt, non quæsívimus partem de possessiónibus, sed te elégimus portiónem. Reliquísti ergo ómnia, Petre, quæ habébas; unde habes quod habére te dicis? Surgit claudus, et sono tui sermónis erígitur: sanitátem áliis donas, qui indigébas ipse tuæ salútis auxílio. Reliquísti ergo quæ habébas, et cepísti, quæ non habébas. Christus tibi pórtio, Christus tibi posséssio est: illíus nomen tibi muníficum, illíus nomen tibi est fructuósum, illíus nomen tibi tribúta depéndit, et bona tribúta, non pecúniæ, sed grátiæ.

\begin{tightlist}{}
\item \textred{℟.} Isti sunt, qui vivéntes in carne, plantavérunt Ecclésiam sánguine suo: non sunt de terris córpora eórum separáta, \textred{*} Quorum mérita sunt in cælo glorióse cumuláta. \textred{℣.} Sancti per fidem vicérunt regna, operáti sunt justítiam, adépti sunt repromissiónes. \textred{Q}uorum.
\end{tightlist}

\subsubsection*{\center \textred{Lectio viij.}}
\lettrine{Q}uicúmque relíquerit patrem, aut matrem, aut domum, aut fratres, aut soróres, céntuplum accípiet in hoc sǽculo, et vitam ætérnam possidébit. Multos quidem in hoc sǽculo, qui sua paupéribus donavérunt, uberióre thesáuro locupletátos fuísse cognóvimus. Est pulchérrimus locus, ut étiam ad misericórdiam provocéntur. Sed non debent a Dómino sæculárem hujúsmodi exígere mercédem, nec speráre quæ mundi sunt; sed magis, quia iis, qui ómnia relíquerint, Deus pórtio est. Ipse est útique merces perfécta virtútum, qui non céntupli enumeratióne, sed perféctæ plenitúdinis æstimatióne censétur. Ego, inquit, Deus tuus sum. Non dixit: Ero, sed: Jam sum, jam inhábito, jam possídeo.

\begin{tightlist}{}
\item \textred{℟.} Isti sunt viri sancti, quos elégit Dóminus in caritáte non ficta, et dedit illis glóriam sempitérnam: \textred{*} Quorum doctrína fulget Ecclésia, ut sole luna. \textred{℣.} In omnem terram exívit sonus eórum, et in fines orbis terræ verba eórum. \textred{Q}uorum.
\end{tightlist}

\subsubsection*{\center \textred{Lectio ix.}}
\lettrine{H}abes, homo, inter multa propósita, étiam hujúsmodi portiónem. Propósuit tibi Dóminus in portióne possessiónes, in portióne aurum, in portióne argéntum, in portióne honóres, in portióne nobilitátem, propósuit étiam in portióne seípsum. Habes ígitur plúrimas portiónes; élige quam putas; non te númerus pertúrbet, sed éxcitet grátia; non te labor avértat, sed fructus invítet.

\begin{tightlist}{}
\item \textred{℟.} Cives Apostolórum et doméstici Dei advenérunt hódie, \textred{*} Portántes facem et illuminántes pátriam, dare pacem géntibus et liberáre pópulum Dómini. \textred{℣.} Audíte preces súpplicum, vitæ ætérnæ poscéntes prǽmia, qui fertis in dextris manípulos justítiæ, quique gaudéntes advenístis hódie. \textred{P}ortántes. Glória Patri. \textred{P}ortántes.
\item Te Deum laudámus. 6.
\item \textred{Sacerdotalis ℣.} Dedísti hereditátem. \textred{℟.} Timéntibus nomen tuum Dómine.
\end{tightlist}

\subsubsection*{\center \textsc{AD LAUDES.}}
\begin{tightlist}{}
\item \textred{℣.} In omnem terram exívit sonus eórum. \textred{℟.} Et in fines orbis terræ verba eórum.
\item \textred{Ad Bened. Aña.} Gloriósi príncipes terræ, quómodo in vita sua dilexérunt se, ita et in morte non sunt separáti.
\end{tightlist}

\begin{center}\textcolor{red}{\rule{\linewidth}{0.5pt}} \end{center}
\clearpage
\section*{\center \textsc{SUFFRAGIA CONSUETA}}
\subsection*{\center \textsc{Seu Commemorationes Communes}}

\fancyhead[C]{\textsc{Commemorationes Communes - In Adventu}}

\begin{tightlist}{}
\item \textred{Dicendæ in fine Vesperarum et Laudum ab Octava Epiphaniæ usque ad Dominicam Passionis exclusive, et ab Octava Corporis Christi usque ad Vigiliam Nativitatis Domini exclusive, in omnibus Dominicis, in quibus non fit Officium Duplex, in Feriis, comprehensis etiam Vigiliis, in Festis Semiduplicibus, Simplicibus, et Simplicissimis extra quascumque Octavas celebratis, juxta Rubricas Generales, Rubr. 41. prout dictæ commemorationes pro suo respectivo tempore inferius distributæ et designatæ sunt.}
\end{tightlist}

\begin{center}\textcolor{red}{\rule{\linewidth}{0.5pt}} \end{center}

\subsection*{\center \textsc{IN ADVENTU}}
\begin{tightlist}{}
\item \textred{Exceptis utrisque Vesperis et Laudibus primæ Dominicæ, et Vesp. in quibus dicitur Aña.} O. \textred{ac Laudibus Feriæ quartæ Quatuor Temporum.}
\end{tightlist}

\subsubsection*{\center \textsc{DE B. VIRGINE}}
\begin{tightlist}{}
\item \textred{Quando Officium diei non est de B. V. etiamsi in Choro recitetur ejus Officium parvum.}
\item \textred{Ad Vesperas et Laudes. Aña.} Spíritus sanctus in te descéndet, María, ne tímeas; habébis in útero Fílium Dei, allelúja.
\item \textred{℣.} Ave María, grátia plena. \textred{℟.} Dóminus tecum.
\item Orémus. \hfill \textred{Oratio.}
\end{tightlist}
\lettrine{D}eus qui de beátæ Maríæ Vírginis útero, Verbum tuum, Angelo nuntiánte, carnem suscípere voluísti: præsta supplícibus tuis; ut, qui vere eam Genitrícem Dei crédimus, ejus apud te intercessiónibus adjuvémur. Per eúmdem Christum Dóminum nostrum. Amen.

\subsubsection*{\center \textsc{DE OMNIBUS SANCTIS.}}
\begin{tightlist}{}
\item \textred{Aña.} Ecce Dóminus véniet, et omnes Sancti ejus cum eo: et erit in die illa lux magna, allelúja.
\item \textred{℣.} Ecce apparébit Dóminus super nubem cándidam. \textred{℟.} Et cum eo Sanctórum míllia.
\item Orémus. \hfill \textred{Oratio.}
\end{tightlist}
\lettrine{C}onsciéntias nostras, quǽsumus Dómine, visitándo purífica: ut véniens Jesus Christus Fílius tuus Dóminus noster, cum ómnibus Sanctis, parátam sibi in nobis invéniat mansiónem. Qui tecum vivit et regnat in unitáte Spíritus Sancti Deus.

\begin{center}\textcolor{red}{\rule{\linewidth}{0.5pt}} \end{center}

\subsubsection*{\center \textsc{DE SS. SACRAMENTO}}
\begin{tightlist}{}
\item \textred{Fit comm. ab Oct. Epiph. usque ad Quadrag. et ab Oct. Corp. Christi usque ad Adventum; solis Feriis quartis ad Vesperas, et Feriis quintis ad Laudes et Vesperas, nisi fit Officium de Venerabili Sacramento, hoc modo:}
\item \textred{Feria quarta ad Vesp., nisi præcedat Festum Duplex, Aña.} O quam suávis est, Dómine, Spíritus tuus: qui, ut dulcédinem tuam in fílios demonstráres, pane suavíssimo de cœlo prǽstito, esuriéntes reples bonis, fastidiósos dívites dimíttens inánes.
\item \textred{℣.} Panem de cœlo præstitísti eis. \textred{℟.} Omne delectaméntum in se habéntem.
\item \textred{Feria quinta ad Laudes, Aña.} Ego sum panis vivus, qui de cœlo descéndi: si quis manducáverit ex hoc pane, vivet in ætérnum, allelúja.
\item \textred{℣.} Pósuit fines tuos pacem. \textred{℟.} Et ádipe fruménti sátiat te.
\item \textred{Feria quinta ad Vesperas, nisi sequatur Festum Duplex, Aña.} O sacrum convívium, in quo Christus súmitur: recólitur memória passiónis ejus: mens implétur grátia, et futúræ glóriæ nobis pignus datur, allelúja.
\item \textred{℣.} Panem de cœlo præstitísti eis. \textred{℟.} Omne delectaméntum in se habéntem.
\item Orémus. \hfill \textred{Oratio.}
\end{tightlist}
\lettrine{D}eus, qui nobis sub Sacraménto mirábili passiónis tuæ memóriam reliquísti: tríbue quǽsumus, ita nos córporis et sánguinis tui sacra mystéria venerári; ut redemptiónis tuæ fructum in nobis júgiter sentiámus. Qui vivis et regnas in sǽcula sæculórum. Amen.

\subsubsection*{\center \textsc{DE RESURRECTIONE}}
\fancyhead[C]{\textsc{Commemorationes Communes - Per Annum}}
\begin{tightlist}{}
\item \textred{Fit com. ab Oct. Corp. Christi usque ad Adventum, Dominicis diebus ad Laudes tantum, quando de Dominica agitur, etiam si sit infra Octav. et hæc præcedere debet omnes alias commem. tam occurrentes proprias, quam communes.}
\item \textred{Aña.} Et valde mane una Sabbatórum véniunt ad monuméntum, orto jam sole, allelúja.
\item \textred{℣.} Surréxit Dóminus de sepúlcro. \textred{℟.} Qui pro nobis pepéndit in ligno.
\item Orémus. \hfill \textred{Oratio.}
\end{tightlist}
\lettrine{P}ræsta quǽsumus omnípotens Deus, ut qui Resurrectiónis Domínicæ commemoratiónem ágimus, ereptiónis nostræ lætítiam suscípere mereámur. Per eúmdem Christum Dóminum nostrum. Amen.

\subsubsection*{\center \textsc{DE CRUCE}}
\begin{tightlist}{}
\item \textred{Ferialibus diebus tantum, (comprehensis etiam Vigiliis) ab Oct. Epiph. usque ad Dom Passionis, et ab Oct. Corp. Christi usque ad Adventum.}
\item \textred{Ad Vesp. et Laud. Aña.} Per signum Crucis de inimícis nostris líbera nos, Deus noster.
\item \textred{℣.} Omnis terra adóret te Deus, et psallat tibi. \textred{℟.} Psalmum dicat nómini tuo Dómine.
\item Orémus. \hfill \textred{Oratio.}
\end{tightlist}
\lettrine{P}erpétua nos quǽsumus Dómine pace custódi: quos per lignum sanctæ Crucis redímere dignátus es. Qui vivis et regnas in sǽcula sæculórum. Amen.

\begin{center}\textcolor{red}{\rule{\linewidth}{0.5pt}} \end{center}

\subsection*{\center \textsc{POST EPIPHANIAM USQUE AD ADVENTUM}}
\begin{tightlist}{}
\item \textred{Ab Oct. Epiph. usque ad Domin. Passionis, et ab Oct. Corp. Christi usque ad Adventum, fiunt sequentes commemorationes.}
\end{tightlist}

\subsubsection*{\center \textsc{DE TITULARI}}
\begin{center}\textred{Vel Patrono Ecclesiæ fit commemoratio ante vel post comm. de Sanctis Ordinis Nostri, juxta ejus dignitatem.}\end{center}

\subsubsection*{\center \textsc{DE BEATA VIRGINE}}
\begin{tightlist}{}
\item \textred{Quando Officium diei non est de B.V. etiamsi in Choro recitetur ejus Officium parvum.}
\item \textred{Ad Vesperas et Laudes, Aña.} Sancta María, succúrre míseris, juva pusillánimes, réfove flébiles: ora pro pópulo, intérveni pro clero, intercéde pro devóto femíneo sexu; séntiant omnes tuum juvámen, quicúmque célebrant tuam sanctam Commemoratiónem.
\item \textred{Ab Oct. Epiph. usque ad Purificationem dicitur:}
\item \textred{℣.} Post partum, Virgo, invioláta permansísti. \textred{℟.} Dei Génitrix intercéde pro nobis.
\item Orémus. \hfill \textred{Oratio.}
\end{tightlist}
\lettrine{D}eus, qui salútis ætérnæ, beátæ Maríæ virginitáte fecúnda, humáno géneri prǽmia præstitísti: tríbue quǽsumus; ut ipsam pro nobis intercédere sentiámus, per quam merúimus auctórem vitæ suscípere, Dóminum nostrum Jesum Christum Fílium tuum. Amen.

\begin{tightlist}{}
\item \textred{A Purificatione usque ad Adventum:}
\item \textred{℣.} Dignáre me laudáre te, Virgo sacráta. \textred{℟.} Da mihi virtútem contra hostes tuos.
\item Orémus. \hfill \textred{Oratio.}
\end{tightlist}
\lettrine{C}oncéde miséricors Deus fragilitáti nostræ præsídium; ut qui sanctæ Dei Genitrícis memóriam ágimus, intercessiónis ejus auxílio, a nostris iniquitátibus resurgámus. Per eúmdem Christum Dóminum nostrum. Amen.

\subsubsection*{\center \textsc{DE S. JOSEPH}}
\begin{tightlist}{}
\item \textred{Fit post commem. B.V. Mariæ et ante quoscumque alios Sanctos, exceptis Angelis et S. Joanne Baptista.}
\item \textred{Ad Vesperas, Aña.} Ecce fidélis servus et prudens, quem constítuit Dóminus super famíliam suam.
\item \textred{℣.} Glória et divítiæ in domo ejus. \textred{℟.} Et justítia ejus manet in sǽculum sǽculi.
\item \textred{Ad Laudes, Aña.} Ipse Jesus erat incípiens quasi annórum trigínta ut putabátur fílius Joseph.
\item \textred{℣.} Os justi meditábitur sapiéntiam. \textred{℟.} Et lingua ejus loquétur judícium.
\item Orémus. \hfill \textred{Oratio.}
\end{tightlist}
\lettrine{D}eus, qui ineffábili providéntia beátum Joseph sanctíssimæ Genitrícis tuæ sponsum elígere dignátus es: præsta quǽsumus, ut quem protectórem venerámur in terris intercessórem habére mereámur in cœlis. Qui vivis et regnas in sǽcula sæculórum. Amen.

\subsubsection*{\center \textsc{DE SANCTIS}}
\begin{center}\textsc{Nostri Ordinis.}\end{center}
\begin{tightlist}{}
\item \textred{Ad Vesperas et Laudes, Aña.} Christi pia grátia Sanctos sublimávit, quos beátæ Vírginis Ordo propagávit. Nos eórum méritis pétimus adjuvári, atque suis précibus Deo commendári.
\item \textred{℣.} Sancti tui, Dómine, benedícent te. \textred{℟.} Glóriam regni tui dicent, et poténtiam tuam loquéntur.
\item Orémus. \hfill \textred{Oratio.}
\end{tightlist}
\lettrine{O}mnípotens et clementíssime Deus, qui Montis Carméli Ordinem gloriosíssimæ Vírginis Maríæ, Genitrícis Fílii tui Dómini nostri Jesu Christi, sacráto título insignítum, sanctórum tuórum Patris nostri Elíæ et Eliséi Prophetárum, Angeli et Anastásii Mártyrum, Cyrílli et Albérti Confessórum, Euphrásiæ, Terésiæ, et Maríæ Magdalénæ Vírginum, et aliórum plurimórum méritis decorásti: tríbue nobis quǽsumus; ut per eórum suffrágia, ab instántibus malis ánimæ et córporis liberáti, ad te verum Carméli vérticem gaudéntes perveníre valeámus. Per eúmdem Christum Dóminum nostrum. Amen.

\subsubsection*{\center \textsc{DE OMNIBUS SANCTIS.}}
\begin{tightlist}{}
\item \textred{Ad Vesperas et Laudes, Aña.} Omnes elécti Dei, nostri memorámini ante Deum, ut vestris précibus adjúti mereámur vobis adjúngi.
\item \textred{℣.} Oráte pro nobis omnes Sancti Dei. \textred{℟.} Ut digni efficiámur promissiónibus Christi.
\item Orémus. \hfill \textred{Oratio.}
\end{tightlist}
\lettrine{O}mnium Sanctórum tuórum, quǽsumus Dómine, intercessióne placátus, et véniam nobis delictórum tríbue, et remédia sempitérna concéde. Per Christum Dóminum nostrum. Amen.

\subsubsection*{\center \textsc{DE PACE.}}
\begin{tightlist}{}
\item \textred{Ad Vesperas et Laudes, Aña.} Da Pacem, Dómine, in diébus nostris: quia non est álius qui pugnet pro nobis nisi tu Deus noster.
\item \textred{℣.} Fiat pax in virtúte tua. \textred{℟.} Et abundántia in túrribus tuis.
\item Orémus. \hfill \textred{Oratio.}
\end{tightlist}
\lettrine{D}eus, a quo sancta desidéria, recta consília, et justa sunt ópera: da servis tuis illam, quam mundus dare non potest, pacem; ut et corda nostra mandátis tuis dédita, et hóstium subláta formídine, témpora sint tua protectióne tranquílla. Per Dóminum nostrum.

\begin{center}\textcolor{red}{\rule{\linewidth}{0.5pt}} \end{center}

\subsection*{\center \textsc{TEMPORE PASCHALI}}
\fancyhead[C]{\textsc{Commemorationes Communes - T.P.}}
\begin{tightlist}{}
\item \textred{Solum fit commem. sequens de Cruce, quæ dicatur etiam in Dominicis, et Festis Semiduplicibus, Simplicibus, et Simplicissimis, extra Octavas celebratis.}
\item \textred{Ad Vesperas et Laudes, Aña.} Crucifíxus surréxit a mórtuis, et redémit nos Sánguine suo, allelúja.
\item \textred{℣.} Dícite in natiónibus. \textred{℟.} Quia Dóminus regnávit a ligno, allelúja.
\item Orémus. \hfill \textred{Oratio.}
\end{tightlist}
\lettrine{D}eus, qui pro nobis Fílium tuum Crucis patíbulum subíre voluísti, ut inimíci a nobis expélleres potestátem; concéde nobis fámulis tuis, ut Resurrectiónis grátiam consequámur. Per eúmdem Dóminum nostrum.

\begin{center}\textcolor{red}{\rule{\linewidth}{0.5pt}} \end{center}

\end{multicols}
\vfill
\begin{center}
\textcolor{red}{\rule{0.8\textwidth}{0.5pt}}
\end{center}

\cleardoublepage

\section*{\center \textsc{PRECES VESPERTINÆ}}
\fancyhead[C]{\textsc{Preces Vespertinæ OCD}}
\subsection*{\center \textsc{Ad usum Carmelitarum Discalceatorum}}
\subsubsection*{\center \textred{Post examen conscientiæ vel communem disciplinam quotidie recitandæ, præside incipiente.}}

\begin{multicols}{2}
\subsubsection*{\center \textsc{De Beata Virgine}}
\begin{tightlist}{}
\item \textred{Ant.} Ait Dóminus Deus * ad serpentem: Inimicítias ponam inter te et mulíerem, et semen tuum et semen illíus: Ipsa cónteret caput tuum.
\item {\small \textred{Et dicitur ter:} Ave, María.}
\item \textred{℣.} Tota pulchra es, María. \textred{℟.} Et mácula originális non est in te.
\item Orémus. \hfill \textred{Oratio}
\end{tightlist}
\lettrine{D}eus, qui per immaculátam Vírginis Conceptiónem dignum Fílio tuo habitáculum præparásti: quǽsumus; ut, qui ex morte ejúsdem Fílii tui prævísa, eam ab omni labe præservásti, nos quoque mundos ejus intercessióne ad te perveníre concédas.

\lettrine{E}cclésiæ tuæ, quǽsumus, Dómine, preces placátus admítte: ut, destrúctis adversitátibus et erróribus univérsis, secúra tibi sérviat libertáte.

\lettrine{D}eus, ómnium fidélium pastor et rector, fámulum tuum N., quem pastórem Ecclésiæ tuæ præésse voluísti, propítius réspice: da ei, quǽsumus, verbo et exémplo quibus præest profícere; ut ad vitam, una cum grege sibi crédito, pervéniat sempitérnam. Per Christum Dóminum nostrum.\\
\textred{℟.} Amen.

\begin{center}\textcolor{red}{\rule{\linewidth}{0.5pt}} \end{center}

\subsubsection*{\center \textsc{De S. P. N. Joseph}}
\begin{tightlist}{}
\item \textred{Ant.} Joseph, fili David, * noli timére accípere Maríam cónjugem tuam: quod enim in ea natum est, de Spíritu Sancto est: páriet autem Fílium et vocábis nomen ejus Jesum.
\item \textred{℣.} Ora pro nobis, sancte Pater Joseph. \textred{℟.} Ut digni efficiámur promissiónibus Christi.
\item Orémus. \hfill \textred{Oratio}
\end{tightlist}
\lettrine{S}anctíssimæ Genitrícis tuæ sponsi, quǽsumus, Dómine, méritis adjuvémur: ut, quod possibílitas nostra non óbtinet, ejus nobis intercessióne donétur: Qui vivis et regnas in sǽcula sæculórum. \textred{℟.} Amen.

\begin{center}\textcolor{red}{\rule{\linewidth}{0.5pt}} \end{center}

\subsection*{\center \textsc{\textred{Officia Mensualia}}}
\subsubsection*{\center \textred{Dicenda singulis mensibus}}

\begin{center}\textcolor{red}{\rule{\linewidth}{0.5pt}} \end{center}

\subsubsection*{\center \textsc{De S. P. N. Elia}}
\begin{tightlist}{}
\item \textred{Ant.} Sancte Prophéta Dei Elía, * Carmelitárum Dux et Pater: intercéde pro nostra omniúmque salúte.
\item \textred{℣.} Ora pro nobis, sancte Pater Elía. \textred{℟.} Ut digni efficiámur promissiónibus Christi.
\item Orémus. \hfill \textred{Oratio}
\end{tightlist}
\lettrine{P}ræsta, quǽsumus, omnípotens Deus: ut, qui beátum Elíam Prophétam tuum et Patrem nostrum, ígneo curru mirabíliter elevátum esse crédimus; eódem interveniénte ad cæléstia sublevémur, et Sanctórum tuórum consórtio gaudeámus. Per Christum Dóminum nostrum.\\ 
\textred{℟.} Amen.

\begin{center}\textcolor{red}{\rule{\linewidth}{0.5pt}} \end{center}

\subsubsection*{\center \textsc{De S. M. N. Teresia}}
\begin{tightlist}{}
\item \textred{Ant.} Sancta Mater Terésia, * réspice de cælo, et vide, et vísita víneam istam: et pérfice eam, quam plantávit déxtera tua.
\item \textred{℣.} Ora pro nobis, sancta Mater Terésia. \textred{℟.} Ut digni efficiámur promissiónibus Christi.
\item Orémus. \hfill \textred{Oratio}
\end{tightlist}
\lettrine{E}xáudi nos, Deus salutáris noster: ut sicut de beátæ Terésiæ Vírginis tuæ et Matris nostræ Commemoratióne gaudémus; ita cæléstis ejus doctrínæ pábulo nutriámur, et piæ devotiónis erudiámur afféctu. Per Christum Dóminum nostrum. \textred{℟.} Amen.

\begin{center}\textcolor{red}{\rule{\linewidth}{0.5pt}} \end{center}

\subsubsection*{\center \textsc{De S. P. N. Joanne}}
\begin{tightlist}{}
\item \textred{Ant.} Fílii, * confortámini, et viríliter ágite in lege: quia in ipsa gloriósi éritis.
\item \textred{℣.} Ora pro nobis, sancte Pater Joánnes. \textred{℟.} Ut digni efficiámur promissiónibus Christi.
\item Orémus. \hfill \textred{Oratio}
\end{tightlist}
\lettrine{D}eus, qui sanctum Joánnem Confessórem tum atque Doctórem, Patrem nostrum, perféctæ sui abnegatiónis, et crucis amatórem exímium effecísti: concéde; ut ejus imitatióni júgiter inhæréntes, glóriam assequámur ætérnam. Per Christum Dóminum nostrum. \textred{℟.} Amen.

\begin{center}\textcolor{red}{\rule{\linewidth}{0.5pt}} \end{center}
\end{multicols}

\end{document}